% Landau Geometry of Computational Resources -- MSCS submission source.
% Compile with pdfLaTeX until references stabilize (normally three passes).
% Self-contained: bibliography and diagrams are embedded; no shell escape.
% Journal: Mathematical Structures in Computer Science.

\documentclass[11pt,a4paper]{article}
\usepackage[T1]{fontenc}
\usepackage[utf8]{inputenc}
\usepackage{lmodern}
\usepackage[margin=29mm]{geometry}
\usepackage{amsmath,amssymb,amsthm,mathtools,mathrsfs}
\usepackage{microtype}
\usepackage[authoryear,round]{natbib}
\setcitestyle{aysep={},yysep={,},notesep={, }}
\usepackage{enumitem,booktabs,array,longtable}
\usepackage{tikz-cd}
\usepackage[colorlinks=true,linkcolor=blue!45!black,citecolor=blue!45!black,urlcolor=blue!45!black]{hyperref}
\usepackage[nameinlink,noabbrev,capitalise]{cleveref}
\hypersetup{pdftitle={Landau Geometry of Computational Resources},pdfauthor={Conor Sheridan}}
\setlist{itemsep=2pt,topsep=5pt}
\setlist[enumerate,1]{label=(\arabic*),ref=\arabic*}
\setlist[enumerate,2]{label=(\alph*),ref=\theenumi(\alph*)}
\setlength{\emergencystretch}{2em}
\numberwithin{equation}{section}
\newtheorem{theorem}{Theorem}[section]
\newtheorem{proposition}[theorem]{Proposition}
\newtheorem{lemma}[theorem]{Lemma}
\newtheorem{corollary}[theorem]{Corollary}
\theoremstyle{definition}
\newtheorem{definition}[theorem]{Definition}
\newtheorem{example}[theorem]{Example}
\theoremstyle{remark}
\newtheorem{remark}[theorem]{Remark}
\theoremstyle{plain}
% Computational domains, classified partiality, and observations.
\newcommand{\E}{\mathcal E}
\newcommand{\N}{\mathbb N}
\newcommand{\QQ}{\mathbb Q}
\newcommand{\OO}{\mathcal O}
\newcommand{\Pow}{\mathcal P}
\newcommand{\Idl}{\operatorname{Idl}}
\newcommand{\Sc}{\operatorname{Sc}}
\newcommand{\Pt}{\operatorname{Pt}_{\Sigma}}
\newcommand{\ED}{\mathbf{EDcpo}_{\Sigma}}
\newcommand{\AS}{\mathcal A_{\Sigma}}
\newcommand{\J}{J_{\Sigma}}
\newcommand{\F}{F_{\Sigma}}
\newcommand{\LL}{L_{\Sigma}}
\newcommand{\LC}{L_{\Sigma,C}}
\newcommand{\Eff}{\mathbf{Eff}}
\newcommand{\id}{\operatorname{id}}
\newcommand{\Id}{\operatorname{Id}}
\newcommand{\ev}{\operatorname{ev}}
\newcommand{\way}{\mathrel{\ll}}
\newcommand{\incl}{\hookrightarrow}
\newcommand{\onto}{\twoheadrightarrow}
\newcommand{\pmap}{\rightharpoonup}
\newcommand{\llbracket}{\mathopen{[\![}}
\newcommand{\rrbracket}{\mathclose{]\!]}}
% Q is the nonnegative extended lower-real domain.
\newcommand{\Q}{\mathsf Q}
\newcommand{\Qeff}{\Q_{\Sigma}}
\newcommand{\Lan}{\operatorname{Lan}}
\newcommand{\JQ}{J^{\Q}_{\Sigma}}
\newcommand{\Sub}{\operatorname{Sub}}
\newcommand{\Loc}{\operatorname{Loc}}
\newcommand{\eps}{\varepsilon}
\setcounter{tocdepth}{2}
\makeatletter
\renewcommand*\l@section{\@dottedtocline{1}{0em}{1.7em}}
\makeatother
\newcommand{\Pos}{\operatorname{Pos}}
\newcommand{\cl}{\operatorname{cl}}
\newcommand{\pair}[2]{\langle #1,#2\rangle}
\newcommand{\bind}{\mathbin{\mathrm{bind}}}


% Continuous linear structure and its budget augmentation.
\newcommand{\LinC}{\mathsf{Lin}^{\J}_S}
\newcommand{\BudC}{\mathsf{Bud}^{\J}_S}
\newcommand{\FL}{\mathsf F}
\newcommand{\GL}{\mathsf G}
\newcommand{\LLin}{\mathsf L}
\newcommand{\pres}{\mathsf{Pres}_{\J}}
\newcommand{\aug}{\operatorname{aug}}
\tikzcdset{arrow style=tikz,diagrams={>=stealth}}


% Landau notation: \QQ denotes rationals,
% \Q denotes nonnegative extended lower reals, and \Qeff denotes effective scalars.
\newcommand{\Qplus}{{\QQ_{\ge0}}}
\newcommand{\A}{{\mathscr A_+}}
\newcommand{\C}{{\LinC}}
\newcommand{\Dcat}{{\mathscr A_{\pm}}}
\newcommand{\Pcat}{\mathscr P}
\newcommand{\Ext}{\mathsf E_S}
\newcommand{\Gr}{\operatorname{Gr}}

\newcommand{\Emb}{\mathsf{Emb}}
\newcommand{\Wit}{\mathsf{Wit}}
\newcommand{\Feas}{\mathsf{Feas}}
\newcommand{\DCx}{\mathsf{C}}
\newcommand{\Sym}{\operatorname{Sym}}
\newcommand{\Mod}{\operatorname{Mod}}
\newcommand{\CAlg}{\operatorname{CAlg}}
\newcommand{\Hom}{\operatorname{Hom}}
\newcommand{\Der}{\operatorname{Der}}
\newcommand{\Ev}{{\operatorname{Ev}}}
\newcommand{\im}{\operatorname{im}}
\newcommand{\abs}[1]{\lvert #1\rvert}
\newcommand{\CostClass}{\mathcal C_{\mathrm{cost}}}
\newcommand{\MComp}{\operatorname{MComp}}
\newcommand{\Weil}{\mathsf W}
\newcommand{\Geo}{\mathcal G_B}
\newcommand{\Rline}{\mathcal R_B}
\newcommand{\Jac}{\operatorname{Jac}}
\newcommand{\Ann}{\operatorname{Ann}}
\newcommand{\gr}{\operatorname{gr}}
\newcommand{\coeq}{\operatorname{coeq}}
\newcommand{\Lamlin}{\Lambda_{\mathrm{lin}}}
\newcommand{\Gamlin}{\Gamma_{\mathrm{lin}}}
\newcommand{\Shape}{\mathfrak I_W}
\newcommand{\InfShape}{\mathfrak I_{\mathrm{inf}}}
\newcommand{\AffB}{\mathsf{Aff}^{\mathrm{fp}}_B}
\newcommand{\InfB}{\mathsf I_B}
\newcommand{\BaseEt}{\mathscr X_{B,\mathrm{et}}}
\newcommand{\ThickEt}{\mathscr H_{B,\mathrm{et}}}
\newcommand{\StEt}{\operatorname{St}_{\mathrm{et}}}
\newcommand{\ShEt}{\operatorname{Sh}_{\mathrm{et}}}
\newcommand{\oeff}{o_{\mathrm{ev}}}
\newcommand{\onn}{o_{\neg\neg}}
\newcommand{\Cplx}{\mathfrak C_{\Ev}}
\newcommand{\Up}{\operatorname{Up}}
\newcommand{\Sh}{\operatorname{Sh}}
\theoremstyle{definition}
\newtheorem{construction}[theorem]{Construction}
\crefname{construction}{Construction}{Constructions}
\Crefname{construction}{Construction}{Constructions}
\theoremstyle{plain}
\newcolumntype{L}[1]{>{\raggedright\arraybackslash}p{#1}}
\title{Landau Geometry of Computational Resources\\
\large Effective semantics, formal geometry, and crystalline resource calculus}
\author{Conor Sheridan\\
  {\normalsize\href{mailto:siridconor@gmail.com}{\texttt{siridconor@gmail.com}}}}
\date{October 2026}

\begin{document}
\maketitle
\begin{abstract}
We construct foundations for a differential geometry of computational
resources in effective domain and realizability settings. Exact rational
cost profiles give a ring of bounded coefficients and an algebra of
eventual germs. Embedded resource modules represent asymptotic budgets;
a coherent locale classifies their families, and measured execution
supplies the implementations witnessing feasibility. Compactness and
supercompactness characterize finite dominating portfolios and optimal
implementations.

Polynomial event valuations connect execution to differential geometry:
their measured Taylor coefficients agree with geometric jets. Embedded
budgets have formal deformation classifiers, and finite systems of budget
and valuation constraints have intrinsic lifting obstructions. Overhead
localization acts on these obstructions through their annihilators.

The primary infinitesimal geometry represents formal principal parts and
completed diagonal cochains. An explicit contraction compares the latter
with de Rham forms; integrable connections yield stratifications and
crystalline jet coalgebras. The comparisons retain resource parameters,
exact coefficients, and effective finite observations. A measured family
exhibits polynomial bounds at every finite order without one polynomial
budget for the complete coefficient family; its generating function also
provides a nonzero logarithmic de Rham class and crystalline transport.
Event instrumentations exhibit curvature, while logarithmic resource
connections distinguish asymptotic pole orders.
\end{abstract}

\medskip
\noindent\textbf{Keywords:} resource semantics; asymptotic complexity;
effective domains; categorical semantics; synthetic differential geometry;
deformation theory.
\par\noindent\textbf{MSC 2020:} Primary 68Q55; Secondary 18F10, 14B10.

\section*{Introduction}
\addcontentsline{toc}{section}{Introduction}
A differential geometry of computational resources must retain both the
implementations that realize a bound and the coefficients through which
that bound varies. We construct such a geometry from measured execution
and asymptotic resource modules. Its formal neighbourhoods, differential
forms, connections, and jets are related by explicit comparison maps;
their finite coefficients remain tied to the underlying measurements.

The starting point is the relation $f=O(g)$. For rational profiles on a
filtered input object, let $K$ be the algebra of eventual germs and
$B=O(1)$ its ring of bounded coefficients. Bounded division identifies
$f=O(g)$ with membership in the embedded module $Bg$. The embedding
records the costs admitted by a budget, while the exact coefficients
support addition, multiplication, and differentiation. Growth ideals
extend principal budgets to families with several or indefinitely many
bounds.

\paragraph{Four connected contributions.}
\begin{enumerate}
\item \emph{Resource spaces and computational witnesses.}
The growth lattice $L$ constructs a coherent locale $\Cplx$ with
\[
 \Loc_\E(Z,\Cplx)\cong\Sub_{\Delta_ZB}(\Delta_ZK).
\]
A generalized point therefore specifies a family of embedded budgets.
A measured span $V\xleftarrow{p}P\xrightarrow{\Phi}K$ supplies the
implementations satisfying them. Its witness sheaf is a left Kan
extension; its feasible open is supercompact exactly when an optimal
implementation exists, and compact exactly when a finite dominating
portfolio exists
(\cref{thm:locale-modules,thm:universal-cost,prop:cat-implementation-kan,thm:optimization-compactness}).

\item \emph{Measured variation and resource deformation.}
Folding an execution history with truncated Taylor-valued event weights
computes the geometric jet of its cost family
(\cref{thm:measured-jet-adequacy}). For a principal marked budget $M$,
finite logarithm identifies its deformation functor with the formal
module determined by $T_M=\Hom_B(M,K/M)$. A feasible valuation $F(x)$
and a moving budget have linearized incidence
\[
 (v,\theta)\longmapsto qDF_x(v)-\theta(F(x)),\qquad q:K\onto K/M.
\]
Its kernel describes tangent directions and differences between lifts;
its cokernel contains successive lifting obstructions. The construction includes finite budget diagrams
and their overhead localization
(\cref{thm:linearized-deformations,thm:moving-budget-obstructions,prop:formal-flag-obstruction,thm:gauge-obstruction-criterion}).

\item \emph{Formal and de Rham resource geometry.}
The primary geometry is an internal \'etale site of nilpotent arrows.
Representability identifies its infinitesimal function objects with
principal parts and completed diagonal cochains. A finite equivariant
contraction compares those cochains with algebraic de Rham forms,
including localized coordinate cells. This constructs the differential
theory over the exact resource coefficients and their classified
parameters
(\cref{thm:etale-endpoint-cohesion,thm:primary-synthetic-cochains,thm:de-rham-comparison,thm:localized-formal-derham}).

\item \emph{Crystalline transport and effective coefficients.}
Integrable connections are equivalent to formal stratifications.
Principal-parts correspondences construct jets and differential
operators; their completed system is the primary jet comonad, with
algebraic stratifications giving its crystalline coalgebras. Effective
finite observations and exact Taylor coefficients agree with these
geometric constructions
(\cref{thm:connection-stratification,thm:completed-jet-comparison,prop:crystalline-jet-coaction,prop:formal-finite-output-tracking}).
\end{enumerate}

\paragraph{Resource questions through the geometry.}
Finite-order feasibility asks which variations preserve a budget;
completion asks whether compatible finite variations admit a common
resource bound. De Rham forms express infinitesimal variation, and
connections test its compatibility with transport. Crystalline jets
collect the response coefficients and specify how they re-centre and
compose. These questions concern the same measured families.

For example, a uniformly constructed family of finite programs has
weighted costs $F_r(z)=\sum_{j=0}^r(n+1)^jz^j$. Every finite coefficient
list has a polynomial budget, but the compatible formal family has no
single polynomial growth degree. Its generating function gives a
logarithmic de Rham cocycle, an integrable connection, and crystalline
transport
(\cref{ex:measured-formal-family,prop:measured-logarithmic-transport,ex:measured-crystalline-family}).
Two event instrumentations detect order through a mixed jet and its
curvature (\cref{prop:instrumentation-curvature}). Coupled budgets show
how an obstruction can persist when each budget separately moves
(\cref{ex:coupled-budget-obstruction}), and logarithmic residues
distinguish asymptotic pole orders
(\cref{cor:logarithmic-residue-invariance}).

\paragraph{Organization and mathematical settings.}
Sections~\ref{sec:foundation}--\ref{sec:feasibility-moduli} construct
the effective coefficient setting, classified resources, and measured
feasibility. Sections~\ref{sec:relative-geometry}--\ref{sec:resource-deformations}
connect history valuations to differential coefficients and constrained
resource deformations. Sections~\ref{sec:primary-geometry}--\ref{sec:jets}
construct the geometric function objects, completion, de Rham theory,
and crystalline jets, developing the measured examples through those
comparisons. Section~\ref{sec:instances} gives numerical execution,
size bounds, certificates, and logarithmic resource connections.
Table~\ref{tab:settings} records the coefficient settings.

\begin{table}[htbp]
\centering
\caption{Objects and their roles in the construction.}
\label{tab:settings}
\begin{tabular}{@{}L{.23\textwidth}L{.34\textwidth}L{.35\textwidth}@{}}
\toprule
Notation & Setting & Role\\
\midrule
$\E,\ \Sigma\incl\Omega$ &
Effective or relative realizability; finite-hit truth values &
Computational names, partial returns and positive observations\\
$\J,\ \AS$ &
Selected ideal completion and its algebras &
Admitted directed suprema and recursive execution\\
$S=\Qeff$ &
Effective nonnegative extended lower reals &
Continuous quantitative interpretation\\
$K,\ B,\ \mathfrak m$ &
Exact rational germs, bounded germs, vanishing germs &
Signed resource coefficients; $O(g)=Bg$, $o(g)=\mathfrak m g$\\
$L,\ \Cplx$ &
Growth lattice and budget locale &
Ordinary families of marked resources\\
$\mathscr S=\Sh_\E(\Cplx)$ &
Budget sheaves, equivalently $[L,\E]$ &
Universal budget and implementation witnesses\\
$\ThickEt,\ \mathcal O_{\mathrm{inf}}$ &
Primary \'etale infinitesimal-arrow sheaves &
Represented functions, formal nerve, de Rham comparison and jets\\
$\Geo,\ j_BM$ &
Based projective Weil diagrams &
Finite resource coefficients and marked deformations\\
\bottomrule
\end{tabular}
\end{table}

\paragraph{Conventions used throughout.}
\label{sec:global-conventions}
Internal existence means a regular image. Formulas on local
representatives define maps on quotients after their kernel-pair
equations have been proved. Finite projectivity means a displayed
finite-free retract, as specified in Section~\ref{sec:foundation}.
Semantic incidence has values in $\Omega$; positive observation
requires its stated factorization through $\Sigma$. A formal
coefficient family is an internal compatible family of finite stages;
a uniform computational name for it is additional data. Computational
environments and size maps are retained in implementation comparisons.
The based geometry supplies finite module coefficients; primary
representability supplies the function objects used for formal cochains.

\paragraph{Relation to published foundations and complexity theory.}
The computational construction draws on effective domain theory
\citep{Smyth1977EffectiveDomains,AbramskyJung1994DomainTheory}
and realizability
\citep{Hyland1982EffectiveTopos,BirkedalOosten2002,VanOosten2008Realizability}.
Within that setting we define the positive-observation selection
of ideal completion and prove the continuous presentation theorem
needed for the resource algebra. The latter adapts the method of
\citep{JungMoshierVickers2008PresentingDcpos} to internal selected
ideals. We also compare the selected-domain recursion with
natural-number-indexed execution. The distinction between synthetic
completeness notions and ordinary execution is discussed in
\citep{VanOostenSimpson2000SyntheticDomains}.

The complexity interpretation connects three established strands.
Blum's machine-independent measures
\citep{Blum1967MachineIndependentComplexity} provide the numerical
comparison: our first-hit construction verifies the two measure
properties, and speedup gives examples of failure of optimality.
Schellekens's complexity spaces
\citep{Schellekens1995SmythCompletion} relate topological complexity
to domain-theoretic completion. Resource monoids and
resource-bounded realizability are developed in
\citep{DalLagoHofmann2011RealizabilityComplexity}.
Here the classifier records embedded asymptotic modules and their
budget families, together with measured implementation witnesses.
The optimality criteria and resource deformations are formulated
for these objects.

The cost-sensitive adequacy method of
\citet{NiuSterlingHarper2024CostSensitiveSDT} relates abstract cost
semantics to computational execution. Their partial cost monad and
the cost--behaviour distinction in
\citep{NiuSterlingGrodinHarper2022Calf,GrodinNiuSterlingHarper2024Decalf}
provide the relevant execution-semantic antecedents. Here retained
histories give domain-and-value adequacy and valuation-independent
erasure; \cref{thm:measured-jet-adequacy} further identifies Taylor-valued
instrumentation with geometric cost jets. These are comparisons for
the constructed execution model and its polynomial valuations.
Passing from runs to size bounds is related to cost-recurrence
extraction \citep{DannerLicataRamyaa2015CostSemantics}.
Representation-sensitive operator complexity
\citep{KawamuraCook2012Operators} motivates retaining input names and
the measured overhead of changing them.

The resource algebra also connects with formal treatments of
$O$-notation \citep{AvigadDonnelly2004ONotation} and filter-based
asymptotic verification \citep{GueneauChargueraudPottier2018Fistful}.
The latter provides a comparison for both the eventual filters
and the finite resource certificates.
Asymptotic gauges and moderate rings
\citep{GiordanoLuperiBaglini2016AsymptoticGauges}, together with
their functorial changes of scale
\citep{LuperiBagliniGiordano2017CategoryColombeau}, are compared
directly in \cref{prop:asymptotic-gauge-comparison}.

Growth and derivations also occur in the asymptotic differential
algebra of
\citep{AschenbrennerVanDenDriesVanDerHoeven2017Transseries}.
The present germ rings retain zero divisors and are not treated
as ordered fields. Finally, algorithm portfolios
\citep{GomesSelman2001AlgorithmPortfolios} provide the computational
comparison for our finite-portfolio property, which is characterized
by compactness of the feasible open.


\paragraph{Geometric antecedents and the resource-specific comparisons.}
Kock's algebraic synthetic differential geometry \citep{Kock}
provides the comprehensive and Weil identities used here.
The symmetric-algebra and universal-derivation constructions belong
to the differential and codifferential categorical framework
\citep{BLWON,BCS}. The contribution in Section~\ref{sec:relative-geometry}
is their construction over the exact Landau coefficients and their
comparison with the retained quantitative presentations.
The primary infinitesimal-arrow site supplies the endpoint
adjunctions and represented function objects needed for that comparison.

Principal parts, algebraic de Rham theory and crystalline
stratifications have the classical foundations
\citep{GrothendieckDeRham,BerthelotOgus,Stacks}.
Sections~\ref{sec:formal-completion}--\ref{sec:jets}
identify their coefficient objects in the primary resource geometry
and construct a finite, equivariant contraction compatible with
completion and displayed projective retracts.
The construction of jets from infinitesimal shape and their
coalgebra interpretation has a direct synthetic precedent in
\citet{KS}. Our comparison specifies the
algebraic endpoint model, the finite principal-parts coefficients,
the effective-epimorphism domain of coalgebra descent, and the
crystalline algebraic sector within the full jet-coalgebra category.
The marked-resource deformation and moving-incidence calculations
of Section~\ref{sec:resource-deformations} use these exact
coefficient structures to describe variations of asymptotic budgets.

\tableofcontents

\section{Effective domains and exact quantitative presentations}
\label{sec:foundation}

The later resource constructions need three interfaces: partial and
recursive execution, exact data with a continuous quantitative
interpretation, and positive computational observations. This section
builds them in the realizability setting. The selected ideal completion
supplies the continuous operations used by execution. Exact positive
presentations then connect that setting to signed resource algebra,
while the observation and name constructions specify how its objects
are presented computationally.

\subsection{Published ambient models and internal reasoning}
We work in the effective topos and its relative-realizability variants
\citep{Hyland1982EffectiveTopos,BirkedalOosten2002,VanOosten2008Realizability},
using effective domain theory as the order-theoretic background
\citep{Smyth1977EffectiveDomains,AbramskyJung1994DomainTheory}.
The selected completion and continuous presentations required here are
constructed below. The supporting presentation and coding arguments
are given in \cref{app:presentations,app:names-codes}.

Write $\E$ for the ambient elementary topos and $\N$ for its natural
numbers object. Internal subsets, images, and quotients are those of $\E$.
An internal family includes its free parameters: $x:Z\to X$ is a
generalized element, and substitution is pullback in the corresponding
slice. Effective quotients allow a map to descend exactly when it
coequalizes the kernel pair of its presentation. A total functional
relation determines a unique arrow, because its projection is both
monic and regular epic.

The numerical realization is $\Eff=\operatorname{RT}(K_1)$, with the
usual partial-recursive application. Its oracle versions use $K_1^Z$.
For supplied function names we also use relative realizability
$\operatorname{RT}(A_\#,A)$: names belong to the ambient partial
combinatory algebra $A$, while uniform trackers belong to its elementary
subalgebra $A_\#$. This includes the usual function-realizability and
Kleene--Vesley presentations. The exact-completion presentation by
partitioned assemblies \citep[Section 10.6, especially Theorem 10.6.7]{Birkedal2000TypesComputability}
provides the projective covers used in
\cref{app:names-codes}. A partitioned assembly consists of a carrier
and its naming map into $A$; a morphism has a single permitted tracker
on all valid names.

\begin{definition}[Finite-hit truth values and overtness]
\label{def:standalone-sigma}
The Rosolini object $i:\Sigma\incl\Omega$ is the image of
\[
 h:2^\N\longrightarrow\Omega,\qquad
 h(\alpha)\ \Longleftrightarrow\ \exists n:\N\;\alpha(n)=1.
\]
A $\Sigma$-classified partial map is a span whose domain mono is
classified by $\Sigma$. An object $X$ is $\Sigma$-overt when
existential quantification restricts uniformly to
$\exists_X:\Sigma^X\to\Sigma$.
\end{definition}

\begin{proposition}[Dominance and uniform natural-number search]
\label{prop:standalone-dominance}
In the effective and relative-realizability models just described,
$\Sigma$ is closed under dependent conjunction, and $\N$ is
$\Sigma$-overt. Thus $\Sigma$ is a dominance and classified partial
maps compose. Finite products, finite sums, regular-epimorphic images,
and $\Sigma$-classified subobjects of overt objects are overt.
\end{proposition}
\begin{proof}\leavevmode
\begin{enumerate}[leftmargin=*,align=left]
\item \textbf{Construct finite truth operations.}
Constant streams and finite interleaving give false, true, finite
conjunction, and finite disjunction.

\item \textbf{Prove closure under dependent conjunction.}
To check dependent conjunction,
pull two successive classified monos back to a partitioned projective
cover $P$ of their ambient object. Lift the first classifier through
$2^\N\onto\Sigma$, obtaining a total tracked stream $\alpha_p$.
On its hit assembly
$H_\alpha=\{(p,n):\alpha_p(n)=1\}$, lift the second classifier to
a total tracked stream $\beta_{p,n}$. This is possible on the displayed
partitioned subassembly. The stream
\[
 \gamma_p(\langle n,k\rangle)=
 \begin{cases}
 \beta_{p,n}(k),&\alpha_p(n)=1,\\
 0,&\alpha_p(n)=0
 \end{cases}
\]
is total: its first branch is evaluated only on a valid hit name.
Its hit image is the composite subobject. Its classifier therefore
coequalizes the kernel pair of $P$ and descends through
$i:\Sigma\incl\Omega$. This proves dominance in every parameter context.

\item \textbf{Construct uniform natural-number search.}
For $\exists_\N$, cover $\Sigma^\N$ projectively and lift evaluation on
the product of that cover with $\N$ to streams $\alpha_{p,n}(k)$.
The single total stream
$\gamma_p(\langle n,k\rangle)=\alpha_{p,n}(k)$ represents the desired
existential. Its classifier descends for the same reason. These
constructions use the represented pairing, finite tests, and guarded
evaluation, so their trackers remain in $A_\#$ on arbitrary supplied
ambient names.

\item \textbf{Preserve overtness under finite constructions.}
Finite sums use disjunction; products and dependent sums use iterated
restricted existential quantification.

\item \textbf{Descend overtness along epimorphisms.}
For an epic $q:X\onto Y$,
$\exists_yu(y)\leftrightarrow\exists_xu(qx)$ proves the image assertion.

\item \textbf{Restrict overtness and verify coherence.}
For a classified subobject, use dependent conjunction before the
existential over its ambient overt object. The corresponding
substitution, Frobenius, and Fubini equations hold in $\Omega$ and
are reflected through the monomorphism $i$.
\end{enumerate}
\end{proof}

The subsequent internal constructions apply in the Rosolini-dominant
recursion setting thereby realized; they retain arbitrary parameter
contexts. In particular, a pointwise external list of semidecisions
is not substituted for a map into $\Sigma^\N$.
The dominance-based partial-element construction is standard
\citep{EscardoKnapp2017}; explicitly,
\begin{equation}\label{eq:standalone-lift}
 \LL Y=\sum_{p:\Sigma}Y^p.
\end{equation}
Its unit has true extent. Multiplication sends a nested partial element
with extents $p$ and $q(u)$ to extent $\exists u:p\,q(u)$ and the
corresponding value. Strength retains an unused parameter. These
formulas represent classified spans and prove the strong, commutative
monad laws by reassociation of dependent extents.

\subsection{The selected ideal completion}
Ordinary ideal completion supplies directed suprema, but the computational
construction must also track positive observations of those suprema.
We therefore select ideals by their complete profile of positive upper
tests. For an internal poset $P$, an ideal is an inhabited directed lower
subset; $\Idl P$ denotes the object of all such ideals. Put
\[
 U_\Sigma P=\Pos(\E)(P,\Sigma),\qquad
 \chi_I(u)=\exists p:P\,[I(p)\wedge u(p)].
\]
\begin{definition}[Admitted ideal]\label{def:standalone-admission}
Define $\J P$ by the pullback
\[
 \J P=\Idl P\times_{\Omega^{U_\Sigma P}}\Sigma^{U_\Sigma P}.
\]
Thus admission is uniform factorization of the complete positive
upper-observation profile. The underlying ideal membership can remain
$\Omega$-valued.
\end{definition}


The defining factorization is the cartesian square
\[
\begin{tikzcd}[column sep=large,row sep=large]
 \J P \arrow[r] \arrow[d,hook]
   & \Sigma^{U_\Sigma P} \arrow[d,hook]\\
 \Idl P \arrow[r,"{I\mapsto\chi_I}"']
   & \Omega^{U_\Sigma P}.
\end{tikzcd}
\]
The lower arrow is the complete upper-test profile; admission
is its uniform factorization through $\Sigma$.


\begin{theorem}[Selected completion and its algebra category]
\label{thm:standalone-domain}
The formulas
\[
 \eta_P(p)=\downarrow p,\qquad
 (\J f)(I)=\downarrow f(I),\qquad
 \mu_P(\mathcal I)=\bigcup_{I\in\mathcal I}I
\]
define a locally monotone monad on internal posets, with
$\J\eta_P\le\eta_{\J P}$ and a monad inclusion $\J\incl\Idl$.
It preserves finite products and disjoint finite sums.
Its algebra category $\AS=\operatorname{Alg}(\J)$ is cartesian closed;
$[A,B]_\Sigma$ is the poset of maps preserving admitted suprema,
with pointwise order and pointwise admitted suprema.
\end{theorem}
\begin{proof}\leavevmode
\begin{enumerate}[leftmargin=*,align=left]
\item \textbf{Admit the monad operations.}
For principal ideals $\chi_{\downarrow p}(u)=u(p)$. For lower images,
$\chi_{\downarrow f(I)}(u)=\chi_I(u f)$. For multiplication define
$\widehat u(I)=\chi_I(u)$, a monotone map $\J P\to\Sigma$.
Then
\[
 \chi_{\mu_P(\mathcal I)}(u)=\chi_{\mathcal I}(\widehat u)\in\Sigma.
\]
This proves closure of all three operations.

\item \textbf{Verify the monad laws and KZ structure.}
Their laws are the
lower-image and union identities in ordinary ideal completion, reflected
through the monic inclusion. If $K\in(\J\eta_P)(I)$, then
$K\subseteq\downarrow p\subseteq I$ for some $p\in I$; this proves
the displayed KZ inequality. Its algebra equations identify an algebra
map $a:\J A\to A$ with the supremum of each admitted ideal, and its
morphisms with maps preserving these suprema.

\item \textbf{Preserve finite products.}
An ideal of $P\times Q$ is the product of its two lower-image
projections: directedness supplies a common bound for any two
coordinate witnesses. Conversely
\[
 \chi_{I\times J}(u)
 =\chi_I\bigl(p\mapsto\chi_J(q\mapsto u(p,q))\bigr)
\]
proves admission of that product. The terminal case follows from
the unique ideal of $1$.

\item \textbf{Preserve disjoint finite sums.}
An inhabited directed ideal of $P+Q$
lies in a unique summand; the local summand choice descends by
uniqueness. Extending a test by false on the other summand proves
admission in both directions. The empty poset has no ideals.

\item \textbf{Construct the exponential algebra.}
Let $H\subseteq B^A$ be the object of monotone supremum-preserving
maps and $I\in\J H$. Define
$f_I(x)=\sup_{h\in I}h(x)$. For $J\in\J A$, the two iterated suprema
over $I$ and $J$ coincide: both are the supremum of the admitted
lower image of $I\times J$ under evaluation. Thus $f_I$ preserves
the supremum of $J$, and the pointwise operation makes $H$ an algebra.

\item \textbf{Verify evaluation and currying.}
The same calculation proves continuity of evaluation.
For $k:C\times A\to B$ continuous, the ideal
$\downarrow c\times J$ proves continuity of $k(c,-)$ in an arbitrary
context $c:C$. The ideals $I\times\downarrow x$ prove continuity
of its transpose $C\to H$. The ordinary evaluation and transposition
identities restrict, giving the exponential universal property.

\item \textbf{Retain parameter dependence.}
Every construction used internal images, pullbacks, or uniquely
specified suprema, so its equations retain their free parameters.
\end{enumerate}
\end{proof}


\paragraph{The Kock--Z\"oberlein structure.}
The inequality $\J\eta\le\eta\J$ identifies the selected completion
as a Kock--Z\"oberlein monad on the locally ordered category of internal
posets.  Its algebra structure is characterized by the adjunction
\[
 a_A:\J A\rightleftarrows A:\eta_A,\qquad
 a_A\dashv\eta_A,\qquad a_A\eta_A=\id_A.
\]
Indeed $a_A(I)\le x$ holds exactly when $I\subseteq\downarrow x$.
This characterizes each admitted supremum by its universal property.
The algebra and algebra-morphism equations are therefore the following
commutative diagrams:
\begin{equation}\label{eq:cat-kz-algebras}
\begin{tikzcd}[column sep=large,row sep=large]
 \J\J A \arrow[r,"{\mu_A}"] \arrow[d,"{\J a_A}"']
   & \J A \arrow[d,"{a_A}"]\\
 \J A \arrow[r,"{a_A}"'] & A
\end{tikzcd}
\qquad
\begin{tikzcd}[column sep=large,row sep=large]
 \J A \arrow[r,"{\J f}"] \arrow[d,"{a_A}"']
   & \J B \arrow[d,"{a_B}"]\\
 A \arrow[r,"{f}"'] & B .
\end{tikzcd}
\end{equation}
The free-algebra adjunction gives
$\AS(\J P,A)\cong\Pos(\E)(P,U_{\AS}A)$.
It is the extension principle used below: an operation specified on
generators extends uniquely when it respects the admitted suprema.
The presentation theorem will apply this principle to linear and
algebraic operations.


A discrete internal object $X$ has only singleton inhabited directed
ideals. Their unique elements define its algebra $\operatorname{Disc}X$.
Every underlying map from it to an algebra preserves those ideals;
hence $\operatorname{Disc}\dashv U_{\AS}$.
A preorder is treated by its effective antisymmetrization before applying
$\J$; monotone maps and the displayed ideal formulas descend.

\begin{proposition}[Lifting and least solutions]\label{prop:standalone-recursion}
The partiality monad $\LL$ lifts to a strong commutative monad on $\AS$.
For every pointed algebra $A$,
\[
 \operatorname{lfp}_A:[A,A]_\Sigma\longrightarrow A,\qquad
 f\longmapsto\bigvee_{n:\N}f^n(\bot)
\]
is an algebra morphism selecting the least fixed point. The construction
commutes with parameters and strict intertwining maps.
\end{proposition}
\begin{proof}\leavevmode
\begin{enumerate}[leftmargin=*,align=left]
\item \textbf{Construct the distributive-law map.}
Order partial elements by extension and increase of their values.
For $I\in\J(\LL P)$, its extent is
$p_I=\chi_I(\operatorname{def})\in\Sigma$. Under this extent form the
lower ideal $V_I$ generated by the defined values of members of $I$.
It is inhabited and directed. For $u:P\to\Sigma$ monotone, the
positive extension
$\widetilde u(z)=\exists d:\operatorname{def}(z)\,u(z(d))$
is $\Sigma$-valued by dominance, and
$\chi_{V_I}(u)=\chi_I(\widetilde u)$ under $p_I$.
Consequently $(p_I,V_I)$ defines
$\lambda_P:\J\LL P\to\LL\J P$.

\item \textbf{Verify the distributive-law identities.}
Principal ideals and total partial elements give the two unit
identities. Flattening either nested ideals or nested extents gives
the same existential extent and the same ideal of defined values;
these are the two multiplication identities. Thus $\lambda$ is a
distributive law, and $\LL a\lambda_A$ is the lifted algebra structure.

\item \textbf{Lift strength and commutativity.}
Strength is continuous because, once one directed member is defined,
the defined members are cofinal among both relevant coordinate
families. Its equations, and commutativity, follow on extents and values.

\item \textbf{Construct the least fixed point.}
An increasing sequence $(x_n)$ generates an admitted lower ideal:
its pairing with $u$ is $\exists_nu(x_n)$, using
\cref{prop:standalone-dominance}. This constructs the displayed
least fixed point. Continuity gives the fixed-point equation;
induction bounds all its iterates by any pre-fixed point.

\item \textbf{Prove continuity of least fixed points.}
For an admitted ideal $I$ of endomorphisms put
$g=\sup I$ and $y=\sup_{h\in I}\operatorname{lfp}(h)$.
Given $f,h\in I$, a common upper member $k$ gives
$f(\operatorname{lfp}(h))\le k(\operatorname{lfp}(k))\le y$.
Continuity and the pointwise supremum give $g(y)\le y$.
Leastness and monotonicity yield
$\operatorname{lfp}(g)=y$, proving the algebra-morphism assertion.

\item \textbf{Verify naturality.}
Strict intertwining identifies the zeroth iterates and then all
successors; preservation of their supremum proves naturality.
\end{enumerate}
\end{proof}

\begin{lemma}[Exact finite return]\label{lem:standalone-return}
For $f:X\to\LL(Y+X)$, let $D_Y,D_X$ be its exit and continuation
domains, with source $c_X:D_X\incl X$, continuation $n:D_X\to X$,
and exit value $v:D_Y\to Y$.
Define $R_0=D_Y$ and $R_{k+1}=D_X\times_XR_k$, with source
$\rho_{k+1}=c_X\pi_1$ and return value $v_{k+1}=v_k\pi_2$.
The map $\sum_{k:\N}R_k\to X$ is a $\Sigma$-classified mono.
Its returned partial map is the least solution of
$a(x)=f(x)\bind[\eta,a]$ (with the usual Kleisli ordering).
For algebra morphisms $f$, it agrees with
\cref{prop:standalone-recursion}.
\end{lemma}
\begin{proof}\leavevmode
\begin{enumerate}[leftmargin=*,align=left]
\item \textbf{Construct the classified return domain.}
Pullback and composition make each $\rho_k$ a classified mono.
Exit and continuation are disjoint, and
$R_{k+1}\times_XR_{l+1}\cong
D_X\times_X(R_k\times_XR_l)$.
Induction proves disjointness of distinct lengths. Their single
internally indexed dependent sum is therefore monic; its classifier
is $\exists k\,\rho_k(x)$ and lies in $\Sigma$.

\item \textbf{Identify the least solution.}
The $j$th bottom approximant has domain $\sum_{k<j}R_k$ and
value $v_k$ on its $k$th component. The union of these domains
gives the asserted partial map. Each solution agrees with $v_k$
there by induction, so this map is least.

\item \textbf{Compare continuous recursion and parameters.}
For continuous $f$,
the lifted supremum has the same extent and stabilized value.
All recursions and images commute with parameters.
\end{enumerate}
\end{proof}

The completion now provides the domains on which the computational
operations act. To relate their semantic elements to effective data,
we next specify the available observations and the names that present
the resulting effective points.

\subsection{Observations, effective points, and names}
\label{sec:standalone-observations}
An observation structure on a poset is an internal family of upper
predicates containing the $\Sigma$-constants and closed under finite
meets and overt-indexed unions. Its generated closure
$\cl_\Sigma(G)$ is the intersection, inside the appropriate power
object, of all such structures containing $G$.
The definition is bounded: overt subfamilies are subobjects of the
given predicate object, and overtness says that their existential
map factors through $\Sigma$.
A family in the intersection is mapped to an overt family in each
member, proving closure of the intersection.
Inverse image preserves generating observations and all these
operations.

An observed dcpo $(D,\OO_D)$ has $\OO_D\incl\Sc(D)$ an observation
structure of Scott opens. Its maps are Scott-continuous maps whose
inverse images preserve the specified observations; their category
is $\ED$. Define
\begin{equation}\label{eq:standalone-points}
 \Pt D=D\times_{\Omega^{\OO_D}}\Sigma^{\OO_D}.
\end{equation}
It has the uniform evaluation $\operatorname{Obs}_D:\Pt D\times\OO_D
\to\Sigma$.


The effective-point construction and its evaluation fit into
\[
\begin{tikzcd}[column sep=large,row sep=large]
 \Pt D \arrow[r] \arrow[d,hook]
   & \Sigma^{\OO_D} \arrow[d,hook]\\
 D \arrow[r,"{\mathrm{ev}^{\flat}}"'] & \Omega^{\OO_D}.
\end{tikzcd}
\]
Transposing the upper map gives the positive evaluation.


\begin{lemma}[Testing generators in context]\label{lem:standalone-tests}
When $\OO_D=\cl_\Sigma(G)$, a parameterized $x:Z\to D$ factors
through $\Pt D$ exactly when its evaluation on $Z\times G$
factors through $\Sigma$. The factorization is unique.
Effective points are functorial and preserve the limits constructed
with componentwise directed suprema and projection-generated observations.
\end{lemma}
\begin{proof}\leavevmode
\begin{enumerate}[leftmargin=*,align=left]
\item \textbf{Extend generator tests to all observations.}
In context $z:Z$, the upper predicates whose value at $x(z)$
lies in $\Sigma$ contain the $\Sigma$-constants and are closed
under finite meets. An overt union of such predicates has value
in $\Sigma$ by restricted existential quantification. This is
an observation structure in that context, hence contains
$\cl_\Sigma(G)$ when it contains $G$.

\item \textbf{Assemble the unique positive evaluation.}
Monicity of $\Sigma\incl\Omega$
and unique choice assemble the joint evaluation map
$Z\times\OO_D\to\Sigma$; transposition gives
\eqref{eq:standalone-points}. The reverse implication is restriction.

\item \textbf{Construct the functorial action.}
For $F:D\to D'$ use
$\operatorname{Obs}_{D'}(F x,U)=\operatorname{Obs}_D(x,F^*U)$.
Uniqueness gives the functor laws.

\item \textbf{Identify effective points of limits.}
At a limit, the generator
condition says precisely that every projected point is effective,
uniformly in the index; the pullback universal property supplies
the limit comparison.
\end{enumerate}
\end{proof}

For a poset $P$ the canonical observed free completion is
\[
 \F P=(\Idl P,\{u^\dagger:u\in U_\Sigma P\}),\qquad
 u^\dagger(I)=\exists p\,[I(p)\wedge u(p)].
\]
Restriction to principal ideals recovers $u$; dagger preserves finite
meets and overt unions, and its predicates are Scott open.
Thus the displayed family is an observation structure and
$\Pt(\F P)=\J P$ by the two defining pullbacks.
For a discrete $C$, $\Idl C\cong C$, whence
$\Pt\F(C_{\mathrm{disc}})\cong C$; its observations are exactly
the maps $C\to\Sigma$.

For completeness the semantic comparison on retracts is explicit.
A $\J$-Kleisli map $r:P\to\J Q$ extends to
$r^\#: \Idl P\to\Idl Q$, $I\mapsto\bigcup_{p\in I}r(p)$.
The inverse image of $u^\dagger$ is $v^\dagger$, where
$v(p)=\chi_{r(p)}(u)\in\Sigma$. Conversely an observed map between
these free completions restricts on principal ideals to an admitted
ideal profile and is recovered by directed extension.
Hence the free observed category is the Kleisli category of $\J$,
equivalently the category of free $\J$-algebras.
Splitting an idempotent by its fixed subobject on both sides extends
this comparison to their retracts. This is the precise comparison
used when an algebra map is represented by a semantic observed map;
a general observed domain retains its separate observation data.

In the realizability models, \cref{prop:standalone-name-lifting}
constructs regular-epimorphic name presentations $q_D:C_D\onto\Pt D$
and tracked implementations. A partial implementation retains both
the returned value equation and equality of its domain with the
pullback of the semantic domain. Names, their sizes, and their costs
are not identified merely because their denotations agree.
For represented operators this distinction is central in the
complexity theory of analysis \citep{KawamuraCook2012Operators}.

\subsection{Lower-real scalars and continuous presentations}
Resource measurements require continuous nonnegative quantities as
well as exact rational profiles. Lower cuts provide the continuous
scalars, with rational thresholds as their observations.
Let $\Q$ be the poset of rounded lower cuts of $\QQ$ containing all
negative rationals. Zero is the negative cut and infinity is the whole
rational object; suprema are unions with the negative cut.
For a cut $x$ put $R(x)=\{0\}\cup\{r\in\QQ_{>0}:r<x\}$.
This inhabited directed family gives
$x=\bigvee_{r\in R(x)}\underline r$. Define
\[
 x+y=\bigvee_{r\in R(x),\,s\in R(y)}\underline{r+s},\qquad
 xy=\bigvee_{r\in R(x),\,s\in R(y)}\underline{rs}.
\]
In particular $0+y=y$ and $0\cdot\infty=0$.
Rational interpolation identifies a positive strict threshold of
these joins with a strict threshold below one of the displayed
rational sums or products. Expanding nested operations and using
common larger rational approximants therefore reduces associativity,
distributivity, and commutativity to the rational identities.
This proof uses the always available zero approximant and makes
no decision whether a general cut is zero. The same finite-witness
calculation proves Scott continuity and arbitrary-join preservation
of multiplication in each variable (including the empty join).

Define $S=\Qeff$ by requiring that the rational strict-threshold profile
of a cut factor uniformly through $\Sigma$.
Rationals and finite rational arrays are overt: encode their numerator
and positive denominator by natural numbers and use the epic-image
part of \cref{prop:standalone-dominance}.
The threshold formulas therefore show that $S$ is closed under the
semiring operations. If $I\in\J S$, the threshold of its supremum is
\[
 q<\sup I\ \Longleftrightarrow\
 \exists x\,[I(x)\wedge q<x]=\chi_I(x\mapsto(q<x)),
\]
so $S$ is a $\J$-algebra and its operations are continuous.
Positivity $\sigma(x)=(0<x)$ preserves multiplicative tensor and
arbitrary semantic joins. Its crisp section is
$\tau(p)=\bigvee\{1:p\}$, with $\sigma\tau=\id$; it preserves
multiplication but is not claimed to preserve arithmetic addition.
The effective version restricts to $p:\Sigma$. Thus the change from
quantitative comparison to positive comparison used below follows
from these explicit threshold maps.

\begin{theorem}[Admitted presentations]\label{thm:standalone-presentation}
For a preorder $P$ and an admitted cover relation
$\mathcal U\incl P\times\J P$, there is a generated algebra
$\pres(P,\mathcal U)$ with a monotone generator map $j$ and
the parameter-uniform restriction isomorphism
\begin{equation}\label{eq:admitted-presentation-interface}
\begin{aligned}
 \AS(\pres(P,\mathcal U),A)\cong
 \{f:P\to U_{\AS}A:\;&f\text{ monotone},\\
 &f(p)\le a_A(\J f(I))
       \text{ for every }(p,I)\in\mathcal U\}.
\end{aligned}
\end{equation}
Finite products of the generated algebras are generated under admitted
suprema by tuples of generators. Finite-arity operations preserving
the covers separately extend continuously, and their equations and
inequalities extend from generators.
\end{theorem}
\begin{proof}\leavevmode
\begin{enumerate}[leftmargin=*,align=left]
\item \textbf{Construct the generated algebra.}
The full construction and extension argument are given in
\cref{app:presentations}. Briefly, form the complete lattice of
cover-closed lower subsets of $P$, and take its least subobject
containing the closed principal ideals and closed under intrinsic
admitted suprema.

\item \textbf{Extend maps and prove uniqueness.}
Extension is constructed in the complete lattice
of admitted-closed lower subsets of the target algebra and then
restricted to the embedded target. Equality and inequality loci
prove uniqueness and order reflection.

\item \textbf{Establish product generation and parameter coherence.}
Two successive section
arguments prove finite-product generation. These are internal
power-object constructions, so the universal property holds with
parameters. This is the selected-ideal version of the dcpo
presentation method of \citep{JungMoshierVickers2008PresentingDcpos};
the selection and parameter arguments are included here.
\end{enumerate}
\end{proof}

Write $\C=\LinC$ for continuous $S$-semimodule objects in $\AS$
and continuous linear maps.
\begin{proposition}[Continuous linear constructions]\label{prop:standalone-linear}
The category $\C$ has free continuous semimodules, a symmetric
monoidal closed tensor representing continuous balanced bilinear
maps, and presented quotient modules. Their constructions and
coherence equations are uniform in parameters.
\end{proposition}
\begin{proof}\leavevmode
\begin{enumerate}[leftmargin=*,align=left]
\item \textbf{Construct free continuous semimodules.}
For $A\in\AS$, form finite formal $S$-linear combinations of elements
of $A$, quotient by the semimodule equations, and impose the least
compatible preorder containing the orders of $A$ and $S$.
Add covers
\[
 [\sup I]\triangleleft\J[-](I)\quad(I\in\J A),\qquad
 (\sup K)v\triangleleft\J(s\mapsto sv)(K)\quad(K\in\J S),
\]
and close them under finite one-hole module contexts. The presentation
theorem gives $\LLin_S A$. Maps out of it are exactly continuous maps
$A\to UN$ extended linearly.

\item \textbf{Construct the balanced tensor.}
For $M\otimes_S^{\J} N$ use generators $m\otimes n$, the balanced
bilinear equations, and continuity covers in each input.
The same theorem gives its bilinear universal property.

\item \textbf{Construct the internal hom.}
The subobject of $[UM,UN]_\Sigma$ cut out by linearity is closed
under admitted pointwise suprema, by interchange of directed
suprema; its pointwise operations give the continuous internal hom.
Currying and evaluation give closedness.

\item \textbf{Verify monoidal coherence.}
The unit is $S$, with inverse to its action given by $m\mapsto1\otimes m$.
The two associative tensor constructions represent the same
trilinear maps. The symmetry and unit comparisons have the
corresponding generator formulas, and all coherence diagrams agree
on generators.

\item \textbf{Construct quotient modules.}
Quotients add the required equations as inequalities
in both directions to the same presentation.
\end{enumerate}
\end{proof}


Throughout, an internally \emph{finite-projective module} means a
\emph{displayed finite-free retract}: specified finite maps
$M\xrightarrow{s}B^n\xrightarrow{r}M$ with $rs=\id_M$, or its
idempotent matrix $sr$. Such a presentation constructs a finite dual
basis, dualizability, and flatness. Indeed
$\Hom_B(M,N)\cong M^\vee\otimes_BN$, where $M^\vee=\Hom_B(M,B)$;
this functor is exact because $M^\vee$ is a finite-free retract.
This is the finite-projective property used below. In an arbitrary
topos it is not identified with unrestricted lifting against every
regular epimorphism of internal modules. Standard finite-projective
and flat-colimit calculations are cited in \citep{Stacks}, with the
internal finite-witness and displayed-retract arguments supplied where
needed.

The continuous module construction gives the linear objects.
Augmentations now record their quantitative grades, and the following
adjunctions relate them to ordinary computational contexts. These are
also the interfaces used by the resource-certificate interpretation.

\subsection{The linear--nonlinear adjunctions}
Put $\BudC=\C/S$, whose objects have a continuous linear
augmentation $\aug_M:M\to S$ and whose arrows preserve it.
Its tensor augmentation is
$\aug_{M\otimes N}(m\otimes n)=\aug_M(m)\aug_N(n)$.
The unit $(S,\id)$ is terminal.
The internal hom has the explicit presentation
\[
 [M,N]_{\BudC}
 =\{(r,f)\in S\times\underline{\Hom}_{\C}(M,N):
       \aug_N f=r\,\aug_M\},\qquad
 \aug(r,f)=r.
\]
It is a continuous submodule. The balance equations and ordinary
linear currying identify maps into it with augmentation-preserving
maps out of the tensor.

\begin{proposition}[Free augmented modules]\label{prop:standalone-lnl}
There is a strong symmetric monoidal adjunction
\[
 \begin{gathered}
 \FL_S:\AS\rightleftarrows\BudC:\GL_S,\\
 \FL_S A=(\LLin_S A,\aug_A),\quad \aug_A([a])=1,\qquad
 \GL_S M=\{m:\aug_M(m)=1\}.
 \end{gathered}
\]
Composing with $\operatorname{Disc}\dashv U_{\AS}$ gives
$F_d=\FL_S\operatorname{Disc}\dashv G_d=U_{\AS}\GL_S$.
Both constructions have their canonical parameter-substitution
comparisons.
\end{proposition}
\begin{proof}\leavevmode
\begin{enumerate}[leftmargin=*,align=left]
\item \textbf{Construct the augmentation.}
The constant map $A\to S$ with value $1$ is continuous,
so extends uniquely to $\aug_A$.

\item \textbf{Restrict the free-module adjunction.}
The fibre $\aug_M^{-1}(1)$ is closed under admitted suprema,
since augmentation preserves them. The free-module universal
property restricts to
\[
 \BudC(\FL_S A,M)\cong\AS(A,\GL_S M).
\]

\item \textbf{Establish strong monoidality.}
The maps sending $[a]\otimes[b]$ to $[(a,b)]$ and conversely
are obtained from the bilinear and free universal properties.
They are inverse on generators; product generation proves they
are inverse everywhere and proves all monoidal coherence equations.
The unit comparison is $\FL_S1\cong(S,\id)$.

\item \textbf{Pass to ordinary contexts.}
Composition of adjunctions gives the second assertion, and
\cref{prop:standalone-parameters} proves substitution coherence.
\end{enumerate}
\end{proof}
Write $!_c=\FL_S\GL_S$ and $!_d=F_dG_d$.
Ordinary dependent contexts are interpreted in the self-indexed
category of $\E$: a family over $\Gamma$ is an object of
$\E/\Gamma$, substitution is pullback, $\Sigma$-types are
composition, $\Pi$-types are the dependent right adjoints, and
identity types are diagonals. Linear families use the displayed
augmented-module construction in the same parameter slices.
Choose the equivalent split presentation by families; the
canonical Beck--Chevalley maps are then its substitution
comparisons. The finite syntax and checking interpretation
needed for resource certificates are given in
\cref{app:certificates}.


\begin{proposition}[Comparison of the two budget exponentials]
\label{prop:dependent-exponential-comparison}
The counit $\epsilon^{\mathrm{disc}}_A:\operatorname{Disc}U_{\AS}A\to A$
constructs a natural comonad morphism
\[
 \lambda_M=\FL_S(\epsilon^{\mathrm{disc}}_{\GL_SM}):
 !_dM\longrightarrow !_cM.
\]
It commutes with parameter substitution. Thus the computational LNL
structure and the ordinary-context dependent LNL structure have a
specified comparison, while retaining their respective domains.
\end{proposition}
\begin{proof}\leavevmode
\begin{enumerate}[leftmargin=*,align=left]
\item \textbf{Construct the comparison map.}
The displayed counit is the identity on the underlying elements and is
continuous from the discrete algebra. Applying the free functor gives
$\lambda$.

\item \textbf{Compare the counits.}
The counit of the composite adjunction is the computational
counit composed with $\lambda$, so
$\epsilon^c\lambda=\epsilon^d$.

\item \textbf{Compare the comultiplications.}
The comultiplications satisfy
\[
 \delta^c_M\lambda_M
 =(!_c\lambda_M)\lambda_{!_dM}\delta^d_M.
\]
On a free generator both sides insert that same generator twice; the
triangle identities of $\operatorname{Disc}\dashv U_{\AS}$ and
$\FL_S\dashv\GL_S$ identify the intervening units. Extension from
generators proves the equation.

\item \textbf{Verify substitution compatibility.}
The same formulas hold in every slice;
\cref{prop:standalone-parameters} identifies their substitution comparisons.
\end{enumerate}
\end{proof}


For a budget morphism $f:M\to N$, the comparison of exponentials
has the naturality square
\[
\begin{tikzcd}[column sep=large,row sep=large]
 !_dM \arrow[r,"{\lambda_M}"] \arrow[d,"{!_df}"']
   & !_cM \arrow[d,"{!_cf}"]\\
 !_dN \arrow[r,"{\lambda_N}"'] & !_cN .
\end{tikzcd}
\]
The comonad-morphism equations proved above also identify its
composites with the two counits and comultiplications.


The two exponentials just constructed organize the ordinary--linear
interface. Later constructions have different roles: the symmetric-algebra
monad $T$ describes polynomial substitution and differentiation, and the
shape operations $\Shape$ and $\InfShape$ describe infinitesimal
neighbourhoods in the based and primary geometries.
The functors of \cref{thm:modality-comparison} compare these structures with
the quantitative layer. Resource certificates use $F_d$ and its
ordinary-context interpretation. Their grades, description lengths, and
execution times are compared through the measured executions and proof
transformations that carry those quantities.

\subsection{Exact quantitative presentations}\label{sec:obstructions}
Differential resource algebra needs signed coefficients, while the
continuous quantitative semantics includes an infinite scalar.
The two obstructions below explain why the signed algebra must be
built from exact data before continuous completion.
Put $\A=\Mod_{\Qplus}(\E)$, $\C=\LinC$, and
$\Dcat=\Mod_{\QQ}(\E)$. Objects of $\A$ are algebraic semimodules,
with no imposed order or cancellativity. This includes both exact
cost profiles and the additive relations needed for universal
differentials.
\begin{proposition}[The augmented category cannot be the ordinary differential base]
The nontrivial monoidal category $\mathsf{Bud}^{\J}_S=\C/S$, with its existing tensor, cannot be made a commutative-monoid-enriched category with bilinear composition and tensor as required by an ordinary differential category.
\end{proposition}
\begin{proof}\leavevmode
\begin{enumerate}[leftmargin=*,align=left]
\item \textbf{Identify the terminal tensor unit.}
Its tensor unit $I=(S,\id)$ is terminal: the unique arrow $(M,a)\to I$ is $a$.

\item \textbf{Derive collapse of an additive enrichment.}
In any proposed additive enrichment, the zero endomorphism of $I$ is therefore $\id_I$. Bilinearity of tensor gives
\[
 \id_M=\id_M\otimes\id_I=\id_M\otimes0=0
\]
after the unit identifications. Every arrow is then zero, contradicting nontriviality.

\item \textbf{Exhibit the augmentation mismatch.}
For the existing pointwise addition the failure is already visible: augmentation-preserving $f,g$ satisfy $a_N(f+g)=2a_M$, and the zero linear map generally does not preserve augmentation.
\end{enumerate}
\end{proof}

\begin{proposition}[The completed scalars have no nonzero signed envelope]
Every unital semiring homomorphism $S\to R'$ to a ring has zero target. Moreover, the additive group completion of the underlying $\Qplus$-semimodule of any $S$-semimodule is zero.
\end{proposition}
\begin{proof}\leavevmode
\begin{enumerate}[leftmargin=*,align=left]
\item \textbf{Exclude a nonzero ring target.}
The effective infinite scalar belongs to $S$, and $\infty+1=\infty$. Cancellation in $R'$ forces $1=0$.

\item \textbf{Collapse additive group completion.}
For an $S$-semimodule, $\infty m+m=\infty m$, so the image of every $m$ in its additive group completion is zero.
\end{enumerate}
\end{proof}

We therefore give the exact positive presentation two realizations:
continuous scalar extension and signed completion. The common exact
source supplies their comparison.

\subsubsection{Continuous scalar extension and signed completion}\label{sec:scalar-span}
\paragraph{Construction of the left-hand functor}
The rational scalar map $\iota:\Qplus\to S$ sends a rational to its principal lower cut. Write $U:\C\to\A$ for restriction of scalars followed by forgetting order and admitted suprema.

\begin{construction}\label{constr:extension}
For $M\in\A$, form the algebraic $S$-semimodule generated by symbols $s\otimes m$, subject to additivity in each variable and
\[
 s\otimes(rm)=s\iota(r)\otimes m\qquad(r\in\Qplus).
\]
Give it the least compatible preorder containing the scalar order of $S$. For every $K\in\J S$ and every algebraic element $v$, impose the cover
\[
 (\sup K)v\ \triangleleft\ \J(s\mapsto sv)(K).
\]
Close these covers under finite one-hole module contexts. Apply \cref{thm:standalone-presentation} to $\pres(P,C)$, including posetal reflection of the preorder. Denote the result by $\Ext M$.
\end{construction}

\begin{proposition}\label{prop:extension-adj}
There is a parameter-uniform adjunction $\Ext\dashv U$ with
\[
 \C(\Ext M,N)\cong\A(M,UN).
\]
The functor $\Ext$ is strong symmetric monoidal:
\[
 \Ext\Qplus\cong S,\qquad
 \Ext(M\otimes_{\Qplus}M')\cong\Ext M\otimes_S^{\J}\Ext M'.
\]
\end{proposition}
\begin{proof}\leavevmode
\begin{enumerate}[leftmargin=*,align=left]
\item \textbf{Construct the scalar-extension adjunction.}
An $\Qplus$-linear map $f:M\to UN$ extends algebraically by $s\otimes m\mapsto sf(m)$. It respects the generated preorder and the displayed covers by continuity of the scalar action on $N$. The admitted-presentation theorem gives a unique continuous linear extension. Restriction to $1\otimes m$ is inverse to extension. Pointwise inequalities and free parameters are preserved by the same theorem.

\item \textbf{Identify the balanced bilinear maps.}
For the tensor assertion, a continuous balanced bilinear map out of $\Ext M\times\Ext M'$ is uniquely determined by an $\Qplus$-bilinear map out of $M\times M'$.

\item \textbf{Extend separately and evaluate.}
To extend such a map, first extend in $M'$; the resulting map into the continuous internal hom is $\Qplus$-linear in $M$. Extend that map in $M$, then evaluate. The internal hom and the parameter-uniform adjunction justify continuity in the other variable.

\item \textbf{Identify tensors and verify coherence.}
Consequently the two displayed tensor objects represent the same bilinear functor. Both unit objects represent selection of an element of $N$, giving $\Ext\Qplus\cong S$. All coherence diagrams agree on elementary generators, so uniqueness proves coherence.
\end{enumerate}
\end{proof}

The extension preserves addition and zero on morphisms:
$\Ext(f+g)=\Ext f+\Ext g$ and $\Ext0=0$. The equalities hold on
exact generators and hence on the generated domain. Here $M$ is
algebraic. If a presentation also specifies order or continuity,
its covers are imposed by a further presentation or retained in the
map $\rho$ of \eqref{eq:presentations}.

\paragraph{Construction of the right-hand functor}
\begin{proposition}\label{prop:group-completion}
Additive group completion defines a strong symmetric monoidal left adjoint
\[
 \Gr:\A\longrightarrow\Dcat.
\]
It takes $\Qplus$ to $\QQ$ and induces group completion of commutative semiring objects and their modules.
\end{proposition}
\begin{proof}\leavevmode
\begin{enumerate}[leftmargin=*,align=left]
\item \textbf{Construct the group-completion quotient.}
Construct $\Gr M$ from pairs $(m,n)$ modulo
\[
 (m,n)\sim(m',n')\quad\Longleftrightarrow\quad
 \exists k\;m+n'+k=m'+n+k.
\]
Addition is componentwise, and $(m,n)$ represents $m-n$.

\item \textbf{Extend scalars and establish the adjunction.}
The $\Qplus$-action extends uniquely to a $\QQ$-action, since the group completion of $\Qplus$ is $\QQ$. Maps to a $\QQ$-module correspond exactly to $\Qplus$-linear maps from $M$ to its underlying monoid.

\item \textbf{Prove additivity and tensor compatibility.}
The functor is additive by its formula on differences. An $\Qplus$-bilinear map $M\times N$ to an additive group extends uniquely to a $\QQ$-bilinear map $\Gr M\times\Gr N$. On differences the extension is
\[
 (m-n,p-q)\longmapsto b(m,p)-b(m,q)-b(n,p)+b(n,q).
\]
This proves $\Gr(M\otimes_\Qplus N)\cong\Gr M\otimes_\QQ\Gr N$ and its coherence.

\item \textbf{Extend algebra and module structures.}
Multiplication and module actions extend by the same bilinear formula. All quotients are internal effective quotients in $\E$.
\end{enumerate}
\end{proof}

For a commutative $\Qplus$-algebra $A$, the same proof gives functors on coefficient fibres
\[
 \Mod_A(\E)\longrightarrow\Mod_{\Ext A}(\C),\qquad
 \Mod_A(\E)\longrightarrow\Mod_{\Gr A}(\E).
\]
They preserve relative tensor products. This follows either from their balanced-map universal properties or by applying the left adjoints to the defining balanced coequalizers. Continuous relative tensors and quotient modules are constructed by the same admitted presentations with the coefficient algebra in place of $S$.


\subsubsection{The category of retained presentations}\label{sec:presentations}
We now package an exact module together with its quantitative
interpretation. Keeping both in one object lets a map carry the exact
data needed for signed realization and the continuous data needed for
quantitative semantics. Define
\begin{equation}\label{eq:presentations}
 \Pcat=(\Ext\downarrow\id_\C).
\end{equation}
An object $(M,N,\rho)$ consists of an exact positive module $M$,
a quantitative module $N$, and a continuous linear map
$\rho:\Ext M\to N$. A morphism $(f,g)$ satisfies
$g\rho=\rho'\Ext f$.

A morphism must preserve that interpretation. Its defining equation
is the commuting square
\begin{equation}\label{eq:cat-presentation-square}
\begin{tikzcd}[column sep=large,row sep=large]
 \Ext M \arrow[r,"{\rho}"] \arrow[d,"{\Ext f}"']
   & N \arrow[d,"{g}"]\\
 \Ext M' \arrow[r,"{\rho'}"'] & N' .
\end{tikzcd}
\end{equation}
The adjunction $\Ext\dashv U$ identifies this comma category with
$(\id_{\A}\downarrow U)$: the transpose of $\rho$ is the exact
interpretation map $M\to UN$.  Under this identification the same
morphism is the commuting square
\[
\begin{tikzcd}[column sep=large,row sep=large]
 M \arrow[r,"{\rho^\flat}"] \arrow[d,"{f}"']
   & UN \arrow[d,"{Ug}"]\\
 M' \arrow[r,"{(\rho')^\flat}"'] & UN' .
\end{tikzcd}
\]
The exact component also has a signed realization. Together the two
realization functors form the span
\begin{equation}\label{eq:cat-realization-span}
\begin{tikzcd}[column sep=large]
 \C & \Pcat \arrow[l,"{\pi_{\C}}"']
       \arrow[r,"{\Gr\pi_{\A}}"] & \Dcat .
\end{tikzcd}
\end{equation}
Its common source retains the exact leg needed for signed realization.
The differential comparisons of \cref{thm:modality-comparison,thm:comma}
will equip both projections with their compatible algebra modalities
and deriving transformations.
Addition is componentwise. Strong monoidality of $\Ext$ defines
the tensor componentwise, with the evident comparison map.
Its unit is $(\Qplus,S,\id)$ under $\Ext\Qplus\cong S$.
These definitions make $\Pcat$ a commutative-monoid-enriched
symmetric monoidal category with finite biproducts (a semiadditive
category). The codifferential axioms use this commutative-monoid enrichment.

The functor $M\mapsto(M,\Ext M,\id)$ is fully faithful and strong
monoidal. If $U$ is the right adjoint of $\Ext$, then
$R_\C(N)=(UN,N,\epsilon_N)$ is a fully faithful right adjoint
to quantitative projection. Indeed a map to $R_\C(N')$ is
determined by its quantitative component, whose adjoint is
the required exact component. This right adjoint is lax
monoidal. Its signed component vanishes by the infinity
calculation in \cref{sec:obstructions}; nonzero signed data
are retained in the chosen exact presentation.

\subsubsection{Raw cost profiles have a faithful quantitative interpretation}
For an input object $X$, use its discrete admitted domain. The function object $S^X$ is a continuous $S$-semimodule under pointwise operations. Exact profiles have a monomorphism
\[
 \Qplus^X\longrightarrow U(S^X),\qquad f\longmapsto\underline f.
\]
Adjunction gives a canonical quantitative realization
\begin{equation}\label{eq:raw-realization}
 \Ext(\Qplus^X)\longrightarrow S^X.
\end{equation}
Because the composite with the exact-generator map is monic, the exact-generator map into $\Ext(\Qplus^X)$ is monic as well. Any exact submodule of raw profiles has the same argument, by restriction. Input substitution is defined on these raw profiles for every map of inputs.

\subsubsection{Exact presentation of free augmented modules}
For $A\in\AS$, let $L_\Qplus(UA)$ be the algebraic free rational-positive module on its underlying object, and let $L_S^{\J}A$ be the free continuous module of \cref{prop:standalone-linear}. Its generator map induces
\[
 \Ext L_\Qplus(UA)\longrightarrow L_S^{\J}A.
\]
This is a presentation object. Its exact augmentation sums rational coefficients, and its quantitative augmentation is the free augmentation of \cref{prop:standalone-lnl}. These are compatible. The construction is strong monoidal on cartesian products of $A$, because both free-module comparisons send a pair of generators to their tensor. Projection gives $\FL_S$. The computational LNL adjunction and $!_c$ therefore have this specified exact lift of the free functor. Precomposition with $\operatorname{Disc}$ gives the corresponding exact lift of $F_d$; its comparison with the computational lift is induced by the same discrete counit as \cref{prop:dependent-exponential-comparison}.


We now have a computational category with recursive execution, an
observation and name interface, and exact presentations that admit both
quantitative and signed realizations. The next section applies the
exact side to cost profiles. Eventual comparison will turn those
profiles into the modules used as resource budgets.

\section{Landau algebra and observed resource modules}
\label{sec:landau-algebra}

An asymptotic budget should record which cost profiles it contains and
how those profiles combine. We obtain this algebraic description by
identifying Landau domination with multiplication by bounded
coefficients. The construction proceeds from exact profiles to eventual
germs and then to growth classes. Keeping these stages distinct lets
us recover both coarse bounds and finer information about coefficients
and errors. The resulting embedded modules will be the objects
classified in the next section.
\subsection{Filters and eventual quotients}
For $X\in\E$, form the object $\operatorname{Filt}(X)$ of predicates
on subobjects closed upward and under finite intersection and
containing $X$. The proper-filter subobject imposes
$\neg\Ev(\varnothing)$. This is a construction inside the power
objects of $\E$, so the following results are families over that
subobject. Write $\Ev(P)$ for membership of a predicate $P$ in
the displayed filter.

A filtered map $u:(Y,\Ev_Y)\to(X,\Ev_X)$ is a map satisfying
$\Ev_X(P)\Rightarrow\Ev_Y(u^{-1}P)$ for every subobject $P$.
Substitution proves the identity and composition laws.

\begin{proposition}[Germs and substitution]\label{prop:germs}
The relations $\Ev(f=g)$ on rational and nonnegative rational
profiles are effective equivalence relations compatible with
addition and multiplication. They construct internal algebras
\[
 K^+=(\Qplus)^X/{=_{\Ev}},\qquad K=\QQ^X/{=_{\Ev}}.
\]
Filtered maps induce algebra homomorphisms by precomposition.
\end{proposition}
\begin{proof}\leavevmode
\begin{enumerate}[leftmargin=*,align=left]
\item \textbf{Verify the eventual equivalence relation.}
Reflexivity uses $\Ev(X)$, symmetry uses symmetry of equality,
and transitivity uses intersection and upward closure.

\item \textbf{Descend the algebraic operations.}
Intersect the extents of two representative equalities to
verify each operation. The effective quotient therefore inherits
the operations and their finite equations.

\item \textbf{Construct filtered substitution.}
The defining implication
for a filtered map preserves representative equality and every
operation. Identity and composition hold on representatives.
\end{enumerate}
\end{proof}

For $X=\N$, the cofinite filter is
\[
 \Ev(P)\equiv\exists N\,\forall n\ge N\,P(n).
\]
For finite words, length gives
$\Ev(P)\equiv\exists N\,\forall x\,(\lvert x\rvert\ge N\Rightarrow P(x))$.
These filters are proper because a purported eventual empty extent,
with cutoff $N$, is contradicted respectively by $n=N$ and by a word
of length $N$. Thus the displayed formulas construct actual inhabitants
of the proper-filter subobject. For a general size map the proper-filter
fibre is retained as part of the context. Arbitrary input maps always
act on raw profiles; a constant input map need not act on cofinite germs.

\subsection{Exact, eventual, and growth information}
The successive comparisons forget different information.
The inclusion $(\Qplus)^X\incl S^X$ preserves exact pointwise profiles.
Eventual equality then identifies profiles that agree on one eventual
extent, and mutual domination makes the further growth quotient.
For example, $n+1$, $2(n+1)$, and $n+2$ have the same growth class
but distinct exact germs. Leading coefficients and small relative
errors must therefore be compared before passing to growth classes.

\subsection{Landau comparison as quantitative enrichment}
\label{sec:enrichment}
A growth comparison records whether a bounded multiplier exists.
Before passing to that predicate, we can retain the size of the
permitted multipliers in a quantitative hom. For nonnegative
rational germs define
\begin{equation}\label{eq:landau-hom}
 h_{\Ev}(f,g)=
 \bigvee\{\underline r:r\in\Qplus,\ \Ev(rf\le g)\}\in\Q.
\end{equation}
The join is an internal lower cut: its strict rational thresholds
are the existentially witnessed thresholds of the displayed
principal cuts. It is a semantic scalar; an effective name is
additional represented data.

\begin{theorem}[Enriched Landau comparison]\label{thm:enrichment}
The homs \eqref{eq:landau-hom} define a $\Q$-enriched category.
They satisfy
\[
 0<h(f,g)\iff \exists a\in\N_{\ge1}\,\Ev(f\le ag),
 \qquad
 h(f,g)=\infty\iff \forall k\ge1\,\Ev(kf\le g).
\]
\end{theorem}
\begin{proof}\leavevmode
\begin{enumerate}[leftmargin=*,align=left]
\item \textbf{Verify well-defined enriched composition.}
Eventual equality of either argument preserves each witnessed
inequality after finite intersection. The witness $r=1$ gives
$1\le h(f,f)$. If $\Ev(rf\le g)$ and $\Ev(sg\le k)$, their
intersection witnesses $\Ev(rsf\le k)$.
Multiplication of lower reals preserves joins, hence
$h(f,g)h(g,k)\le h(f,k)$.

\item \textbf{Identify positive comparison.}
A positive threshold of the join supplies a rational witness
$r>0$ with $\Ev(rf\le g)$. Choose an integer $a$ with $ar\ge1$.
This gives $\Ev(f\le ag)$. Conversely such an inequality supplies
$r=1/a$.

\item \textbf{Identify infinite comparison.}
Unboundedness of the cut supplies a witness $r>k$ for
each positive integer $k$, yielding $\Ev(kf\le g)$.
The integer witnesses in the other direction make the cut
unbounded.
\end{enumerate}
\end{proof}


\paragraph{Change of enrichment and separation.}
Positivity is a morphism of commutative quantales
\[
 \sigma:(\Q,\cdot,1)\longrightarrow(\Omega,\wedge,\top),
 \qquad \sigma(q)=(0<q):
\]
it preserves arbitrary joins and multiplication by the rational
threshold formulas of \cref{sec:foundation}.  Change of enrichment
along $\sigma$ sends $h(f,g)$ to the comparison predicate
$\sigma h(f,g)=\llbracket f=O(g)\rrbracket$.  Transitivity is the
image of enriched composition:
\[
\begin{tikzcd}[column sep=large,row sep=large]
 h(f,g)\,h(g,k) \arrow[r,phantom,"{\le}"]
   \arrow[d,mapsto,"{\sigma}"']
   & h(f,k) \arrow[d,mapsto,"{\sigma}"]\\
 \llbracket f=O(g)\rrbracket\wedge\llbracket g=O(k)\rrbracket
   \arrow[r,phantom,"{\le}"]
   & \llbracket f=O(k)\rrbracket .
\end{tikzcd}
\]
Taking positivity retains the domination predicate and forgets the
numerical value of the quantitative hom. Quotienting mutual domination
then gives the growth poset. These predicates take values in $\Omega$;
the later admission construction selects their positive-observation
profiles.


Write $f\preceq g$ for the positive comparison and
$\mathsf L_{\Ev}$ for its quotient by mutual comparison.
Its join is $[f+g]$, its meet is $[\min(f,g)]$, and its bottom is
$[0]$. Common upper bounds on $f,g$ give an upper bound on
$f+g$ after adding their constants; common lower bounds give
one for their minimum after taking the larger constant.
The distributive identities follow from the pointwise identities
for minimum and maximum, with $f+g$ growth-equivalent to
$\max(f,g)$. Multiplication is monotone and distributes over joins.

\begin{definition}[Admitted resources]\label{def:admitted}
The resource object is $\mathsf{Res}_{\Ev}=\J\mathsf L_{\Ev}$.
Its members are the growth ideals selected by
\cref{def:standalone-admission}. The full semantic ideal object is
$\Idl(\mathsf L_{\Ev})$.
\end{definition}
The positivity theorem above identifies $\sigma h(f,g)$ with
the displayed growth comparison. Applying the pullback of
\cref{def:standalone-admission} to that order gives its admitted ideals.
The full semantic ideal object retains ideals whose positive test
profile need not factor uniformly through $\Sigma$.

For signed arguments, comparison uses $h(\abs f,\abs g)$.
The hom $h(\abs{f-g},\abs g)$ retains the exact difference used
in asymptotic equivalence. On nonnegative constant ratios
$f=cg$, with $g$ eventually nonzero and $c>0$ rational,
one has $h(f,g)=1/c$. Thus the enrichment retains numerical
comparison information that positivity alone forgets.

\subsection{Landau calculus}
\label{sec:landau}
\subsubsection{Definitions and their quantifiers}
For signed germs put
\begin{equation}\label{eq:domination}
 f\preceq g\quad\Longleftrightarrow\quad
 \exists a\in\N_{\ge1}\,\Ev(\abs f\le a\abs g).
\end{equation}
Positive rational constants give the same predicate by rational
rounding. The notation $f=O(g)$ is shorthand for membership
$f\in O(g)$, and similarly for the other class notations.

\begin{definition}\label{def:landau}
Define
\[
\begin{aligned}
 f\in O(g)&\iff f\preceq g,&
 f\in\Omega(g)&\iff g\preceq f,\\
 f\in\Theta(g)&\iff f\preceq g\ \wedge\ g\preceq f,&
 f\in o(g)&\iff\forall k\ge1\,\Ev(k\abs f\le\abs g),\\
 f\in\omega(g)&\iff g\in o(f),&
 f\sim g&\iff f-g\in o(g).
\end{aligned}
\]
We use $f\asymp g$ as a synonym for $f\in\Theta(g)$.
The symbol $\way$ has its approximation-theoretic
meaning; Landau domination uses $\preceq$.
For nonnegative costs $\Omega$ is the eventual lower-bound
convention of \citep{Knuth}; $\sim$ is asymptotic equality in
the sense of relative error \citep{DLMF}.
\end{definition}

These division-free predicates include zero-support profiles.
For instance $o(0)=O(0)=\{0\}$ as germs. The object
$K_{\mathrm{nz}}=\{g:\Ev(g\ne0)\}$ is constructed by its
displayed predicate. On this object rational division, defined
arbitrarily as zero off the support, gives the usual statements:
$f\in o(g)$ exactly when $\abs{f/g}$ tends to zero, and
$f\sim g$ exactly when $f/g$ tends to one. The support condition
makes these formulas independent of the off-support extension.

The relations $O,\Omega,\Theta,o,\omega$ are recovered from the
two directed enriched homs by positivity, mutual positivity,
and the value $\infty$. Asymptotic equivalence is recovered
from $h(\abs{f-g},\abs g)=\infty$.

\subsubsection{Arithmetic and replacement laws}
\begin{proposition}[Landau arithmetic]\label{prop:landau-arithmetic}
Domination is a preorder and direct smallness is transitive.
Each of $O,\Omega,\Theta,o,\omega$ is unchanged by replacing
either argument by a $\Theta$-equivalent germ. For germs
$f,g,h,k$ one has
\[
\begin{aligned}
 O(g)+O(h)&\subseteq O(\abs g+\abs h),&
 O(g)O(h)&\subseteq O(gh),\\
 o(g)+o(h)&\subseteq o(\abs g+\abs h),&
 o(g)O(h)&\subseteq o(gh),\\
 O(g)o(h)&\subseteq o(gh).
\end{aligned}
\]
Furthermore $f\sim g$ is an equivalence relation, implies
$f\in\Theta(g)$, and is preserved by multiplication of two
asymptotically equivalent pairs. For nonnegative profiles,
$f\sim g$ and $h\sim k$ imply $f+h\sim g+k$.
\end{proposition}
\begin{proof}\leavevmode
\begin{enumerate}[leftmargin=*,align=left]
\item \textbf{Prove transitivity and replacement.}
Multiply domination constants and intersect their extents.
To compose smallness with domination by $a$, use the smallness
inequality at $ka$. To compose two smallness relations, use
smallness of one and the $k=1$ bound of the other. These
arguments also prove replacement under two-sided domination.

\item \textbf{Establish the sum and product estimates.}
For sums, use the two smallness inequalities at $2k$;
for products, multiply the inequalities after intersecting
their extents. Absolute values satisfy the needed triangle
and product identities pointwise.

\item \textbf{Deduce tightness from asymptotic equality.}
If $f-g\in o(g)$, the
inequality at $k=2$ gives eventually
$\abs g/2\le\abs f\le3\abs g/2$. Hence $f\in\Theta(g)$.

\item \textbf{Prove equivalence of asymptotic equality.}
Replacement and $\abs{f-g}=\abs{g-f}$ give symmetry.
For transitivity write $f-h=(f-g)+(g-h)$, replace the
reference budget $g$ by $h$, and use the sum law.
Reflexivity is zero error.

\item \textbf{Preserve asymptotic equality under products.}
Write $f=g+r$ and $h=k+s$ with $r\in o(g)$ and $s\in o(k)$.
Then $fh-gk=rk+gs+rs\in o(gk)$ by the product and sum laws.

\item \textbf{Preserve asymptotic equality under nonnegative sums.}
For nonnegative sums the error is $r+s$ and the reference
budget is exactly $g+k$, giving the assertion. For signed
sums the preceding estimate uses $\abs g+\abs k$; cancellation
is retained in the actual reference profile.
\end{enumerate}
\end{proof}

The definition of $\sim$ uses the signed difference, not merely
the two absolute values. For example $-g$ and $g$ are
$\Theta$-equivalent, while their relative error is $-2$ on
the eventually nonzero locus.

\subsubsection{Logical refinements and uniformity}
The alternative double-negation relation is retained explicitly:
\[
 f\in\onn(g)\iff f\preceq g\ \wedge\
 \forall k\ge1\,\neg\neg\Ev(k\abs f\le\abs g).
\]
The direct relation is also denoted $\oeff(g)$ when comparison
with an explicitly witnessed eventual presentation is useful;
throughout this paper plain $o$ means this direct relation.
We have $o(g)\subseteq\onn(g)\subseteq O(g)$.
Finite conjunction is stable under double negation, so the
arithmetic and replacement arguments above give the analogous
closure statements for $\onn$. Quantifiers are left in their
displayed order.

For cofinite input, a chosen modulus is an element of
\[
 \operatorname{Mod}_o(f,g)=
 \{m:\N_{\ge1}\to\N:
   \forall k\,\forall n\ge m(k),\ k\abs{f(n)}\le\abs{g(n)}\}.
\]
It maps to the direct smallness predicate. A represented
computable modulus carries its implementation as well.
The modulus object and its represented points record these two levels of data.

For $f,g:A\times\N\to\QQ$, uniform domination is the predicate
\[
 \exists c,N\;\forall a\in A\;\forall n\ge N,\quad
 \abs{f(a,n)}\le c\abs{g(a,n)}.
\]
Pointwise-in-$A$ domination has $\forall a$ before $\exists c,N$.
Uniform smallness has $\forall k\,\exists N\,\forall a\,\forall n\ge N$.
All the arithmetic proofs preserve uniformity because only
finitely many bounds are combined at each tolerance.

\subsubsection{Oscillation and infinitely-often bounds}
On cofinite numerical input define
\[
 f\in\Omega_{\mathrm{io}}(g)\iff
 \exists c\in\QQ_{>0}\,\forall N\,\exists n\ge N,\quad
 \abs{f(n)}\ge c\abs{g(n)}.
\]
For signed one-sided statements replace the last inequality
by $f(n)\ge c\abs{g(n)}$ or $f(n)\le-c\abs{g(n)}$,
giving $\Omega_{\mathrm{io},+}$ and $\Omega_{\mathrm{io},-}$.
These definitions use positive witnesses of arbitrarily late
inputs. They are not defined by negating smallness.

Let $g(n)=n+1$ and let $f(n)=n+1$ on even inputs and $1$ on odd
inputs. Then $f\in O(g)\cap\Omega_{\mathrm{io}}(g)$, while
$f\notin o(g)$ and $f\notin\Theta(g)$. Each failure has explicit
arbitrarily large parity witnesses. If instead $f(n)=(n+1)^2$
on even inputs and $1$ on odd inputs, then
$f\notin O(g)$ and $f\notin\omega(g)$.
Thus lower and strict lower bounds are actual predicates,
not complements of upper bounds. The eventual/infinitely-often
distinction is also the distinction emphasized in \citep{Knuth}.

Filtered substitution preserves every direct eventual comparison,
and the arithmetic proofs commute with it. An isomorphism of
filtered inputs reflects the comparisons by its inverse.
Infinitely-often comparisons use their own displayed witness
transport and are not silently inferred from preservation of
eventual filters.

The comparison laws can now be expressed inside the germ algebra.
Bounded coefficients will describe $O$-bounds, vanishing coefficients
will describe $o$-bounds, and their units will describe the corresponding
equivalence relations. This converts the calculus of inequalities into
module membership and group actions.

\subsection{Bounded coefficients, vanishing errors, and units}
\label{sec:coefficients}
\subsubsection{Positive coefficients and bounded division}
Within the exact positive germ algebra put
\[
 B^+=\{b\in K^+:b\in O(1)\}.
\]
Bounds add and multiply, so $B^+$ is a unital $\Qplus$-algebra.
It embeds into $K^+$. Both algebras are cancellative: equality
$f+h=g+h$ on an eventual extent implies $f=g$ on the same extent.
For a profile $g$ let $e_g$ denote the indicator germ of $g\ne0$.
Rational equality is decidable, so this is an actual bounded
profile with $e_g^2=e_g$. Write $z_g=1-e_g$ on the signed side.

\begin{lemma}[Positive bounded division]\label{lem:positive-division}
For $f,g\in K^+$, $f=O(g)$ if and only if $f=bg$ for some $b\in B^+$. If $e_g$ is the indicator germ of $g\ne0$, then
\[
 B^+g\cong B^+e_g,\qquad
 B^+g+B^+h=B^+(g+h),\qquad
 B^+g\cap B^+h=B^+\min(g,h).
\]
\end{lemma}
\begin{proof}\leavevmode
\begin{enumerate}[leftmargin=*,align=left]
\item \textbf{Construct the bounded multiplier.}
Define $b(x)=f(x)/g(x)$ where $g(x)\ne0$ and $b(x)=0$ otherwise. Rational equality is decidable. An eventual bound $f\le ag$ makes $b$ eventually bounded and gives $f=bg$ as germs. The converse is multiplication of the bound.

\item \textbf{Identify the support module.}
Multiplication by $g$ gives the support isomorphism, with inverse given by the same division on its support.

\item \textbf{Compute the module sum.}
The two generators lie in $B^+(g+h)$ and their sum lies in their module sum.

\item \textbf{Compute the intersection and descend representatives.}
For intersections, combine the two eventual bounds with their maximum and apply bounded division to $\min(g,h)$. These arguments are local on representatives and descend through the effective quotient.
\end{enumerate}
\end{proof}


\begin{proposition}[Signed coefficients]\label{prop:signed-coefficients}
There are canonical algebra isomorphisms
\[
 \Gr K^+\cong K,\qquad
 \Gr B^+\cong B=\{b\in K:b\in O(1)\}.
\]
The algebra $B$ is reduced and $K$ is its total quotient algebra.
For signed germs $f,g$,
\[
 f\in O(g)\iff f=bg\text{ for some }b\in B.
\]
\end{proposition}
\begin{proof}\leavevmode
\begin{enumerate}[leftmargin=*,align=left]
\item \textbf{Identify the signed coefficient rings.}
Cancellativity embeds the group completions into $K$.
Every signed rational profile has the canonical decomposition
$f=f_+-f_-$, where $f_+=\max(f,0)$ and
$f_-=\max(-f,0)$. Boundedness of $\abs f$ bounds both parts,
proving both isomorphisms.

\item \textbf{Prove signed bounded division.}
Division $f/g$ on the support of $g$,
extended by zero elsewhere, proves signed bounded division
just as in \cref{lem:positive-division}.

\item \textbf{Prove reducedness.}
If $b^r=0$ for a positive integer $r$, rational arithmetic gives
$b=0$ on the same eventual extent, proving reducedness.

\item \textbf{Compute annihilators and regular elements.}
The annihilator of $b$ in $B$ is $Bz_b$: an annihilating
multiplier is zero on the support of $b$, and conversely.
Consequently a regular $b$ has $z_b=0$ and is eventually
nonzero. Its rational reciprocal is an inverse in $K$.

\item \textbf{Construct the total quotient ring.}
For arbitrary $f\in K$, the germ $s=(1+\abs f)^{-1}$ is regular
and bounded, and $sf$ is bounded. Thus every $f$ is a fraction
of bounded germs with regular denominator. Equality of these
fractions is reflected in $K$ by regularity, giving the
localization universal property.
\end{enumerate}
\end{proof}

\subsubsection{Units and the vanishing ideal}
\begin{lemma}[Bounded inverse]\label{lem:units}
For $u\in B$, the following are equivalent:
\[
 u\in B^\times;\qquad
 \exists a\in\N_{\ge1}\,\Ev(1\le a\abs u).
\]
Thus the units are exactly the bounded germs bounded away from
zero. In particular $B^\times=\Theta(1)$.
\end{lemma}
\begin{proof}\leavevmode
\begin{enumerate}[leftmargin=*,align=left]
\item \textbf{Extract a lower bound from an inverse.}
If $uv=1$ and $\abs v\le a$ eventually, then $1\le a\abs u$
on the intersection.

\item \textbf{Construct and descend the reciprocal.}
Conversely define $v=1/u$ where $u\ne0$
and $v=0$ elsewhere. The displayed bound gives $\abs v\le a$
and $uv=1$ eventually. These equations descend to the germ
ring. The last assertion includes the already specified
boundedness of $u$.
\end{enumerate}
\end{proof}

Put $\mathfrak m=o(1)$ and $\mathfrak m_{\neg\neg}=\onn(1)$.
The sum inequalities at $2k$ and the product inequality at
$ka$, for a bound $a$ on a multiplier, prove that both are
ideals of $B$. No infinite intersection of eventual extents
is used.

\begin{theorem}[Vanishing errors and the Jacobson radical]
\label{thm:jacobson}
Define the unit-theoretic Jacobson radical by
\[
 \Jac(B)=\{a\in B:\forall b\in B,\ 1-ab\in B^\times\}.
\]
Then
\[
 \Jac(B)=\mathfrak m.
\]
In particular $1+\mathfrak m$ is a subgroup of $B^\times$.
\end{theorem}
\begin{proof}\leavevmode
\begin{enumerate}[leftmargin=*,align=left]
\item \textbf{Show that vanishing germs lie in the radical.}
Let $a\in\mathfrak m$ and let $\abs b\le c$ eventually.
The smallness bound at $2c$ gives $\abs{ab}\le1/2$ eventually,
so $\abs{1-ab}\ge1/2$. The bounded-inverse lemma makes $1-ab$
a unit.

\item \textbf{Construct a cutoff multiplier.}
Conversely, fix $a\in\Jac(B)$ and $k\ge1$. On a representative
define
\[
 b_k(x)=
 \begin{cases}
  a(x)^{-1},&k\abs{a(x)}>1,\\
  0,&k\abs{a(x)}\le1.
 \end{cases}
\]
This is bounded by $k$ and descends to a bounded germ.

\item \textbf{Deduce vanishing from radical membership.}
The element $q_k=1-ab_k$ is the indicator of
$k\abs a\le1$, hence idempotent. It is invertible by the
definition of $\Jac(B)$. Multiplying $q_k^2=q_k$ by its inverse
gives $q_k=1$, which is precisely $\Ev(k\abs a\le1)$.
This holds for each $k$, proving $a\in\mathfrak m$.

\item \textbf{Verify the subgroup of units.}
For $a,b\in\mathfrak m$, one has
$(1+a)(1+b)=1+(a+b+ab)$ with $a+b+ab\in\mathfrak m$.
Also $(1+a)^{-1}-1=-a(1+a)^{-1}\in\mathfrak m$.
These formulas prove the subgroup assertion.
\end{enumerate}
\end{proof}

The theorem uses a definition of the radical by units and finite
rational tests. The quotient
$B/\mathfrak m$ retains nontrivial idempotents. For example, on
cofinite numerical input the parity indicator and its complement
are neither zero nor vanishing: the tolerance $1/2$ and
arbitrarily late inputs of each parity prove this directly.
Its class is therefore an idempotent distinct from zero and one.
The residue algebra is not treated as a field.

\begin{theorem}[Multiplier and orbit formulas]\label{thm:orbits}
For every signed germ $g$,
\[
\begin{aligned}
 O(g)&=Bg,& o(g)&=\mathfrak m g,&
 \onn(g)&=\mathfrak m_{\neg\neg}g,\\
 \Theta(g)&=B^\times g,&
 \{f:f\sim g\}&=g+\mathfrak m g=(1+\mathfrak m)g.
\end{aligned}
\]
Consequently
\[
 f\in\Omega(g)\iff Bg\subseteq Bf,\qquad
 f\in\omega(g)\iff g\in\mathfrak m f.
\]
\end{theorem}
\begin{proof}\leavevmode
\begin{enumerate}[leftmargin=*,align=left]
\item \textbf{Identify the bounded orbit.}
Bounded division proves the first formula.

\item \textbf{Identify the vanishing orbit.}
For $f\in o(g)$
use the same ratio $a=f/g$ on the support of $g$, zero
elsewhere. Every smallness inequality implies
$k\abs a\le1$ on its eventual extent, while the $k=1$
inequality ensures $f=ag$ as germs. Thus $a\in\mathfrak m$.
Conversely multiply the inequalities for $a$ by $\abs g$.

\item \textbf{Identify the double-negated orbit.}
For $\onn$ the explicit domination supplies the germ equation
and boundedness; applying double negation to each implication
gives the remaining bounds.

\item \textbf{Identify the supported unit orbit.}
If $f\in\Theta(g)$, the two domination inequalities imply that
$f$ and $g$ have the same support eventually. Define $u=f/g$
where both are nonzero and $u=1$ elsewhere. On the intersection
of these eventual extents it is bounded, its reciprocal is
bounded, and $f=ug$. The bounded-inverse lemma gives $u\in B^\times$.
Conversely a unit and its bounded inverse give both domination
bounds. This argument includes $g=0$.

\item \textbf{Identify the asymptotic unit orbit.}
Finally $f\sim g$ means $f-g\in\mathfrak m g$, proving the
coset and orbit descriptions.

\item \textbf{Recover the lower-bound relations.}
The two reversed comparisons
follow immediately from the definitions and bounded division.
\end{enumerate}
\end{proof}

For nonnegative profiles the unit orbit may be restricted to
eventually positive units. The construction of $u$ in the proof
has this property. Signed absolute-value comparison also records
sign-changing units, while the subgroup $1+\mathfrak m$ retains
the relative sign in asymptotic equivalence.

\subsection{Resource modules and admitted growth ideals}
\label{sec:resources}
Principal budgets have the form $Bg$. To include budgets with several
or indefinitely many bounds, we collect growth classes into ideals.
The following correspondence identifies these ideals with embedded
$B$-submodules of $K$, so the marking records exactly which costs a
budget contains. For $I\in\Idl(\mathsf L_{\Ev})$ define
\[
 M_I^+=\{f\in K^+:[f]\in I\},\qquad
 M_I=\{f\in K:[\abs f]\in I\}.
\]
Here an ideal is inhabited, lower, and directed; equivalently
it contains zero and is closed under finite joins.

\begin{theorem}[Positive resource representation]\label{thm:positive-rep}
There is an order isomorphism
\[
 \Idl(\mathsf L_\Ev)\cong\operatorname{Sub}_{B^+}(K^+),\qquad
 I\longmapsto M_I^+=\bigcup_{[g]\in I}B^+g.
\]
It preserves sums, intersections, and module products. Multiplication gives
\[
 M_I^+\otimes_{B^+}M_J^+\cong M_{IJ}^+.
\]
\end{theorem}
\begin{proof}\leavevmode
\begin{enumerate}[leftmargin=*,align=left]
\item \textbf{Identify ideals with positive resource modules.}
A directed union of the displayed principal modules is a submodule. Conversely, for a submodule $M\subseteq K^+$, its growth classes form a lower set by bounded division and are directed by $f+g$. Both constructions are inverse.

\item \textbf{Preserve sums, intersections, and products.}
The formulas in Lemma~\ref{lem:positive-division} prove preservation of sums and intersections; finite products give the product statement.

\item \textbf{Construct the principal tensor isomorphism.}
Each $B^+e$ is a retract of $B^+$ by multiplication with the
idempotent $e$. For principal resources, the tensor identification is
\[
 B^+g\otimes_{B^+}B^+h
 \cong B^+e_g\otimes_{B^+}B^+e_h
 \cong B^+(e_ge_h)\cong B^+gh.
\]
The middle map is multiplication; the inverse sends
$b e_ge_h$ to $be_g\otimes e_h$.

\item \textbf{Verify naturality under inclusion.}
If one principal module
is contained in another, these maps agree after inclusion into
$K^+$, so they form a natural isomorphism of the principal-stage
diagrams.

\item \textbf{Extend the tensor isomorphism to all resources.}
General resources are their internal directed colimits.
The balanced-map universal property makes tensor preserve these
colimits in each variable. Passing to the colimits gives the general
identification.
\end{enumerate}
\end{proof}

\begin{theorem}[Signed resource lattice]\label{thm:signed-rep}
The assignments
\[
 I\longmapsto M_I,\qquad
 M\longmapsto\{[g]:g\ge0,\ g\in M\}
\]
are inverse order isomorphisms
$\Idl(\mathsf L_{\Ev})\cong\Sub_B(K)$.
Ideal joins correspond to module sums, meets to intersections,
and ideal products to products inside $K$.
For every $g$,
\[
 \Ann_B(g)=Bz_g,\qquad Bg\cong B/Bz_g\cong Be_g,
\]
and
\[
 Bg+Bh=B(\abs g+\abs h),\qquad
 Bg\cap Bh=B\min(\abs g,\abs h).
\]
\end{theorem}
\begin{proof}\leavevmode
\begin{enumerate}[leftmargin=*,align=left]
\item \textbf{Pass between signed modules and growth classes.}
For every $f\in K$, the profile $\operatorname{sgn}(f)$ is bounded
and hence defines an element of $B$. The identities
\[
 |f|=\operatorname{sgn}(f)\,f,\qquad
 f=\operatorname{sgn}(f)\,|f|
\]
show that $f\in M$ if and only if $|f|\in M$ for every
$B$-submodule $M\subseteq K$.
Its nonnegative growth classes are lower by bounded division
and directed by addition. Conversely the directed union
of principal budgets is a submodule. These constructions are
inverse on elements, giving the order isomorphism.

\item \textbf{Identify principal supports.}
The annihilator was calculated in
\cref{prop:signed-coefficients}; the idempotent splitting gives
the principal isomorphism.

\item \textbf{Establish the lattice formulas.}
The sum and intersection arguments
of \cref{lem:positive-division} apply after absolute value.
Finite sums and products describe the general lattice operations.
\end{enumerate}
\end{proof}

\begin{theorem}[Conservative signed resource realization]\label{thm:signed-resources}
There are canonical isomorphisms
\[
 \Gr K^+\cong K=\QQ^X/{=_{\Ev}},\qquad
 \Gr B^+\cong B=\{b\in K:\abs b=O(1)\},\qquad
 \Gr M_I^+\cong M_I\subseteq K.
\]
On marked resource submodules these comparisons preserve and reflect inclusions and membership, and preserve sums, intersections, products, and tensor products.
\end{theorem}
\begin{proof}\leavevmode
\begin{enumerate}[leftmargin=*,align=left]
\item \textbf{Identify the positive part after group completion.}
Positive germs are cancellative: equality $f+h=g+h$ on an eventual set gives $f=g$ on that set. Their group completion therefore embeds into the signed germ algebra. Every signed rational profile splits canonically as
\[
 f=f_+-f_-,\qquad f_+=\max(f,0),\quad f_-=\max(-f,0).
\]
This proves the first isomorphism.

\item \textbf{Identify the signed coefficient ring.}
If $\abs f$ is bounded, so are $f_+$ and $f_-$, proving the second.

\item \textbf{Identify each signed resource module.}
If $f\in M_I$, then $f_\pm\le\abs f$ and their growth belongs to $I$. Conversely the difference of two members of $M_I^+$ belongs to $M_I$ because their sum dominates its absolute value. This proves the third, with the same ambient marking.

\item \textbf{Reflect inclusion and membership.}
A positive germ belongs to $\Gr M_I^+$ inside $K$ exactly when its growth belongs to $I$. Thus inclusions and membership are reflected.

\item \textbf{Transfer the lattice and tensor formulas.}
Theorems~\ref{thm:positive-rep} and~\ref{thm:signed-rep} identify the same ideal lattice, so sums and intersections agree. Product and tensor comparisons follow from Theorem~\ref{thm:positive-rep} and strong monoidality on coefficient fibres.
\end{enumerate}
\end{proof}


\begin{corollary}[Monoidal resource representation]
\label{cor:cat-monoidal-resources}
Let $\Emb_B(K)$ have objects marked submodules $i_M:M\incl K$
and maps $f:M\to N$ satisfying $i_Nf=i_M$.  The ideal-to-module
correspondence is a strong symmetric monoidal equivalence
\[
 (\Idl L,\star,\downarrow[1])
   \simeq(\Emb_B(K),\otimes_B,B).
\]
The tensor marking is multiplication into $K$; it is monic by the
resource tensor theorem.  Joins correspond to embedded module sums,
meets to intersections, and directed joins to directed unions.
\end{corollary}
\begin{proof}\leavevmode
\begin{enumerate}[leftmargin=*,align=left]
\item \textbf{Establish the equivalence and tensor comparison.}
The underlying equivalence is \cref{thm:signed-rep}.
Its tensor comparison is the already constructed multiplication
$M_I\otimes_BM_J\cong M_{I\star J}$, and its unit comparison is
$M_{\downarrow[1]}=B$.

\item \textbf{Verify monoidal coherence.}
Both paths in
\[
\begin{tikzcd}[column sep=large,row sep=large]
 (M_I\otimes_BM_J)\otimes_BM_H
   \arrow[r,"{\cong}"] \arrow[d,"{\cong}"']
   & M_{I\star J}\otimes_BM_H \arrow[d,"{\cong}"]\\
 M_I\otimes_B(M_J\otimes_BM_H)
   \arrow[r,"{\cong}"']
   & M_{I\star J\star H}
\end{tikzcd}
\]
send $(f\otimes g)\otimes h$ to $fgh$ in $K$; this proves
associativity coherence.  The unit and symmetry diagrams follow
from $1f=f$ and $fg=gf$.  Their markings determine the comparison
maps uniquely.

\item \textbf{Preserve lattice operations and directed colimits.}
The lattice assertions are
\cref{thm:signed-rep}; the directed-colimit assertion follows from
the membership formula defining $M_I$.
\end{enumerate}
\end{proof}


Restricting to $\J\mathsf L_\Ev$ selects the resources whose positive
test profiles are admitted. Natural-valued costs embed conservatively:
their domination formula is unchanged by rational inclusion.

\begin{theorem}[Constructive finite-generator splitting]
\label{thm:finite-generator-splitting}
Every finitely generated submodule $N\subseteq B^n$ has a displayed
finite-free retract, hence a finite idempotent-matrix presentation.
More precisely, if $w_1,\ldots,w_m\in B^n$ generate $N$ and
$p:E\to N$ is a $B$-linear map with displayed lifts
$v_i\in E$ satisfying $p(v_i)=w_i$, then these finite data
construct a $B$-linear section $s:N\to E$ of $p$.
All operations commute with substitution of coefficient and
generator parameters.
\end{theorem}
\begin{proof}\leavevmode
\begin{enumerate}[leftmargin=*,align=left]
\item \textbf{Initialize the induction.}
Induct on $n$. For $n=0$ the zero module has its unique section.

\item \textbf{Split a module with one coordinate.}
For $n=1$ put $h=\sum_i|w_i|$, so $N=Bh\cong Be_h$ by
bounded division. Write $q(c,h)=c/h$ on the support of $h$,
zero elsewhere, and define
\[
 s(c)=q(c,h)\sum_i\operatorname{sgn}(w_i)v_i\qquad(c\in Bh).
\]
The quotient is a well-defined $B$-linear function of $c\in Bh$:
if $c=bh$, then $q(c,h)=be_h$. Since
$\sum_i\operatorname{sgn}(w_i)w_i=h$, we have
$ps(c)=q(c,h)h=c$.

\item \textbf{Split the final coordinate.}
For $n>1$ let $\pi:N\to B$ be the final-coordinate map,
$a_i=\pi(w_i)$, and $h=\sum_i|a_i|$. The rank-one formula,
using the same displayed lifts $v_i$, constructs
$j:Bh\to E$ with $(\pi p)j=\id_{Bh}$. Put
\[
 q_0=\id_E-j\pi p,\qquad q_0^2=q_0,\qquad \pi p q_0=0.
\]

\item \textbf{Generate the kernel of the final coordinate.}
The elements $w_i'=p q_0(v_i)$ have final coordinate zero and
generate $N'=\ker(\pi:N\to Bh)$. Indeed, for $w\in N'$, generation
of $N$ gives locally $w=\sum_i b_iw_i$. Then
\[
 \sum_i b_iw_i'
 =\sum_i b_i(w_i-pj\pi w_i)
 =w-pj\pi w=w.
\]

\item \textbf{Apply the induction hypothesis.}
The corestriction $p q_0:E\to N'$ has displayed lifts $v_i$
of these generators. Induction gives $s':N'\to E$ with
$p q_0s'=\operatorname{incl}_{N'}$. Put $\bar s'=q_0s'$,
so $p\bar s'=\operatorname{incl}_{N'}$.

\item \textbf{Reconstruct a section of the original surjection.}
For $w\in N$ set
\[
 s(w)=j(\pi w)+\bar s'\bigl(w-pj(\pi w)\bigr).
\]
The second argument lies in $N'$, and direct substitution gives
$ps(w)=pj(\pi w)+w-pj(\pi w)=w$.

\item \textbf{Obtain finite projectivity and parameter stability.}
Taking $E=B^m$ and $p$ to be the canonical generating
surjection gives $ps=\id_N$ and the idempotent $sp$ on $B^m$,
whose image is $N$. Every displayed operation is finite support
division, finite arithmetic, or induction, and the equations hold
uniformly in parameters.
\end{enumerate}
\end{proof}

\begin{corollary}[The algebra needed by the later geometry]
Every signed embedded resource is a directed union of finite projective principal modules. It is flat over $B$, and multiplication identifies $M\otimes_BN$ with $MN\subseteq K$. The algebra $K$ is the total quotient algebra of $B$; $B$ is reduced and semihereditary.
\end{corollary}
\begin{proof}\leavevmode
\begin{enumerate}[leftmargin=*,align=left]
\item \textbf{Identify principal projective modules.}
For signed $g$, bounded division and its support idempotent give
$Bg\cong Be_g$, an explicit finite-projective summand of $B$.

\item \textbf{Express resources as directed unions.}
For any finite family $f_1,\ldots,f_n\in M$, the element
$h=\sum_i|f_i|\in M$ satisfies $f_i\in Bh$, so the principal
submodules form an internal directed system with union $M$.

\item \textbf{Deduce flatness.}
Every finite relation is witnessed at a common principal stage,
and this stage is flat. The finite-witness filtered-colimit
argument is the internal form of the standard exactness
calculation \citep[Tag~05UT]{Stacks}.

\item \textbf{Establish the tensor formulas.}
The principal tensor multiplication isomorphisms commute with
stage inclusions, and tensor preserves directed colimits.
Their colimit therefore gives $M\otimes_BN\cong MN\subseteq K$.

\item \textbf{Verify reducedness.}
Reducedness follows pointwise from rational arithmetic on the same eventual set.

\item \textbf{Identify the total quotient ring.}
For any $f\in K$, the regular bounded germ $s=(1+\abs f)^{-1}$ satisfies $sf\in B$. Conversely the zero-locus indicator of a regular $b\in B$ is zero; its pointwise reciprocal, set to zero off its support, is an inverse in $K$. This proves the total-quotient description.

\item \textbf{Deduce semihereditariness.}
Every finitely generated ideal of $B$ is one of the projective principal modules just described, proving semihereditariness.
\end{enumerate}
\end{proof}


In particular $f\in\Theta(g)$ says that $f$ and $g$ generate
the same marked principal module. Asymptotic equivalence is
equality with $g$ in the quotient
$Bg/\mathfrak m Bg$. The lower-bound incidences $\Omega(g)$ and $\omega(g)$ exclude zero when $g\ne0$.

The lattice of all resources is a frame, since it is the
ideal completion of the distributive growth lattice.
The admitted subobject is selected by $\J$.
For numerical costs the passage to rational resources is the
following explicitly typed comparison.
\begin{proposition}[Conservative rational extension of growth ideals]
\label{prop:natural-rational-growth}
Let $\mathsf L^{\N}_{\Ev}$ be the growth poset of nonnegative
integer-valued profiles and let
$\ell:\mathsf L^{\N}_{\Ev}\to\mathsf L_{\Ev}$ be induced by
$\N\hookrightarrow\Qplus$. It is order-reflecting. For an ideal
$I\in\Idl(\mathsf L^{\N}_{\Ev})$, its rational lower image is
\[
 \ell_!I=\{q\in\mathsf L_{\Ev}:\exists n\in I,\ q\le\ell(n)\}.
\]
Then, for a natural-valued profile $f$,
\[
 [f]_{\N}\in I\quad\Longleftrightarrow\quad
 f\in M_{\ell_!I}.
\]
The map $\J\ell$ carries admitted ideals to admitted ideals and
reflects their inclusion order; its underlying ideal map is $\ell_!$.
\end{proposition}
\begin{proof}\leavevmode
\begin{enumerate}[leftmargin=*,align=left]
\item \textbf{Compare rational and integer bounds.}
A rational eventual domination coefficient is bounded above by an
integer coefficient; hence rational domination between integer
profiles reflects integer domination. Therefore
\[
 \ell([f]_{\N})\in\ell_!I
 \ \Longleftrightarrow\
 \exists n\in I\;([f]_{\N}\le n)
 \ \Longleftrightarrow\ [f]_{\N}\in I.
\]
By definition of $M_{\ell_!I}$ this is the displayed comparison.

\item \textbf{Reflect admitted resources and their operations.}
Functoriality of $\J$ gives admission of the lower image. Moreover,
for every $f$,
$[f]_{\N}\in I\iff \ell([f]_{\N})\in(\J\ell)(I)$;
these tests recover inclusion of ideals and prove order reflection.
All statements are internal and uniform in the filter and input.
\end{enumerate}
\end{proof}

\subsubsection{The observation domain of computational resources}
\label{sec:resource-observations}
The module classification describes resource membership semantically.
We next specify which observations of a resource are positive
computational tests. Put $L=\mathsf L_{\Ev}$. Its free observed domain
from \cref{sec:standalone-observations} has the type and observation family
\[
 \F L=(\Idl L,\OO_{\F L}),\qquad
 \OO_{\F L}=\{u^\dagger:u\in\operatorname{Mon}(L,\Sigma)\},
\]
where
\[
 u^\dagger(I)=\exists a:L\;[I(a)\wedge u(a)].
\]
Its effective-point object is $\Pt(\F L)=\J L$. Thus the uniform
positive pairing is a map
\begin{equation}\label{eq:resource-positive-pairing}
 \chi:\J L\times\operatorname{Mon}(L,\Sigma)\to\Sigma.
\end{equation}
By contrast, the marked-resource membership and measured incidences have
types
\[
 \mu:\Idl L\times K\to\Omega,\qquad
 \mu(I,f)\iff[\abs f]\in I,
 \qquad (I,e)\mapsto\mu(I,\Phi(e)):\J L\times P\to\Omega.
\]


\begin{proposition}[Effective transport of resource observations]
\label{prop:effective-resource-transport}
A monotone map $r:L\to L'$ induces an effective Scott map
$\F r:\F L\to\F L'$ and its admitted-point map
$\J r:\J L\to\J L'$. Explicitly,
\[
 (\F r)(I)=\downarrow r(I),\qquad
 \chi_{\J r(I)}(u)=\chi_I(u\circ r).
\]
\end{proposition}
\begin{proof}\leavevmode
\begin{enumerate}[leftmargin=*,align=left]
\item \textbf{Construct the continuous ideal map.}
The lower image is an inhabited directed ideal, and its formula
commutes with directed unions.

\item \textbf{Transport observations and profiles.}
For $u:L'\to\Sigma$ monotone,
\[
 \exists b\,[b\in\downarrow r(I)\wedge u(b)]
 \iff\exists a\,[a\in I\wedge u(r(a))].
\]
The forward implication uses $b\le r(a)$ and monotonicity; the reverse
uses $b=r(a)$. The inverse image of $u^\dagger$ is therefore
$(u\circ r)^\dagger$, proving observation preservation.

\item \textbf{Transport resources and verify composition.}
For admitted
$I$ the right side factors uniformly through $\Sigma$, proving the
admitted-point assertion. Composition is literal substitution of tests.
\end{enumerate}
\end{proof}

\begin{proposition}[Admitted membership is not generally positive]
\label{prop:nonpositive-budget}
In the effective numerical fibre with the cofinite input filter, the
predicate that a uniformly given exact profile lies in the principal
budget $\downarrow[1]$ need not be $\Sigma$-classified.
\end{proposition}
\begin{proof}\leavevmode
\begin{enumerate}[leftmargin=*,align=left]
\item \textbf{Construct the clocked family of germs.}
The bounded numerical interpreter of \cref{prop:standalone-numerical-cover,prop:numerical-clock} gives a Boolean test
$H(e,e,n)$ for halting within $n$ updates. Define the primitive-recursive
family
\[
 r_e(n)=\begin{cases}1,&\neg H(e,e,n),\\n+1,&H(e,e,n).\end{cases}
\]

\item \textbf{Characterize boundedness of the family.}
If the diagonal computation never halts, $r_e$ is constantly one. If
it halts, $r_e$ is eventually $n+1$ and cannot be eventually bounded.

\item \textbf{Deduce failure of semidecidability.}
Hence $r_e=O(1)$ is exactly the complement of diagonal halting. A
$\Sigma$ factorization of this membership predicate would semidecide
that complement, contradicting \cref{lem:standalone-halting}.

\item \textbf{Verify admission of the principal resource.}
The principal ideal $\downarrow[1]$ is admitted by the unit of $\J$.
\end{enumerate}
\end{proof}

The pairing \eqref{eq:resource-positive-pairing} supplies the positive resource observations. Matrix equations over germ coefficients are interpreted in $\Omega$; finite certificates have the checking and search maps constructed in \cref{app:certificates}.

\subsection{Gauge localization, soft notation, and coefficient descent}
\label{sec:transport}
\subsubsection{Comparison with asymptotic gauges}
The asymptotic gauges of
\citep{GiordanoLuperiBaglini2016AsymptoticGauges} define moderate
rings by domination by a member of a scale family.
The following comparison uses rational-valued profiles and the
same eventual filter as our germ algebra. It concerns moderate
elements before a negligible-element quotient.

\begin{definition}[Rational asymptotic gauge]
A rational asymptotic gauge is a family $\mathcal G$ of exact
rational germs with a member tending to $+\infty$, such that
for $g,h\in\mathcal G$ there are $k,l\in\mathcal G$ with
$gh=O(k)$ and $|g|+|h|=O(l)$.
Rational rescaling is absorbed by the $O$-relation.
Its generated growth ideal is
$I_{\mathcal G}=\downarrow\{[|g|]:g\in\mathcal G\}$.
\end{definition}

\begin{proposition}[Gauge-moderate germs]\label{prop:asymptotic-gauge-comparison}
The family $I_{\mathcal G}$ is an inhabited directed growth ideal,
and its marked resource module is
\[
 M_{I_{\mathcal G}}
   =\{f\in K:\exists g\in\mathcal G,\ f=O(g)\}
   =\bigcup_{g\in\mathcal G}B|g|.
\]
It is a unital $B$-subalgebra of $K$. Thus rational moderate
profiles modulo eventual equality give exactly this marked algebra.
If $\mathcal G$ and $\mathcal H$ dominate one another memberwise,
their ideals and marked algebras coincide.
\end{proposition}
\begin{proof}\leavevmode
\begin{enumerate}[leftmargin=*,align=left]
\item \textbf{Verify directedness of the moderate resource.}
The divergent member makes $\mathcal G$ inhabited and dominates
$1$. The sum condition provides a common upper growth bound
for any two generators, so their lower closure is directed.

\item \textbf{Identify its resource module.}
Membership in $M_I$ means precisely that the growth of the
absolute value belongs to $I$; bounded division identifies this
with membership in $B|g|$ for some $g$. This proves both equalities.

\item \textbf{Construct the moderate subalgebra.}
The sum and product conditions respectively close the displayed
union under addition and multiplication; bounded multiplication
and negation preserve domination. The divergent member puts
$1$, hence all of $B$, in this union.

\item \textbf{Descend multiplication to eventual classes.}
Eventual equality preserves every domination condition, proving
the quotient statement. Memberwise domination gives inclusions
of the generated lower ideals in both directions.
\end{enumerate}
\end{proof}

A general resource module need not have the gauge's multiplicative
closure: on cofinite numerical input $Bn$ contains $n$ but not
$n^2$. Also $B$ itself contains no divergent germ. These facts
distinguish the full resource classifier from its gauge-generated
algebra sector. Admission of $I_{\mathcal G}$ is still the uniform
positive-pairing property of \cref{def:standalone-admission};
it is not inferred solely from the algebraic gauge conditions.
Colombeau constructions additionally quotient moderate elements
by negligible ones, whereas $K$ uses eventual equality.
The comparison above therefore does not identify those quotient
algebras with $K$ or with its Weil thickenings.
Functorial changes of asymptotic gauges are studied in
\citep{LuperiBagliniGiordano2017CategoryColombeau}; the input and
coefficient transports below retain their specified filter and
module maps.


\subsubsection{Overhead gauges}
A chosen family of overheads changes which multipliers count as
bounded. We encode that family before changing the coefficient ring.
The gauge object is the subobject of $\Pow((\Qplus)^X)$ consisting
of multiplicative families $\mathcal P$ of profiles $p\ge1$ containing
$1$. Its unit and finite products are part of this definition.
For example, powers of the numerical bitlength profile below form
such a family by NNO recursion. Its downward growth closure is
closed under finite joins, since $p+q\le2pq$. Define
\[
 f\in O_{\mathcal P}(g)\iff
 \exists p\in\mathcal P,\quad f\in O(pg),
\]
and define $\Omega_{\mathcal P}$ by reversal and
$\Theta_{\mathcal P}$ by mutual comparison.
The corresponding coefficient ring inside $K$ is
\[
 B_{\mathcal P}=\bigcup_{p\in\mathcal P}Bp.
\]
The displayed directed union is a ring: products multiply
gauges, and sums use the gauge $pq$.

\begin{lemma}[Admitted multiplication of growth ideals]
\label{lem:admitted-growth-product}
Let $L=\mathsf L_{\Ev}$ with its monotone multiplication
$m:L\times L\to L$. For $P,I\in\J L$ put
$P\star I=\J m(P\times I)$. Then $P\star I\in\J L$;
these products are associative and commutative with principal unit
$\downarrow[1]$. For every positive upper predicate $u:L\to\Sigma$,
\[
 \chi_{P\star I}(u)=
 \chi_P\bigl(p\mapsto\chi_I(x\mapsto u(px))\bigr).
\]
Consequently an admitted multiplicatively closed growth ideal
$P$ containing $[1]$ induces an inflationary idempotent operation
$I\mapsto P\star I$ on admitted ideals.
\end{lemma}
\begin{proof}\leavevmode
\begin{enumerate}[leftmargin=*,align=left]
\item \textbf{Construct admitted multiplication.}
Product preservation for $\J$ makes $P\times I$ an admitted ideal
of $L\times L$; functoriality then admits its lower image under $m$.

\item \textbf{Compute the multiplication formula.}
Expanding the lower-image pairing gives the displayed iterated
existential. Alternatively the inner pairing belongs uniformly to
$\Sigma$ by admission of $I$, is monotone in $p$, and is therefore a
positive upper test admitted by $P$.

\item \textbf{Verify algebraic laws and saturation.}
Associativity, commutativity and
the unit follow from the corresponding monoid identities on $L$ by
functoriality. If $P$ contains $[1]$ and is closed under products,
then $P\star P=P$: one inclusion is multiplicative closure and the
other uses the unit. Thus $I\le P\star I$ and
$P\star(P\star I)=P\star I$.
\end{enumerate}
\end{proof}

\begin{theorem}[Gauge localization]\label{thm:gauge}
Let $D=\{p^{-1}:p\in\mathcal P\}\subseteq B$.
Then
\begin{equation}\label{eq:gauges}
 B_{\mathcal P}\cong D^{-1}B,\qquad
 O_{\mathcal P}(g)=B_{\mathcal P}g,\qquad
 B_{\mathcal P}\otimes_BM_I\cong M_{\mathcal P I}.
\end{equation}
The same construction on exact positive coefficients commutes
with $\Gr$. The induced resource saturation is an inflationary,
idempotent, finite-meet-preserving operation.
\end{theorem}
\begin{proof}\leavevmode
\begin{enumerate}[leftmargin=*,align=left]
\item \textbf{Identify the gauge localization.}
Each $p^{-1}$ is bounded, regular, and invertible in $K$.
A fraction $b/p^{-1}$ is $bp$, and every element of the
displayed union has this form. This proves the localization
isomorphism and its universal property.

\item \textbf{Compute localized resource modules and tensors.}
Bounded division gives
the second formula. The tensor-product theorem for marked
resources, or its balanced-map proof on principal cells,
gives the third, with marking $a\otimes m\mapsto am$.

\item \textbf{Verify compatibility with group completion.}
Positive fractions have the same presentation and their
differences give precisely the signed fractions, proving
compatibility with $\Gr$.

\item \textbf{Prove saturation and finite-meet preservation.}
Containment of $1$ makes saturation inflationary, and
multiplicative closure makes it idempotent. Localization is
flat: a finite relation of fractions is cleared by one product
of denominators and is then a relation before localization.
Thus it preserves the kernel representing an intersection
of two submodules. This proves finite-meet preservation.

\item \textbf{Restrict saturation to admitted resources.}
For an admitted multiplicative growth ideal, the induced operation
on admitted resources is precisely the $\J$-product of
\cref{lem:admitted-growth-product}. Its outputs are admitted by that
lemma; finite-meet preservation on full resource modules is the
localization calculation above.
\end{enumerate}
\end{proof}

Equivalently an admitted multiplicative ideal $P$ containing
$[1]$ has coefficient algebra $M_P$. Its denominators can be
taken as $(1+g)^{-1}$ for $[g]\in P$, $g\ge0$:
\[
 f=\frac{f/(1+g)}{(1+g)^{-1}}\quad\text{whenever }f\in O(g).
\]
The product of denominators for $g,h$ is the denominator for
$g+h+gh$. This proves the same localization directly from
the ideal presentation.

\subsubsection{Soft Landau notation}
On numerical input put
$\lambda(n)=1+\operatorname{bitlength}(n+1)$.
This is an exact computable integer profile, comparable to
$1+\log_2(n+1)$ by fixed positive constants. Its finite powers
form a gauge. Define
\[
 \widetilde O(g)=\bigcup_{k\ge0}O(\lambda^k g),\quad
 f\in\widetilde\Omega(g)\iff g\in\widetilde O(f),\quad
 \widetilde\Theta=\widetilde O\cap\widetilde\Omega.
\]
Then $\widetilde O(g)=B[\lambda]g$, where
$B[\lambda]=B[(\lambda^{-1})^{-1}]$ inside $K$.
In particular $f\in\widetilde\Theta(g)$ means equality of
their principal modules after this localization.
The profile uses finite bit-length computation, so no logarithm
operation on arbitrary lower-real names is needed.

\begin{proposition}[The admitted polylogarithmic gauge]
\label{prop:admitted-soft-gauge}
In the realizability models of \cref{sec:foundation}, the growth ideal
\[
 P_\lambda=\downarrow\{[\lambda^k]:k\in\N\}
\]
is $\J$-admitted, contains $[1]$, and is closed under growth
multiplication and finite joins.  Consequently
$I\mapsto P_\lambda\star I$ is a constructed inflationary,
idempotent soft-resource operation on $\J\mathsf L_{\Ev}$.
\end{proposition}
\begin{proof}\leavevmode
\begin{enumerate}[leftmargin=*,align=left]
\item \textbf{Prove admission of the power ideal.}
The powers are increasing because $\lambda\ge1$, and their
lower closure is an ideal.  For a positive upper predicate
$u:\mathsf L_{\Ev}\to\Sigma$,
\[
 \chi_{P_\lambda}(u)
 =\exists k\in\N\,u([\lambda^k])\in\Sigma
\]
by overtness of the NNO, uniformly in $u$.

\item \textbf{Verify multiplicative closure and induced saturation.}
The unit is the
zeroth power; multiplication adds exponents.  Also
$[\lambda^j]\vee[\lambda^k]\le[\lambda^{j+k}]$.
The admitted-growth-product lemma now constructs the stated
operation and proves its laws.
\end{enumerate}
\end{proof}

\begin{proposition}[Vanishing under overhead inversion]
\label{prop:gauge-vanishing}
If $s\in\mathfrak m$ is inverted by a gauge localization,
then $\mathfrak m B_{\mathcal P}=B_{\mathcal P}$.
In particular this occurs when $s=t$, $t(n)=1/(n+1)$,
is inverted, and when $s=\lambda^{-1}$ is inverted.
\end{proposition}
\begin{proof}\leavevmode
\begin{enumerate}[leftmargin=*,align=left]
\item \textbf{Show that inversion makes the ideal contain the unit.}
The element $1=s^{-1}s$ belongs to the extended ideal.

\item \textbf{Verify the stated vanishing profiles.}
The profiles $t$ and $\lambda^{-1}$ satisfy each vanishing
bound by their explicit unbounded denominators.
\end{enumerate}
\end{proof}
Thus gauge-relative domination is calculated in the localized
ring, while original small-error incidences remain recorded
in the exact presentation over $B$.
The extended ideal $\mathfrak m B_{\mathcal P}$ is not
redefined to be the vanishing ideal of the original numerical
filter. This formula specifies exactly how its marking changes.

\subsection{Eventual quotients and continuous completion}
\begin{lemma}[Admitted rational approximation]\label{lem:admitted-rationals}
In the realizability models of \cref{sec:foundation}, put
$r_n=(n+1)/(n+2)$ and
$I=\downarrow\{\underline{r_n}:n\in\N\}\subseteq S$.
Then $I\in\J S$ and $\sup I=1$.
\end{lemma}
\begin{proof}\leavevmode
\begin{enumerate}[leftmargin=*,align=left]
\item \textbf{Prove admission by a countable presentation.}
The family is inhabited and increasing, hence its lower closure is
an ideal.  For every monotone $u:S\to\Sigma$,
\[
 \chi_I(u)=\exists s\in I\,u(s)
 \quad\Longleftrightarrow\quad\exists n\in\N\,u(\underline{r_n}).
\]
The right-hand existential factors uniformly through $\Sigma$
by \cref{prop:standalone-dominance}.  Hence $I$ passes the defining $\J$ pairing test.

\item \textbf{Compute the supremum.}
Every $r_n\le1$, and for each positive rational $q<1$ the
Archimedean property of rational arithmetic gives an $n$ with
$q<r_n$.  Equality of the resulting lower cuts gives $\sup I=1$.
\end{enumerate}
\end{proof}

\begin{proposition}\label{prop:collapse}
In the realizability models of \cref{sec:foundation}, let $X=\N$ carry the cofinite filter, let $t(n)=1/(n+1)$, and let $B^+$ be the semiring of eventually bounded nonnegative rational germs.  In $B^+$, $1+t\ne1$, but their images in $\Ext B^+$ are equal.
\end{proposition}
\begin{proof}\leavevmode
\begin{enumerate}[leftmargin=*,align=left]
\item \textbf{Separate the germs before completion.}
Pointwise $t>0$, so the germ equality is false.

\item \textbf{Construct the rational comparison.}
Let $x$ be the image of $1+t$ and let $e$ be the image of $1$. For each rational $0\le r<1$, the inequality $r(1+t)\le1$ holds eventually. The nonnegative rational profile
\[
 c_r=\max\{1-r(1+t),0\}
\]
satisfies $r(1+t)+c_r=1$ as a germ. Therefore $rx\le e$ in the positive order of $\Ext B^+$.

\item \textbf{Use scalar continuity to identify the images.}
The admitted ideal of \cref{lem:admitted-rationals} has supremum $1$, so scalar continuity gives $x=\sup_n r_nx\le e$. Conversely $e\le x$ by positivity of the image of $t$. Antisymmetry gives $x=e$.
\end{enumerate}
\end{proof}

Continuous completion and eventual equality therefore have different
effects on a cost profile. The construction keeps the exact profile
and forms its eventual relations there; the presentation object in
$\Pcat$ records the quantitative interpretation alongside it.
In particular, \eqref{eq:raw-realization} remains available before the
eventual quotient.

The resource algebra now supplies embedded modules, their tensor and
lattice operations, and the coefficient localizations associated with
overhead gauges. We next classify these resources in families, so that
one construction can express a budget varying over a parameter space.


\section{The universal structured resource space}
\label{sec:complexity-locale}

The preceding section describes individual budgets as embedded modules.
To vary a budget over parameters, we need a classifier carrying a
universal family. The growth order already determines its topology:
a basic open expresses availability of one growth class.
For $L=\mathsf L_{\Ev}$, form the frame of upper predicates
\[
 \Up(L)=\{U:L\to\Omega:\forall a\le b,\ U(a)\Rightarrow U(b)\}.
\]
Write $\Cplx$ for the internal locale with
$\mathcal O(\Cplx)=\Up(L)$. We use the standard internal locale and
localic-sheaf constructions \citep{Vickers,MacLaneMoerdijk}.
The following results identify its generalized points with marked
resource families and construct their universal module. Thus the
geometry is determined by growth comparison itself.

\subsection{Ideal classification and the Scott frame}
\begin{theorem}[Canonical classifying frame]
\label{thm:classifying-frame}
For every internal frame $H$, frame maps
$\Up(L)\to H$ are naturally equivalent to maps $I:L\to H$
satisfying
\begin{equation}\label{eq:frame-ideals}
 I(0)=\top,\qquad
 a\le b\Rightarrow I(b)\le I(a),\qquad
 I(a)\wedge I(b)=I(a\vee b).
\end{equation}
Consequently
\[
 \operatorname{Pt}_{\E}(\Cplx)\cong\Idl(L),\qquad
 \mathcal O(\Cplx)\cong\Sc(\Idl(L)).
\]
The principal opens $U_a=\{b:a\le b\}$ are supercompact,
are closed under finite intersections, and generate the frame.
Thus $\Cplx$ is a coherent internal locale.
\end{theorem}
\begin{proof}\leavevmode
\begin{enumerate}[leftmargin=*,align=left]
\item \textbf{Extract an ideal from a frame map.}
Every upper predicate has the canonical expansion
$U=\bigvee_{a\in U}U_a$. Because $L$ has bottom and
finite joins,
\[
 U_0=\top,\qquad U_a\wedge U_b=U_{a\vee b}.
\]
Given a frame map $\theta:\Up(L)\to H$, put
$I_\theta(a)=\theta(U_a)$. The displayed identities and
$U_b\subseteq U_a$ for $a\le b$ prove
\eqref{eq:frame-ideals}.

\item \textbf{Construct a frame map from an ideal.}
Conversely, given $I$ with those equations, define
\[
 \theta_I(U)=\bigvee_{a\in U}I(a).
\]
Arbitrary-join preservation is immediate. Frame distributivity gives
\[
 \theta_I(U)\wedge\theta_I(V)
 =\bigvee_{a\in U,b\in V}I(a\vee b)
 =\theta_I(U\cap V).
\]
For the last equality, $a\vee b\in U\cap V$ gives one
inequality, while each $c\in U\cap V$ is obtained with $a=b=c$
for the other. Preservation of top follows from $I(0)=\top$.

\item \textbf{Verify the classification and its points.}
The inverse equations are
\[
 \theta_{I_\theta}(U)=\bigvee_{a\in U}\theta(U_a)=\theta(U),
 \qquad I_{\theta_I}(a)=\bigvee_{b\ge a}I(b)=I(a).
\]
For $H=\Omega$, conditions \eqref{eq:frame-ideals} say exactly
that $I$ is an inhabited directed lower ideal: bottom belongs
to it, and the join of two members is a member. This gives
the asserted internal point comparison.

\item \textbf{Identify the Scott frame.}
For the Scott-frame comparison, map an upper predicate $U$ to
\[
 \widehat U(I)\iff\exists a\in I\;U(a).
\]
It is upper and preserves inhabited directed joins of ideals,
which are unions. Conversely send a Scott-open $W$ to
$u_W(a)=W(\downarrow a)$. Each ideal is the directed union
$I=\bigvee_{a\in I}\downarrow a$, so Scott openness gives
$W(I)\iff\exists a\in I\,W(\downarrow a)$.
Also $\widehat U(\downarrow a)\iff U(a)$ by upperness.
These are the two inverse equations.

\item \textbf{Prove supercompactness and coherence.}
Finally, if
$U_a\le\bigvee_jV_j$, evaluating at $a$ gives locally a
witness $j$ with $a\in V_j$; upperness gives $U_a\le V_j$.
This proves supercompactness internally. The principal opens
form a finite-meet-stable generating family, including top,
so the locale is coherent.
\end{enumerate}
\end{proof}

\subsection{Embedded modules and the universal resource family}
The ideal classification now extends to families of embedded modules.
For an internal locale $Z$, write
$p_Z:\Sh_{\E}(Z)\to\E$ for its localic geometric morphism and
$\Delta_Z=p_Z^*$ for inverse image. A budget over $Z$ is a module
subobject in $\Sh_{\E}(Z)$; this allows its membership conditions
to vary locally with the parameter.

The family-valued moduli problem is the contravariant indexed poset
\[
 \mathcal B:\Loc_\E^{\mathrm{op}}\longrightarrow\mathbf{Pos},
 \qquad
 \mathcal B(Z)=\Sub_{\Delta_ZB}(\Delta_ZK),\qquad
 \mathcal B(f)=f^*.
\]
Each object retains its embedding into the ambient module.
The next theorem represents this indexed poset by $\Cplx$.
Its naturality says that reindexing a budget family is composition
with its classifying map, and hence determines the universal family.


\begin{theorem}[Localic resource-module classification]
\label{thm:locale-modules}
There is a natural order isomorphism
\begin{equation}\label{eq:locale-modules}
 \operatorname{Loc}_{\E}(Z,\Cplx)
 \ \cong\ \Sub_{\Delta_ZB}(\Delta_ZK),
\end{equation}
where the right-hand side is the poset of
$\Delta_ZB$-submodules of the fixed ambient module $\Delta_ZK$.
Under it, the identity of $\Cplx$ classifies a canonical marked
module
\begin{equation}\label{eq:universal-module}
 \mathscr M\incl\Delta_{\Cplx}K
 \qquad\text{in }\Sh_{\E}(\Cplx).
\end{equation}
For a generalized ideal point $I$, its pullback is $M_I\incl K$.
\end{theorem}
\begin{proof}\leavevmode
\begin{enumerate}[leftmargin=*,align=left]
\item \textbf{Express subobjects by open-valued membership.}
The localic subobject-classifier formula
$p_{Z*}\Omega_{\Sh_{\E}(Z)}\cong\mathcal O(Z)$ and the
adjunction $\Delta_Z\dashv p_{Z*}$ give
\[
 \Sub_{\Sh_{\E}(Z)}(\Delta_ZK)
 \cong\E(K,\mathcal O(Z)).
\]
Thus a subobject has an open-valued membership function
$\lambda:K\to\mathcal O(Z)$. It is a module precisely when
the equalities and implications expressing membership of zero,
addition, negatives, and multiplication by $B$ hold in this
frame-valued logic.

\item \textbf{Descend membership to growth classes.}
The sign multiplier $\operatorname{sgn}(f)\in B$ satisfies
$\abs f=\operatorname{sgn}(f)f$ and
$f=\operatorname{sgn}(f)\abs f$; hence
$\lambda(f)=\lambda(\abs f)$. If $f,g\ge0$ have equal
growth classes, the bounded-division theorem supplies
multipliers in both directions, giving $\lambda(f)=\lambda(g)$.
Therefore positive membership descends across the effective
growth quotient $K^+\onto L$ to a map
$I_\lambda:L\to\mathcal O(Z)$.

\item \textbf{Verify the ideal equations.}
It satisfies $I_\lambda(0)=\top$. Bounded division gives
lowerness. On positive elements, module closure gives
$\lambda(f)\wedge\lambda(g)\le\lambda(f+g)$,
and both $f,g$ are bounded multiples of $f+g$, giving
\[
 \lambda(f)\wedge\lambda(g)=\lambda(f+g).
\]
Since $[f+g]=[f]\vee[g]$, these are precisely
\eqref{eq:frame-ideals}.

\item \textbf{Recover the module from its ideal.}
Conversely an
$\mathcal O(Z)$-valued ideal $I$ defines
$\lambda_I(f)=I([\abs f])$. The triangle inequality,
$\abs{bf}\preceq\abs f$ for $b\in B$, and
\eqref{eq:frame-ideals} prove closure under all module
operations. The two constructions are inverse on their
membership predicates.

\item \textbf{Classify modules by locale maps.}
By \cref{thm:classifying-frame}, the resulting
$\mathcal O(Z)$-valued ideals are precisely frame maps
$\Up(L)\to\mathcal O(Z)$, equivalently locale maps
$Z\to\Cplx$. Every step is natural in $Z$ by pullback of
membership predicates.

\item \textbf{Identify the universal module and its fibres.}
At $Z=\Cplx$, the identity frame
map defines \eqref{eq:universal-module}, with membership
of $f\in K$ given by the open
\begin{equation}\label{eq:universal-membership}
  \llbracket f\in\mathscr M\rrbracket=U_{[\abs f]}.
\end{equation}
Evaluating this at an ideal point $I$ gives
$[\abs f]\in I$, exactly the membership predicate of
$M_I$ from \cref{sec:resources}. This proves the fibre assertion.
\end{enumerate}
\end{proof}


For $b:Z\to\Cplx$, the classified family is
\begin{equation}\label{eq:cat-classified-budget-pullback}
 \mathscr M_b=b^*\mathscr M\incl\Delta_ZK,
 \qquad \mathscr M_{b\circ f}=f^*\mathscr M_b.
\end{equation}
Thus the category of marked budget families and cartesian maps is
$\Loc_\E/\Cplx$, with universal family corresponding to
$\id_{\Cplx}$.

\subsection{Generalized points as families of budgets}
\label{sec:budget-functor-of-points}
\label{app:complexity-functor-of-points}
\label{app:complexity-generalized-points}

The classification can be used directly to evaluate a resource
condition on a varying family. Retain the indexed budget functor
$\mathcal B$ and the coordinate identities
\[
 U_a\wedge U_b=U_{a\vee b},\qquad U_0=\top,\qquad
 U_a(I)\iff a\in I,\qquad M_{\downarrow[g]}=Bg=O(g).
\]
Write $I\models U$ for evaluation of an open $U$ at an ideal point
$I$. This notation expresses an internal truth value.

The budget functor of points is the classification
\begin{equation}\label{eq:app-budget-functor}
 h_{\Cplx}(Z)=\Loc_\E(Z,\Cplx)
 \cong\Sub_{\Delta_ZB}(\Delta_ZK)
\end{equation}
of \cref{thm:locale-modules}. The following example evaluates a
resource condition on a nonconstant family.

\paragraph{A nonconstant family of budgets.}
To exhibit a nonconstant budget family, take the Sierpi\'nski locale
$Z=\mathbb S$ itself, and write $u\in\OO(Z)$ for its distinguished
proper open.  Let $a_0=[1]$ and $a_1=[n+1]$; thus $a_0<a_1$ in the
numerical cofinite growth order.  Construct the following
$\OO(Z)$-valued ideal on $L$:
\begin{equation}\label{eq:app-generic-budget-family}
 I_Z(c)=
 \llbracket c\le a_0\rrbracket_Z
 \ \vee\
 \bigl(u\wedge\llbracket c\le a_1\rrbracket_Z\bigr),
 \qquad c\in L.
\end{equation}
Here $\llbracket P\rrbracket_Z$ denotes the constant open of $Z$
classified by the internal truth value $P$; it is not assumed that
arbitrary order relations on $L$ are decidable.  Since $a_0\le a_1$,
the expression in \eqref{eq:app-generic-budget-family} sends $0$ to
top, is lower in $c$, and satisfies
$I_Z(c)\wedge I_Z(d)=I_Z(c\vee d)$.  Indeed the component with no
factor $u$ requires both $c,d\le a_0$, whereas all remaining terms
require $u$ and both $c,d\le a_1$.  By
\cref{thm:classifying-frame,thm:locale-modules} it defines a
generalized point $b:Z\to\Cplx$.  Its budget is constant-bounded at
the closed specialization and linear-bounded on the distinguished
open $u$.

\begin{equation}\label{eq:app-nonconstant-probe}
 b^*U_{a_1}
 =\llbracket a_1\le a_0\rrbracket_Z
   \vee\bigl(u\wedge\llbracket a_1\le a_1\rrbracket_Z\bigr)
 =u.
\end{equation}
Thus a resource condition can hold on a proper open of the parameter
locale.  The family-valued classification records this varying locus,
whereas an individual ideal point records only its specialization.
In \cref{app:complexity-size-example} the same family will be probed
by a measured numerical implementation.

\begin{theorem}[Flat universal resource module]
\label{thm:universal-flat}
The universal module $\mathscr M$ is an internal directed colimit
of finite projective principal modules and is flat over
$\Delta_{\Cplx}B$.
\end{theorem}
\begin{proof}\leavevmode
\begin{enumerate}[leftmargin=*,align=left]
\item \textbf{Construct projective principal stages.}
Work in the internal logic of $\Sh_{\E}(\Cplx)$.
The bounded-division maps and support idempotents are obtained by
pullback of their parameter-uniform constructions in
\cref{lem:positive-division,prop:signed-coefficients}.
Every local element $f\in\mathscr M$ generates
$Bf\cong Be_f$, a split summand of the rank-one free
module. Its inclusion into $\mathscr M$ is the multiplication
map with this marking.

\item \textbf{Express the universal module as a filtered colimit.}
Principal submodules are directed:
$f,g\in\mathscr M$ have common upper principal stage
$B(\abs f+\abs g)\subseteq\mathscr M$.
They cover $\mathscr M$, because $f\in Bf$. Their internal
filtered colimit is therefore $\mathscr M$.

\item \textbf{Deduce flatness from finite relations.}
Each principal stage is projective and flat. A finite module
relation in this colimit is, locally, represented at finitely
many stages; directedness supplies one common principal stage,
where projectivity resolves it. This finite-equation proof of
filtered-colimit flatness applies in the internal module category
(cf. \citep[Tag~05UT]{Stacks}), proving flatness.
\end{enumerate}
\end{proof}

\subsection{The ringed resource topos and its diagram presentation}
\label{sec:structured-resources}
The universal family now equips the budget classifier with coefficients.
Put $\mathscr S=\Sh_\E(\Cplx)$, $\bar B=\Delta_{\Cplx}B$ and
$\bar K=\Delta_{\Cplx}K$. The structured resource space is
$(\mathscr S,\bar B)$ with its marked module $\mathscr M\incl\bar K$.
A cartesian map $(Z,N)\to(Z',N')$ is a locale map $f:Z\to Z'$
with $N=f^*N'$ as marked submodules. By
\cref{thm:locale-modules}, these families form a category equivalent
to $\Loc_\E/\Cplx$. Allowing $N\subseteq f^*N'$ instead gives the
indexed poset of resources.

The marking fixes the image in the ambient module. In particular,
$i\alpha=i$ forces a marked automorphism $\alpha$ to be the identity.
Choices of a presentation can still have change-of-frame maps.
The following diagram description makes both the coefficients and
their budget transitions explicit.

\begin{theorem}[Resource sheaves as budget diagrams]
\label{thm:budget-diagrams}
There is an equivalence of topoi over $\E$
\begin{equation}\label{eq:budget-diagrams}
 \Sh_\E(\Cplx)\simeq[L,\E],\qquad
 F\longmapsto\bigl(a\longmapsto F(U_a)\bigr).
\end{equation}
The right side is the category of internal covariant diagrams on
the internal poset $L$. Under this equivalence $\bar B,\bar K$
are constant diagrams and
\[
 \mathscr M(a)=M_a:=M_{\downarrow a},\qquad
 (a\le b)\longmapsto(M_a\incl M_b).
\]
At a generalized ideal point $I$, inverse image is
$I^*F\cong\varinjlim_{a\in I}F(a)$.
\end{theorem}
\begin{proof}\leavevmode
\begin{enumerate}[leftmargin=*,align=left]
\item \textbf{Construct a sheaf from a budget diagram.}
The principal opens form a basis, and $a\le b$ gives the restriction
$F(U_a)\to F(U_b)$. Supercompactness says that a covering sieve on
this basis contains the identity internally; thus no additional
matching condition is imposed on a diagram. Explicitly, extend a
diagram $D$ by
\[
 \widetilde D(U)=\{(x_a)_{a\in U}:x_a\in D(a),\
            D(a\le b)x_a=x_b\text{ for }a\le b\text{ in }U\}.
\]

\item \textbf{Verify gluing and the inverse equivalence.}
These compatible families glue uniquely across unions of opens.
On $U_a$, the element $a$ is initial in its indexing poset, so
$\widetilde D(U_a)\cong D(a)$. Principal decomposition of an open
proves the other inverse equation. All constructions take place
in the internal logic, including in parameter slices.

\item \textbf{Identify constant and universal module diagrams.}
Constant diagrams represent $\Delta_{\Cplx}$, and
\eqref{eq:universal-membership} gives the stated module diagram.

\item \textbf{Compute the fibre at an ideal.}
The inverse image at $I$ is the colimit over its principal
neighbourhoods; $I$ is inhabited and directed. Equivalently this
colimit preserves finite limits by taking a common stage for a
finite tuple and its equations, and sends $U_b$ to the truth value
$b\in I$. It is therefore the inverse image classified by $I$.
\end{enumerate}
\end{proof}


Evaluation at a generalized ideal point is compatible with the
diagram presentation:
\[
\begin{tikzcd}[column sep=large,row sep=large]
 \Sh_\E(\Cplx) \arrow[r,"{\simeq}"] \arrow[d,"{I^*}"']
   & {[L,\E]} \arrow[d,"{{\varinjlim_{a\in I}}}"]\\
 \E \arrow[r,equal] & \E .
\end{tikzcd}
\]
The commutativity is the canonical natural isomorphism of inverse-image
functors, including their action on marked modules and witness sheaves.


This equivalence preserves the observation data specified in
\cref{sec:resource-observations}: the effective points are $\J L$
with their positive-test pairing. Membership in a principal open
$U_a$ remains the corresponding $\Omega$-valued resource condition.

The individual resource operations now lift to universal families.
Tensor multiplication, sums, and intersections must preserve the
markings as well as the underlying modules. The next results construct
these operations on the classifier and compare them with admission.

\subsection{Multiplication, lattice structure, and admission}
\begin{theorem}[Localic resource multiplication]
\label{thm:localic-multiplication}
There is a commutative monoid structure
\[
 \mu:\Cplx\times\Cplx\to\Cplx,
 \qquad e:1\to\Cplx,
\]
whose generalized-point formulas are
$\mu(I,J)=I\star J=\downarrow\{ab:a\in I,b\in J\}$ and
$e=\downarrow[1]$. On principal opens,
\begin{equation}\label{eq:localic-product-open}
 \mu^*U_g
 =\bigvee_{g\le ab}
    (\pi_1^*U_a\wedge\pi_2^*U_b).
\end{equation}
For the universal module there is a canonical coherent
isomorphism of marked modules over $\Cplx\times\Cplx$:
\begin{equation}\label{eq:universal-tensor}
 \mu^*\mathscr M
 \ \cong\ \pi_1^*\mathscr M
     \otimes_{\Delta B}\pi_2^*\mathscr M.
\end{equation}
Finite lattice joins and meets of resources likewise define
locale maps, represented by sums and intersections of their
universal modules.
\end{theorem}
\begin{proof}\leavevmode
\begin{enumerate}[leftmargin=*,align=left]
\item \textbf{Construct the product ideal and locale map.}
On $\Cplx\times\Cplx$, let $I_1,I_2$ denote the two
universal frame-valued ideals. Define
\[
 (I_1\star I_2)(g)
 =\bigvee_{g\le ab}\bigl(I_1(a)\wedge I_2(b)\bigr).
\]
This is inhabited because both ideals contain zero, and is
lower by construction. For two witnesses $a_1b_1,a_2b_2$,
the join witnesses $a=a_1\vee a_2$ and $b=b_1\vee b_2$
lie in the respective ideals and dominate both products.
The ideal-classification theorem therefore supplies a locale
map $\mu$ with exactly the inverse-image formula
\eqref{eq:localic-product-open}.

\item \textbf{Verify the monoid laws.}
Multiplication of growth
classes proves associativity and commutativity; multiplication
by $[1]$ proves the unit laws. Uniqueness in the classifying
property turns these into equations of locale maps.

\item \textbf{Construct the principal tensor comparison.}
The pullback $\mu^*\mathscr M$ is the resource module of the
product ideal. Multiplication in $\Delta K$ supplies a balanced
map from $\pi_1^*\mathscr M\times\pi_2^*\mathscr M$ into it.
At principal stages it induces the isomorphisms
\[
 Bf\otimes_BBg\cong Be_f\otimes_BBe_g
 \cong Be_fe_g\cong Bfg.
\]
Their inverses are the support-idempotent inverse maps of
\cref{thm:signed-rep}; their naturality under inclusion into
$K$ was verified in \cref{sec:resources}.

\item \textbf{Extend the comparison and verify coherence.}
Tensor preserves
the internally directed colimits of these stages. Passing to
the colimits proves \eqref{eq:universal-tensor}.
Both associativity composites take $f\otimes g\otimes h$ to
$fgh$, and the unit and symmetry composites have the analogous
formulas. Equality on principal stages proves coherence.

\item \textbf{Construct the lattice operations.}
For meets, the inverse image of a principal open is
$U_g\mapsto\pi_1^*U_g\wedge\pi_2^*U_g$.
For joins, it is
\[
 U_g\longmapsto
 \bigvee_{g\le a\vee b}
  (\pi_1^*U_a\wedge\pi_2^*U_b).
\]
These formulas classify intersection and sum of resource modules
by \cref{thm:signed-rep}; preservation of frame operations follows
from the classifying ideal property, proving the final assertion.
\end{enumerate}
\end{proof}


The universal tensor isomorphism preserves the markings:
\[
\begin{tikzcd}[column sep=large,row sep=large]
 \pi_1^*\mathscr M\otimes_{\Delta B}\pi_2^*\mathscr M
   \arrow[r,"{\cong}"] \arrow[d]
   & \mu^*\mathscr M \arrow[d,hook]\\
 \Delta K\otimes_{\Delta B}\Delta K
   \arrow[r,"{\mathrm{mult}}"'] & \Delta K .
\end{tikzcd}
\]
All objects in this square live over $\Cplx\times\Cplx$.
Associativity of this square is the family-valued version of the
monoidal coherence in \cref{cor:cat-monoidal-resources}.


The effective points of this locale are selected by the \emph{existing}
admission doctrine, not by declaring every principal open effective.
\begin{proposition}[Admitted points and positive observation]
\label{prop:localic-admission}
The comparison $\operatorname{Pt}_{\E}(\Cplx)\cong\Idl(L)$ identifies
$\J L$ with the subobject of points $I$ for which
\[
 \chi_I(u)=\exists a:L\,[I(a)\wedge u(a)]\in\Sigma
 \quad\text{uniformly in monotone }u:L\to\Sigma.
\]
Localic multiplication restricts on this subobject to the
admitted product $\J m$ of \cref{lem:admitted-growth-product}.
The canonical map from the discrete locale on $\J L$ to
$\Cplx$ pulls the universal module back to $I\mapsto M_I$.
\end{proposition}
\begin{proof}\leavevmode
\begin{enumerate}[leftmargin=*,align=left]
\item \textbf{Identify the admitted points.}
The first assertion is exactly the defining pullback of the
selected monad of \cref{def:standalone-admission}; the uniform factorization
is unchanged under the internal point isomorphism.

\item \textbf{Prove closure under multiplication.}
For an upper positive predicate $u$, the admitted-product
calculation reads
\[
 \chi_{I\star J}(u)
 =\chi_I\bigl(a\mapsto\chi_J(b\mapsto u(ab))\bigr).
\]
The inner pairing is uniformly $\Sigma$-valued and monotone;
the outer pairing is $\Sigma$-valued by admission. Thus the
product is admitted and agrees with $\J m$ by its
principal-weight calculation.

\item \textbf{Identify the pulled-back universal family.}
The canonical map of discrete
locales is the transpose of the inclusion
$\J L\incl\operatorname{Pt}_{\E}(\Cplx)$. Pulling
\eqref{eq:universal-membership} back along it gives
$I(a)$ on the principal growth class $a$, proving the
family statement.  The exact localic relation between
admitted and semantic resources is calculated in
\cref{prop:admitted-dense}.
\end{enumerate}
\end{proof}

\begin{proposition}[Admitted points are locale-dense]
\label{prop:admitted-dense}
Let $a:\operatorname{Disc}(\J L)\to\Cplx$ be the canonical
locale map from admitted points.  Its inverse-image map is
\[
 a^*:\Up(L)\longrightarrow\Omega^{\J L},
 \qquad (a^*U)(I)=\exists x\in I\,U(x).
\]
It preserves arbitrary joins and finite meets, and is
order-reflecting.  Consequently the localic image nucleus
$a_*a^*$ is the identity on $\Up(L)$: the localic image is
all of $\Cplx$.
\end{proposition}
\begin{proof}\leavevmode
\begin{enumerate}[leftmargin=*,align=left]
\item \textbf{Verify preservation of frame operations.}
Joins commute with existential quantification.  For finite
meets, two witnesses $x,y\in I$ have their join $x\vee y\in I$,
and upperness of $U,V$ gives the common witness; the reverse
inclusion follows by projecting a witness.  This proves that
$a^*$ is a frame homomorphism.

\item \textbf{Prove order reflection using principal ideals.}
Every principal ideal $\downarrow x$ is $\J$-admitted: its
pairing with a positive upper test $u$ is $u(x)$, as in
\cref{sec:foundation}.  Evaluation on it gives
\[
 (a^*U)(\downarrow x)\iff U(x).
\]
Hence $a^*U\le a^*V$ implies $U\le V$, so $a^*$ is
order-reflecting.

\item \textbf{Establish density by the frame adjunction.}
Since $a^*$ preserves arbitrary joins, it
has a right adjoint $a_*$.  The adjunction unit gives
$U\le a_*a^*U$, while the counit gives
$a^*a_*a^*U\le a^*U$; order reflection yields
$a_*a^*U\le U$.  Thus $a_*a^*=\id$.  This proves density of the admitted points.
\end{enumerate}
\end{proof}

Finite diagrams of the classifier parameterize compatible budgets for compiler maps and resource deformations.

\subsection{Classifying finite diagrams of budgets}
\label{sec:budget-diagrams-moduli}
\begin{theorem}[Finite resource-diagram classifiers]
\label{thm:budget-flags}
The equalizer
\[
 \operatorname{Flag}_2=\operatorname{Eq}
   (\vee,\pi_2:\Cplx^2\rightrightarrows\Cplx)
\]
classifies pairs of marked modules $M_0\subseteq M_1$ over every
localic parameter. More generally, a finite diagram of inclusions
between resource modules has a classifier obtained by finite
limits from products of $\Cplx$. Finite equations and inclusions
between resource expressions formed using sums, intersections,
and products have classifiers by the same construction.
\end{theorem}
\begin{proof}\leavevmode
\begin{enumerate}[leftmargin=*,align=left]
\item \textbf{Classify a single inclusion.}
A map to $\Cplx^2$ classifies two marked modules. Factoring
through the equalizer says $M_0+M_1=M_1$, equivalently
$M_0\subseteq M_1$, by \cref{thm:localic-multiplication}.

\item \textbf{Classify finite diagrams and resource equations.}
For a finite graph of inclusions, impose this equation for each
edge and take their common equalizer. Commutativity is automatic
for inclusions into the fixed ambient module. Resource expressions
are interpreted by the already constructed locale operations;
their equations are equalizers, and $R\subseteq S$ is the
equation $R+S=S$. These constructions commute with parameter
pullback because their stated universal properties do.
\end{enumerate}
\end{proof}

For a monotone growth transformation $r:L\to L'$, the inverse-image
frame map $r^{-1}$ constructs $s_r:\Cplx\to\Cplx'$.
The same construction gives a locale
\begin{equation}\label{eq:compiler-budget-classifier}
 \mathfrak R_r=\operatorname{Eq}
 \bigl(\vee\circ(s_r\pi_1,\pi_2),\pi_2:
                 \Cplx\times\Cplx'\rightrightarrows\Cplx'\bigr),
\end{equation}
classifying $\downarrow r(I)\subseteq J$. An implemented compiler
will act over this budget-compatibility locale in
\cref{sec:implementation-sheaf}, with its code supplied separately.

\subsection{Gauge saturation as a classifying sublocale}
\label{sec:gauge-sublocale}
Coefficient localization closes a resource under chosen overheads.
We now identify the families fixed by this operation and their
classifying sublocale. Let $P\in\Idl(L)$ be a multiplicatively closed
ideal containing $[1]$, regarded as a point of the object of such
ideals. Its coefficient ring is the localization $B_P=M_P\subseteq K$
of \cref{sec:transport}. Define
\[
 S_P:\Idl(L)\longrightarrow\Idl(L),\qquad I\longmapsto P\star I.
\]


\begin{theorem}[Gauge saturation and the localized moduli problem]
\label{thm:gauge-sublocale}
The map $S_P$ is Scott-continuous, inflationary and idempotent.
Consequently its inverse-image operation
\[
 j_P=S_P^{-1}:\OO(\Cplx)\longrightarrow\OO(\Cplx)
\]
is a nucleus.  It is determined on the principal opens by
\begin{equation}\label{eq:gauge-nucleus}
 j_P(U_g)=\bigvee_{\substack{p\in P,\ a\in L\\g\le pa}}U_a.
\end{equation}
The nucleus constructs a sublocale
$\iota_P:\Cplx^{(P)}\incl\Cplx$ satisfying, naturally in every
internal locale $Z$,
\begin{equation}\label{eq:localized-moduli-classification}
 \operatorname{Loc}_{\E}(Z,\Cplx^{(P)})
 \cong\operatorname{Sub}_{\Delta_ZB_P}(\Delta_ZK).
\end{equation}
Its points are precisely the $P$-saturated ideals $I=P\star I$,
and its universal marked module is $\iota_P^*\mathscr M$ with
its canonical $\Delta_{\Cplx^{(P)}}B_P$-action.  Resource
multiplication restricts to a commutative monoid on $\Cplx^{(P)}$
with unit $P$, and its universal marked tensor comparison is
balanced over $B_P$.
\end{theorem}
\begin{proof}\leavevmode
\begin{enumerate}[leftmargin=*,align=left]
\item \textbf{Construct the gauge nucleus.}
Products of ideals distribute over directed unions, so $S_P$
preserves directed joins.  The inequalities $I\le P\star I$ and
$P\star(P\star I)=I\star(P\star P)=P\star I$ follow from the
unit and product closure of $P$.  Under the Scott-frame
identification in \cref{thm:classifying-frame}, inverse image
along $S_P$ preserves arbitrary joins and finite meets.  Since
$S_P$ is inflationary and idempotent, its inverse image is also
inflationary and idempotent; hence $j_P$ is a nucleus.

\item \textbf{Compute the nucleus formula.}
Expanding $g\in P\star I$ gives
\[
 g\in P\star I\iff\exists p\in P\,\exists a\in I\;(g\le pa),
\]
which proves \eqref{eq:gauge-nucleus}.

\item \textbf{Characterize factorization through the sublocale.}
A classifying map $f:Z\to\Cplx$ corresponds to an open-valued
ideal $I_f$ and its marked module $M_{I_f}\subseteq\Delta_ZK$.
It factors through the nucleus sublocale precisely when
$f^*j_P(U)=f^*U$ for every open $U$.  Since the $U_g$ generate,
this is equivalent to
\[
 (P\star I_f)(g)=I_f(g)\quad\text{for all }g\in L.
\]

\item \textbf{Identify localized modules and their universal family.}
The module representation and the localization theorem identify
this equation with $B_PM_{I_f}=M_{I_f}$.  Since $1\in B_P$,
this is exactly closure under the $B_P$-action in $\Delta_ZK$.
The two classifications therefore restrict to the bijection
\eqref{eq:localized-moduli-classification}; pullback proves its
naturality and supplies the stated universal module.

\item \textbf{Restrict multiplication and its unit.}
If $I,J$ are saturated, associativity gives
$P\star(I\star J)=(P\star I)\star J=I\star J$, so
multiplication restricts.  Its unit is $S_P(\downarrow[1])=P$.

\item \textbf{Restrict the coherent tensor comparison.}
As $B_P\otimes_BB_P\cong B_P$ by localization, for two
$B_P$-modules the balanced tensors over $B$ and $B_P$ agree.
The universal tensor comparison of \cref{thm:localic-multiplication}
therefore restricts to one balanced over $B_P$.  Its associativity,
unit and symmetry follow from multiplication of marked elements.
\end{enumerate}
\end{proof}


\begin{proposition}[Reflection into overhead-stable resources]
\label{prop:cat-gauge-reflection}
Extension of marked modules and restriction of their coefficient action
give an adjunction with fully faithful right adjoint
\[
 B_P(-):\Emb_B(K)\rightleftarrows\Emb_{B_P}(K):\mathrm{incl}.
\]
Its unit is the marked inclusion $M\incl B_PM$, and its counit is
the identity on a $B_P$-submodule.  Under
\cref{cor:cat-monoidal-resources}, the induced idempotent monad
is $S_P(I)=P\star I$.  This reflection preserves resource products
with their localized unit, and finite intersections.
\end{proposition}
\begin{proof}\leavevmode
\begin{enumerate}[leftmargin=*,align=left]
\item \textbf{Prove the reflection adjunction.}
For an embedded $B_P$-module $N$, closure under $B_P$ gives
\[
 B_PM\subseteq N\quad\Longleftrightarrow\quad M\subseteq N.
\]
This is the adjunction, including uniqueness of its marked maps.

\item \textbf{Verify idempotence and the comparison formulas.}
The multiplication isomorphism $B_P\otimes_BB_P\cong B_P$ gives
idempotence.  The tensor and intersection comparisons are the
resource product and flat localization formulas of \cref{thm:gauge}.

\item \textbf{Identify the commuting classification square.}
In particular the square of functors
\[
\begin{tikzcd}[column sep=large,row sep=large]
 \Idl L \arrow[r,"{\simeq}"] \arrow[d,"{S_P}"']
   & \Emb_B(K) \arrow[d,"{B_P(-)}"]\\
 (\Idl L)^{S_P} \arrow[r,"{\simeq}"']
   & \Emb_{B_P}(K)
\end{tikzcd}
\]
commutes up to the canonical marked isomorphism.  Here the lower
left object denotes the fixed ideals, with their inherited maps.
\end{enumerate}
\end{proof}
The adjunction identifies the same operation in three forms:
extension to $B_P$-coefficients, reflection into overhead-stable budgets,
and the classifying nucleus $j_P=S_P^{-1}$. Its unit embeds each
budget into its saturation, compatibly with the marking:
\[
\begin{tikzcd}[column sep=large]
 M \arrow[r,hook] \arrow[dr,hook]
   & B_PM \arrow[d,hook]\\
 & K
\end{tikzcd}
\qquad
\begin{tikzcd}[column sep=large,row sep=large]
 B_P\otimes_BM \arrow[r,"{\cong}"] \arrow[d]
   & B_PM \arrow[d,hook]\\
 B_P\otimes_BK \arrow[r,"{\cong}"'] & K .
\end{tikzcd}
\]


\begin{corollary}[Functorial overhead geometries]
\label{cor:nested-gauges}
For multiplicative ideals $P\subseteq Q$ containing $[1]$,
there is a canonical inclusion
$\Cplx^{(Q)}\incl\Cplx^{(P)}\incl\Cplx$, compatible with
multiplication and the localized universal marked modules.
For numerical cofinite input, the explicitly admitted bitlength
ideal $P_\lambda$ of \cref{prop:admitted-soft-gauge} constructs
$\Cplx^{(P_\lambda)}$, which classifies budgets closed under
polylogarithmic overhead.
\end{corollary}
\begin{proof}\leavevmode
\begin{enumerate}[leftmargin=*,align=left]
\item \textbf{Derive the inclusion of fixed-resource locales.}
The equations $Q\star P=Q$ and $Q\star Q=Q$ yield
$S_QS_P=S_Q$.  Every $Q$-fixed ideal is therefore $P$-fixed,
which gives the sublocale inclusion by the defining nuclei.

\item \textbf{Verify coefficient and tensor compatibility.}
Compatibility of coefficient actions and tensors follows from
$B_P\subseteq B_Q\subseteq K$.

\item \textbf{Specialize to the polylogarithmic gauge.}
For $P_\lambda$ use the
explicitly constructed admitted multiplicative ideal and
$B_{P_\lambda}=B[\lambda]$ of \cref{sec:transport}.
\end{enumerate}
\end{proof}

\begin{theorem}[The localized growth order]
\label{thm:localized-growth-order}
For the multiplicative resource ideal $P$, define a preorder on $L$ by
\begin{equation}\label{eq:soft-growth-preorder}
 a\preceq_P b\quad\Longleftrightarrow\quad
 \exists p\in P\;(a\le pb).
\end{equation}
Let $L/P$ be its separated quotient and
$q_P:L\to L/P$ the quotient map.  Finite joins descend to $L/P$,
and there are natural isomorphisms
\begin{equation}\label{eq:soft-locale-identification}
 \Idl(L/P)\cong\{I\in\Idl(L):P\star I=I\},\qquad
 \OO(\Cplx^{(P)})\cong\Up(L/P).
\end{equation}
In particular the localized principal opens
$U^{(P)}_{q_P(a)}$ are supercompact and form a
finite-meet-stable basis of $\Cplx^{(P)}$.
\end{theorem}
\begin{proof}\leavevmode
\begin{enumerate}[leftmargin=*,align=left]
\item \textbf{Verify the localized preorder.}
Reflexivity uses $[1]\in P$.  If $a\le pb$ and $b\le qc$,
then $a\le pqc$ with $pq\in P$, proving transitivity.

\item \textbf{Descend finite joins to the quotient.}
To descend joins, replace $p,q$ by $p\vee[1],q\vee[1]\in P$:
their product dominates both, so
$a\preceq_P b$ and $c\preceq_P d$ imply
$a\vee c\preceq_P b\vee d$.  The quotient is therefore a
finite-join semilattice with bottom $q_P(0)$.

\item \textbf{Identify saturated ideals.}
A lower ideal $I$ of $L$ is saturated precisely when it is
lower also for $\preceq_P$: if $b\in I$ and $a\le pb$,
saturation puts $pb$ and hence $a$ in $I$; the converse
follows by setting $a=pb$.  Since $q_P$ preserves finite
joins, inverse image along $q_P$ identifies ideals of $L/P$
with exactly these saturated ideals.

\item \textbf{Compute the nucleus on upper predicates.}
For an upper predicate $U$ on $L$, evaluate its nucleus image
at $a$ using the principal ideal $\downarrow a$:
\[
 (j_PU)(a)\iff\exists p\in P\;U(pa).
\]

\item \textbf{Identify fixed opens and prove coherence.}
The equation $j_PU=U$ is equivalent to $U$ being upward
closed also under $\preceq_P$.  Indeed, $pa\preceq_Pa$
and hence upward closure for $\preceq_P$ makes
$U(pa)\Rightarrow U(a)$; conversely, if $a\le pb$ and
$U(a)$, original upperness gives $U(pb)$, and the fixed-point
equation gives $U(b)$.  Such predicates are precisely the
pullbacks of upper predicates on the separated quotient.
This proves the second isomorphism in
\eqref{eq:soft-locale-identification}.  Applying
\cref{thm:classifying-frame} to the finite-join semilattice
$L/P$ proves the principal-open and coherence assertions.
\end{enumerate}
\end{proof}

\subsection{The multiplicative quantale and its support spectrum}
\label{sec:resource-quantale}
The resource product also defines a spectrum. Calculating it will
distinguish the support information carried by multiplication from the
growth information carried by $\Cplx$.
Write $\mathfrak Q=\Idl(L)$, ordered by inclusion, with resource product
$\star$ and unit $\downarrow[1]$. This product is distinct from the
finite meet of the Scott frame $\OO(\Cplx)$.
Let $\operatorname{Idem}(B)=\{e\in B:e^2=e\}$, with Boolean operations
$e\wedge f=ef$, $e\vee f=e+f-ef$, and $\neg e=1-e$.

\begin{theorem}[Quantale and constructive multiplicative spectrum]
\label{thm:resource-spectrum}
The complete lattice $\mathfrak Q$ is a commutative unital
quantale: multiplication preserves arbitrary joins in each variable.
Its unital multiplicative frame-valued characters, that is, maps
$\chi:\mathfrak Q\to H$ to an internal frame $H$ satisfying
\[
 \chi\bigl(\bigvee_iI_i\bigr)=\bigvee_i\chi(I_i),\quad
 \chi(I\star J)=\chi(I)\wedge\chi(J),\quad
 \chi(\downarrow[1])=\top,
\]
are classified by the Stone locale of $\operatorname{Idem}(B)$:
\begin{equation}\label{eq:support-spectrum-frame}
 \OO(\operatorname{Sp}_{\otimes}(\mathfrak Q))
 \cong\Idl(\operatorname{Idem}(B)).
\end{equation}
More precisely every such character satisfies
\begin{equation}\label{eq:spectrum-support-only}
 \chi(\downarrow[g])=\chi(\downarrow[e_g]),
 \qquad
 \chi(I)=\bigvee_{[g]\in I}\chi(\downarrow[e_g]).
\end{equation}
Thus the unital multiplicative spectrum sees support idempotents,
whereas the original classifying locale retains growth comparisons.
\end{theorem}
\begin{proof}\leavevmode
\begin{enumerate}[leftmargin=*,align=left]
\item \textbf{Verify the quantale operations.}
The join of resource ideals is generated by finite joins of their
members.  Multiplying such finite joins distributes over their
products in $L$, giving
\[
 I\star\bigvee_jJ_j=\bigvee_j(I\star J_j).
\]
The zero ideal is absorbing and $\downarrow[1]$ is a unit;
associativity and commutativity are those of profile multiplication.

\item \textbf{Extract the principal character equations.}
A join-preserving $\chi$ is determined by its values
$d_g=\chi(\downarrow[g])$, because every ideal is the
join of its principal ideals.  Finite joins and resource products
imply
\[
 d_0=\bot,\qquad d_1=\top,\qquad
 d_{a\vee b}=d_a\vee d_b,\qquad d_{ab}=d_a\wedge d_b.
\]

\item \textbf{Reduce characters to support idempotents.}
For a nonnegative germ $g$ let $g^\dagger$ be its rational
reciprocal on its nonzero support and zero elsewhere.  It is a
germ in $K^+$ with $gg^\dagger=e_g$ and $ge_g=g$.
The multiplicative equations give
\[
 d_g=d_g\wedge d_{e_g}\le d_{e_g},\qquad
 d_{e_g}=d_g\wedge d_{g^\dagger}\le d_g.
\]
Hence $d_g=d_{e_g}$.  These idempotents are invariant under
growth equivalence, and $e_{g+h}=e_g\vee e_h$ and
$e_{gh}=e_ge_h$.  The $d_e$ therefore form a Boolean-algebra
homomorphism $\operatorname{Idem}(B)\to H$ (the images of
$e$ and $1-e$ are complements).

\item \textbf{Recover a character from Boolean support data.}
Conversely any such map
$\varphi$ defines a quantale character by
$\chi(I)=\bigvee_{[g]\in I}\varphi(e_g)$: finite joins handle
ideal generation, and frame distributivity handles products.
These constructions are inverse.

\item \textbf{Classify the support spectrum.}
Finally, the ideal completion of a Boolean algebra is the free
frame on its bounded-distributive-lattice operations: a map
$\varphi:E\to H$ preserving finite meets, joins and bounds
extends uniquely to
$\widehat\varphi(J)=\bigvee_{e\in J}\varphi(e)$.
Its inverse is restriction to principal ideals.  Taking
$E=\operatorname{Idem}(B)$ proves the frame classification
\eqref{eq:support-spectrum-frame} by its universal property.  This is the point-free
support/Pierce spectrum; the general quantalic spectrum
perspective is discussed in \citep{ManuellSpectrum}.
\end{enumerate}
\end{proof}

For cofinite numerical input, the parity indicator yields a
nontrivial complemented open in the support spectrum.  The support spectrum records these idempotent decompositions; $\Cplx$ records growth budgets.

\subsection{Filtered-input substitution and coherence}
\label{sec:input-substitution}
For a filtered map $u:(Y,\Ev_Y)\to(X,\Ev_X)$, substitution
gives $B_X\to B_Y$ and $K_X\to K_Y$.
On a principal resource,
\[
 B_Y\otimes_{B_X}B_Xg
 \cong B_Ye_{g\circ u}\cong B_Y(g\circ u).
\]
Taking directed colimits proves
\begin{equation}\label{eq:input-basechange}
 B_Y\otimes_{B_X}M_I\cong M_{\downarrow u^*I}.
\end{equation}
The map sends $b\otimes m$ to $b(m\circ u)$.

Both composite comparisons for successive substitutions or
localizations have this same formula on elementary tensors.
Their equality proves unit and associativity coherence.
Raw profiles still reindex along every input map, as in
\eqref{eq:raw-realization}.


\begin{proposition}[Natural base change of classifying resources]
\label{prop:locale-input-change}
For a filtered map
$u:(Y,\Ev_Y)\to(X,\Ev_X)$, write
$\varphi_u:L_X\to L_Y$ for precomposition of growth profiles.
The frame map $\varphi_u^{-1}:\Up(L_Y)\to\Up(L_X)$ defines
$s_u:\mathfrak C_X\to\mathfrak C_Y$, acting on ideal points by
$I\mapsto\downarrow\varphi_u(I)$.  There is a natural
isomorphism of marked modules over $\mathfrak C_X$,
\begin{equation}\label{eq:locale-input-basechange}
 \Delta_{\mathfrak C_X}B_Y
 \otimes_{\Delta_{\mathfrak C_X}B_X}\mathscr M_X
 \ \cong\ s_u^*\mathscr M_Y,
\end{equation}
where the marking on the left is multiplication after
substitution into $\Delta_{\mathfrak C_X}K_Y$.
These isomorphisms satisfy the identity and composition laws.
\end{proposition}
\begin{proof}\leavevmode
\begin{enumerate}[leftmargin=*,align=left]
\item \textbf{Compute the locale inverse image.}
The inverse-image formula preserves arbitrary unions and finite
intersections.  On the principal opens of $\mathfrak C_Y$ it is
\[
 s_u^*U_h=\bigvee_{\substack{g\in L_X\\h\le\varphi_u(g)}}U_g.
\]

\item \textbf{Construct the comparison on principal resources.}
Work internally in $\Sh_{\E}(\mathfrak C_X)$.  Principal
resources generate $\mathscr M_X$ locally and form a directed
family, by \cref{thm:universal-flat}.  For such a local generator
$g$, support idempotents give the explicit marked comparison
\[
 B_Y\otimes_{B_X}B_Xg
 \cong B_Ye_{g\circ u}\cong B_Y(g\circ u)
 \hookrightarrow K_Y.
\]
The map takes $b\otimes m$ to $b(m\circ u)$.

\item \textbf{Prove injectivity of the total marking.}
To verify injectivity after passage to the directed union, let
$z=\sum_i b_i\otimes m_i$ be a finite tensor expression that
maps to zero in $K_Y$.  Locally the finitely many $m_i$ are
contained in a common principal stage: the principal budget
generated by $\sum_i|m_i|$.  The displayed principal-stage
comparison is an isomorphism onto a submodule of $K_Y$,
so $z=0$ at that stage.  Kernel membership is local, making
the total marking monic.

\item \textbf{Identify its image with the pulled-back module.}
For surjectivity onto $s_u^*\mathscr M_Y$, evaluate the membership
predicate of $h\in K_Y$:
\[
 \begin{aligned}
 \llbracket h\in s_u^*\mathscr M_Y\rrbracket
 &=s_u^*U_{[|h|]}\\
 &=\bigvee_{\substack{g\in L_X\\{}[|h|]\le\varphi_u(g)}}U_g.
 \end{aligned}
\]
Thus membership locally supplies $g\in\mathscr M_X$ for which
$h\in B_Y(g\circ u)$.  Bounded division produces the multiplier
$b\in B_Y$ with $h=b(g\circ u)$, represented by $b\otimes g$.
The image of the monic total marking is therefore precisely
$s_u^*\mathscr M_Y$, proving the marked-module isomorphism.

\item \textbf{Verify composition coherence.}
For successive filtered maps, both comparison composites send
an elementary tensor to the same iterated substituted product.
Equality on the directed principal stages proves composition
coherence, and the identity case is immediate.
\end{enumerate}
\end{proof}


Filtered substitution intertwines resource representation and scalar
extension by the natural comparison
\[
\begin{tikzcd}[column sep=large,row sep=large]
 \Idl(L_X) \arrow[r,"{\simeq}"] \arrow[d,"{\downarrow\varphi_u(-)}"']
   & \Emb_{B_X}(K_X) \arrow[d,"{B_Y\otimes_{B_X}-}"]\\
 \Idl(L_Y) \arrow[r,"{\simeq}"'] & \Emb_{B_Y}(K_Y).
\end{tikzcd}
\]
The right-hand marking is induced by $K_X\to K_Y$.
The theorem proves its required monicity on these resources.


\begin{proposition}[Extension and restriction of input budgets]
\label{prop:input-budget-adjunction}
For the filtered substitution $u:Y\to X$, the growth map
$\varphi_u:L_X\to L_Y$ preserves $0$ and finite joins. There is
a locale map $t_u:\mathfrak C_Y\to\mathfrak C_X$ with
\[
 t_u^*U_a^X=U_{\varphi_u(a)}^Y,\qquad
 t_u(J)=\varphi_u^{-1}(J),\qquad
 s_u(I)\subseteq J\ \Longleftrightarrow\ I\subseteq t_u(J).
\]
In the diagram presentations,
\begin{equation}\label{eq:input-kan}
 s_u^*G=G\circ\varphi_u,\qquad
 (t_u^*F)(b)=\varinjlim_{\varphi_u(a)\le b}F(a).
\end{equation}
The coefficient map $B_X\to B_Y$ equips $t_u$ with a ringed-topos
map. Its module pullback has a canonical monomorphism into
$\mathscr M_Y$; at $J$ its marked image is
$M_{s_ut_u(J)}\subseteq M_J$.
\end{proposition}
\begin{proof}\leavevmode
\begin{enumerate}[leftmargin=*,align=left]
\item \textbf{Construct the reverse locale map and ideal adjunction.}
Precomposition preserves zero and sums of profiles, hence finite
joins of growth classes. The assignment on principal opens
preserves their finite-meet relations, so
\cref{thm:classifying-frame} constructs $t_u$. Its point formula
and the displayed ideal adjunction follow by membership.

\item \textbf{Identify the diagram inverse images.}
Precomposition of diagrams sends the subterminal $U_b$ to
$\varphi_u^{-1}(U_b)$, proving the first formula in
\eqref{eq:input-kan}. For the second, the comma poset
$\{a:\varphi_u(a)\le b\}$ contains $0$ and is directed by
joins. Its colimit preserves finite limits by the finite-witness
argument in \cref{thm:budget-diagrams}. The resulting left Kan
extension preserves colimits and sends $U_a$ to
$U_{\varphi_u(a)}$, so it is exactly $t_u^*$.

\item \textbf{Construct the ringed pullback.}
It sends the constant coefficient ring $B_X$ to the constant
ring $B_X$; the substitution map to $B_Y$ defines the ringed
structure. Module pullback is
$\bar B_Y\otimes_{t_u^*\bar B_X}t_u^*(-)$.

\item \textbf{Compute the resource comparison and its fibres.}
On a principal stage $a$, the marked scalar-extension comparison
of \cref{prop:locale-input-change} identifies its value with
$M_{\varphi_u(a)}$. At budget $b$, these form a directed union
inside $M_b$, proving monicity. At an ideal $J$ this union has
growth ideal $\downarrow\varphi_u(\varphi_u^{-1}J)$.

\item \textbf{Verify coherence and observation transport.}
Identity and composition follow on elementary tensors and
principal opens. The forward map $s_u$ retains the effective
observation transport of \cref{prop:effective-resource-transport};
no assertion that $t_u$ preserves admission is needed or made.
\end{enumerate}
\end{proof}
\subsection{Finite principal-open descent}
To use the resource coefficients on affine charts, we need to glue
modules and their maps across localizations. For a commutative
coefficient algebra $A$, finite lists of generators and equations
construct the category of finitely presented $A$-algebras.
A finite principal cover is displayed by $s_1,\ldots,s_r$ and a
unit-ideal witness $1=\sum_i a_i s_i$. The next result supplies
the gluing calculation used later for differential coefficients.
\begin{proposition}[Module gluing]\label{prop:descent}
For every $A$-module $M$ such a cover gives an equalizer
\[
 M\longrightarrow\prod_i M_{s_i}
 \rightrightarrows\prod_{i,j}M_{s_is_j}.
\]
These equalizers commute with restriction and refinement.
\end{proposition}
\begin{proof}\leavevmode
\begin{enumerate}[leftmargin=*,align=left]
\item \textbf{Construct unit combinations of denominator powers.}
This is the finite standard-open \v{C}ech equalizer
(see \citep[Tag~01X9]{Stacks}). Its denominator calculation
is finite, and hence applies internally. First, for every $d\ge1$,
$1=\sum_i c_i s_i^d$ follows by expanding
$(\sum_i a_i s_i)^{r(d-1)+1}$ and grouping each monomial by
an index whose exponent is at least $d$.

\item \textbf{Prove injectivity of the equalizer map.}
If $m$ maps to zero in
every $M_{s_i}$, choose a common $d$ with $s_i^d m=0$.
The unit combination yields $m=0$, proving injectivity.

\item \textbf{Choose compatible denominator representatives.}
For compatible local elements choose a common denominator exponent
$n$ and representatives $x_i=m_i/s_i^n$ with $m_i\in M$.
Compatibility on every overlap gives a common $q$ such that
\[
 (s_is_j)^q(s_j^n m_i-s_i^n m_j)=0\qquad(i,j).
\]

\item \textbf{Construct the global glued element.}
Put $d=n+q$ and $\widetilde m_i=s_i^q m_i$. Then
$s_j^d\widetilde m_i=s_i^d\widetilde m_j$. Choose
$1=\sum_i c_i s_i^d$ by the unit-combination calculation above and set
$m=\sum_i c_i\widetilde m_i$. For each $j$,
\[
 s_j^d m=\sum_i c_i s_i^d\widetilde m_j=\widetilde m_j,
\]
so $m/1=x_j$ in $M_{s_j}$.

\item \textbf{Verify uniqueness and refinement coherence.}
This constructs the gluing map;
injectivity proves uniqueness, restriction, and refinement
coherence. All common exponents arise from finite maxima.
\end{enumerate}
\end{proof}


The gluing condition is the equalizer diagram
\[
\begin{tikzcd}[column sep=large]
 M \arrow[r]
   & {\prod_iM_{s_i}}
       \arrow[r,shift left=.6ex]
       \arrow[r,shift right=.6ex]
   & {\prod_{i,j}M_{s_is_j}} .
\end{tikzcd}
\]
The parallel arrows restrict to the two factors of each overlap;
refinement compares the equalizers through localization.

We now have a classifier of budget families, its marked coefficient
module, and compatible resource operations and transports. The
remaining computational question is which of these budgets contain
the cost of an implementation of a given problem. Answering it
requires the measured execution construction that follows.


\section{Effective implementations and universal incidence}
\label{sec:measurements}
The budget classifier tells us what a resource constraint is. To ask
whether a problem satisfies it, we must supply implementations and
measure their execution. We first construct returning histories and
their resource valuations, then use finite descriptions and input names
to fix the measured implementation family. Its denotation and cost maps
form the span from which universal feasibility is obtained.

The partiality and cost construction is related to cost-sensitive
synthetic domain semantics \citep{NiuSterlingHarper2024CostSensitiveSDT}.
Here the exact-return argument establishes the required comparison in
the selected effective-domain setting. The subsequent sheaf
construction carries the resulting descriptions and observation maps
over the budget classifier.


\subsection{Finite words in the computational category}
Histories must record finitely many events in their execution order.
We construct their word objects in $\AS$, the computational category
of \cref{sec:foundation}. The classified domains of partial execution
remain subobjects in $\E$.

\begin{lemma}[Internally indexed disjoint sums]\label{lem:word-sums}
Every object of $\E$ with its discrete order has a canonical $\J$-algebra.
For an internally given family $(P_n)_{n:\N}$ of posets, with the disjoint
order on its dependent sum, there is an isomorphism
\[
 \J\Bigl(\sum_{n:\N}P_n\Bigr)\cong\sum_{n:\N}\J P_n
\]
compatible with principal ideals and ideal unions. Thus an internal
family of computational algebras has a componentwise disjoint-sum algebra.
\end{lemma}
\begin{proof}\leavevmode
\begin{enumerate}[leftmargin=*,align=left]
\item \textbf{Identify admitted ideals in a discrete poset.}
An inhabited directed ideal in a discrete poset has exactly one element:
two members with a common upper bound are equal. Unique existence gives
its element and identifies $\J A$ with $A$.

\item \textbf{Identify the unique summand and its admitted ideal.}
In a disjoint sum, members with a common upper bound have the same index.
Inhabitedness and directedness therefore specify a unique $n$. Restrict
the ideal to $P_n$. To prove its admission, extend each monotone test
$P_n\to\Sigma$ by false on the other summands, using decidable equality
on $\N$. Conversely, every monotone positive test on the sum restricts
to a test on the chosen fibre. These constructions identify the pairing
predicates and are inverse.

\item \textbf{Verify monad and algebra compatibility.}
Principal ideals and unions retain the unique
index, proving the monad compatibility and hence the algebra assertion.
The sum is one internally constructed family, not an external countable
coproduct.
\end{enumerate}
\end{proof}

\begin{proposition}[Computational finite words]\label{prop:finite-words}
For $A\in\AS$, put
\[
 A^*=\sum_{n:\N}A^{[n]},\qquad [n]=\{j:\N:j<n\},
\]
ordered componentwise at each fixed length. This is a computational
algebra, and empty word, singleton, and concatenation are algebra morphisms
\[
 []:1\to A^*,\qquad [-]:A\to A^*,\qquad
 {+\!+}:A^*\times A^*\to A^*.
\]
It is the free monoid on $A$ in $\AS$. In particular, for a monoid
$(C,e,\mu)$ in $\AS$ and an algebra morphism $w:A\to C$, its unique
monoid extension $\overline w:A^*\to C$ is constructed by
\[
 \overline w([])=e,\qquad
 \overline w(a::s)=\mu(w(a),\overline w(s)).
\]
\end{proposition}
\begin{proof}\leavevmode
\begin{enumerate}[leftmargin=*,align=left]
\item \textbf{Construct the word algebra.}
Finite powers carry the product algebra of \cref{thm:standalone-domain}, uniformly by recursion
on the length. Apply \cref{lem:word-sums}.

\item \textbf{Prove continuity of concatenation and the unit.}
An admitted ideal of a pair
of words has fixed lengths in both components, so preservation of its
supremum by concatenation reduces to the finite-product algebra formula.
The empty word and singleton satisfy the same equation.

\item \textbf{Construct the continuous fold and prove uniqueness.}
At each length
the displayed fold is a finite composite of algebra morphisms. The
disjoint-sum algebra therefore makes it an algebra morphism on all words.
Induction proves preservation of concatenation, and proves uniqueness
from the empty and singleton clauses. These constructions retain all
free parameters.
\end{enumerate}
\end{proof}

\subsection{Resource-valued partiality}
A returning computation now carries both a value and a resource.
Sequential execution combines resources using a monoid; commutativity
is unnecessary because the execution order is retained.
Let $(C,e,\mu)$ be an internal monoid in $\AS$, with algebra morphisms
$e:1\to C$ and $\mu:C\times C\to C$, and write $cd=\mu(c,d)$.
Define
\[
 \LC A=\LL(C\times A),\qquad
 \LC u=\LL(\id_C\times u),\qquad
 \eta^C_A(a)=\eta_{C\times A}(e,a).
\]
For $u\in\LC A$ and $g:A\to\LC B$, define resource-valued bind by
\begin{equation}\label{eq:cost-bind}
 u\bind_C g
 =u\bind_\Sigma\bigl((c,a)\mapsto
       g(a)\bind_\Sigma((d,b)\mapsto\eta_{C\times B}(cd,b))\bigr).
\end{equation}
Here $\bind_\Sigma$ is the partiality bind of \eqref{eq:standalone-lift}. The strength has type
\[
 X\times\LC A\longrightarrow\LC(X\times A),\qquad
 (x,u)\longmapsto
 u\bind_\Sigma((c,a)\mapsto\eta(c,(x,a))).
\]

\begin{theorem}[Constructed resource partiality]\label{thm:cost-partiality}
These maps make $\LC$ a strong monad on $\AS$. Cost erasure
\[
 \epsilon^C_A=\LL\pi_A:\LC A\to\LL A
\]
is a strong monad morphism. A monoid morphism $h:C\to C'$ induces the
strong monad morphism $\LL(h\times\id_A):\LC A\to L_{\Sigma,C'}A$.
For commutative $C$, the constructed monad is commutative.
\end{theorem}
\begin{proof}\leavevmode
\begin{enumerate}[leftmargin=*,align=left]
\item \textbf{Verify continuity of the resource operations.}
All operations are composites of the lifted partiality operations,
strength, products, and $e,\mu$, so they are algebra morphisms.

\item \textbf{Compute composite domains and costs.}
Write a partial value as $(p,(c,a))$ in context $p:\Sigma$, and the next
partial value as $(q,(d,b))$ in that context. Its composite has extent
$\exists t:p\;q(t)$ and value $(cd,b)$. Dominance classifies this extent.

\item \textbf{Verify associativity, units, and strength.}
Three successive computations have the same composite domain under both
associations; their values differ only by $(cd)r=c(dr)$. The unit equations
are $ec=c=ce$. Strength retains the same unused parameter throughout
these domains and values, so the strength equations of \cref{prop:standalone-recursion} prove its
laws.

\item \textbf{Verify erasure and change-of-cost morphisms.}
Erasure forgets precisely the accumulated monoid value. Applying
$h$ to a product gives $h(cd)=h(c)h(d)$ and $h(e)=e$, proving the two
monad-morphism claims.

\item \textbf{Prove commutativity for independent computations.}
For independent calls the two orders have the same
domain by commutativity of $\LL$ and costs $cd$ and $dc$; commutativity of
$C$ proves the final assertion.
\end{enumerate}
\end{proof}

The notation $L_{\Sigma,C}$ distinguishes resource-valued partiality
from the symmetric-algebra monad $T$ and the infinitesimal-shape modality
$\Shape$. The monoids $(\N,0,+)$, $(\N,0,\max)$, and their product are
constructed using the discrete algebra of \cref{lem:word-sums}.
They accumulate declared event costs, recorded peak occupancy, and joint
time--space values, respectively.

The resource-valued bind describes finite sequential composition.
Recursive execution requires the same account of costs on every
returning branch. We obtain it by attaching resource values to the
finite return objects of the selected-domain recursion construction.

\subsection{Weighted exact returns and iteration}
For an equation morphism $f:X\to\LC(Y+X)$ in $\AS$, the map
\[
 \Psi:[X,\LC(Y+X)]_\Sigma\times[X,\LC Y]_\Sigma
       \longrightarrow[X,\LC Y]_\Sigma,
 \quad\Psi(f,a)(x)=f(x)\bind_C[\eta^C_Y,a]
\]
is an algebra morphism. The second function object is pointed by the
everywhere undefined map. The parameterized least-fixed-point operation
of \cref{prop:standalone-recursion} therefore constructs
\[
 f^{\ddagger_C}=\operatorname{lfp}(\Psi(f,-)):X\to\LC Y.
\]
The operation $f\mapsto f^{\ddagger_C}$ is itself an algebra morphism
between the displayed function objects.

To calculate it, represent $f$ by the partial-map span
\[
 X\xleftarrow{m}D\xrightarrow{o}C\times(Y+X),\qquad D=D_Y+D_X.
\]
Its branches have maps in $\E$
\[
 \begin{array}{lll}
 c_Y:D_Y\incl X,&w_Y:D_Y\to C,&v:D_Y\to Y,\\
 c_X:D_X\incl X,&w_X:D_X\to C,&n:D_X\to X.
 \end{array}
\]
Both source monos are $\Sigma$-classified. Retain exactly the return
objects of \cref{lem:standalone-return}, and attach the indicated valuations:
\begin{align*}
 R_0&=D_Y,&R_{k+1}&=D_X\times_{n,X,\rho_k}R_k,\\
 \rho_0&=c_Y,&\rho_{k+1}&=c_X\pi_1,\\
 v_0&=v,&v_{k+1}&=v_k\pi_2,\\
 a_0&=w_Y,&a_{k+1}&=\mu(w_X\pi_1,a_k\pi_2).
\end{align*}


\begin{theorem}[Resource-valued finite return]\label{thm:weighted-return}
The domain of $f^{\ddagger_C}$ is the $\Sigma$-classified mono
\[
 D_f=\sum_{k:\N}R_k\incl X,
\]
and its value on $R_k$ is $(a_k,v_k)$. Consequently
\begin{equation}\label{eq:erase-weighted-return}
 \epsilon^C_Y f^{\ddagger_C}
 =(\epsilon^C_{Y+X}f)^\ddagger.
\end{equation}
These constructions and equations commute with substitution of parameters.
\end{theorem}
\begin{proof}\leavevmode
\begin{enumerate}[leftmargin=*,align=left]
\item \textbf{Construct the classified finite-return domain.}
\Cref{lem:standalone-return} applies to the erased equation.
Pullback and composition prove that each $\rho_k$ is $\Sigma$-classified.
Disjointness of exit and continuation, followed by induction, proves
disjointness of distinct return lengths. Hence the dependent sum is monic.
Its classifier $\exists k:\N\;\rho_k(x)$ lies in $\Sigma$ by
\cref{prop:standalone-dominance}.

\item \textbf{Identify the finite approximants and their costs.}
Starting $\Psi(f,-)$ at bottom gives $b_0=\bot$ and
$b_{j+1}=\Psi(f,b_j)$. Induction shows that $b_j$ has domain
$\sum_{k<j}R_k$ and value $(a_k,v_k)$ on $R_k$. The return branch
contributes $w_Ye=w_Y$; the continuation branch contributes the current
cost followed by the recursive cost. Once defined, the value remains
unchanged in later approximants.

\item \textbf{Recover the partial map and parameter coherence.}
Their admitted supremum is therefore
the displayed partial map, proving the first claim in $\AS$.
Erasure leaves precisely the original return span, including its domain.
All pullbacks, finite recursions, and classified existential images retain
their parameters, as does the cost recursion.
\end{enumerate}
\end{proof}

Equivalently, define
\[
 \widetilde f:C\times X\to\LL((C\times Y)+(C\times X))
\]
by executing $f(x)$ and sending $(d,\mathsf{inl}\,y)$ to
$\mathsf{inl}(cd,y)$ and $(d,\mathsf{inr}\,x')$ to
$\mathsf{inr}(cd,x')$. The constructed iteration satisfies
\[
 f^{\ddagger_C}=\widetilde f^\ddagger\pair{e!}{\id_X}.
\]
Indeed a return in $R_k$ started with accumulator $c$ ends with $ca_k$,
by induction on $k$. This identifies the two constructions. The return length is the projection
$D_f\to\N$; the predicates $R_k$ need not be Boolean. Unit cost on every
call gives $k+1$ on $R_k$. Unit cost on continuation and zero on return
gives $k$; we specify the convention in each concrete measurement.

\begin{theorem}[Resource-valued Elgot structure]\label{thm:weighted-elgot}
Write $g\diamond_C f$ for Kleisli composition and
$g\oplus_C h=[\LC(\mathsf{inl})g,\LC(\mathsf{inr})h]$.
For $f:X\to\LC(Y+X)$, $g:Y\to\LC Z$, and
$h:X\to\LC((Y+X)+X)$, the constructed operation satisfies
\begin{align}
 f^{\ddagger_C}&=[\eta^C_Y,f^{\ddagger_C}]\diamond_C f,
 \label{eq:weighted-fix}\\
 g\diamond_C f^{\ddagger_C}
 &=((g\oplus_C\eta^C_X)\diamond_C f)^{\ddagger_C},
 \label{eq:weighted-nat}\\
 (h^{\ddagger_C})^{\ddagger_C}
 &=(\LC[\id_{Y+X},\mathsf{inr}]h)^{\ddagger_C}.
 \label{eq:weighted-codiag}
\end{align}
For $u:Z\to X$ and $k:Z\to\LC(Y+Z)$,
\[
 fu=\LC(\id_Y+u)k\quad\Longrightarrow\quad f^{\ddagger_C}u=k^{\ddagger_C}.
\]
Thus $\LC$ is a complete Elgot monad. Monoid morphisms commute with
this iteration. Coproduct feedback is the algebra morphism
\begin{equation}\label{eq:weighted-trace}
 \operatorname{Tr}^{U}_{+,C}(f)
 =[\eta^C_Y,(f\iota_U)^{\ddagger_C}]\diamond_C f\iota_X
 \quad(f:X+U\to\LC(Y+U)).
\end{equation}
It satisfies the symmetric coproduct-trace equations.
\end{theorem}
\begin{proof}\leavevmode
\begin{enumerate}[leftmargin=*,align=left]
\item \textbf{Verify the fixed-point equation.}
Exposing the first transition proves the fixed-point equation.

\item \textbf{Verify naturality by matching histories.}
For
naturality, both sides execute the same returning $f$-history and then
$g$. Their partial domains agree, including failure of $g$, and their
ordered products append the same $g$-cost.

\item \textbf{Verify codiagonal by cutting and flattening histories.}
For codiagonal, tag each
continuation by its inner or outer occurrence of $X$, and cut a finite
history at the outer tags. This gives exactly the successive inner
returns of the left side. Flattening the blocks is inverse. Each original
$h$-call occurs once in its original order, so associativity and the unit
law identify the costs.

\item \textbf{Transport return domains and costs.}
For uniformity, the commuting equation transports transitions along $u$
with identical tags, costs, and outputs. Induction identifies each exact
return domain with the corresponding pullback and identifies its cost.
The classified union proves the equation.

\item \textbf{Verify change-of-cost compatibility and continuity.}
A monoid morphism preserves
the displayed cost recurrence, which proves its iteration compatibility.
The operation is an algebra morphism by the construction using $\Psi$.

\item \textbf{Verify the trace equations.}
Formula~\eqref{eq:weighted-trace} first enters the equation and then
iterates the interface branch. For each trace equation, matching the
finite-return components retains every primitive call in order: naturality
adjoins the same entry and exit calls, dinaturality moves the interface
cut, and vanishing cuts or flattens the same blocks. The cost products
are consequently equal. Superposing leaves the unused component
unchanged. The symmetry in yanking is a pure map followed by $\eta^C$,
so its cost is the monoid unit. These arguments prove all trace equations,
including equality of classified domains, in every parameter context.
\end{enumerate}
\end{proof}


\paragraph{Functorial change of resource monoid.}
Changing a resource valuation should commute with recursive execution.
The monoid-morphism clause of the theorem gives exactly this
compatibility: the construction is a functor to strong complete Elgot
monads and iteration-preserving strong monad morphisms,
\[
 C\longmapsto\LC,\qquad
 h:C\to C'\longmapsto\alpha_h=L_\Sigma(h\times-).
\]
Functoriality follows from functoriality of $L_\Sigma$; preservation
of iteration is the monoid-morphism assertion just proved.
In particular the following square commutes:
\begin{equation}\label{eq:cat-elgot-naturality}
\begin{tikzcd}[column sep=large,row sep=large]
 {[X,L_{\Sigma,C}(Y+X)]_\Sigma}
   \arrow[r,"{\ddagger_C}"] \arrow[d,"{(\alpha_h)_*}"']
   & {[X,L_{\Sigma,C}Y]_\Sigma} \arrow[d,"{(\alpha_h)_*}"]\\
 {[X,L_{\Sigma,C'}(Y+X)]_\Sigma}
   \arrow[r,"{\ddagger_{C'}}"']
   & {[X,L_{\Sigma,C'}Y]_\Sigma}.
\end{tikzcd}
\end{equation}
Here $(-)_*$ means postcomposition, and all four objects are the
constructed computational function objects.  Erasure is the case of
the unique monoid morphism $C\to1$.


\subsection{Histories and their resource valuations}
We can measure several resources from the same execution by recording
its events before applying a valuation. For $A\in\AS$ and
$f:X\to\LL(A\times(Y+X))$, replace each event by its singleton
word and execute in the free monoid:
\[
 f_* =\LL([-]\times\id_{Y+X})f,\qquad
 H_f=f_*^{\ddagger_{A^*}}:X\to\LL(A^*\times Y).
\]
This is a constructed returning-history map with the original domain.


The event weight $w:A\to C$ extends through the free-monoid unit:
\begin{equation}\label{eq:cat-history-universality}
\begin{tikzcd}[column sep=large,row sep=large]
 A \arrow[r,"{{[-]}}"] \arrow[dr,"{w}"']
   & A^* \arrow[d,dashed,"{\overline w}"]\\
 & C .
\end{tikzcd}
\end{equation}
The dashed arrow is the unique monoid morphism extending the event
weight. It induces
$\alpha_{\overline w}:L_{\Sigma,A^*}\to L_{\Sigma,C}$.
Thus \cref{eq:cat-elgot-naturality} says that folding the returning
history gives the same result as accumulating its weights during
execution. The following theorem records this comparison.


\begin{theorem}[Finite-history valuation]\label{thm:history-valuation}
For an algebra morphism $w:A\to C$ and its monoid extension
$\overline w:A^*\to C$,
\begin{equation}\label{eq:history-fold}
 \LL(\overline w\times\id_Y)H_f
 =\bigl(\LL(w\times\id_{Y+X})f\bigr)^{\ddagger_C}.
\end{equation}
\end{theorem}
\begin{proof}\leavevmode
\begin{enumerate}[leftmargin=*,align=left]
\item \textbf{Identify the word and cost recurrences.}
The word on $R_0$ is the singleton return label. On $R_{k+1}$ it is the
continuation label followed by the word on $R_k$. The fold recurrence
is precisely the recurrence for $a_k$ in \cref{thm:weighted-return}.

\item \textbf{Compare values on every return summand.}
Induction identifies the values on every summand of the common domain.
\end{enumerate}
\end{proof}


State events are recorded by a typed history map. Given $b:X\to\LL(Y+X)$, strength constructs
\[
 X\to\LL((X\times(Y+X))\times(Y+X)),\qquad
 x\mapsto b(x)\bind_\Sigma(z\mapsto\eta((x,z),z)).
\]
Its words record actual state events. A finite word may contain domain
values or supplied names, so validity of an arbitrary proposed history
is not thereby decidable. For a total algebra morphism $\nu:A^*\to C$,
$\LL(\nu\times\id)H_f$ is a measured execution even if $\nu$ is not a
monoid morphism. Only a monoid fold inherits the sequential product law
of \eqref{eq:history-fold}. In particular, peak occupancy uses the maximum
of recorded occupancy values, including the initial and final states
when relevant; allocation increments alone do not give that measurement.

\begin{proposition}[Universal classification of event weights]
\label{prop:universal-history-weights}
For an event algebra $A\in\AS$ and an internal monoid
$(C,e,\mu)$ in $\AS$, the computational exponential
$\mathcal W_C=[A,C]_\Sigma$ represents all continuous event
weights $w:A\to C$.  There is a canonical, parameter-uniform
valuation map
\[
 \operatorname{fold}_C:\mathcal W_C\times A^*\to C
\]
whose value on $(w,[a_1,\ldots,a_n])$ is the ordered finite
product $w(a_1)\cdots w(a_n)$, with empty product $e$.
This represents monoid morphisms from $A^*$ to $C$ in every
parameter context; its pullback at $w$ is the history valuation
of \cref{thm:history-valuation}.
\end{proposition}
\begin{proof}\leavevmode
\begin{enumerate}[leftmargin=*,align=left]
\item \textbf{Construct the parametrized fold.}
Cartesian closure of $\AS$ constructs $\mathcal W_C$ and its
algebra-morphism evaluation
$\ev:\mathcal W_C\times A\to C$.  The finite-word object is
an NNO-indexed disjoint sum of finite powers.  Iterate
$\mu$ and $\ev$ on each finite-power summand, retaining the
parameter $w$.  The formulas give the displayed map from the
disjoint sum and preserve admitted suprema componentwise by
\cref{lem:word-sums,prop:finite-words}.

\item \textbf{Verify the monoid equations.}
Induction on word
length proves
\[
 \operatorname{fold}_C(w,[]) =e,\qquad
 \operatorname{fold}_C(w,s{+\!+}t)
 =\operatorname{fold}_C(w,s)\operatorname{fold}_C(w,t).
\]

\item \textbf{Prove the universal classification of weights.}
A monoid map from $A^*$ is determined by its value on
singleton words and hence by an arrow $A\to C$.  Restriction
and the displayed fold are inverse by the same induction,
uniformly in free parameters.  Evaluating at a fixed $w$
recovers its unique monoid extension.  Thus a measured
history valuation is a point of this classification object.  A chosen
resource interpretation still selects which point is measured.
\end{enumerate}
\end{proof}

The execution operations are now available. To form a complexity
judgment, we must specify which descriptions and primitive operations
are being measured. Fixing that presentation also fixes what is charged
by an event valuation and what is supplied by the environment.

\subsection{Implementation descriptions}
\label{sec:retained-execution}
\begin{construction}[Fixed computational environments]
\label{constr:fixed-environment}
The finite description construction of \cref{constr:standalone-syntax} has typed contexts
of supplied name and operation parameters. For such a context, form
its environment object $\Theta$ by the indicated finite products and
dependent sums of parameter carriers. Varying a finite context uses
the corresponding internally constructed sum of these objects. A
description in an environment refers to its supplied parameters by
their typed slots. Closed constructor data and implemented programs
are still constructed by their finite syntax. Every unrestricted opaque
literal is recorded in the environment; it is not an extra choice made
while taking an implementation fibre.

Pairing a well-typed description with its environment constructs
$E_{X,Y}\to\Theta$. Interpretation is a morphism over $\Theta$, obtained
by the finite recursion of \cref{lem:standalone-syntax}. The primitive signature and its selected execution descriptions
are retained in this context; by
\cref{prop:universal-history-weights}, selecting a history
valuation is precisely selecting an event-weight point of a
constructed computational exponential. We work internally at a generic
$\theta:\Theta$, suppressing $\theta$ in the formulas. Thus $E_{X,Y}$
below means $E_{X,Y,\theta}$, and an implementation quantifier fixes
$\theta$. Products, function objects, images, and quantifiers in this
construction are those of $\E/\Theta$.

An implemented substitution of environments precomposes these
interpretations. A substitution which itself performs charged work is
represented by its execution description and contributes its history;
ordinary reindexing alone supplies no assertion that this work is free.
The numerical instance has the fixed register signature and no oracle
parameters. The relative instance retains its declared oracle in
$\Theta$. Program-generated function values are obtained by evaluating
their generating descriptions, with that evaluation recorded.
\end{construction}

Use the finite typed descriptions with supplied literal parameters
constructed in \cref{constr:standalone-syntax,lem:standalone-syntax}.
Fix a typed fragment of its implemented signature. Its primitive
operations have types $b_j:U_j\to\LL V_j$, and the fragment includes
the products, sums, composition, and execution operations it uses.
The well-typed subobject of $\mathsf{Tree}_{\mathcal S}$ carries
these descriptions, their interpretation, and their name presentation.
Literal carriers remain supplied parameters; a numerical enumeration
of them is not part of the construction.

For the finitely many primitive operation tags form the computational
event object $A=\sum_j(U_j\times V_j)$. Define the logged primitive by
\[
 b_j^H(u)=b_j(u)\bind_\Sigma
       \bigl(v\mapsto\eta([\iota_j(u,v)],v)\bigr)
       :U_j\to\LL(A^*\times V_j).
\]
Pure structural maps are followed by $\eta^{A^*}$, so have empty word.
Interpret composition by $\diamond_{A^*}$, coproduct branching by
copairing, and recursive execution by $\ddagger_{A^*}$ or
\eqref{eq:weighted-trace}. A specified sequential pair of computations
$f:U\to L_{\Sigma,A^*}V$, $g:U\to L_{\Sigma,A^*}W$ has interpretation
\[
 u\mapsto f(u)\bind_{A^*}
       \bigl(v\mapsto g(u)\bind_{A^*}(w\mapsto\eta^{A^*}(v,w))\bigr).
\]
Thus its evaluation order is retained in the word. These clauses are
operations on the same finite descriptions; no quotient by unweighted
semantic equality is performed.

The substitution here is substitution in the description's typed
operation parameters. For a context of arrow parameters
\[
 G=\prod_i[U_i,L_{\Sigma,A^*}V_i]_\Sigma,
\]
interpretation of a description $t:X\to Y$ has type
$\llbracket t\rrbracket_H:G\to[X,L_{\Sigma,A^*}Y]_\Sigma$.
For substitutions $\sigma_i$ from a parameter context $G'$, put
$\sigma_H=\langle\llbracket\sigma_i\rrbracket_H\rangle_i:G'\to G$.
The substitution equation is
\[
 \llbracket t[\sigma]\rrbracket_H
 =\llbracket t\rrbracket_H\sigma_H.
\]
Ordinary data parameters are retained in the same constructor maps;
their evaluation uses the stated sequencing convention.

\begin{proposition}[Interpretation of retained descriptions]
\label{prop:logged-descriptions}
For each input and output pair in the fragment, structural recursion
constructs maps in $\E$
\[
 \llbracket-\rrbracket_H:E_{X,Y}\to[X,\LL(A^*\times Y)]_\Sigma,
 \qquad
 \operatorname{ev}_H:E_{X,Y}\times X\to\LL(A^*\times Y).
\]
They commute with the supplied finite-tree substitutions. Erasing words
recovers the original interpretation of every description, including
its classified partial domain. A fixed algebra morphism $\nu:A^*\to\N$
constructs its measured evaluator
\begin{equation}\label{eq:measured-evaluator}
 \operatorname{ev}_\nu
 =\LL(\nu\times\id_Y)\operatorname{ev}_H:
 E_{X,Y}\times X\to\LL(\N\times Y).
\end{equation}
\end{proposition}
\begin{proof}\leavevmode
\begin{enumerate}[leftmargin=*,align=left]
\item \textbf{Interpret the typed constructors.}
Each constructor clause above is a morphism on the appropriate
computational function objects: composition and pairing use cartesian
closure and strength; iteration and feedback use
\cref{thm:weighted-elgot}. Interpret a well-typed finite tree by its
constructors, retaining its literal parameters. Induction on its typing
derivation gives the displayed function-object type. This is the finite
structural recursion on the tree carrier of \cref{constr:standalone-syntax}.

\item \textbf{Verify substitution by structural induction.}
Substitution at a variable evaluates the substituted tree, substitution
at a literal retains its interpretation, and a compound constructor
applies the same operation to the substituted children. Induction proves
the substitution equation.

\item \textbf{Verify erasure and the evaluation maps.}
Erasure commutes with the primitive clauses,
unit, bind, strength, and iteration by
\cref{thm:cost-partiality,thm:weighted-return}; induction therefore gives
the original domain and value. Evaluation and postcomposition give the
last two maps.
\end{enumerate}
\end{proof}

The signature, valuation, and environment $\theta$ remain fixed when
we quantify over $E_{X,Y,\theta}$. This specifies what each event
measures. An opaque function literal is an oracle for an invocation-cost
measurement; measuring its internal work instead requires its execution
description and the events of that description. The numerical instance
constructs the configuration update in
\cref{constr:numerical-execution}. Projective lifting supplies a map on
names and its domain comparison; the history of an execution supplies
its cost.

Fixing the environment is essential to the complexity interpretation.
If arbitrary functions could be varied as literals while application
had unit cost, the description $x\mapsto f(x)$ would give every such
denotation a one-call bound. Here an opaque oracle $f$ is fixed by
$\theta$. To measure its internal execution we substitute its supplied
description and charge the resulting events. The language and the
environment together determine the implementations under comparison.

The word fold is stable under the cut and concatenation operations in
the weighted trace proof. Counting outer loop iterations alone is not:
regrouping two primitive calls into one body changes that count. Nor
does equality after the unweighted denotation quotient imply equal
cost. Resource-preserving transformations preserve the recorded word,
or carry a proved comparison of its selected valuation.

\subsection{Measured name presentations}
\label{sec:measured-names}
A represented input may have several names, and execution on those
names may have different costs. We therefore carry the name maps
through the measurement. In the generic fixed environment, retain
the input and output presentations
\[
 q_X:C_X\onto X_0=\Pt D_X,\qquad
 q_Y:C_Y\onto Y_0=\Pt D_Y.
\]
Use the projective presentations and their implemented structural maps
constructed in \cref{prop:standalone-name-lifting}, with partitioned refinements after pullback. Each
primitive retains its actual partial name map and the equality of its
domain with the pullback of the semantic domain. Finite products and
sums of these named sorts provide the typed execution fragment. The
discrete admitted algebra on a name carrier allows the finite-word and
weighted-execution constructions to be applied to these actual maps.
Let $\widehat A$ be their event object: an invocation event records the
primitive tag and its input and output names; an unfolded primitive
instead records its displayed internal execution events.

\begin{theorem}[Measured name-presentation comparison]
\label{thm:measured-name-presentation}
Finite recursion on the retained implementation descriptions constructs
\begin{equation}\label{eq:named-history-evaluator}
 \widehat{\operatorname{ev}}_H:
 E\times C_X\longrightarrow\LL(\widehat A^*\times C_Y).
\end{equation}
Write $\operatorname{ev}^{\mathrm{sem}}:E\times X_0\pmap Y_0$ for
the interpretation of \cref{lem:standalone-syntax} on effective points. Then, as
partial maps,
\begin{equation}\label{eq:measured-name-adequacy}
 q_Y\,\operatorname{erase}\widehat{\operatorname{ev}}_H
 =\operatorname{ev}^{\mathrm{sem}}(\id_E\times q_X),\qquad
 \widehat D=(\id_E\times q_X)^*D_{\mathrm{sem}}.
\end{equation}
For the fixed implemented total valuation
$\widehat\nu:\widehat A^*\to\N$, let $P\incl E$ classify total
named executions. The returned values and measured histories construct
\begin{equation}\label{eq:named-measured-image}
 p_0:P\to Y_0^{X_0},\qquad
 \widehat\Phi_0:P\to\N^{C_X},\qquad
 P\xrightarrow{p\,\onto}V\incl Y_0^{X_0}.
\end{equation}
The cost descends to $\N^{X_0}$ exactly on the subobject of $P$ where
\begin{equation}\label{eq:cost-name-descent}
 q_Xc=q_Xc'\quad\Longrightarrow\quad
 \widehat\Phi_0(e)(c)=\widehat\Phi_0(e)(c')
 \qquad(c,c':C_X).
\end{equation}
All constructions commute with substitution of the retained environment
and finite-description parameters.
\end{theorem}
\begin{proof}\leavevmode
\begin{enumerate}[leftmargin=*,align=left]
\item \textbf{Interpret the measured name primitives.}
Log the actual name primitives and apply the clauses of
\cref{prop:logged-descriptions}. Their selected events determine the
measurement; lifting a semantic arrow alone has assigned no runtime.

\item \textbf{Verify adequacy for finite constructors.}
Erasure of a primitive is its retained name implementation, whose
domain-and-value square is that of \cref{prop:standalone-name-lifting}. Composition pulls back
the second domain along the first returned name; the two primitive
adequacy squares identify this pullback with the semantic composite
domain. Its returned value satisfies the composite square. Coproduct
branching and the stated sequential pairing give the same conclusion
by their finite case and bind clauses.

\item \textbf{Verify adequacy on exact return domains.}
For iteration, induction on the number of transitions identifies each
exact finite-return domain with the pullback of the corresponding
semantic return domain, and identifies the returned values.

\item \textbf{Pass to total return domains and substitution.}
The union
of these domains is the classified return domain on both sides by
\cref{thm:weighted-return}. Pullback and the natural-number existential of
\cref{prop:standalone-dominance} identify those unions. This proves
\eqref{eq:measured-name-adequacy}, including domains, for every finite
description. The same constructor induction proves substitution.

\item \textbf{Construct the totality locus and descend values.}
Write $\widehat d:E\times C_X\to\Sigma$ for the return classifier.
Define $P(e)\iff\forall c:C_X\;\widehat d(e,c)$. Since $q_X$ is a
regular epimorphism, the domain equation identifies this with
$\forall x:X_0\;d_{\mathrm{sem}}(e,x)$. On $P\times C_X$ the evaluator
is total. Its returned name, followed by $q_Y$, is invariant on the
kernel pair of $\id_P\times q_X$, by value adequacy. Effective quotient
descent therefore gives the unique map $P\times X_0\to Y_0$.

\item \textbf{Construct the implementation image and measured cost.}
Transposition and regular image factorization give $p_0,p,V$.
Folding its retained history gives the total map
$P\times C_X\to\N$ and hence $\widehat\Phi_0$.

\item \textbf{Characterize cost descent and parameter coherence.}
Finally form the subobject of $P$ defined by
\eqref{eq:cost-name-descent}. On that subobject the measured map
coequalizes the same kernel pair and thus descends uniquely to $X_0$.
Conversely every descended cost coequalizes that pair. Pullback stability
of these quotients, dependent-product Beck--Chevalley, and the finite
interpretation equations prove parameter coherence.
\end{enumerate}
\end{proof}


The return denotation descends by its extensionality equation.
Cost descends precisely when it is constant on names of the same
input. For a fixed total implementation $e$, the coequalizer
criterion for $q_X$ expresses that condition:
\begin{equation}\label{eq:cat-cost-descent}
\begin{tikzcd}[column sep=large,row sep=large]
 C_X\times_{X_0}C_X
   \arrow[r,shift left=.6ex,"{\pi_1}"]
   \arrow[r,shift right=.6ex,"{\pi_2}"']
   & C_X \arrow[r,two heads,"{q_X}"]
       \arrow[d,"{\widehat\Phi_0(e)}"']
   & X_0 \arrow[dl,dashed,"{\Phi_0^{\mathrm{ext}}(e)}"]\\
 & \N &
\end{tikzcd}
\end{equation}
The dashed map exists uniquely when the two cost composites from the
kernel pair agree. The value leg has its own descent equation,
\cref{eq:measured-name-adequacy}. Until the cost equation also holds,
the name carrier remains the domain of the measured profile.


\begin{example}[Value descent without cost descent]
Take $q_X:\N\times\N\onto\N$, $q_X(n,k)=n$. In the retained
signature a description can perform $k$ implemented idle invocations
and return $0$. For their counting valuation its cost is $k$.
Its erased name map $(n,k)\mapsto0$ has true domain and commutes
with $q_X$ and the identity presentation of the output. It therefore
implements the constant semantic primitive; unfolding this particular
implementation records its name-dependent waiting history.
The returned value descends to $\N$, but the cost fails
\eqref{eq:cost-name-descent}. This is a name-level measurement, not
an extensional function of $n$.
\end{example}

\subsection{Total implementations and the measured image}
We now collect the total implementations of the fixed measured
evaluator and separate their denotations from their costs.
First take actual execution input and return carriers $X,Y$.
Write $E=E_{X,Y}$, or use the numerical family of
\cref{constr:numerical-execution}, with evaluator
\eqref{eq:measured-evaluator}. If $d:E\times X\to\Sigma$ classifies
its domain, totality defines the subobject
\[
 P\incl E,\qquad P(e)\iff\forall x:X\;d(e,x).
\]
Totality is an $\Omega$-valued predicate. On a returning execution write
\begin{equation}\label{eq:history-measurement}
 H_e(x)\in A^*,\qquad T_{\nu,e}(x)=\nu(H_e(x)).
\end{equation}
Both expressions have exactly the return domain of the represented
execution.

\begin{proposition}[Measured implementation image]\label{prop:measured-image}
The construction gives canonical arrows
\[
 p_0:P\to Y^X,\qquad\Phi_0:P\to\N^X,
 \qquad P\xrightarrow{p\,\onto}V\incl Y^X,
 \quad V=\im(p_0).
\]
Their uncurried values are the returned value and the selected finite
resource measurement. The cost map has domain $P$.
\end{proposition}
\begin{proof}\leavevmode
\begin{enumerate}[leftmargin=*,align=left]
\item \textbf{Restrict the evaluator to its totality locus.}
Restriction of the evaluator to $P\times X$ has true extent, so its
partial-map span is a unique total map $P\times X\to\N\times Y$.

\item \textbf{Construct the implementation image and cost map.}
Transposing its two projections gives $\Phi_0,p_0$. Regular image
factorization in $\E$ gives $p$ and $V$.
\end{enumerate}
\end{proof}

For represented inputs, use \cref{thm:measured-name-presentation}
instead: the denotation lies in $Y_0^{X_0}$ and the raw measurement in
$\N^{C_X}$. From now on $X$ in a cost profile denotes the actual
execution input carrier, so $X=C_X$ in this case. The notation $p:P\onto V$, $\Phi_0:P\to\N^X$ covers both constructions. All these objects remain over the fixed
environment $\theta$ of \cref{constr:fixed-environment}.

The order of this construction matters: the language, environment,
and event valuation are fixed before the denotation image is formed.
Otherwise arbitrary clocks or unit-cost semantic primitives would
trivialize the measurement. For example, the clock constantly equal
to true with return value $g(x)$ has first hit zero for every total
$g:X\to Y$. The concrete numerical clocks of
\cref{prop:numerical-clock} instead come from the specified transition
system. The general return theorem \cref{thm:weighted-return} does
not require those clocks.

For Landau comparison, rational inclusion and the eventual quotient of \cref{prop:germs}
give the exact cost map
\[
 \Phi:P\xrightarrow{\Phi_0}\N^X\longrightarrow\Qplus^X
      \xrightarrow{q_+}K^+.
\]
Each measurement is fixed before taking the existential over its
implementation fibre. Joint resources use a product valuation on the
same histories and retain the same implementation witness.

The proper-filter context used in $K$ is a context on this actual input
carrier $X$. Bounds specified by an extensional input or by an input-size
map are pulled back to $X$. See \cref{prop:size-indexed-bounds}.


\paragraph{The measured implementation span.}
For a fixed computational environment, a total implementation
determines both a denotation and a measured cost. These assignments
have the same source and form the span in $\E/\Theta$
\begin{equation}\label{eq:cat-measured-span}
\begin{tikzcd}[column sep=large]
 V & P \arrow[l,two heads,"{p}"'] \arrow[r,"{\Phi}"] & K .
\end{tikzcd}
\end{equation}
Its apex retains the total implementations. To impose a budget,
we will pull its marked module back along $\Phi$, then compose with
$p$ to obtain the family of implementations meeting that budget.
Taking the image will give the feasible denotations.

At the raw-profile level the cost endpoint is $\N^X$.
Rational inclusion and the germ quotient change it to $K$ while
preserving the apex. The pair of legs also factors through its
regular image:
\[
\begin{tikzcd}[column sep=large]
 P \arrow[r,two heads] \arrow[rr,bend left=25,"{{(p,\Phi)}}"]
   & \mathscr R_{p,\Phi} \arrow[r,hook] & V\times K .
\end{tikzcd}
\]
The relation records which denotation--cost pairs occur.
The span carries the witnesses of those pairs, including distinct
implementations with the same denotation and cost. The following
slice construction preserves this information until the feasibility
image is taken.

\begin{proposition}[Witness transport and its feasibility image]
\label{prop:cat-witness-transport}
The measured span induces the slice functor
\[
 \Wit_{p,\Phi}=p_!\Phi^*:\E/K\longrightarrow\E/V,\qquad
 \Wit_{p,\Phi}(M\to K)=(P\times_KM\to V).
\]
Here $p_!$ is dependent sum in the slice.  On marked subobjects,
taking the image gives the monotone map
\[
 \Feas_{p,\Phi}=\exists_p\Phi^*:
       \Sub(K)\longrightarrow\Sub(V).
\]
Both constructions commute with substitution of environment parameters.
A commuting regular-epimorphic refinement of the apex induces a
regular-epimorphic map of witness families and preserves their
feasibility images.
\end{proposition}
\begin{proof}\leavevmode
\begin{enumerate}[leftmargin=*,align=left]
\item \textbf{Construct witness transport and its feasibility image.}
The pullback and subsequent projection are
\begin{equation}\label{eq:cat-witness-pullback}
\begin{tikzcd}[column sep=large,row sep=large]
 P\times_KM \arrow[r] \arrow[d]
   & M \arrow[d]\\
 P \arrow[r,"{\Phi}"'] \arrow[d,"{p}"']
   & K\\
 V &
\end{tikzcd}
\end{equation}
and define $p_!\Phi^*$.  Their regular image is
$\exists_p\Phi^*M$ by the definition of existential image.

\item \textbf{Verify parameter substitution.}
Beck--Chevalley for pullbacks and stability of regular images give
parameter substitution.

\item \textbf{Prove invariance under a covering presentation.}
If $r:P'\onto P$ satisfies
$p'=pr$ and $\Phi'=\Phi r$, the square
\[
\begin{tikzcd}[column sep=large,row sep=large]
 P'\times_KM \arrow[r,two heads] \arrow[d]
   & P\times_KM \arrow[d]\\
 P' \arrow[r,two heads,"{r}"'] & P
\end{tikzcd}
\]
is a pullback.  Its upper arrow is therefore regular epic.
The composites to $V$ have the same image.
\end{enumerate}
\end{proof}

For simultaneous measurements $\Phi_T:P\to K_T$ and
$\Phi_S:P\to K_S$, the right leg is their pairing and the
joint witness family is the pullback
\[
\begin{tikzcd}[column sep=large,row sep=large]
 Z_{M_T,M_S} \arrow[r] \arrow[d]
   & M_T\times M_S \arrow[d,hook]\\
 P \arrow[r,"{{(\Phi_T,\Phi_S)}}"']
   & K_T\times K_S .
\end{tikzcd}
\]
The pullback imposes both constraints on the same implementation.
Its image in $V$ is therefore simultaneous feasibility. This fixes
the order of the operations: combine the resource conditions first,
then quantify over their shared witness.


\subsection{Universal implementation incidence and its fibres}
\label{sec:complexity}
\label{sec:universal-incidence}
The measured span can now be applied to the universal budget family.
Write its fixed-environment maps as $p:P\onto V$,
$\Phi:P\to K^+\incl K$, and $\kappa:P\to L$; with the parameter
displayed they are $p_\theta:P_\theta\onto V_\theta$ and
$\Phi_\theta$. For represented inputs, these are the maps of
\cref{thm:measured-name-presentation}, with germs formed on the
name carrier.
In the ringed resource topos define
\begin{equation}\label{eq:universal-incidence}
 \mathscr Z=\Delta_{\Cplx}P
       \times_{\bar K}\mathscr M,
 \qquad \mathscr F=\im(\mathscr Z\to\Delta_{\Cplx}V).
\end{equation}


This applies \cref{prop:cat-witness-transport} in $\mathscr S$.
The cartesian square selects the implementations meeting the
universal budget, and the regular image records their feasible
denotations:
\begin{equation}\label{eq:cat-universal-incidence-square}
\begin{tikzcd}[column sep=large,row sep=large]
 \mathscr Z \arrow[r] \arrow[d,hook]
   & \mathscr M \arrow[d,hook]\\
 \Delta P \arrow[r,"{\Delta\Phi}"'] \arrow[d,"{\Delta p}"']
   & \bar K\\
 \Delta V &
\end{tikzcd}
\qquad
\begin{tikzcd}[column sep=large]
 \mathscr Z \arrow[r,two heads] \arrow[rr,bend right=25]
   & \mathscr F \arrow[r,hook] & \Delta V .
\end{tikzcd}
\end{equation}
 At an ideal point $I$, write
\begin{equation}\label{eq:incidence}
 Z_I=I^*\mathscr Z=P\times_KM_I,\qquad
 \CostClass(I)=I^*\mathscr F=\im(Z_I\to V).
\end{equation}
This defines semantic budget fibres; restricting $I$ to $\J L$
gives the admitted computational instances.

\begin{theorem}[Classified computational upper bounds]
\label{thm:universal-cost}
The incidence and its image are the budget diagrams
\[
 \mathscr Z(a)=\{e:P:\kappa(e)\le a\},\qquad
 \mathscr F(a)=\{v:V:\exists e\in P_v,\ \kappa(e)\le a\}.
\]
The open-valued classifier of $\mathscr F\incl\Delta V$ is
\begin{equation}\label{eq:classified-cost-open}
 \mathcal U_v=\bigvee_{e\in P_v}U_{\kappa(e)}.
\end{equation}
At every ideal point,
\begin{equation}\label{eq:classified-cost-fibre}
 I^*\mathscr F=\CostClass(I)
 =\{v:\exists e\in P_v,\ \kappa(e)\in I\}.
\end{equation}
Regular-epimorphic refinements of implementation covers preserving
denotation and measured cost leave this image unchanged.
\end{theorem}
\begin{proof}\leavevmode
\begin{enumerate}[leftmargin=*,align=left]
\item \textbf{Compute incidence and feasibility at each budget.}
Apply \cref{thm:budget-diagrams} to the pullback in
\eqref{eq:universal-incidence}. Images are computed at each
budget, giving the two diagrams. Membership of $e$ has extent
$U_{\kappa(e)}$; existential quantification along $p$ gives
\eqref{eq:classified-cost-open}.

\item \textbf{Compute the fibres.}
Geometric inverse image preserves
pullbacks and regular images, proving the fibre formula.

\item \textbf{Prove invariance under presentation refinement.}
For a refinement $r:P'\onto P$, $\exists_r\top=\top$;
composition of images proves the invariance.
\end{enumerate}
\end{proof}

The budget fibres satisfy
$\CostClass(I)=\bigvee_{a\in I}\CostClass(\downarrow a)$
by interchanging their two existential quantifiers. Their incidence
has type $\J L\times V\to\Omega$, with the separately specified
positive observations of \cref{sec:resource-observations}.

\subsection{Complexity judgments as subobject adjunctions}
\label{sec:complexity-quantifiers}
An upper bound asks for an implementation whose cost meets the bound.
A lower bound constrains every implementation of the same denotation.
The adjunctions $\exists_p\dashv p^*\dashv\forall_p$ express these
different quantifiers over the fixed implementation family.
Their classifiers are $\Omega$-valued; positive observations retain
the separately specified meaning of \cref{sec:resource-observations}.
For a nonnegative bound $g$, form
\[
 A_g=\{e:P:\Phi(e)\in Bg\},\qquad
 D_g=\{e:P:g\in B\Phi(e)\},
\]
and define
\begin{equation}\label{eq:complexity-adjoints}
 \operatorname{Upper}_g=\exists_pA_g,\qquad
 \operatorname{Lower}_g=\forall_pD_g,\qquad
 \operatorname{Tight}_g=\operatorname{Upper}_g\wedge
                       \operatorname{Lower}_g.
\end{equation}
Their internal formulas are respectively
$\exists e\in P_v\,\Phi(e)\in O(g)$ and
$\forall e\in P_v\,\Phi(e)\in\Omega(g)$.
With the environment displayed, the first is
\[
 \operatorname{Upper}^{\theta}_g(v)
 \iff\exists e\in P_{\theta,v}\;\Phi_\theta(e)\in Bg,
\]
and the lower and tight judgments use that same $P_{\theta,v}$.
The lower-bound constants and eventual extents occur inside
$\forall e$; the variant with uniform constants is a different
dependent-quantifier expression.

\begin{proposition}[Complexity bounds as inclusions of feasibility opens]
\label{prop:bounds-as-opens}
For $g\in K^+$ and $a=[g]$, the upper, lower and tight judgments
have the geometric descriptions
\begin{equation}\label{eq:bounds-as-opens}
 \begin{aligned}
 \operatorname{Upper}_g(v)&\iff U_a\subseteq\mathcal U_v,\\
 \operatorname{Lower}_g(v)&\iff\mathcal U_v\subseteq U_a,\\
 \operatorname{Tight}_g(v)&\iff\mathcal U_v=U_a.
 \end{aligned}
\end{equation}
These equivalences are internal and uniform in the original
denotation, bound and computational-environment parameters.
\end{proposition}
\begin{proof}\leavevmode
\begin{enumerate}[leftmargin=*,align=left]
\item \textbf{Reflect growth order through principal opens.}
The principal opens satisfy
$U_b\subseteq U_a\iff a\le b$: evaluate an inclusion at $b$
for one direction, and use upperness for the other.

\item \textbf{Characterize the upper-bound judgment.}
If $\kappa(e)\le a$, then
$U_a\subseteq U_{\kappa(e)}\subseteq\mathcal U_v$.
Conversely, evaluation of $U_a\subseteq\mathcal U_v$ at $a$
and \eqref{eq:classified-cost-open} give
$\exists e\in P_v\;\kappa(e)\le a$, precisely the upper judgment.
Equivalently, supercompactness of $U_a$ extracts a member of the
canonical implementation cover.

\item \textbf{Characterize the lower and tight judgments.}
For the lower judgment,
\[
 \begin{aligned}
 \mathcal U_v\subseteq U_a
 &\iff\forall e\in P_v\;U_{\kappa(e)}\subseteq U_a\\
 &\iff\forall e\in P_v\;a\le\kappa(e)
 \iff\operatorname{Lower}_g(v).
 \end{aligned}
\]
Combining the two comparisons proves the tight statement.
\end{enumerate}
\end{proof}

The feasibility open therefore determines every principal upper,
lower and tight judgment.  A tight bound is the assertion that this
open is the corresponding principal open.  The implementation cover
will identify principality with an attained optimal growth class in
\cref{thm:optimization-compactness,thm:optimal-sections}.
The lower statement in \eqref{eq:bounds-as-opens} is an inclusion
between constructed opens; its proof retains the universal
implementation quantifier and its individual domination witnesses.

\begin{proposition}[Matching bounds and shared witnesses]
\label{prop:matching}
A tight bound supplies an implementation with cost in $\Theta(g)$.
On the eventually nonzero comparison locus, a lower bound excludes
an implementation in $o(g)$. An existential implementation in
$\Theta(g)$ does not imply the universal lower bound.
For two incidence subobjects $A,D\incl P$, the shared-witness
classification is $\exists_p(A\wedge D)$, which lies below
$(\exists_pA)\wedge(\exists_pD)$; equality is not automatic.
\end{proposition}
\begin{proof}\leavevmode
\begin{enumerate}[leftmargin=*,align=left]
\item \textbf{Combine bounds using a shared implementation.}
Frobenius reciprocity and the counit of $p^*\dashv\forall_p$ give
\[
 (\exists_pA_g)\wedge(\forall_pD_g)
 =\exists_p(A_g\wedge p^*\forall_pD_g)
 \le\exists_p(A_g\wedge D_g).
\]
Mutual domination gives $\Theta(g)$.

\item \textbf{Exclude simultaneous smallness and a nonzero lower bound.}
If a small implementation
also satisfies $g\preceq\Phi(e)$ with constant $a$, smallness
at $2a$ gives eventually $2|g|\le|g|$, contradicting the
eventually nonzero extent and properness of the filter.

\item \textbf{Distinguish an implementation cost from a universal bound.}
A constant-output function has both an immediate-return
implementation of bounded cost and one scanning every input
letter, of cost $\Theta(|x|+1)$. The latter does not make that
bound a universal lower bound.

\item \textbf{Establish the shared-witness inclusion.}
Finally monotonicity of images
gives the shared-witness inclusion; taking the two images
separately permits different implementations.
\end{enumerate}
\end{proof}

All these predicates are defined in $\E$ with their displayed
parameters. Constant-diagram transport preserves their particular
quantifications: a map of constant objects has the same fibre at
every budget and Weil stage, and a constant predicate has the same
universal clause on that fibre. Parameter substitution inside a
topos obeys the dependent-product Beck--Chevalley law, by its
right-adjoint universal property.
For a cover refinement $r:P'\onto P$, both
$\exists_r r^*$ and $\forall_r r^*$ are the identity on
subobjects of $P$, so pulled-back upper and lower judgments
are invariant. This follows from the inhabited fibres of the regular epi.

The open $\mathcal U_v$ is Scott open by its principal cover.
Universal resource assertions need not be: in the cofinite
numerical fibre put $g_k(n)=(n+1)^k$ and
$I_m=\downarrow[g_m]$. The assertion $\forall k\,[g_k]\in I$
holds at $\bigcup_m I_m$ but at no $I_m$, since
$[g_{m+1}]\not\le[g_m]$. Universal judgments therefore retain
their dependent-product type instead of being identified with
feasibility opens.
\subsection{The implementation sheaf and compiler morphisms}
\label{sec:implementation-sheaf}
\begin{proposition}[Substitution of fixed environments]
\label{prop:environment-base-change}
For a parameter substitution $s:\Theta'\to\Theta$, reindexing the
retained execution construction gives the same measured evaluator as
interpreting its substituted descriptions. It carries $P,V,p,\Phi$,
their implementation incidence and feasibility image, and their upper
and lower judgments to the corresponding constructions over $\Theta'$.
The implementation-sheaf, compiler, and optimality constructions below
are likewise performed in this fixed-environment family.
\end{proposition}
\begin{proof}\leavevmode
\begin{enumerate}[leftmargin=*,align=left]
\item \textbf{Transport interpretation through environment substitution.}
The interpretation statement is the finite substitution equation of
\cref{prop:logged-descriptions,thm:measured-name-presentation}.

\item \textbf{Identify totality, denotations, and cost profiles.}
Pullback between parameter slices preserves finite limits and regular
images. The dependent-product Beck--Chevalley equation for the input
projection identifies the totality subobjects; exponentiation in the
slice identifies their transposed denotations and cost profiles.
The quotient and image comparisons then identify $p,V,\Phi$.

\item \textbf{Transport feasibility and lower judgments.}
The same pullback and image equations identify incidence and upper
feasibility. Beck--Chevalley for $p$ identifies the universal lower
judgment.

\item \textbf{Verify parameter coherence of the remaining constructions.}
All budget and coefficient constructions retain these
parameters; the sheaf constructions below have the same pointwise
formulas in each slice. In particular an implementation quantifier is
never replaced by an existential over environments. The displayed
comparisons satisfy identity and composition by their universal
properties.
\end{enumerate}
\end{proof}


In particular the denotation objects and their implementation covers
form a cartesian square under $s:\Theta'\to\Theta$:
\[
\begin{tikzcd}[column sep=large,row sep=large]
 P_{\Theta'} \arrow[r] \arrow[d,two heads,"{p_{\Theta'}}"']
   & P_{\Theta} \arrow[d,two heads,"{p_{\Theta}}"]\\
 V_{\Theta'} \arrow[r] & V_{\Theta}.
\end{tikzcd}
\]
Pullback of the measured cost leg gives the compatible incidence square;
Beck--Chevalley identifies its existential and universal judgments.


Fixing a denotation in the universal incidence leaves a family of
implementations indexed by budgets. We now describe this witness
sheaf explicitly. Its image is the feasible open, while its sections
retain the choices of implementation needed for the later optimality
criterion.

\begin{theorem}[The universal sheaf of feasible implementations]
\label{thm:implementation-sheaf}
For $v:V$ the pullback $\mathscr W_v$ of $\mathscr Z$ along
$\Delta\{v\}\to\Delta V$ is the dependent coproduct
\begin{equation}\label{eq:implementation-sheaf}
 \mathscr W_v\cong
 \coprod_{e\in P_v}U_{\kappa(e)}
 \qquad\text{in }\Sh_{\E}(\Cplx).
\end{equation}
Its image over $\Cplx$ is $\mathcal U_v$, and its
restriction to that open is a regular epimorphism
$\mathscr W_v\onto\mathcal U_v$.
By the sheaf--local-homeomorphism correspondence, this
constructs an internal \'etale locale $E_v\to\Cplx$
whose fibre over a resource ideal $I$ is
$\{e\in P_v:\kappa(e)\in I\}$.
\end{theorem}
\begin{proof}\leavevmode
\begin{enumerate}[leftmargin=*,align=left]
\item \textbf{Identify the implementation sheaf by membership.}
Over an implementation $e:P_v$, the defining pullback
$\mathscr Z=\Delta P\times_{\Delta K}\mathscr M$
restricts to the membership subterminal
$\llbracket\Phi(e)\in\mathscr M\rrbracket
 =U_{\kappa(e)}$.  Dependent disjoint union over $P_v$
gives \eqref{eq:implementation-sheaf}.

\item \textbf{Compute feasibility and the regular epimorphism.}
The image of its
projection is the join of these membership opens, exactly
$\mathcal U_v$ by \cref{thm:universal-cost}.  Factoring
through the image gives the regular epi.

\item \textbf{Construct the associated locale and its fibres.}
The standard
internal sheaf--\'etale-locale equivalence for a localic
sheaf topos \citep{Vickers,MacLaneMoerdijk} constructs $E_v$;
pullback along $I$ computes the displayed fibre.
\end{enumerate}
\end{proof}


Under \cref{thm:budget-diagrams}, this is the internal diagram
\[
 \mathsf{Impl}_v:L\to\E,\qquad
 \mathsf{Impl}_v(a)=\{e\in P_v:\kappa(e)\le a\},
\]
whose transition maps retain the same implementation as the budget
increases. Its ideal fibre is the directed colimit over $a\in I$.

Each implementation is available at its own cost and at every larger
budget. Extending this generating data along the cost map produces
the entire witness diagram. The following universal property says
that a transformation out of the diagram is determined by what it
does to those generators.

\begin{proposition}[The free budget extension of an implementation family]
\label{prop:cat-implementation-kan}
Write $\kappa_v:\operatorname{Disc}(P_v)\to L$ for the cost functor
from the discrete internal category of implementations.  The witness
diagram is the left Kan extension of its terminal diagram:
\[
 \mathsf{Impl}_v\cong\Lan_{\kappa_v}\mathbf1,\qquad
 \mathsf{Impl}_v(a)
   \cong\coprod_{e\in P_v}L(\kappa_v(e),a).
\]
Here $L(b,a)$ is the subterminal object expressing $b\le a$.
For every internal diagram $D:L\to\E$ there is a natural isomorphism
\begin{equation}\label{eq:cat-implementation-kan-property}
 \operatorname{Nat}(\mathsf{Impl}_v,D)
 \cong\prod_{e\in P_v}D(\kappa_v(e)).
\end{equation}
\end{proposition}
\begin{proof}\leavevmode
\begin{enumerate}[leftmargin=*,align=left]
\item \textbf{Identify the budget-indexed implementation family.}
At $a$, the displayed dependent coproduct has exactly the elements
$e\in P_v$ with $\kappa_v(e)\le a$, and its transition maps retain
$e$.

\item \textbf{Extract data at each implementation's cost.}
For a transformation $\alpha$, evaluate $\alpha$ at the element
$e$ of $\mathsf{Impl}_v(\kappa_v(e))$ to obtain
$d_e\in D(\kappa_v(e))$.

\item \textbf{Construct the inverse and prove the universal property.}
Conversely transport $d_e$ along
$D(\kappa_v(e)\le a)$ whenever $e$ is feasible at $a$.
These operations are inverse by naturality.  They are internal
dependent operations, so require no selection from unrelated fibres.
This proves the Kan-extension universal property.
\end{enumerate}
\end{proof}
The Kan-extension formula retains each implementation across all budgets
that admit it; its regular image is the feasible open.

\subsection{A problem probes the locale by an open-valued observable}
\label{app:complexity-feasibility-probes}

The witness sheaf describes implementations at each budget.
Its support answers the corresponding question about the problem:
on which budget families does an implementation exist?
For $v:V$, use the fibre $P_v$ of the measured span
\eqref{eq:cat-measured-span}, with witness diagram
$\Lan_{\kappa_v}\mathbf1$. Its support is
\begin{equation}\label{eq:app-problem-open}
 \mathcal U_v=\bigvee_{e\in P_v}U_{\kappa(e)}
 \quad\in\OO(\Cplx).
\end{equation}
Thus $I\models\mathcal U_v$ if and only if
$\exists e\in P_v\;\Phi(e)\in M_I$.  This is the same cost
incidence as in \cref{thm:universal-cost}, with neither names nor
implementation costs quotiented by denotation.

\begin{proposition}[The functor-of-points feasibility probe]
\label{prop:app-feasibility-probe}
For each represented denotation $v:V$, there is a locale map
$\chi_v:\Cplx\to\mathbb S$ to the Sierpi\'nski locale, uniquely
classified by $\mathcal U_v$.  On every test locale $Z$ it induces
a natural transformation
\begin{equation}\label{eq:app-natural-probe}
 h_{\Cplx}(Z)\longrightarrow\OO(Z),
 \qquad b\longmapsto b^*\mathcal U_v.
\end{equation}
For the budget family $\mathscr M_b\subseteq\Delta_ZK$, the open
$b^*\mathcal U_v$ is the regular-image locus of the pullback
\begin{equation}\label{eq:app-family-incidence}
 \Delta_ZP_v\times_{\Delta_ZK}\mathscr M_b
 \longrightarrow 1_Z,
\end{equation}
where the first map to $\Delta_ZK$ is $\Delta_Z\Phi$.
Consequently the functor of points of the open sublocale
$\mathcal U_v\incl\Cplx$ is the subfunctor
\begin{equation}\label{eq:app-feasible-subfunctor}
 h_{\mathcal U_v}(Z)
 \cong\{b:Z\to\Cplx:b^*\mathcal U_v=\top_Z\}
 \ \incl h_{\Cplx}(Z).
\end{equation}
For an ideal point $I$, the specialization is
\[
 \chi_v(I)=\top\quad\Longleftrightarrow\quad
 v\in\CostClass(I)
 \quad\Longleftrightarrow\quad
 \exists e\in P_v\;\kappa(e)\in I.
\]
\end{proposition}
\begin{proof}\leavevmode
\begin{enumerate}[leftmargin=*,align=left]
\item \textbf{Classify the feasibility open and verify naturality.}
The Sierpi\'nski classification of opens gives the unique $\chi_v$.
Under the natural identification
$\Loc_\E(Z,\mathbb S)\cong\OO(Z)$, composition with $\chi_v$
is exactly \eqref{eq:app-natural-probe}.  For a map $f:Z'\to Z$,
$(bf)^*\mathcal U_v=f^*(b^*\mathcal U_v)$, proving naturality.

\item \textbf{Pull back the witness family and its image.}
The universal incidence is the pullback
$\mathscr Z=\Delta P\times_{\bar K}\mathscr M$.
Pulling its $v$-fibre back along $b$ gives
\eqref{eq:app-family-incidence}.  Geometric inverse image preserves
pullbacks and regular images, so the image open is $b^*\mathcal U_v$.

\item \textbf{Characterize feasible factorization and ideal fibres.}
An open immersion $\mathcal U_v\incl\Cplx$ has the asserted
factorization property for locale maps, proving
\eqref{eq:app-feasible-subfunctor}.  The last equivalences are the
fibre formula in \cref{thm:universal-cost}.
\end{enumerate}
\end{proof}


For $f:Z'\to Z$, the problem observable respects the parameter triangle
\[
\begin{tikzcd}[column sep=large,row sep=large]
 Z' \arrow[r,"{f}"] \arrow[dr,"{bf}"']
   & Z \arrow[d,"{b}"]\\
 & \Cplx \arrow[d,"{\chi_v}"]\\
 & \mathbb S .
\end{tikzcd}
\]
Its open is $(bf)^*\mathcal U_v=f^*(b^*\mathcal U_v)$.
The pullback witness family supplies its local implementation witnesses.


Equivalently, the Yoneda correspondence identifies the probe with
\[
 \operatorname{Nat}(h_{\Cplx},h_{\mathbb S})
 \ \cong\ \Loc_\E(\Cplx,\mathbb S)
 \ \cong\ \OO(\Cplx),
\]
where $h_{\mathbb S}(Z)=\OO(Z)$.  In this precise sense
\emph{problems probe the complexity locale}: the generalized points
are budget families, while a problem gives an open-valued observable
on those families.  By \cref{thm:feasibility-classifier}, the same
observable is the point $c(v)$ of the constructed exponential locale
$\mathfrak F=\mathbb S^{\Cplx}$, with
$c:\operatorname{Disc}(V)\to\mathfrak F$ natural in the denotation
parameter.

The witness sheaf remains distinct from its feasible open.  By
\cref{thm:implementation-sheaf}, it has the presentation
\[
 \mathscr W_v\cong\coprod_{e\in P_v}U_{\kappa(e)},
 \qquad \mathscr W_v\onto\mathcal U_v.
\]
After pullback along $b:Z\to\Cplx$, the equation
$b^*\mathcal U_v=\top_Z$ says that the witness projection is
locally inhabited over $Z$. A section supplies the stronger datum
of compatible implementation witnesses throughout the family.
The distinction between these two statements will govern the
optimality result.

\subsection{A constant problem and redundant programs}
\label{app:complexity-constant-example}

The distinction between a program's cost and a problem's complexity
already appears for a constant function. Work on $X=\N$ with its
cofinite filter. Select an implemented unit-valued idle primitive
$\mathsf a$ that leaves the result unchanged, and count its recorded
invocations with $w(\mathsf a)=1$ and empty-word cost $0$.
This is the history-fold measurement of
\cref{prop:finite-words,prop:universal-history-weights}.

Let $v:\N\to\N$ be the constant-zero function.  The pure
constant-return description $e_0$ has an empty event word on every
input, and therefore
\[
 p(e_0)=v,\qquad H_{e_0}(n)=[],\qquad
 T_{e_0}(n)=0,\qquad\kappa(e_0)=[0].
\]
A second finite description $e_1$ performs $n$ consecutive calls to
$\mathsf a$ and then returns zero.  The iteration is
NNO recursion applied to the implemented primitive; the log equation
and finite induction give
\[
 p(e_1)=v,\qquad H_{e_1}(n)=\mathsf a^n,\qquad
 T_{e_1}(n)=n,\qquad\kappa(e_1)=[n].
\]
Here $\mathsf a^n$ is the finite word of charged idle events.

The first program contributes $U_{[0]}=\top$ to the feasibility
open.  Since every growth ideal contains its bottom element, it
follows that
\begin{equation}\label{eq:app-constant-open}
 \mathcal U_v=\top_{\Cplx}.
\end{equation}
The additional contribution $U_{[n]}\subseteq U_{[0]}$ leaves this
open unchanged. Thus the problem is feasible at the zero-cost budget
although $e_1$ has linear event cost. For the corresponding bounded
and linear machine-update example, see \cref{prop:matching}.

We next compare witness families by translating their implementations.
The Kan-extension property reduces the semantic comparison to the
cost of each translated generator. A finite compiler description
then supplies the implementation of that translation.

\subsection{Implemented natural transformations}

\begin{proposition}[Lax cost squares classify witness transformations]
\label{prop:cat-lax-compiler}
For a monotone budget transformation $r:L_v\to L_w$ there is a
natural isomorphism
\begin{equation}\label{eq:cat-lax-compiler}
 \operatorname{Nat}(\mathsf{Impl}_v,\mathsf{Impl}_w\circ r)
 \cong
 \{\gamma:P_v\to P_w:\kappa_w\gamma\le r\kappa_v\}.
\end{equation}
These correspondences preserve identities and composition.
\end{proposition}
\begin{proof}\leavevmode
\begin{enumerate}[leftmargin=*,align=left]
\item \textbf{Apply the Kan-extension classification.}
Apply \cref{eq:cat-implementation-kan-property} with
$D=\mathsf{Impl}_w\circ r$.  Its right side is the dependent
product of the objects
$\{d\in P_w:\kappa_w(d)\le r(\kappa_v(e))\}$, which is precisely
the object of displayed maps $\gamma$.

\item \textbf{Identify the action and its composition laws.}
The transformation sends
$e$ at every feasible budget to $\gamma(e)$.
The two descriptions therefore have the same identity and composition.
\end{enumerate}
\end{proof}
The compiler's cost bound is the inequality in the lax square
\begin{equation}\label{eq:cat-lax-cost-square}
\begin{tikzcd}[column sep=large,row sep=large]
 P_v \arrow[r,"{\gamma}"] \arrow[d,"{\kappa_v}"']
   \arrow[dr,phantom,"{\scriptstyle\le}" description]
   & P_w \arrow[d,"{\kappa_w}"]\\
 L_v \arrow[r,"{r}"'] & L_w .
\end{tikzcd}
\qquad \kappa_w\gamma\le r\kappa_v.
\end{equation}
Monotonicity transports the bound to every larger budget. The resulting
components send the same implementation to the same compiled code, which
proves naturality. We now retain finite-code realizations of these maps.

\begin{construction}[The category of measured compiler comparisons]
\label{constr:measured-compilers}
Work at the fixed environment of \cref{constr:fixed-environment}.
Source and target signatures may share that environment; a comparison
using an implemented change of environment retains the specified map
and is formed after the corresponding pullback. Its execution cost is
obtained from that change's retained description.
Use the data extension of the implementation signature in
\cref{sec:retained-execution}, retaining its code carriers as sorts
as in \cref{constr:standalone-syntax}. For two code carriers
$E_v,E_w$, let $T_{v,w}$ be their constructed carrier of well-typed
finite transformation descriptions, with interpreted partial map
\[
 \operatorname{ev}^{\mathrm{tr}}_{v,w}:T_{v,w}\times E_v\to\LL E_w.
\]
This interpretation is the finite-tree construction of \cref{lem:standalone-syntax}; its
identity and sequential-composition constructors give actual descriptions
of the identity and composite code maps. In the numerical instance
$T_{v,w}=\N$ with its partial numerical evaluator on program codes;
the identity and composition wrappers are those of \cref{lem:standalone-syntax,prop:standalone-numerical-cover} and
\cref{constr:numerical-execution}.
For the measured fibre inclusions $i_v:P_v\incl E_v$,
$i_w:P_w\incl E_w$, form the
subobject
\[
 \operatorname{CompCode}(P_v,P_w)=
 \{t:T_{v,w}:\forall e:P_v\;\exists d:P_w,\
    \operatorname{ev}^{\mathrm{tr}}_{v,w}(t,i_v e)=\eta_{E_w}(i_wd)\}.
\]
The condition means that the partial domain is true and its returned
value belongs to $P_w$; it is an $\Omega$-valued condition. Restriction
to this true domain and transposition construct
\[
 \gamma:\operatorname{CompCode}(P_v,P_w)\to(P_w)^{P_v}.
\]
Define the raw bound subobject
\[
 R(v,w)=\left\{(t,r):\begin{array}{l}
 t\in\operatorname{CompCode}(P_v,P_w),\quad
 r\in\operatorname{Mon}(L_v,L_w),\\
 \forall e:P_v,\quad\kappa_w(\gamma_t e)\le r(\kappa_v(e))
 \end{array}\right\}.
\]
The hom object of measured comparisons is its regular image
\begin{equation}\label{eq:measured-compilers-object}
 \MComp(v,w)=\im\left(
 R(v,w)\xrightarrow{(\gamma,\pi_2)}
 (P_w)^{P_v}\times\operatorname{Mon}(L_v,L_w)\right).
\end{equation}
Thus a comparison retains an actual code-output map and its resource
transformation, with local evidence of a finite implementation. Its
identity and composition are
\[
 1_v=(\id_{P_v},\id_{L_v}),\qquad
 (\delta,s)\circ(\gamma,r)=(\delta\gamma,sr).
\]
They land in the image because the finite constructors implement identity
and composition. Their bounds follow from
\[
 \kappa_z(\delta\gamma e)\le s(\kappa_w(\gamma e))
                          \le sr(\kappa_v(e)).
\]
Associativity and units are equations of interpreted functions and
monotone maps, proving the internal category laws for every retained
internal family of objects. The codes themselves remain a presentation
of this category; raw composition wrappers need not be equal as syntax.
Nor are two maps identified merely because their outputs have equal
denotations: distinct codes in $P_w$ can have different costs. The
quotient in \eqref{eq:measured-compilers-object} uses equality in
$(P_w)^{P_v}$, not equality after $p_w:P_w\to V_w$.
\end{construction}

\begin{theorem}[Compiler action on the implementation moduli]
\label{thm:compiler-sheaves}
For a measured comparison $(\gamma,r):v\to w$ in
$\MComp$, the implementation map $\gamma:P_v\to P_w$ and monotone
resource transformation $r:L_v\to L_w$ satisfy
\begin{equation}\label{eq:compiler-moduli-bound}
 \kappa_w(\gamma e)\le r(\kappa_v(e))\qquad(e\in P_v).
\end{equation}
The growth transformation induces a locale map
$s_r:\mathfrak C_v\to\mathfrak C_w$, acting on ideals by
$I\mapsto\downarrow r(I)$.  The compiler induces a morphism
of witness sheaves over $s_r$,
\[
 \mathscr W_v\longrightarrow s_r^*\mathscr W_w,
\]
and its regular image gives
\begin{equation}\label{eq:compiler-open}
 \mathcal U_v\le s_r^*\mathcal U_w.
\end{equation}
Identity and composed comparisons in $\MComp$ act by identity and
composition of these morphisms. This action is independent of the
choice of finite code representing the same interpreted comparison
in \cref{constr:measured-compilers}.
\end{theorem}
\begin{proof}\leavevmode
\begin{enumerate}[leftmargin=*,align=left]
\item \textbf{Construct the compiler transformation.}
Under \cref{thm:budget-diagrams}, $s_r^*$ is precomposition with $r$.
The lax comparison therefore gives the required witness transformation
by \cref{prop:cat-lax-compiler}.  Its explicit component at budget $a$
sends $e$ to $\gamma e$, with
\[
 \kappa_w(\gamma e)\le r(\kappa_v(e))\le r(a),
\]
so this is a natural transformation
$\mathsf{Impl}_v\to\mathsf{Impl}_w\circ r$.

\item \textbf{Transport feasibility and verify coherence.}
Taking images proves \eqref{eq:compiler-open}. Both composites
send an implementation to the same composite code and use the
same composed monotone bound, proving identity, composition,
and independence of the finite description of the interpreted map.
\end{enumerate}
\end{proof}


For a generalized budget $b:Z\to\mathfrak C_v$, the same compiler
therefore gives
\[
 b^*\mathscr W_v\longrightarrow(s_rb)^*\mathscr W_w,
 \qquad b^*\mathcal U_v\le(s_rb)^*\mathcal U_w.
\]
These comparisons commute with further substitution of parameter
locales because geometric inverse image preserves the witness map
and its image.  The family statement retains the actual code map
$\gamma$; an inclusion of feasibility opens alone supplies no compiler.

Over the locale $\mathfrak R_r$ of
\eqref{eq:compiler-budget-classifier}, this comparison further
gives $\pi_1^*\mathscr W_v\to\pi_2^*\mathscr W_w$.
Indeed $\kappa_v(e)\in I$ implies
$\kappa_w(\gamma e)\in\downarrow r(I)\subseteq J$.
The map still retains $\gamma e$; it is not obtained by choosing
a witness from the target feasibility image. Observation transport
is the equation
$\chi_{\J r(I)}(u)=\chi_I(u\circ r)$ of
\cref{prop:effective-resource-transport}, and implemented name
maps for $\gamma$ are supplied by the finite compiler descriptions above.
The budget-name operation is a separate constructed implementation,
as follows.

\begin{proposition}[Named transport of admitted budgets]
\label{prop:named-budget-transport}
The observed map $\F r$ has the admitted-point map $\J r$ and, on the
projective name presentations of \cref{prop:standalone-name-lifting}, an implementation
\begin{equation}\label{eq:named-budget-transport}
 \widehat{\J r}:C_L\to C_{L'},\qquad
 q_{L'}\widehat{\J r}=\J r\,q_L.
\end{equation}
For a family of monotone $r$, this construction is uniform after
retaining a projective name for that family parameter. The compiler
code for $\gamma$ and this budget tracker jointly present the action
on named implementations and admitted budgets.
\end{proposition}
\begin{proof}\leavevmode
\begin{enumerate}[leftmargin=*,align=left]
\item \textbf{Construct the observation-preserving budget map.}
Put $R=\operatorname{Mon}(L,L')$. The underlying uniform map
\[
 R\times\J L\longrightarrow\J L',\qquad (r,I)\longmapsto\J r(I)
\]
is constructed by the lower-image formula. Its observation equation is
\[
 \chi_{\J r(I)}(u)=\chi_I(u\circ r),
\]
uniformly in $r,I,u$. This is the observation-preserving comparison
required by \cref{prop:standalone-name-lifting}.

\item \textbf{Lift budget transport to retained names.}
Take its constructed projective covers $C_R\onto R$,
$q_L:C_L\onto\J L$, and $q_{L'}:C_{L'}\onto\J L'$.
The product of the displayed partitioned-assembly presentations is
projective. Pull back $q_{L'}$ along the preceding map on
$C_R\times C_L$ and use \cref{prop:standalone-name-lifting}. It gives a tracked lift
$C_R\times C_L\to C_{L'}$, whose commutative square is
\eqref{eq:named-budget-transport} with the parameter name retained.

\item \textbf{Verify parameter substitution on refinements.}
For a dependent parameter pullback, take a partitioned
refinement of the actual input presentation before extracting the
tracker. The same equation holds on that refinement; comparisons of
choices use their common projective refinement from \cref{prop:standalone-name-lifting}.

\item \textbf{Verify coherence after name realization.}
Identity and sequential composition of these retained name maps give
implementations of $\J\id$ and $\J(s r)$; their squares follow by
composition and functoriality of $\J$. Their positive observations
compose by substitution of tests. Different lifts can return different
names for the same budget, so coherence is asserted after the displayed
name-realization maps, not as equality of unrelated trackers.
Together with the finite code-output tracker for $\gamma$, these are
the two components of the named compiler action.
\end{enumerate}
\end{proof}

Neither projective lifting nor the observation equation assigns the
budget tracker's runtime. A runtime comparison for that operation is
obtained by retaining and measuring its execution description. Likewise,
the semantic bound parameter $r$ is supplied data of a comparison; it
is not computed merely by executing the program transformation $\gamma$.

\subsection{Measured changes of input representation}
\label{sec:measured-representation-change}
Keep $X_0,Y_0$ and the environment fixed. Use two input presentations
$q_X:C_X\onto X_0$ and $q'_X:C'_X\onto X_0$. Form the object of
retained conversion descriptions whose total returned-name map
$r:C'_X\to C_X$ satisfies $q_Xr=q'_X$. Its defining predicates are
formed as subobjects, just as for $\operatorname{CompCode}$; its generic
element carries its actual history $H_r$ and selected measurement.
For resource transport, take the subobject on which the two supplied
proper-filter contexts make $r$ a filtered-input map of
\cref{sec:input-substitution}. This retains the eventual-substitution
witness needed to form $r^*:K_X\to K_{C'_X}$.

\begin{proposition}[Conversion overhead from retained histories]
\label{prop:representation-overhead}
Sequential composition constructs an implementation $e^r$ on $C'_X$
with the same extensional denotation as $e$. Its domain is the
pullback of the original domain along $r$. In the retained writer
interpretation its history and additive-time profile are
\begin{equation}\label{eq:representation-overhead}
 H_{e^r}(c')=H_r(c'){+\!+}H_e(r(c')),
 \qquad T_{e^r}=T_r+T_e\circ r.
\end{equation}
On the total-implementation subobject, its exact germ is
$\Phi(e^r)=h+r^*\Phi(e)$, where
$h=[T_r]$. If $\varphi_r:L_X\to L_{C'_X}$ is the induced growth map,
the monotone budget transformation
\[
 a\longmapsto[h]\vee\varphi_r(a)
\]
constructs the corresponding measured compiler comparison.
\end{proposition}
\begin{proof}\leavevmode
\begin{enumerate}[leftmargin=*,align=left]
\item \textbf{Identify the composite domain and output.}
The code-composition constructor first runs the converter and then $e$.
Totality of the converter and the bind domain equation give the stated
domain. The two value-adequacy squares and $q_Xr=q'_X$ identify the
extensional output.

\item \textbf{Compute the exact history and cost.}
Writer bind concatenates the two histories, proving
the first equation; the additive fold proves the second. Passing to
germs gives the exact equality.

\item \textbf{Derive the growth bound and its retained witness.}
For nonnegative profiles, the growth
class of a sum is the join of their classes. Thus the cost of $e^r$ has
class $[h]\vee\varphi_r(\kappa(e))$, and monotonicity gives the
required compiler bound at every larger budget. The conversion and
composition descriptions are retained witnesses of its implementation.
\end{enumerate}
\end{proof}

For a different history valuation the formula is
$\nu(H_r{+\!+}H_e\circ r)$; no additive identity is inferred. In a
numerical compilation every actual dispatch and simulation instruction
belongs to the compiled history. Its bound is computed from that full
history, including any overhead of realizing the writer-level
composition. Opposite representation conversions give opposite measured
comparisons with their own overheads. This supplies the quantitative
comparison between presentations; equality of their unmeasured
denotations alone does not identify their complexity profiles.


The measured span now supplies a witness sheaf and a feasible open
for each represented problem, together with the action of implemented
compilers. The next section studies the feasible open itself.
Its covering properties will say whether one implementation, or a
finite portfolio, accounts for all feasible budgets.

\section{Feasibility moduli and optimality}
\label{sec:feasibility-moduli}

A problem determines an open of feasible budgets. We first classify
these open-valued specifications and then compare their covers with
the implementation witnesses that generate them. One witness
accounting for the whole open is an optimal implementation; finitely
many such witnesses form a dominating portfolio. The same construction
applies to joint resources, provided every constraint is imposed on
one shared implementation.
\subsection{The classifying locale of feasibility specifications}
\label{sec:feasibility-classifier}
The environment parameter remains fixed throughout this section:
$V,P_v$ mean $V_\theta,P_{\theta,v}$. Its classifying maps and
optimality statements are families over that parameter, with the
substitution comparisons of \cref{prop:environment-base-change}.
Let $\mathbb S$ denote the Sierpi\'nski locale, which classifies
opens, and let $K_c(\Cplx)$ be the internal lattice of compact
opens. Its elements are precisely finite unions of principal
opens: compactness extracts a finite subcover from the principal
decomposition, and finite unions of compact opens are compact.
Finite unions and intersections make $K_c(\Cplx)$ a bounded
distributive lattice.

\begin{theorem}[Feasibility specifications are classified]
\label{thm:feasibility-classifier}
The locale $\mathfrak F$ with frame
$\OO(\mathfrak F)=\Up(K_c(\Cplx))$ has the natural universal
property
\begin{equation}\label{eq:feasibility-classifier}
 \Loc_\E(Z,\mathfrak F)\cong\OO(Z\times\Cplx).
\end{equation}
Thus it constructs the exponential $\mathbb S^{\Cplx}$.
Its universal open is
\[
 \mathscr U=\bigvee_{a\in L}V_{U_a}\times U_a
       \ \in\OO(\mathfrak F\times\Cplx),
 \qquad V_k=\{\ell\in K_c(\Cplx):k\subseteq\ell\}.
\]
The measured incidence determines a canonical map
$c:\operatorname{Disc}(V)\to\mathfrak F$ classifying
$v\mapsto\mathcal U_v$.
\end{theorem}
\begin{proof}\leavevmode
\begin{enumerate}[leftmargin=*,align=left]
\item \textbf{Describe opens of the product by monotone functions.}
For $H=\OO(Z)$, the frame of the product $Z\times\Cplx$
is naturally the frame of monotone maps $w:L\to H$.
Indeed a rectangle $z\times U_a$ corresponds to the function
$b\mapsto\bigvee_{a\le b}z$, and every monotone function is
$\bigvee_a w(a)\times U_a$.

\item \textbf{Verify the frame-coproduct universal property.}
For frame maps $f:H\to H'$ and
$h:\Up(L)\to H'$, the unique extension is
$w\mapsto\bigvee_a f(w(a))\wedge h(U_a)$.
Finite meets are preserved using $U_a\cap U_b=U_{a\vee b}$,
and joins and top are preserved by principal decomposition.
This proves the frame-coproduct universal property, including
the inverse formulas.

\item \textbf{Pass from classifying ideals to monotone functions.}
By \cref{thm:classifying-frame}, a map $Z\to\mathfrak F$
is an $H$-valued ideal $j$ on $K_c(\Cplx)$. It determines
$w(a)=j(U_a)$.

\item \textbf{Construct the inverse ideal and verify independence.}
Conversely, for monotone $w$, put
\[
 j\left(\bigcup_{i=1}^n U_{a_i}\right)
       =\bigwedge_{i=1}^n w(a_i).
\]
This is independent of the finite presentation: if
$U_a\subseteq\bigcup_iU_{b_i}$, supercompactness gives
internally an $i$ with $b_i\le a$, hence
$\bigwedge_iw(b_i)\le w(a)$. Applying this to each generator
proves order reversal and independence; the empty union has
value $\top$, and unions become meets.

\item \textbf{Identify the universal open and the classification map.}
The constructions are
inverse and natural in $H$. Their composition proves
\eqref{eq:feasibility-classifier}, and sends the universal ideal
to the displayed open $\mathscr U$.
Apply it to the open relation \eqref{eq:classified-cost-open}
over $\operatorname{Disc}(V)$ to obtain $c$.
\end{enumerate}
\end{proof}


Let $\chi_{\mathscr F}$ classify the universal measured feasibility
family on $\operatorname{Disc}(V)\times\Cplx$.
The exponential transpose $c$ is characterized by
\[
\begin{tikzcd}[column sep=large,row sep=large]
 \operatorname{Disc}(V)\times\Cplx
   \arrow[r,"{c\times\id}"] \arrow[dr,"{\chi_{\mathscr F}}"']
   & \mathfrak F\times\Cplx \arrow[d,"{\mathrm{ev}}"]\\
 & \mathbb S .
\end{tikzcd}
\]
Here $\mathrm{ev}$ classifies the universal open $\mathscr U$.
This triangle relates the problem-indexed observable to the
representing locale of feasibility specifications.


The point $c(v)$ records the feasible specification. Its measured
instance also carries the epimorphism $\mathscr W_v\onto\mathcal U_v$,
which records the implementations. Selecting a section of that
epimorphism is a further question. Effective names and positive
observations retain the structures already constructed; the
classification alone assigns neither to an arbitrary point of
$\mathfrak F$.
The principal opens contributed by implementations cover the feasible
open. If one member covers the whole, its implementation is optimal.
If a finite subfamily covers it, those implementations suffice as a
portfolio. Supercompactness and compactness express these two
possibilities intrinsically.

\subsection{Optimization and topological compactness}
\label{sec:optimality-topology}
The implementation comparison is with abstract complexity measures
\citep{Blum1967MachineIndependentComplexity}; finite portfolios have
also been studied as algorithm-selection objects
\citep{GomesSelman2001AlgorithmPortfolios}. Here the relevant
properties will be expressed by compactness of a feasibility open.
Fix a denotation $v:V$ and write $P_v=\{e:P:p(e)=v\}$.
The universal cost open of \cref{thm:universal-cost} is
$\mathcal U_v=\bigvee_{e\in P_v}U_{\kappa(e)}$.
We use the internal frame definitions: an open $W$ is
\emph{supercompact} if $W\le\bigvee_{i:J}W_i$ implies
$\exists i:J\;(W\le W_i)$, and \emph{compact} if it implies
$\exists n\in\N\,\exists i_1,\ldots,i_n:J\;
 W\le\bigvee_{k=1}^nW_{i_k}$.
The empty family is included in these definitions.

\begin{theorem}[Optimal algorithms and finite portfolios]
\label{thm:optimization-compactness}
The classified feasibility open satisfies
\begin{equation}\label{eq:supercompact-optimal}
 \mathcal U_v\text{ is supercompact}
 \quad\Longleftrightarrow\quad
 \exists e_0\in P_v\ \forall e\in P_v\;
       \kappa(e_0)\le\kappa(e).
\end{equation}
It is compact precisely when a finite asymptotically dominating
portfolio exists:
\begin{equation}\label{eq:compact-portfolio}
 \begin{aligned}
 \mathcal U_v\text{ is compact}\quad\Longleftrightarrow\quad
 &\exists n\in\N\,\exists e_1,\ldots,e_n\in P_v\;\forall e\in P_v\\
 &\qquad\exists k\in\{1,\ldots,n\}\;
                 \kappa(e_k)\le\kappa(e).
 \end{aligned}
\end{equation}
In particular the empty portfolio accounts for the empty
feasibility open.  Both equivalences are internal and stable under
substitution of denotation parameters.
\end{theorem}
\begin{proof}\leavevmode
\begin{enumerate}[leftmargin=*,align=left]
\item \textbf{Reflect growth order through principal opens.}
The principal opens satisfy
\begin{equation}\label{eq:principal-open-reflection}
 U_b\le U_a\quad\Longleftrightarrow\quad a\le b.
\end{equation}
For the forward implication evaluate inclusion at
$\downarrow b$; the reverse implication is upperness.

\item \textbf{Deduce supercompactness from an optimal implementation.}
If an $e_0$ as in \eqref{eq:supercompact-optimal} exists,
\eqref{eq:principal-open-reflection} makes every
$U_{\kappa(e)}$ lie below $U_{\kappa(e_0)}$, so
$\mathcal U_v=U_{\kappa(e_0)}$ is supercompact.

\item \textbf{Extract an optimal implementation.}
Conversely, apply supercompactness to the canonical cover
$\mathcal U_v=\bigvee_{e\in P_v}U_{\kappa(e)}$.
It supplies $e_0$ with $\mathcal U_v\le U_{\kappa(e_0)}$.
For every $e$, the principal-open reflection then yields
$\kappa(e_0)\le\kappa(e)$.

\item \textbf{Extract a finite dominating portfolio.}
Compactness of $\mathcal U_v$ applied to the same cover
produces a finite list $e_1,\ldots,e_n$ with
$\mathcal U_v\le\bigvee_kU_{\kappa(e_k)}$.
For any $e$, supercompactness of $U_{\kappa(e)}$ extracts
one $k$ with $U_{\kappa(e)}\le U_{\kappa(e_k)}$; the
reflection equation gives the displayed domination.

\item \textbf{Deduce compactness from a portfolio.}
Conversely, a finite dominating portfolio makes
$\mathcal U_v=\bigvee_kU_{\kappa(e_k)}$, a finite union
of compact principal opens.  All arguments use finite lists and internal existential elimination.
\end{enumerate}
\end{proof}

\begin{theorem}[Optimal implementations are sections]
\label{thm:optimal-sections}
For $v:V$, there is a natural isomorphism of internal $\E$-objects
\begin{equation}\label{eq:optimal-sections}
 \operatorname{Sect}(\mathscr W_v\to\mathcal U_v)
 \cong\{e_0\in P_v:\forall e\in P_v,\
                         \kappa(e_0)\le\kappa(e)\}.
\end{equation}
Here $\operatorname{Sect}$ represents sections in all parameter
contexts, equivalently $p_{\Cplx*}(\mathscr W_v^{\mathcal U_v})$.
Consequently the witness epimorphism splits internally exactly
when $\mathcal U_v$ is supercompact.
\end{theorem}
\begin{proof}\leavevmode
\begin{enumerate}[leftmargin=*,align=left]
\item \textbf{Show that a section selects one common implementation.}
In \cref{thm:budget-diagrams}, a section assigns a feasible
implementation $e_a$ to every feasible budget $a$, with
$e_a=e_b$ when $a\le b$. Any two feasible budgets have the
feasible common upper bound $a\vee b$; hence all selected
implementations agree. Since $v\in V=\im(p_0)$, the fibre
$P_v$ is internally inhabited. Thus there is a unique common
implementation $e_0$; unique existence descends through the
inhabited parameter cover.

\item \textbf{Prove optimality of the selected implementation.}
At $a=\kappa(e)$
feasibility gives $\kappa(e_0)\le\kappa(e)$.

\item \textbf{Construct the inverse section and verify naturality.}
Conversely such an $e_0$ is available at every feasible budget
and gives a constant compatible section. The two operations
are inverse and commute with parameters. The final assertion
is \cref{thm:optimization-compactness}.
\end{enumerate}
\end{proof}


An optimal implementation is therefore exactly a splitting
\[
\begin{tikzcd}[column sep=large]
 \mathscr W_v \arrow[r,two heads,shift left=.8ex,"{\pi}"]
   & \mathcal U_v \arrow[l,shift left=.8ex,"{s}"]
\end{tikzcd}
\qquad \pi s=\id_{\mathcal U_v}.
\]
A section makes compatible choices at all feasible budgets.
In the diagram presentation, the common upper budget compares
the choices made at $a$ and $b$:
\[
\begin{tikzcd}[column sep=large,row sep=large]
 \mathsf{Impl}_v(a) \arrow[dr]
   && \mathsf{Impl}_v(b) \arrow[dl]\\
 & \mathsf{Impl}_v(a\vee b) &
\end{tikzcd}
\]
Both transition maps retain the implementation code, so compatibility
forces those choices to agree. This is why a section supplies one
optimal witness throughout the feasible open.

The optimality criterion can also be evaluated after allowing chosen
overheads. Pulling the witness family to the gauge sublocale gives
the same questions relative to the localized notion of resource.

\subsection{Gauge-relative feasibility}
\begin{corollary}[Soft optimality and portfolios]
\label{cor:soft-optimization}
Let $\mathcal U_v^{(P)}$ be the restriction of the measured
feasibility open to $\Cplx^{(P)}$.  Then
\[
 \mathcal U_v^{(P)}
   =\bigvee_{e\in P_v}U^{(P)}_{q_P(\kappa(e))}.
\]
It is supercompact precisely when an implementation $e_0$
satisfies $\kappa(e_0)\preceq_P\kappa(e)$ for all $e\in P_v$.
It is compact precisely when finitely many implementations
dominate every implementation in this gauge-relative preorder.
For $P=P_\lambda$, this is optimality and finite-portfolio
completeness for soft-$O$ comparison.
\end{corollary}
\begin{proof}\leavevmode
\begin{enumerate}[leftmargin=*,align=left]
\item \textbf{Compute feasibility in the localized growth order.}
Under \eqref{eq:soft-locale-identification}, membership of a
cost in a saturated budget is membership of its $q_P$-class
in the corresponding ideal of $L/P$.  This gives the displayed
open formula.

\item \textbf{Characterize soft optimality and finite portfolios.}
Principal opens of $L/P$ are supercompact;
applying the proof of \cref{thm:optimization-compactness}
to that principal cover proves both equivalences.  The bitlength
gauge identifies $\preceq_{P_\lambda}$ with soft domination by
\cref{sec:transport}.
\end{enumerate}
\end{proof}

\subsection{Joint-resource moduli and Pareto portfolios}
\label{sec:multi-resource-moduli}
Separate resource bounds may be witnessed by different programs.
Joint feasibility instead requires one program meeting every bound.
Let $L_T,L_S$ be growth lattices for time and space measurements
of the same implementation object, with classifying locales
$\mathfrak C_T,\mathfrak C_S$. The product construction below
retains that common implementation while varying both budgets.

\begin{theorem}[Product moduli and joint feasibility]
\label{thm:multi-resource}
There is a canonical frame isomorphism
\begin{equation}\label{eq:product-growth-frame}
 \OO(\mathfrak C_T\times\mathfrak C_S)
 \cong\Up(L_T\times L_S).
\end{equation}
Each principal rectangle $U_a^T\times U_b^S$ is
supercompact.  For joint measured costs
$\kappa_T:P_v\to L_T$, $\kappa_S:P_v\to L_S$, the
simultaneous feasibility open is
\begin{equation}\label{eq:joint-feasibility}
 \mathcal U_v^{T,S}=
 \bigvee_{e\in P_v}
    (U^T_{\kappa_T(e)}\times U^S_{\kappa_S(e)}).
\end{equation}
It is supercompact exactly when one implementation is
simultaneously asymptotically optimal in both costs.
It is compact exactly when finitely many implementations
jointly dominate every implementation: for each $e$ there
is a member $e_i$ of the finite list satisfying
\[
 \kappa_T(e_i)\le\kappa_T(e),\qquad
 \kappa_S(e_i)\le\kappa_S(e).
\]
The witness sheaf is the coproduct of the principal rectangles
in \eqref{eq:joint-feasibility}; hence both budget tests
refer to the same implementation.
\end{theorem}
\begin{proof}\leavevmode
\begin{enumerate}[leftmargin=*,align=left]
\item \textbf{Construct the frame map on principal rectangles.}
Every upper predicate on $L_T\times L_S$ is the union of
principal rectangles.  Given frame maps
$h_T:\Up(L_T)\to H$ and $h_S:\Up(L_S)\to H$, define
\[
 h(W)=\bigvee_{(a,b)\in W}
          (h_T(U_a^T)\wedge h_S(U_b^S)).
\]

\item \textbf{Verify the product universal property.}
The identity
$(U_a^T\times U_b^S)\wedge(U_c^T\times U_d^S)
 =U_{a\vee c}^T\times U_{b\vee d}^S$
and frame distributivity prove that $h$ preserves finite
meets; principal decomposition proves join preservation.
Restriction along the two projections gives the inverse
bijection.  This is the coproduct universal property of the
frame of the localic product, proving
\eqref{eq:product-growth-frame}.

\item \textbf{Prove supercompactness of principal rectangles.}
Evaluating a covering
family at $(a,b)$ proves supercompactness of each principal
rectangle as in \cref{thm:classifying-frame}.

\item \textbf{Construct joint feasibility and its witness sheaf.}
The simultaneous incidence is the pullback of the two
membership predicates along the diagonal implementation
map, followed by the regular image to the budget product.
The image is exactly \eqref{eq:joint-feasibility}; taking
its implementation fibre gives the stated witness sheaf.

\item \textbf{Characterize joint optimality and portfolios.}
The supercompactness and compactness proofs now repeat
\cref{thm:optimization-compactness} with the product order
and principal rectangles.  Both comparisons are coordinatewise,
so they yield exactly the two displayed cost inequalities.
\end{enumerate}
\end{proof}

In particular, intersecting the two separately existential
feasibility predicates may use different implementations.
Equation~\eqref{eq:joint-feasibility} is the geometric
construction of simultaneous time--space feasibility.


The section characterization of \cref{thm:optimal-sections} also
applies to the product budget order: a section of the joint witness
sheaf is exactly one implementation optimal in every displayed cost.
Its proof uses the componentwise join as the common upper budget.


\subsection{A compact two-resource portfolio}
\label{app:complexity-portfolio-example}

The next example calculates a compact feasible open for which neither
implementation is optimal in both resources. We use a selected
two-description presentation and two event-count valuations, so the
entire implementation fibre can be inspected.
Work on $X=\N$ with its cofinite filter, and fix two implemented
unit-valued idle primitives with distinct tags $\mathsf a,\mathsf b$
and identical extensional effect. Define the product valuation by
\[
 w(\mathsf a)=(1,0),\qquad
 w(\mathsf b)=(0,1),\qquad w([])=(0,0),
\]
and extend it to words by componentwise addition. These are two
points of the history-weight object of
\cref{prop:universal-history-weights}. They count the selected
events; total machine time and peak-space occupancy are separate
measurements.

Finite NNO iteration and multiplication construct two fixed total
descriptions $e_A,e_B$ with the same output, namely the identity
function $v(n)=n$, and respective logged histories
\begin{equation}\label{eq:app-portfolio-histories}
 H_{e_A}(n)=\mathsf a^{n}\,\mathsf b^{n^2},\qquad
 H_{e_B}(n)=\mathsf a^{n^2}\,\mathsf b^{n}.
\end{equation}
The exponents here are finite repetition counts implemented by the
retained arithmetic and iteration constructors.  In each case the
program returns its unchanged input after the logged calls.  The
history fold therefore gives \emph{exact} raw cost pairs
\[
 \bigl(T_{1,A}(n),T_{2,A}(n)\bigr)=(n,n^2),\qquad
 \bigl(T_{1,B}(n),T_{2,B}(n)\bigr)=(n^2,n).
\]
Put $a=[n+1]$ and $b=[(n+1)^2]$ in the numerical growth lattice.
Then $a<b$: the inequality $a\le b$ holds, while $b\not\le a$
is witnessed by arbitrarily large numerical inputs.
The growth pairs of the two implementations are $(a,b)$ and $(b,a)$,
which are incomparable in the coordinatewise order.

Take the two-description subpresentation
$E_{\mathrm{port}}=\{e_A,e_B\}$. Its evaluator is the copairing
of the two total evaluations, its denotation image contains just
$v$, and its implementation object consists of the two displayed
members. The measured-image and incidence constructions therefore
give its complete joint feasible open:
\begin{equation}\label{eq:app-portfolio-open}
 \mathcal U^{\mathrm{port}}_v
 =\bigl(U^{(1)}_{a}\times U^{(2)}_{b}\bigr)
  \ \vee\
  \bigl(U^{(1)}_{b}\times U^{(2)}_{a}\bigr)
 \quad\in\OO(\Cplx^{(1)}\times\Cplx^{(2)}).
\end{equation}
The budget points $(\downarrow a,\downarrow b)$ and
$(\downarrow b,\downarrow a)$ are feasible, whereas
$(\downarrow a,\downarrow a)$ is not feasible in this presentation.
Each program meets one of the smaller bounds, but neither meets
both. Choosing different witnesses for the two resources would
therefore give the wrong joint judgment.


For clarity, put $I_a=\downarrow a$ and $I_b=\downarrow b$.  Direct
evaluation of the two principal rectangles gives the following
complete table for the displayed portfolio:
\[
\begin{array}{c|cc|c}
 \text{joint budget}& e_A\text{ feasible}&e_B\text{ feasible}
       &\text{some implementation feasible}\\ \hline
 (I_a,I_a)&\text{no}&\text{no}&\text{no}\\
 (I_a,I_b)&\text{yes}&\text{no}&\text{yes}\\
 (I_b,I_a)&\text{no}&\text{yes}&\text{yes}\\
 (I_b,I_b)&\text{yes}&\text{yes}&\text{yes}
\end{array}
\]
Each entry tests both resource components of one of the two measured
cost pairs. The table evaluates the complete open of this finite
presentation.


The open in \eqref{eq:app-portfolio-open} is compact, being a union of
two supercompact principal rectangles.  It is not supercompact.  Indeed,
if one of the two rectangles contained the union, it would contain the
other rectangle.  The budget point $(I_b,I_a)$ belongs to the second
rectangle but not the first, while $(I_a,I_b)$ belongs to the first
but not the second.  Thus the displayed two-rectangle cover has no
single covering member.  By
\cref{thm:optimization-compactness,thm:multi-resource}, the two codes
form a finite dominating portfolio, whereas neither code is
simultaneously optimal in both measurements \emph{within this selected
subpresentation}.  This computes both sides of the compactness
criterion on an implemented instance.

In the full description language, this calculated open is a subopen
of $\mathcal U_v^{(1,2)}$. Other implementations can enlarge it:
an uncharged pure identity description makes the full open top for
these idle-event counters. The example thus establishes the
portfolio criterion for its specified two-description fibre.
The product theorem also applies to the time and occupancy
measurements of \cref{constr:numerical-execution}; identifying a
particular frontier there additionally requires the corresponding
universal lower-bound analysis.

\label{app:complexity-probe-scope}
The classifier now describes budget families, their feasible opens,
and their implementation witnesses. Keeping a measured history fixed
and varying its valuation gives the resource families studied next.

\section{Differential calculus of resource constraints}
\label{sec:relative-geometry}
The measured span separates a program's denotation from its resource
valuation. Keeping its execution history fixed while varying event
weights produces a family of resource values. We begin with this
family and its budget equation, then construct the finite coefficient
calculus that differentiates them.

\subsection{Measured valuation families and budget variations}
\label{sec:valuation-geometry-bridge}
For a fixed total implementation $e$ with finitely many measured event
types, let $N_{e,i}$ be its exact event-count profiles. Additive rational
weights give the measured family
\[
 F_e(w)=\sum_i N_{e,i}w_i.
\]
This follows by folding the returning history
(\cref{thm:history-valuation}). Passing to germs and extending
coefficients gives the $B$-linear map $F_e:B^k\to K$ with the same
formula. A polynomial parametrization $w=w(z)$ therefore gives a
polynomial resource law. Its finitely many coefficients lie in the
principal resource generated by the sum of their absolute values,
so it is calculated on a finite-projective resource cell.

For an embedded budget $M\subseteq K$, put $q_M:K\to K/M$.
Feasibility is the equation $q_MF_e(z)=0$. Over the dual numbers,
polynomial expansion gives
\[
 q_MF_e(z+\eps v)=q_MF_e(z)+\eps q_MDF_{e,z}(v).
\]
At a feasible point, the fixed-budget tangent condition is therefore
$q_MDF_{e,z}(v)=0$. If a finite-projective budget also varies, its
tangent class is $\theta\in\Hom_B(M,K/M)$ and the condition becomes
\[
 q_MDF_{e,z}(v)=\theta(F_e(z)).
\]
\Cref{thm:embedded-deformation-tangent,thm:moving-incidence} construct
the budget tangent and prove this equation by comparing coefficients.

\begin{theorem}[Measured Taylor coefficients and geometric jets]
\label{thm:measured-jet-adequacy}
Let $H_e$ be the returning history of a retained implementation, with
finitely many measured event types and polynomial weights
$w_i(z)\in\QQ[z_1,\ldots,z_m]$. Write
\[
 \tau_r:\QQ[z]\longrightarrow
 \QQ[z,\eps_1,\ldots,\eps_m]/(\eps_1,\ldots,\eps_m)^{r+1},
 \qquad p(z)\longmapsto p(z+\eps).
\]
Folding $H_e$ with the weights $\tau_rw_i$ gives the Taylor coefficients
of its measured cost through order $r$:
\begin{equation}\label{eq:measured-jet-fold}
 \operatorname{fold}(\tau_rw,H_e)
 =\tau_r\operatorname{fold}(w,H_e).
\end{equation}
The two executions have the same return domain and returned value.
For a total implementation on the retained input carrier, taking exact
germs identifies these coefficients with the order-$r$ geometric jet
of $F_e(z)=\sum_iN_{e,i}w_i(z)$. The comparisons commute with truncation,
polynomial parameter substitution, and the retained environment maps.
Their finite coefficient operations have uniform implementations in
the order and the displayed polynomial data.
\end{theorem}
\begin{proof}\leavevmode
\begin{enumerate}[leftmargin=*,align=left]
\item \textbf{Compare Taylor substitution with finite history folds.}
The map $\tau_r$ is the algebra homomorphism specified by its displayed
values on generators. In particular it preserves the additive monoid
of weights. On a returning word $[a_1,\ldots,a_k]$, both sides of
\eqref{eq:measured-jet-fold} are
$\sum_{j=1}^k w_{a_j}(z+\eps)$ modulo total degree $r+1$.
Induction on word length proves the equality.

\item \textbf{Extend the comparison to iteration and erasure.}
The finite-return
components of \cref{thm:weighted-return} have this same word fold;
their common classified union proves the equality for iteration.
Valuation changes only the resource component, so erasure and its
domain agree by \cref{thm:cost-partiality,thm:measured-name-presentation}.

\item \textbf{Descend the coefficients to germs.}
At every actual input the equality is coefficientwise in exact
rational polynomials. The germ quotient preserves the finite sums and
products, hence carries it to $K$.

\item \textbf{Identify the geometric jet on a principal cell.}
The finitely many coefficients of
$F_e$ lie in the principal cell generated by the sum of their absolute
values. The polynomial realization and principal-parts comparisons
of \cref{thm:polynomial-lift,prop:jet-coefficients,thm:primary-jet-coefficients}
identify substitution by $z+\eps$ on that cell with its geometric jet.
Thus the geometric coefficients are the measured coefficients.

\item \textbf{Verify substitution and truncation coherence.}
Truncation deletes higher monomials, and polynomial substitution
commutes with the displayed algebra homomorphisms. The history
substitution theorem gives the environment comparison.

\item \textbf{Construct uniform coefficient operations.}
Monomial
enumeration, rational arithmetic, and finite list folds construct the
coefficient operations by NNO recursion; the retained arithmetic name
maps therefore implement them uniformly. The Taylor coefficients of
positive cost valuations may be signed, and are interpreted by the
exact signed coefficient construction.
\end{enumerate}
\end{proof}

The theorem supplies the execution interpretation of the differential
calculus: finite response coefficients are obtained by instrumenting
the same returning history. Its compatible truncations will give the
measured formal families used in the completion and crystalline
comparisons.


\begin{example}[A measured polynomial resource family]
\label{ex:weighted-resource-family}
Work on $\N$ with the cofinite filter and the unit-valued idle events
$\mathsf a,\mathsf b$ of \cref{app:complexity-portfolio-example}.
Finite iteration constructs a fixed total program that emits one
$\mathsf a$ and then $(n+1)^2$ copies of $\mathsf b$, and returns its
input. Assign weights $1$ and $\lambda^2$, respectively. The history
fold gives the exact nonnegative rational cost
\[
 T_\lambda(n)=1+(n+1)^2\lambda^2\qquad(\lambda\in\QQ).
\]
Writing $t(n)=1/(n+1)$, its polynomial coefficient extension is
\begin{equation}\label{eq:running-resource-family}
 F:B\longrightarrow K,\qquad F(\lambda)=1+t^{-2}\lambda^2.
\end{equation}
The base value $F(0)=1$ belongs to $B$. At first order,
$F(d)=1$ in $K[d]/d^2$. At second order the additional coefficient
is $t^{-2}$, whose class in $K/B$ is nonzero. We will use this same
family to calculate fixed, moving, and coupled-budget obstructions in
\cref{ex:second-obstruction,ex:moving-removes-obstruction,ex:coupled-budget-obstruction}.
\end{example}

Free commutative algebras represent the polynomial expressions above.
Their universal derivatives compare the exact, quantitative, and
signed interpretations. Based Weil stages then represent finite
coefficient variations, and module realization carries the budget
equation into the differential geometry.

\subsection{The differential extension and its comparison maps}\label{sec:differential}
\subsubsection{Free continuous commutative algebra}
\begin{theorem}[Free continuous commutative algebra]\label{thm:free-algebra}
The forgetful functor $\CAlg(\C)\to\C$ has a constructed left adjoint. Its underlying monad $T=\Sym_S^{\J}$ is an algebra modality on $\C$.
\end{theorem}
\begin{proof}\leavevmode
\begin{enumerate}[leftmargin=*,align=left]
\item \textbf{Present the finite polynomial algebra.}
For $M\in\C$, begin with finite commutative polynomial expressions over $S$ in symbols $[m]$, imposing
\[
 [m+m']=[m]+[m'],\qquad [sm]=s[m],\qquad[0]=0,
\]
and the commutative unital algebra equations. Give this algebra the least compatible preorder incorporating the orders of $S$ and $M$.

\item \textbf{Construct the continuous algebra.}
Impose covers
\[
 [\sup I]\triangleleft\J[-](I)\quad(I\in\J M),
 \qquad
 \sup K\triangleleft\J\iota_S(K)\quad(K\in\J S),
\]
where $\iota_S$ inserts scalar constants. Close the covers under finite one-hole algebra contexts. \Cref{thm:standalone-presentation,lem:standalone-product-generation} construct a continuous algebra $TM$. Scalar contexts give joint continuity of scalar action, and multiplication contexts give continuous multiplication. Equations extend from generators by product generation.

\item \textbf{Prove the free-algebra universal property.}
A continuous linear map $M\to A$ into a continuous commutative algebra extends algebraically to the finite expressions. It respects every cover by continuity, including contextual covers. It therefore extends uniquely to a continuous algebra homomorphism $TM\to A$. This proves the free-algebra universal property.

\item \textbf{Construct the algebra modality and verify its laws.}
The unit $\eta_M$ inserts a linear generator; $\mu_M:TTM\to TM$ substitutes polynomials. The monad laws follow because both sides have the same action on free generators. The multiplication $m_M:TM\otimes TM\to TM$ and unit $u_M:S\to TM$ are natural, and $\mu_M$ is an algebra map. Thus this is an algebra modality.
\end{enumerate}
\end{proof}

The same finite-expression construction, without admitted covers, gives $T_+=\Sym_\Qplus$ on $\A$ and $T_\QQ=\Sym_\QQ$ on $\Dcat$.

\subsubsection{Universal derivative and its identities}
Finite products in $\C$ are also biproducts: $(m,n)=(m,0)+(0,n)$, and all structure maps are continuous linear. For a continuous algebra $A$ and a continuous $A$-module $N$, form the square-zero extension
\begin{equation}\label{eq:square-zero}
 A\ltimes N=A\oplus N,\qquad
 (a,n)(b,n')=(ab,an'+bn),\qquad 1=(1,0).
\end{equation}
The finite equations prove associativity and commutativity; its operations are continuous because the module action and addition are continuous.

\begin{theorem}[Constructed deriving transformation]\label{thm:derivative}
There is a natural continuous linear map
\[
 d_M:TM\longrightarrow TM\otimes_S M
\]
which is the universal continuous $S$-derivation of the free algebra $TM$. It satisfies the constant, Leibniz, linear, monadic chain, and interchange identities.
\end{theorem}
\begin{proof}\leavevmode
\begin{enumerate}[leftmargin=*,align=left]
\item \textbf{Construct the universal derivation by a square-zero extension.}
Regard $TM\otimes M$ as the free $TM$-module on $M$. The continuous linear map
\[
 M\longrightarrow TM\ltimes(TM\otimes M),\qquad
 m\longmapsto(\eta_Mm,1\otimes m)
\]
extends uniquely to an algebra map. Its first component is $\id_{TM}$ by freeness. Call the second component $d_M$. Equation~\eqref{eq:square-zero} gives
\begin{equation}\begin{aligned}
 d_M(1)&=0, & d_M(ab)&=a\,d_Mb+b\,d_Ma,\\
 d_M(\eta_Mm)&=1\otimes m.\label{eq:d-identities}
\end{aligned}\end{equation}

\item \textbf{Prove the representing property.}
More generally, a continuous linear $v:M\to N$ into any $TM$-module defines an algebra map with generator $(\eta m,v(m))$. This proves
\[
 \Der_S^{\J}(TM,N)\cong\C(M,N)
 \cong\Mod_{TM}(TM\otimes M,N),
\]
so $d_M$ is universal.

\item \textbf{Verify naturality.}
For $f:M\to N$, both $(Tf\otimes f)d_M$ and $d_NTf$ are derivations along $Tf$ and send a generator $m$ to $1\otimes f(m)$. Universality proves naturality.

\item \textbf{Prove the monadic chain identity.}
Use composition in the usual right-to-left order. The chain identity is
\begin{equation}\label{eq:monadic-chain}
 d_M\mu_M
 =(m_M\otimes1_M)(\mu_M\otimes d_M)d_{TM}.
\end{equation}
Both sides are derivations from the free algebra $TTM$, with target module structure through $\mu_M$. At its generator $p\in TM$, the left side is $d_Mp$; the right side is also $d_Mp$. The universal property therefore proves~\eqref{eq:monadic-chain}.

\item \textbf{Prove interchange on monomials.}
On a finite monomial,
\begin{equation}\label{eq:monomial-d}
 d_M(m_1\cdots m_n)
 =\sum_i\left(\prod_{j\ne i}m_j\right)\otimes m_i.
\end{equation}
Applying $d_M$ again to its first factor sums over ordered pairs of distinct positions. Interchanging the two selected positions permutes this finite sum. Hence $(d_M\otimes1_M)d_M$ is invariant under switching its final two factors.

\item \textbf{Extend the identities to the presented domain.}
Both sides are continuous linear maps, so the identity extends from finite expressions to the entire presented domain. This is the interchange identity. The other displayed identities were already proved by square-zero multiplication and hold on completed elements as well.
\end{enumerate}
\end{proof}

These universal properties and identities establish a codifferential
category in the terminology of \citep{BLWON,BCS}. In the admitted
setting, its free continuous algebra and derivative have been
constructed explicitly, so the later realization comparisons can
be checked on their generators.

\subsubsection{Universal differentials of a general algebra}
\begin{proposition}[Universal continuous differentials]\label{prop:kahler}
Every continuous commutative algebra $A$ has a constructed module $\Omega^{1,\J}_{A/S}$ representing continuous $S$-derivations. For a continuous $A$-module $N$,
\[
 \Mod_A(\Omega^{1,\J}_{A/S},N)
 \cong\Der_S^{\J}(A,N)
 \cong\CAlg(\C)_{/A}(A,A\ltimes N).
\]
\end{proposition}
\begin{proof}\leavevmode
\begin{enumerate}[leftmargin=*,align=left]
\item \textbf{Present the module of continuous differentials.}
Start with the free continuous $A$-module $A\otimes_S A$ on the underlying $S$-module of $A$, writing $da=1\otimes a$. Impose the module congruence
\[
 d1=0,\qquad d(ab)=a\,db+b\,da.
\]
Construct its continuous quotient by the admitted-presentation theorem, using equations as inequalities in both directions and imposing the original module's continuity covers.

\item \textbf{Prove the universal bijections.}
Maps from this quotient are exactly continuous $S$-linear maps satisfying the displayed relations. This is the first bijection. Equation~\eqref{eq:square-zero} proves the second.

\item \textbf{Construct relative differentials.}
Relative differentials $\Omega_{A/A_0}$ are constructed by additionally imposing $d(a_0)=0$ for the specified coefficient map $A_0\to A$.
\end{enumerate}
\end{proof}

The algebraic versions over $\Qplus$ and $\QQ$ use the same generators and relations. One must use congruences, not subtract generators, on the semiring side.

\subsubsection{The differential comparisons commute}\label{sec:comparison}
\begin{theorem}[Comparison of modalities and derivations]\label{thm:modality-comparison}
There are natural coherent isomorphisms
\begin{equation}\begin{aligned}
 \Ext(T_+M)&\cong T(\Ext M),\\
 \Gr(T_+M)&\cong T_\QQ(\Gr M).\label{eq:symmetric-comparison}
\end{aligned}\end{equation}
They commute with algebra multiplication, monad unit and multiplication, and the deriving transformations. For an exact commutative $\Qplus$-algebra $A$ there are also natural isomorphisms
\begin{equation}\begin{aligned}
 \Ext\Omega^1_{A/\Qplus}&\cong\Omega^{1,\J}_{\Ext A/S},\\
 \Gr\Omega^1_{A/\Qplus}&\cong\Omega^1_{\Gr A/\QQ},\label{eq:kahler-comparison}
\end{aligned}\end{equation}
and corresponding relative-coefficient comparisons.
\end{theorem}
\begin{proof}\leavevmode
\begin{enumerate}[leftmargin=*,align=left]
\item \textbf{Extend the module adjunction to algebras.}
Strong monoidality equips $\Ext A$ with a continuous algebra
structure. Write $U_+:\CAlg(\C)\to\CAlg(\A)$ for restriction
of scalars on algebras, $U_{\C}:\CAlg(\C)\to\C$ for
forgetting the algebra operations, and $U:\C\to\A$ as above.
The module adjunction restricts to an algebra adjunction:
an exact algebra map $A\to U_+N$ extends to $\Ext A\to N$,
and its multiplication and unit equations hold on exact generators
and hence by admitted generation. Conversely scalar restriction
preserves them.

\item \textbf{Compare the free-algebra modalities.}
Consequently, for $N\in\CAlg(\C)$,
\[
\begin{aligned}
 \CAlg(\C)(\Ext T_+M,N)
 &\cong\CAlg(\A)(T_+M,U_+N)\\
 &\cong\A(M,U U_{\C}N)\\
 &\cong\C(\Ext M,U_{\C}N)\\
 &\cong\CAlg(\C)(T(\Ext M),N).
\end{aligned}
\]
This is the universal property of $T(\Ext M)$. The group-completion comparison has the same proof.

\item \textbf{Verify compatibility of the modality structures.}
The isomorphisms send each exact generator to the corresponding linear generator. The unit, substitution, algebra multiplication, and all their coherence diagrams therefore agree by their universal properties.

\item \textbf{Compare the deriving transformations.}
Both functors preserve finite biproducts and tensor products, and hence square-zero extensions. Under the free-algebra comparison, the image of the algebra map defining $d$ is the algebra map sending a generator to $(\eta m,1\otimes m)$. Uniqueness proves compatibility with $d$, including on the continuous side's admitted closure.

\item \textbf{Compare the representing modules of derivations.}
For general differentials, let $N$ be an $\Ext A$-module. The adjunction on module fibres gives
\[
 \Mod_{\Ext A}(\Ext\Omega_{A/\Qplus}^1,N)
 \cong\Mod_A(\Omega_{A/\Qplus}^1,UN)
 \cong\Der_\Qplus(A,UN).
\]
In the last expression the $A$-action is through $A\to U\Ext A$. A derivation on $A$ is an algebra map $A\to U(\Ext A\ltimes N)$ with prescribed first component. Its adjoint is an algebra section $\Ext A\to\Ext A\ltimes N$, and thus a continuous derivation. This proves the first comparison by representability. Group completion gives the second in the same way.

\item \textbf{Deduce the relative comparisons.}
Requiring vanishing on a coefficient algebra commutes with these correspondences, proving the relative assertions.
\end{enumerate}
\end{proof}


To differentiate a retained presentation, its exact and quantitative
derivatives must have the same interpretation. The free-algebra
comparison supplies this compatibility. With
$\kappa_M:\Ext T_+M\cong T\Ext M$, it is the square
\begin{equation}\label{eq:cat-deriving-comparison}
\begin{tikzcd}[column sep=large,row sep=large]
 \Ext(T_+M) \arrow[r,"{\Ext d_M^+}"] \arrow[d,"{\kappa_M}"']
   & \Ext(T_+M\otimes_{\Qplus}M) \arrow[d,"{\cong}"]\\
 T(\Ext M) \arrow[r,"{d_{\Ext M}}"']
   & T(\Ext M)\otimes_S\Ext M .
\end{tikzcd}
\end{equation}
The corresponding square for $\Gr$ has the same construction:
\[
\begin{tikzcd}[column sep=large,row sep=large]
 \Gr(T_+M) \arrow[r,"{\Gr d_M^+}"] \arrow[d,"{\cong}"']
   & \Gr(T_+M\otimes_{\Qplus}M) \arrow[d,"{\cong}"]\\
 T_\QQ(\Gr M) \arrow[r,"{d_{\Gr M}}"']
   & T_\QQ(\Gr M)\otimes_\QQ\Gr M .
\end{tikzcd}
\]
The squares identify the derivatives under both realizations.
They therefore supply the compatibility needed to define the
derivative on a retained presentation.


Together these universal properties give a span preserving the
differential-algebra structure. The signed leg still acts on the exact
presentation, as required by Section~\ref{sec:obstructions}.


\subsubsection{The differential category of presentations}
\begin{theorem}[Codifferential structure on retained presentations]\label{thm:comma}
The category $\Pcat$ in~\eqref{eq:presentations} is a
commutative-monoid-enriched codifferential category with finite
biproducts. Its tensor and modality are constructed componentwise
using the comparisons of \cref{thm:modality-comparison}. Projection
$\pi_\C:\Pcat\to\C$ and signed realization
$\Gr\pi_\A:\Pcat\to\Dcat$ are strong monoidal and preserve the
algebra modality and deriving transformation.
\end{theorem}
\begin{proof}\leavevmode
\begin{enumerate}[leftmargin=*,align=left]
\item \textbf{Construct the additive monoidal structure.}
An arrow $(M,N,\rho)\to(M',N',\rho')$ is a pair $(f,g)$ with $g\rho=\rho'\Ext f$. Addition and zero are componentwise. Define
\[
 (M,N,\rho)\otimes(M',N',\rho')
 =\bigl(M\otimes M',N\otimes N',(\rho\otimes\rho')\phi^{-1}\bigr),
\]
where $\phi:\Ext M\otimes\Ext M'\cong\Ext(M\otimes M')$. The unit is $(\Qplus,S,\id)$ under the unit comparison. Coherence and bilinearity follow componentwise.

\item \textbf{Define the modality componentwise.}
Write $\kappa_M:\Ext T_+M\cong T\Ext M$. Set
\[
 T_\Pcat(M,N,\rho)=\bigl(T_+M,TN,T\rho\,\kappa_M\bigr).
\]
The modality unit, substitution, multiplication, and deriving map are
respectively the pairs of those maps on $\A$ and $\C$.

\item \textbf{Verify compatibility and the differential laws.}
Finite biproducts, their injections and projections, and the hom-monoid
zeros are componentwise because $\Ext$ preserves them; the defining
compatibility equations are stable under addition. Naturality and the
differential comparison show that the structural maps satisfy the
commuting-square condition for arrows of $\Pcat$. Their laws hold
componentwise. For example, the square for $(d_M,d_N)$ follows by first commuting $\Ext d_M$ with $\kappa_M$, then using naturality of $d$ along $\rho$. This proves the theorem.

\item \textbf{Prove freeness on retained presentations.}
In fact the modality is the free commutative algebra modality: the two free extensions of a compatible pair of maps remain compatible by uniqueness on generators.
\end{enumerate}
\end{proof}


For a presentation $(M,N,\rho)$, the structural map of its lifted
free algebra is the composite in
\[
\begin{tikzcd}[column sep=large]
 \Ext T_+M \arrow[r,"{\kappa_M}","{\cong}"']
   \arrow[rr,bend right=25,"{T\rho\,\kappa_M}"']
   & T\Ext M \arrow[r,"{T\rho}"] & TN .
\end{tikzcd}
\]
The comma-square condition for its derivative is
\[
\begin{tikzcd}[column sep=large,row sep=large]
 \Ext T_+M \arrow[r,"{T\rho\,\kappa_M}"]
   \arrow[d,"{\Ext d_M^+}"']
   & TN \arrow[d,"{d_N}"]\\
 \Ext(T_+M\otimes_{\Qplus}M)
   \arrow[r,"{{(T\rho\,\kappa_M)\otimes\rho}}"']
   & TN\otimes_SN .
\end{tikzcd}
\]
The lower label uses the strong monoidal comparison to identify
its source with $\Ext T_+M\otimes_S\Ext M$.
Thus the derivative is a morphism of retained presentations,
and both realization legs of \cref{eq:cat-realization-span}
preserve it.


\subsection{Finite-projective stages for resource geometry}
\label{sec:weil}
\subsubsection{Rational stages and finite-projective closure}
A rational Weil presentation is a finite-dimensional augmented commutative
$\QQ$-algebra $W=\QQ\oplus\mathfrak n$, specified by a rational basis,
a multiplication table, and $r$ with $\mathfrak n^{r+1}=0$.
Let $C_W=B\otimes_\QQ W$.  The \emph{rational-stage category}
$\Weil_B^{\mathrm{rat}}$ has these objects and all $B$-algebra maps
between them, encoded by finite lists of images satisfying their
multiplication equations.  Tensor product gives coproducts of these
stages, but the category with all $B$-algebra maps is not, in general,
closed under equalizers.

\begin{example}[A coefficient-dependent rational equalizer]
\label{ex:weilstage-rank}
In the cofinite numerical fibre, let $e$ be the parity idempotent,
which is neither zero nor one.  On $C=B[\eps]/\eps^2$ the maps
$u(\eps)=\eps$ and $v(\eps)=e\eps$ have algebraic equalizer
\[
 \operatorname{Eq}(u,v)=B\oplus Be\eps.
\]
Its nilpotent summand is rank one on the support of $e$ and rank zero
on the complementary support.  It is finite projective over $B$, but
is not $B\otimes_\QQ W$ for a fixed rational Weil algebra $W$:
on the two idempotent components the augmentation ranks of such a
rational extension would agree.  Thus the full rational-stage category
cannot serve as the stage category for all coefficient-dependent
finite Weil limits.
\end{example}

\begin{theorem}[Finite-projective Weil-stage closure]
\label{prop:projective-weil}
There is a small internal category $\Weil_B^{\mathrm{proj}}$ whose
objects are augmented commutative $B$-algebras
$C=B\oplus N$ with $N^{r+1}=0$ and a displayed finite-free
retract (idempotent-matrix presentation) of $N$. It contains
$\Weil_B^{\mathrm{rat}}$ as a full subcategory and is closed
under finite limits and Weil tensor coproducts. Each stage is
finite projective in this explicit sense and flat over $B$.
These presentations and closure operations are uniform in parameters.
\end{theorem}
\begin{proof}\leavevmode
\begin{enumerate}[leftmargin=*,align=left]
\item \textbf{Construct the internal category of Weil stages.}
A stage code consists of an integer $n$, an idempotent
$p\in M_n(B)$, a finite bilinear multiplication table on
$N=\operatorname{im}(p)$, and a positive integer $r$ witnessing
$N^{r+1}=0$. Commutativity and associativity are finite
matrix equations. A stage morphism is similarly a finite
matrix representing a multiplication-preserving map between
these augmentation modules. Matrix identity, substitution, and
composition construct the internal category. The validity conditions are matrix equations over $B$, interpreted as in \cref{sec:resource-observations}. The ring is
$B\oplus N$, with multiplication
$(b,x)(c,y)=(bc,by+cx+xy)$.

\item \textbf{Construct coproducts and finite products.}
The tensor coproduct of $B\oplus N$ and $B\oplus N'$ has
underlying module
\[
 B\oplus N\oplus N'\oplus(N\otimes_BN'),
\]
a finite-free retract. The augmentation ideal is nilpotent:
a sufficiently long product contains either more than $r$
factors from the first augmentation ideal or more than $r'$
from the second. The terminal stage is $B$. Products in the
augmented-algebra category are fibre products over $B$, with
augmentation ideal $N\oplus N'$.

\item \textbf{Split the image of the difference map.}
For equalizers take two stage maps $u,v:C\rightrightarrows C'$.
Every such map preserves augmentation: its composite with
$C'\to B$ kills the nilpotent ideal of $C$, because $B$ is
reduced. Thus $d=u-v$ restricts to a $B$-linear map
$d:N\to N'$. Its image $J\subseteq N'$ is generated by
images of the displayed finite generators of $N$. Embed $N'$
as a displayed summand of a finite free module. Apply
\cref{thm:finite-generator-splitting} to the corestriction
$d:N\twoheadrightarrow J$, with those images and their
specified preimages in $N$. It constructs a section
$s:J\to N$ satisfying $ds=\id_J$.

\item \textbf{Construct the equalizer idempotent.}
The endomorphism
\[
 p_N=\id_N-sd:N\longrightarrow N
\]
satisfies
\[
 p_N^2=\id_N-2sd+s(d s)d=p_N,\qquad
 d p_N=d-dsd=0.
\]
For $x\in\ker d$ one has $p_Nx=x$, so
$\operatorname{im}(p_N)=\ker d$. Therefore the algebraic
equalizer is
\[
 \operatorname{Eq}(u,v)=B\oplus\operatorname{im}(p_N),
\]
a displayed finite-free retract of $C$ as a module.

\item \textbf{Verify algebra closure and finite-limit stability.}
It is
closed under multiplication because $u$ and $v$ are algebra
maps; its augmentation ideal inherits nilpotence from $N$.
Finite limits follow from products and equalizers, with the
idempotent presentations retained. All calculations are finite,
compatible with parameter substitution, and use the displayed lifts.
\end{enumerate}
\end{proof}

For the remainder of the paper put
$\Weil_B=\Weil_B^{\mathrm{proj}}$ and
\[
 \Geo=[\Weil_B,\E].
\]
The functor topos is internal and covariant, equivalently the
presheaf topos on the opposite stage category; finite limits,
colimits, and images are computed stagewise.
Write
\[
 \Rline(C)=C,\quad
 j_BM(C)=C\otimes_BM,\quad
 D_A(C)=\Hom_{B\text{-alg}}(A,C)
 \qquad(A,C\in\Weil_B).
\]
The resource line is a ring object and $j_BM$ an $\Rline$-module.
Each $D_A$ is a representable infinitesimal object.

A budget family and an infinitesimal coefficient stage can vary
together. The next construction combines these two parameters, so
each budget carries its realized resource module at every Weil stage.

\subsection{Relative geometry of budgets and Weil coefficients}
\label{sec:relative-probes}
\begin{theorem}[Budget--Weil presentation]
\label{thm:budget-weil}
For the finite-matrix stage presentations constructed above,
\[
 \Weil_{\bar B}^{\mathrm{proj}}
       \cong\Delta_{\Cplx}\Weil_B^{\mathrm{proj}},\qquad
 \mathcal G_{\Cplx}:=[\Weil_{\bar B}^{\mathrm{proj}},\mathscr S]
       \simeq[L\times\Weil_B^{\mathrm{proj}},\E].
\]
At a probe $(a,A)$ the structure ring and universal resource are
\begin{equation}\label{eq:resource-probes}
 \mathcal R_{\bar B}(a,A)=A,\qquad
 \mathfrak M(a,A)=A\otimes_BM_a.
\end{equation}
For $a\le b$ and $u:A\to A'$, the transition sends
$c\otimes m$ to $u(c)\otimes m$ using $M_a\incl M_b$.
\end{theorem}
\begin{proof}\leavevmode
\begin{enumerate}[leftmargin=*,align=left]
\item \textbf{Transport the finite stage presentations.}
A stage code comprises natural numbers, finite coefficient lists,
idempotent and multiplication matrices, and finite equations.
The nilpotence bound is a displayed natural number and its
equations are checked on the finite generators. The object of
codes and the object of arrows are therefore built from finite
limits and natural-number-indexed finite lists of coefficients.
The constant-diagram functor preserves these constructions;
in particular it preserves their equality subobjects. This identifies both
stage categories, including identities and composition.

\item \textbf{Identify budget diagrams of Weil diagrams.}
An internal diagram over the resulting constant category is
equivalently a budget diagram of Weil diagrams: the two actions
commute. Currying gives the asserted equivalence.

\item \textbf{Compute the probe formulas.}
Tensor and
module operations are computed at each probe, giving
\eqref{eq:resource-probes}.
\end{enumerate}
\end{proof}

Write $i_W:\mathscr S\to\mathcal G_{\Cplx}$ for constancy
in the Weil variable and $\Gamma_W Z(a)=Z(a,B)$. The restricted
based-Weil adjunctions are
$i_W\dashv\Gamma_W\dashv i_W\dashv\Gamma_W$.
They calculate reduction on fixed-base probes at each budget.
The finite-hit dominance of $\mathscr S$ is $\Delta_{\Cplx}\Sigma$:
constant diagrams preserve $2^{\N}$ and its finite-hit image,
and composition of the classified monos is checked at each
budget. Applying the same construction in the Weil variable
gives $\Sigma_{\mathcal G_{\Cplx}}=i_W\Delta_{\Cplx}\Sigma$.
For every positive test $u:Z\to\Sigma_{\mathcal G_{\Cplx}}$,
\begin{equation}\label{eq:relative-positive-reduction}
 u=i_W(\Gamma_Wu)\eta_Z.
\end{equation}
This follows from $\Gamma_W\dashv i_W$ exactly as in
\cref{prop:geometric-dominance}. Budget locality is retained;
infinitesimal coefficients are invisible to these positive tests.
For a name cover $q_D:C_D\onto\Pt D$ and an implemented map
$\widehat F:C_D\to C_{D'}$, let $k=i_W\Delta_{\Cplx}$.
The transported cover remains a regular epimorphism and
\[
 (kq_{D'})(k\widehat F)=k(\Pt F)(kq_D),\qquad
 k\operatorname{Obs}_D:k\Pt D\times k\OO_D\to k\Sigma.
\]
These equations are evaluated at each probe as the original
name-realization and observation equations; pointwise finite limits
and images preserve their types. For a partial implementation the
same argument preserves its $\Sigma$-classified domain span and
its value map. Thus the relative geometry retains both domains
and returned values.
\subsection{Reduction, observations, and Weil coefficient calculus}
\label{sec:cohesion}
\subsubsection{Reduction and infinitesimal shape}
Every $B$-algebra map $C=B\oplus N\to B$ kills the nilpotent
augmentation ideal because $B$ is reduced.  It therefore equals
the augmentation.  Thus $B$ is terminal as well as initial
in the full projective Weil-stage category $\Weil_B$.
Write $i:\E\to\Geo$ for constant objects and $\Gamma(F)=F(B)$.

\begin{theorem}[The ambidextrous based-Weil restriction]\label{thm:cohesion}
On the based-Weil functor topos, the functor $i$ is fully faithful
and there is an ambidextrous adjoint quadruple
\[
 i\dashv\Gamma\dashv i\dashv\Gamma.
\]
The endofunctor $\Shape=i\Gamma$ is left exact and idempotent,
with natural maps
\[
 \epsilon_X:\Shape X\to X,\qquad
 \eta_X:X\to\Shape X,\qquad
 \eta_X\epsilon_X=\id_{\Shape X}.
\]
It carries the induced reduction comonad and infinitesimal-shape
monad structures. The associated triple of adjoint idempotent
endofunctors has coincident underlying functor $\Shape$.
\end{theorem}
\begin{proof}\leavevmode
\begin{enumerate}[leftmargin=*,align=left]
\item \textbf{Construct the adjunction from initiality.}
Initiality of $B$ gives
\[
 \operatorname{Nat}(iA,F)\cong\Hom_{\E}(A,F(B)):
\]
a map at $B$ extends at $C$ through $F(B\to C)$.

\item \textbf{Construct the adjunction from terminality.}
Terminality gives
\[
 \operatorname{Nat}(F,iA)\cong\Hom_{\E}(F(B),A):
\]
a map at $B$ extends at $C$ through $F(C\to B)$.

\item \textbf{Verify the adjunction quadruple and full faithfulness.}
Both inverse equations are naturality along those stage maps.
These two adjunctions repeat to give the quadruple.
Since $\Gamma i=\id$, $i$ is fully faithful.

\item \textbf{Identify the units, counits, and modal structures.}
At $C$ the maps $\epsilon_X,\eta_X$ are respectively
$X(B\to C)$ and $X(C\to B)$. Their composite on $X(B)$
is the identity. Evaluation and constant formation preserve
finite limits, giving left exactness; $\Gamma i=\id$ gives
idempotence. The units, counits, and their triangle identities
give the stated monad, comonad, and adjoint structures.
\end{enumerate}
\end{proof}

The ambidexterity here has a specific source: both endpoints of a
based probe reduce to the fixed $B$-point. It describes the based-Weil
restriction of the primary \'etale endpoint geometry of
\cref{thm:etale-endpoint-cohesion}. The comparison in
\cref{thm:completion-etale} places the following Weil coefficient
and microlinearity calculations within that larger geometry.
For the line,
\[
 \Shape\Rline=iB,\qquad
 \eta_{\Rline}(C):C\to B
\]
is augmentation; for every representable probe, $\Shape D_A=iD_A(B)=1$.
Thus reduction extracts the reduced point of each infinitesimal probe.

The reduction $B\to B/\mathfrak m$ has a different type:
it removes numerical errors vanishing along the eventual
filter. The map $\eta_X$ removes formal nilpotent directions.
Both are used with their displayed coefficient objects.

\subsubsection{Computational observations inside the cohesive geometry}
\label{sec:geometric-effectivity}
The functor $i:\E\to\Geo$ preserves finite limits, images, and the NNO
stagewise. Between constant objects it preserves exponentials:
\[
 (iY)^{iX}\cong i(Y^X).
\]
Indeed a natural map $h_C\times iX\to iY$ is determined by its value
on $(\id_C,x)$; naturality along $u:C\to C'$ forces the value on
$(u,x)$ to be the same. This proves the isomorphism in every parameter
context.

\begin{proposition}[Positive observations factor through reduction]
\label{prop:geometric-dominance}
The canonical finite-hit Rosolini truth-value object of $\Geo$ is
\[
 \Sigma_{\Geo}\cong i\Sigma.
\]
It is a dominance. Every positive computational test has the unique
factorization
\[
 u:Z\to\Sigma_{\Geo},\qquad
 u=i(\Gamma u)\eta_Z:Z\to i\Gamma Z\to i\Sigma.
\]
Consequently its pullbacks along the two projections of
$Z\times_{\Shape Z}Z$ are equal.
\end{proposition}
\begin{proof}\leavevmode
\begin{enumerate}[leftmargin=*,align=left]
\item \textbf{Construct the geometric positive truth-value object.}
Constant-subobject classification gives a monomorphism
$i\Omega_{\E}\to\Omega_{\Geo}$: at each stage, evaluation of the
classified natural family on its identity arrow recovers the original
truth value. The preceding exponential calculation gives
$(i2)^{i\N}\cong i(2^\N)$. The finite-hit predicate is the same
natural-number existential at every stage; its classifier is the
constant finite-hit classifier followed by this monomorphism. Images
are stagewise, so its image is $i\Sigma$.

\item \textbf{Factor positive observations through reduction.}
The adjunction $\Gamma\dashv i$ identifies a map $Z\to i\Sigma$ with
its component $\Gamma Z\to\Sigma$ and gives the displayed factorization.
A subobject so classified is the pullback along $\eta_Z$ of its reduced
classified subobject.

\item \textbf{Verify dominance and formal-disk invariance.}
For two composable such subobjects, left exactness
of $\Gamma$ identifies the reduced composite; dominance in $\E$
classifies it, and pullback along $\eta_Z$ recovers the original
composite. Truth and pullback stability follow by the same argument.
Thus $i\Sigma$ is a dominance. The two formal-disk projections have
the same composite with $\eta_Z$, proving the last assertion.
\end{enumerate}
\end{proof}

For $j_BM$, reduction is augmentation
$C\otimes_BM\to M$, since $\Gamma j_BM=M$.
Canonical positive tests therefore observe the reduced value and
are constant along its formal disk. Infinitesimal coefficients
are instead read by the algebraic coefficient calculations below.


Computational names, natural numbers, and denotations are constant
objects of this geometry. The coefficient calculation in
\cref{thm:weil} also expresses their infinitesimal discreteness as
$(iX)^{D_A}\cong iX$. Resource values lie in the nonconstant modules
$j_BM$. These two roles allow an implementation witness to stay
fixed while its resource constraint varies infinitesimally.

\subsubsection{Weil coefficient identity and comparison of stages}
\begin{theorem}[Exact realization and Weil coefficients]
\label{thm:weil}
The functor $j_B:\Mod_B(\E)\to\Mod_{\Rline}(\Geo)$ is exact,
fully faithful, strong symmetric monoidal, and reflects exactness
and isomorphisms.  For every $F\in\Geo$ and projective stage $A$,
\[
 F^{D_A}(C)\cong F(C\otimes_BA),
\]
and therefore
\begin{equation}\label{eq:weil-coefficients}
 (j_BM)^{D_A}(C)
 \cong (C\otimes_BA)\otimes_BM
 \cong (C\otimes_BM)\otimes_BA.
\end{equation}
All comparisons are natural in the stages, their maps, and the
module, and commute with augmentation and Weil substitution.
\end{theorem}
\begin{proof}\leavevmode
\begin{enumerate}[leftmargin=*,align=left]
\item \textbf{Compute exponentials by Weil disks.}
Writing $h_C=\Hom_{B\text{-alg}}(C,-)$, tensor coproduct and
Yoneda give
\[
 h_C\times D_A\cong h_{C\otimes_BA},\qquad
 F^{D_A}(C)=\operatorname{Nat}(h_C\times D_A,F)
 \cong F(C\otimes_BA).
\]
This also proves the displayed coefficient formula.

\item \textbf{Prove exactness of realization.}
Every stage is finite projective and hence flat, so each
$C\otimes_B-$ is exact, making $j_B$ exact stagewise.

\item \textbf{Prove full faithfulness and reflection.}
A natural $\Rline$-linear map $j_BM\to j_BN$ is determined by
its component at $B$: naturality along $B\to C$ determines
its values on $1\otimes m$, and $C$-linearity determines
all $c\otimes m$.  Conversely any $B$-linear map gives such
a natural transformation by scalar extension, proving full
faithfulness.  Evaluation at $B$ reflects exactness and
isomorphisms.

\item \textbf{Verify strong monoidality and coherence.}
The associativity of balanced tensor gives
$j_B(M\otimes_BN)\cong j_BM\otimes_{\Rline}j_BN$ and
$j_BB\cong\Rline$; uniqueness on elementary tensors proves
unit, symmetry, and associativity coherence.
\end{enumerate}
\end{proof}

\begin{proposition}[Restriction to rational Weil probes]
\label{prop:rational-restriction}
The full inclusion
$k:\Weil_B^{\mathrm{rat}}\hookrightarrow\Weil_B^{\mathrm{proj}}$
induces restriction
$k^*:\Geo\to[\Weil_B^{\mathrm{rat}},\E]$ with coherent
isomorphisms
\[
 k^*j_B^{\mathrm{proj}}M\cong j_B^{\mathrm{rat}}M,
 \qquad k^*D_{C_W}^{\mathrm{proj}}\cong D_{C_W}^{\mathrm{rat}},
 \qquad k^*i^{\mathrm{proj}}X\cong i^{\mathrm{rat}}X.
\]
It preserves the rational Weil coefficient and exponential
identities and commutes with evaluation at $B$ and the
resulting infinitesimal-shape operation.
\end{proposition}
\begin{proof}\leavevmode
\begin{enumerate}[leftmargin=*,align=left]
\item \textbf{Compare restriction and Weil exponentiation.}
All three functors have the displayed stagewise formulas.
For rational stages $C_W,C_V$, their tensor coproduct
$C_{W\otimes V}$ is again rational; applying the exponential
calculation of \cref{thm:weil} before or after restriction
therefore gives the same expression
$F(C_W\otimes_BC_V)$.

\item \textbf{Compare reduction and verify coherence.}
Evaluation at $B$ commutes with
restriction, and constant functors do so stagewise.
These identities give the stated shape comparison, while equality
of the elementary-tensor formulas gives all coherence diagrams.
\end{enumerate}
\end{proof}

\subsubsection{Finite Kock--Lawvere identities and microlinearity}
For $A=B[\eps]/\eps^2$, put $D=D_A$.  Yoneda identifies
$D\subseteq\Rline$ with the square-zero elements.
Formula~\eqref{eq:weil-coefficients} gives
\[
 \Rline^D\cong\Rline\times\Rline,\qquad
 (j_BM)^D\cong j_BM\times j_BM,
\]
where $(a,b)$ corresponds to the map $d\mapsto a+db$.
The same formula for every finite-projective probe is the finite
Weil--Kock--Lawvere coefficient identity \citep{Kock}.

\begin{theorem}[Rational and projective microlinearity]
\label{thm:microlinearity}
For every $B$-module $M$, the augmented coefficient functor
\[
 W\longmapsto\bigl((j_BM)^{D_{C_W}}\to j_BM\bigr)
\]
preserves finite limits of rational Weil diagrams with rational
structure maps.  For every embedded resource module $M\subseteq K$,
the functor
\[
 A\longmapsto\bigl((j_BM)^{D_A}\to j_BM\bigr)
\]
preserves finite limits of the full finite-projective Weil-stage
category.  In both statements the codomain is $\Geo/j_BM$ and
the displayed maps are induced by augmentation.
\end{theorem}
\begin{proof}\leavevmode
\begin{enumerate}[leftmargin=*,align=left]
\item \textbf{Preserve rational Weil limits by split exactness.}
Rational Weil algebras with rational algebra homomorphisms have
finite limits: the product over $\QQ$ is an augmented fibre
product, and the equalizer is a rational linear kernel.  The
latter has a basis constructed by Gaussian elimination.  Its
finite exact sequences split over $\QQ$, so tensoring with any
$\QQ$-module preserves these limits.  At a stage $C$ the
coefficient functor is
$W\mapsto(C\otimes_BM)\otimes_\QQ W$ over $C\otimes_BM$;
this proves the first assertion stagewise.

\item \textbf{Establish flatness of resource coefficients.}
For the second, \cref{thm:signed-rep} writes each embedded resource
as an internally directed colimit of finite-projective principal
modules.  Its internal filtered-colimit flatness proof in
\cref{sec:resources} makes $M$ flat.  Each stage $C$ is
finite projective and flat, so $T=C\otimes_BM$ is flat over $B$.

\item \textbf{Preserve projective Weil limits stagewise.}
The formula at $C$ is now $A\mapsto T\otimes_BA$ with
augmentation to $T$.  Finite limits of projective augmented
Weil algebras are computed on the underlying $B$-modules by
fibre products and kernels of differences.  Flatness preserves
these finite module limits, including the augmentation maps.
The resulting comparisons are natural at stages and therefore
are isomorphisms in $\Geo/j_BM$.
\end{enumerate}
\end{proof}

\begin{proposition}[Failure for arbitrary coefficient modules]
\label{prop:nonflat-weil}
For the numerical cofinite ring $B$ and regular nonunit
$t(n)=1/(n+1)$, the augmented coefficient functor of
\cref{thm:microlinearity} need not preserve projective Weil
equalizers for $M=B/tB$.
\end{proposition}
\begin{proof}\leavevmode
\begin{enumerate}[leftmargin=*,align=left]
\item \textbf{Compute the original equalizer.}
On $A=B[\eps]/\eps^2$, take
$u(\eps)=t\eps$ and $v(\eps)=0$.
Regularity of $t$ gives $\operatorname{Eq}(u,v)=B$.

\item \textbf{Exhibit failure of the equalizer comparison.}
After tensoring with $M=B/tB$, the two maps coincide and their
equalizer becomes $M\oplus M\eps$, whereas tensoring the
original equalizer gives only $M$.  The two are not isomorphic
as the indicated augmented equalizer comparison: its image is
the constant-coefficient summand and misses $\eps$.
The embedded resources of \cref{thm:microlinearity} do not have
this obstruction because they are flat.
\end{enumerate}
\end{proof}

At $C=B[\eps]/\eps^2$ the generic infinitesimal $\eps$ is
nonzero: coefficient extraction would otherwise force $1=0$
in $B$, contradicting properness of the filter.  Thus the
constructed model has nontrivial infinitesimal directions.

\subsubsection{Coefficient change and reduced-base cohesion}
Finite Weil coefficient realization commutes with a ring map
$\beta:B\to B'$ on rational probes: for any rational Weil $W$,
\[
 (B'\otimes_\QQ W)\otimes_{B'}(B'\otimes_BM)
 \cong B'\otimes_B((B\otimes_\QQ W)\otimes_BM).
\]
Both sides send an elementary coefficient--Weil--module tensor
to the same expression; the unit and substitution diagrams follow
from tensor coherence.  The adjoint quadruple of
\cref{thm:cohesion} has the further property that the coefficient
stage is terminal, which is a consequence of reducedness rather
than of tensor coherence alone.

\begin{proposition}[A coefficient reduction producing nilpotents]
\label{prop:nonreduced-reduction}
In the cofinite numerical fibre, let $t(n)=1/(n+1)$ and
$z(n)=1/(1+\lfloor\sqrt{n+1}\rfloor)$.  Then
$z\notin tB$ but $z^2\in tB$.  Thus $B'=B/tB$ has a nonzero
square-zero element, and its coefficient stage is not terminal
in the category of $B'$-Weil stages.
\end{proposition}
\begin{proof}\leavevmode
\begin{enumerate}[leftmargin=*,align=left]
\item \textbf{Construct a nonzero square-zero coefficient.}
Write $m=\lfloor\sqrt{n+1}\rfloor$.  The inequalities
\[
 \frac{z(n)^2}{t(n)}=\frac{n+1}{(1+m)^2}\le1,
 \qquad
 \frac{z(n)}{t(n)}=\frac{n+1}{1+m}\ge\frac{m^2}{1+m}
\]
show that $z^2/t$ is bounded but $z/t$ is unbounded along
square inputs.  Hence the class $\bar z=z+tB$ is nonzero and
satisfies $\bar z^{\,2}=0$.

\item \textbf{Show failure of terminality.}
The two $B'$-algebra maps
\[
 B'[\eps]/\eps^2\rightrightarrows B',
 \qquad \eps\longmapsto0,\quad\eps\longmapsto\bar z
\]
are distinct, so $B'$ is not terminal among those stages.
This explicitly separates coefficient-tensor comparisons from
the reduced-base cohesion calculation.
\end{enumerate}
\end{proof}

\subsubsection{Explicit linear adjoints}
Let $\mathscr M_B=\Mod_{\Rline}(\Geo)$ and
$\Gamlin(F)=F(B)$ with its $B$-module structure.

\begin{theorem}[Linear realization adjoints]\label{thm:linear-adjoints}
There is a constructed adjoint triple
\[
 \Lamlin\dashv j_B\dashv\Gamlin,\qquad
 \Lamlin j_B\cong\id\cong\Gamlin j_B.
\]
\end{theorem}
\begin{proof}\leavevmode
\begin{enumerate}[leftmargin=*,align=left]
\item \textbf{Construct the right-adjoint correspondence.}
A $B$-linear $v:M\to F(B)$ extends at $C$ by
$c\otimes m\mapsto cF(B\to C)(v(m))$.
Naturality and linearity prove it is inverse to evaluation,
giving $j_B\dashv\Gamlin$.

\item \textbf{Construct the coend candidate.}
For the left adjoint put $C^\vee=\Hom_B(C,B)$ with
right $C$-action $(\phi c)(x)=\phi(cx)$, and construct
\begin{equation}\label{eq:linear-coend}
 \Lamlin F=
 \coeq\left(
 \bigoplus_{u:C\to C'}(C')^\vee\otimes_C F(C)
 \rightrightarrows
 \bigoplus_C C^\vee\otimes_C F(C)\right).
\end{equation}

\item \textbf{Define the balanced coequalizer maps.}
One map sends $\phi\otimes x$ to $(\phi u)\otimes x$
in the $C$-summand; the other to $\phi\otimes F(u)x$
in the $C'$-summand. The right $C$-action on $(C')^\vee$
uses $u$, so both maps are balanced. The indexing objects
are the object and arrow objects of the internal stage
category; internal free modules and quotients construct
these sums and this coequalizer.

\item \textbf{Establish the stagewise duality correspondence.}
Finite projectivity gives
\[
 \Hom_C(F(C),C\otimes_BM)
 \cong
 \Hom_B(C^\vee\otimes_CF(C),M).
\]
It sends $v$ to
$\phi\otimes x\mapsto(\phi\otimes1)v(x)$.
For a dual basis $(c_i,\phi_i)$ the inverse sends $h$
to $x\mapsto\sum_i c_i\otimes h(\phi_i\otimes x)$.
The dual-basis identities and balancing prove the inverse
equations and $C$-linearity.

\item \textbf{Prove the left adjunction and composite identities.}
Naturality of the family $v$ is exactly equalization of
the two maps in \eqref{eq:linear-coend}.
The coequalizer universal property proves
$\Lamlin\dashv j_B$. Full faithfulness of $j_B$ gives
the two composite identities.
\end{enumerate}
\end{proof}

The coefficient model now lets us express resource operations as
geometric maps. We first identify the coordinate algebra of a
finite-projective module. Its dual appears because coordinates are
linear functions on that module. Universal differentials then
describe the first-order change of those functions.

\subsection{Affine resource coordinates and universal differentials}
\label{sec:geometry}

\subsubsection{Dual coordinate modules}
\begin{proposition}[Finite-cell differential comparison]\label{prop:dual-coordinates}
Let $M$ be a finite projective $B$-module and set $M^\vee=\Hom_B(M,B)$. Its affine coordinate algebra is $\Sym_B(M^\vee)$, and for every $B$-algebra $C$,
\[
 \Hom_{B\text{-alg}}(\Sym_B(M^\vee),C)
 \cong\Hom_B(M^\vee,C)\cong C\otimes_BM.
\]
Under this identification, the universal algebraic derivative gives the coefficient of $\eps$ in evaluation on $C[\eps]/\eps^2$.
\end{proposition}
\begin{proof}\leavevmode
\begin{enumerate}[leftmargin=*,align=left]
\item \textbf{Establish the finite-duality bijections.}
The first bijection is the free-algebra universal property. For a finite dual basis $(m_i,\lambda_i)$, the inverse of the second sends a functional $f:M^\vee\to C$ to $\sum_i f(\lambda_i)\otimes m_i$; the dual-basis identity proves both inverse equations and independence of the presentation.

\item \textbf{Extract the dual-number derivative.}
A polynomial coordinate $p$ satisfies
\[
 p(x+\eps v)=p(x)+\eps\,dp_x(v)
\]
because multiplying factors chooses either no $\eps$ or exactly one $\eps$. This is precisely formula~\eqref{eq:monomial-d}. Addition and relations preserve the equality. Thus the induced map of K\"ahler modules, transposed to tangent vectors, agrees with the dual-number coefficient.
\end{enumerate}
\end{proof}

Positive principal resources $B^+g\cong B^+e_g$ are dualizable, with dual $B^+e_g$. Strong monoidal functors preserve their evaluation and coevaluation maps and hence their duals. Theorem~\ref{thm:modality-comparison} therefore compares these coordinate algebras and their derivatives before and after signed realization.

For a general resource $M=\bigcup_g Bg$, a finite tuple of inputs lies in a common finite projective cell. A finite polynomial resource operation lands in the finite principal cell generated by the sum of the absolute values of its finitely many coefficient-budget products. The comparison above applies on each such cell. Cell inclusions commute with evaluation and derivation by naturality, so these comparisons assemble over the directed resource presentation. This proves the derivative comparison for resource arithmetic and polynomial overhead cells on this directed presentation.


\begin{theorem}[Canonical polynomial lifting to Weil geometry]
\label{thm:polynomial-lift}
Let $E,N$ be finite-projective $B$-modules.  A polynomial resource
map $p:E\to N$, specified by its coordinate-algebra map
\[
 p^*:\Sym_B(N^\vee)\longrightarrow\Sym_B(E^\vee),
\]
induces canonically a geometric map
$\widehat p:j_BE\to j_BN$.  It preserves identities and
polynomial composition, and
\[
 D\widehat p(x,v)=dp_x(v).
\]
These maps and derivatives commute with finite-projective resource
cell inclusions whenever the polynomial expressions and their
coefficient budgets are transported by those inclusions.
\end{theorem}
\begin{proof}\leavevmode
\begin{enumerate}[leftmargin=*,align=left]
\item \textbf{Construct the geometric polynomial map.}
For each Weil stage $C$, finite duality identifies
$C\otimes_BE$ with the $C$-points of
$\Sym_B(E^\vee)$, and similarly for $N$.
Precomposition with $p^*$ therefore gives
$C\otimes_BE\to C\otimes_BN$, naturally in every stage map.
This is the required geometric transformation.

\item \textbf{Verify functoriality.}
Algebra-map
composition establishes functoriality and identity preservation.

\item \textbf{Identify the derivative by finite duality.}
For the dual-number coefficient $x+\eps v$, a coordinate
$\lambda\in E^\vee$ evaluates to
$\lambda(x)+\eps\lambda(v)$.  Expanding the finite polynomial
$p^*(\mu)$ for each $\mu\in N^\vee$ extracts the term containing
exactly one $\eps$, which is $dp_x(v)$ by the universal
Leibniz formula.  Finite duality detects equality of the resulting
$N$-valued coefficients, proving the derivative identity.

\item \textbf{Extend compatibility over principal cells.}
Every finite polynomial expression in an embedded resource uses
finitely many arguments and coefficient budgets.  Their absolute
sum supplies a common finite-projective principal cell by
\cref{thm:signed-rep}; naturality of polynomial evaluation and
the differential comparison proves compatibility over directed
cell presentations.
\end{enumerate}
\end{proof}

\begin{proposition}[Finite polynomial laws on arbitrary modules]
\label{prop:polynomial-laws}
For $B$-modules $H,N$, a finite family of $B$-multilinear maps
$m_r:H^{\otimes_B r}\to N$ for $0\le r\le d$ (with $m_0\in N$)
constructs the polynomial law
$p(x)=\sum_{r=0}^d m_r(x,\ldots,x)$ and a canonical natural
geometric lift $\widehat p:j_BH\to j_BN$, satisfying
\[
 D\widehat p(x,v)=
 \sum_{r=1}^d\sum_{j=1}^r
 \widetilde m_r(x,\ldots,x,\underset{j}{v},x,\ldots,x),
\]
where $\widetilde m_r$ is the scalar-extended multilinear map at
the generalized Weil stage.  On finite-projective modules this
agrees with the coordinate-algebra
construction of \cref{thm:polynomial-lift}.
\end{proposition}
\begin{proof}\leavevmode
\begin{enumerate}[leftmargin=*,align=left]
\item \textbf{Extend the multilinear coefficients to Weil stages.}
For each Weil stage $C$, balanced tensor extension of $m_r$
produces a $C$-multilinear map
$(C\otimes_BH)^{\otimes_Cr}\to C\otimes_BN$.
Evaluating the finite sum diagonally defines
$\widehat p_C:C\otimes_BH\to C\otimes_BN$.

\item \textbf{Verify naturality and extract the derivative.}
Tensor functoriality makes these maps natural in every
$B$-algebra stage morphism.  Evaluating at $x+\eps v$,
the square-zero multiplication removes all terms with two or
more occurrences of $\eps$.  The terms with exactly one
occurrence give the displayed formula, by coefficient uniqueness.

\item \textbf{Compare with the finite-duality construction.}
For finite projective $H,N$, both stagewise constructions are
polynomial evaluations with the same coordinate values, so
finite duality identifies them.
\end{enumerate}
\end{proof}

An arbitrary numerical value family also has a canonical, but
different, realization.  The constant-object/evaluation adjunction
of \cref{thm:cohesion} gives
\[
 \Geo(iX,j_BK)\cong\E(X,K),
 \qquad f\longmapsto(x\mapsto1\otimes f(x)).
\]
The source $iX$ is infinitesimally discrete:
$(iX)^D\cong iX$. This realization therefore preserves the values
of an arbitrary numerical family without differentiating its discrete
input. Polynomial resource variables have a different role:
the lifting theorem realizes their polynomial maps geometrically
and identifies their derivatives.

\subsubsection{Relative differential modules}
For a commutative $B$-algebra $A$, form the free $A$-module
on symbols $da$ and impose additivity,
$d(ba)=b\,da$ for $b\in B$, and
$d(aa')=a\,da'+a'\,da$.
The quotient $\Omega^1_{A/B}$ represents $B$-derivations
by direct evaluation of these generators and relations.
Equivalently a derivation is a section of
$A\ltimes N\to A$ as a $B$-algebra map.
For a finite projective coordinate cell this gives
\[
 \Omega^1_{\Sym_B(M^\vee)/B}
 \cong \Sym_B(M^\vee)\otimes_BM^\vee.
\]
Dual bases identify its evaluation pairing with the tangent
coefficient of \cref{prop:dual-coordinates}.

The absolute module $\Omega^1_{B/\QQ}$ is obtained by fixing
only the rational constants. The comparison
\eqref{eq:kahler-comparison} identifies it with signed
realization of $\Omega^1_{B^+/\Qplus}$.
These constructions retain their coefficient maps.
For instance every relative derivation over $B$ kills a
scale $t\in B$, whereas its absolute differential $dt$ is
the specified generator in $\Omega^1_{B/\QQ}$.

Localization is compatible with relative differentials:
\[
 \Omega^1_{A_s/B}\cong A_s\otimes_A\Omega^1_{A/B}.
\]
To prove this, a derivation extends uniquely by
$d(s^{-1})=-s^{-2}ds$ and the quotient rule
$d(a/s^r)=s^{-r}da-r a s^{-r-1}ds$.
Clearing denominators verifies independence of the fraction,
and the Leibniz law verifies the extension.
Representability proves the isomorphism.
The gluing equalizer of \cref{prop:descent} therefore
also glues these differential modules and their maps.


\subsection{Finite tangent calculus}
\label{sec:finite-tangents}

\begin{proposition}[Differentiation in the realized geometry]\label{prop:geometric-d}
Every map $F:j_BM\to j_BN$ in $\Geo$ has a unique first-order expansion
\[
 F(x+\eps v)=F(x)+\eps\,DF(x,v).
\]
Its derivative is linear in $v$, obeys the chain rule, and agrees with Proposition~\ref{prop:dual-coordinates} on polynomial resource cells.
\end{proposition}
\begin{proof}\leavevmode
\begin{enumerate}[leftmargin=*,align=left]
\item \textbf{Obtain existence and uniqueness of the derivative.}
The dual-number coefficient identity gives existence and uniqueness.

\item \textbf{Prove additivity in the tangent direction.}
For addition in $v$, use the Weil algebra with two generators $\eps_1,\eps_2$ and all their pairwise products zero. Apply $F$ to $x+\eps_1v+\eps_2w$. Projections to each dual-number axis identify its two linear coefficients as $DF(x,v)$ and $DF(x,w)$. The algebra map sending both $\eps_1$ and $\eps_2$ to $\eps$ then proves additivity.

\item \textbf{Prove scalar linearity.}
The stage map $\eps\mapsto a\eps$ proves scalar linearity, including generalized coefficients $a\in C$.

\item \textbf{Prove the chain rule.}
Composing two expansions proves the chain rule by uniqueness of coefficients.

\item \textbf{Compare algebraic and geometric derivatives.}
The finite polynomial calculation in Proposition~\ref{prop:dual-coordinates} gives the comparison.

\item \textbf{Obtain higher symmetry and Taylor expansions.}
Iterated probes and permutation of their generators give the higher symmetry identities; truncated probes give finite Taylor expansions with the rational factorials available.
\end{enumerate}
\end{proof}

For a linear map $L$, its coefficient formula is $DL(x,v)=L(v)$.
For addition and multiplication in the resource line it gives
\[
 D(+)((x,y),(v,w))=v+w,\qquad
 D(\cdot)((x,y),(v,w))=yv+xw.
\]
For a composite the chain rule is
$D(GF)(x,v)=DG(Fx,DF(x,v))$.
The tangent functor is $T_{\mathrm{tan}}X=X^D$ with
projection induced by $0:1\to D$. On realized modules it
is the usual pair of base point and direction. Its zero
section, addition of directions, and scalar action are
the coefficient maps just constructed.

Iterated tangent uses
$\QQ[\eps_1,\eps_2]/(\eps_1^2,\eps_2^2)$.
Its coefficient of $\eps_1\eps_2$ is invariant under
interchanging the two generators. This proves symmetry
of the second base-point derivative. Iteration and
adjacent transpositions prove symmetry of every higher
base-point derivative, with directions held fixed.
These are equations of natural maps, so include
generalized coefficients at every stage.

We can now apply the differential calculus to resource relations.
The algebraic descriptions of $O$, $o$, $\Theta$, and $\sim$ use
multiplier witnesses. Realizing those multiplication maps first
retains the witnesses needed to differentiate the constraints.

\subsection{Realization of the Landau incidences}
Multiplication makes $j_BK$ an $\Rline$-algebra.
Write $U=\Rline^\times$ and $\mathcal M=j_B\mathfrak m$.
The subobject $1+\mathcal M$ is a subgroup of $U$.
To see this at $C=B\oplus N$, augment
$a\in C\otimes_B\mathfrak m$ to $a_0\in\mathfrak m$.
Then $1+a_0$ is a unit of $B$, while $a-a_0$ lies in the
nilpotent augmentation ideal. A finite geometric inverse
makes $1+a$ a unit of $C$, and
$(1+a)^{-1}-1=-a(1+a)^{-1}$ lies in the same extended ideal.

Define geometric binary relations as the images of
\[
\begin{array}{rcl}
 \Rline\times j_BK&\longrightarrow&j_BK\times j_BK,
       \quad (b,g)\longmapsto(bg,g),\\
 \mathcal M\times j_BK&\longrightarrow&j_BK\times j_BK,
       \quad (a,g)\longmapsto(ag,g),\\
 U\times j_BK&\longrightarrow&j_BK\times j_BK,
       \quad (u,g)\longmapsto(ug,g),\\
 \mathcal M\times j_BK&\longrightarrow&j_BK\times j_BK,
       \quad (a,g)\longmapsto((1+a)g,g).
\end{array}
\]
They realize respectively $O,o,\Theta,\sim$.
Reversing the first two gives the realizations of
$\Omega,\omega$. These formulas specify the geometric
relations also on nonconstant points.

\begin{theorem}[Exact Landau incidence]\label{thm:all-incidence}
Pullback of each of these six geometric relations along
$i(K\times K)\to j_BK\times j_BK$ is the corresponding
constant exact Landau relation. This statement is uniform
in exact computational parameters.
\end{theorem}
\begin{proof}\leavevmode
\begin{enumerate}[leftmargin=*,align=left]
\item \textbf{Extend exact witnesses to all Weil stages.}
An exact multiplier, unit, or small multiplier extends to
every stage and witnesses the forward inclusion by
\cref{thm:orbits}.

\item \textbf{Reflect membership through augmentation.}
Conversely, augment a stagewise witness
to $B$. Augmentation takes an arbitrary scalar to a bounded
coefficient, a unit to a unit, and an element of
$C\otimes_B\mathfrak m$ to an element of $\mathfrak m$.
Its equation becomes the exact equation on $f,g$.

\item \textbf{Identify incidence and the reversed comparisons.}
Thus membership is reflected. Images and the resulting pullbacks are computed stagewise. Reversal gives the two lower-bound cases.
\end{enumerate}
\end{proof}

Differentiating a witnessed incidence varies both its value and
its multiplier. The resulting equations describe how the resource
relation is preserved to first order; taking its image recovers
the corresponding relation on values.

\subsection{Differentials of the Landau incidences}
\begin{proposition}[Tangents of multiplier orbits]
\label{prop:orbit-tangents}
For an exact reference germ $g$, let $\widehat\Theta_g$ be the image
of $U=\Rline^\times\to j_BK$, $u\mapsto ug$, and put
$\widehat{\sim}_g=g+j_B(\mathfrak m g)$. At the common point $g$,
\begin{equation}\label{eq:orbit-tangents}
 T_g\widehat\Theta_g\cong j_B(Bg),\qquad
 T_g\widehat{\sim}_g\cong j_B(\mathfrak m g).
\end{equation}
\end{proposition}
\begin{proof}\leavevmode
\begin{enumerate}[leftmargin=*,align=left]
\item \textbf{Extract a tangent direction from a unit witness.}
A tangent in the first image at stage $C$ is a point over
$C[\eps]/(\eps^2)$ reducing to $g$. By the regular-image
presentation, after a regular-epimorphic refinement it has a unit
witness $u=u_0+\eps u_1$. Its value is
$ug=u_0g+\eps u_1g$. Reduction gives $u_0g=g$, and coefficient
extraction places the direction $u_1g$ in $C\otimes_BBg$.

\item \textbf{Descend the direction through the witness cover.}
This membership descends because $C\otimes_BBg\hookrightarrow
C\otimes_BK$ is monic and subobject inclusions are reflected by
regular-epimorphic covers.

\item \textbf{Realize every supported tangent direction.}
Conversely, for every
$b\in C$, the element $1+\eps b$ is a unit with inverse
$1-\eps b$ and has orbit value $g+\eps bg$.
Thus the tangent directions at every stage are exactly
$C\otimes_BBg$, naturally in $C$.

\item \textbf{Compute the small-multiplier tangent.}
For the second locus, translate
$j_B(\mathfrak m g)$ by $g$ and use the dual-number coefficient
identity; its directions are exactly $C\otimes_B\mathfrak m g$.
The construction and inverse are natural in the stage and reflect
equalities of tangent coefficients.
\end{enumerate}
\end{proof}

Lower-bound witnesses also have a differential equation.
In the incidence $g=af$, with $a$ a bounded multiplier,
the tangent equation is
\[
 \delta g=a\,\delta f+f\,\delta a.
\]
For the strict lower-bound incidence the multiplier and its
variation lie in $j_B\mathfrak m$.
This equation is the derivative of the witnessed multiplication
constraint, whose image realizes the lower-bound relation.
Thus the multiplier witness is retained while the differential
equation is formed.

\subsection{The universal differential resource family}
\begin{theorem}[Universal differential resource family]
\label{thm:universal-weil}
The functor
$j_{\bar B}:\Mod_{\bar B}(\mathscr S)
\to\Mod_{\mathcal R_{\bar B}}(\mathcal G_{\Cplx})$,
given at a stage $A$ by $A\otimes_{\bar B}-$, is exact,
fully faithful, and strong symmetric monoidal. The module
$\mathfrak M=j_{\bar B}\mathscr M$ is microlinear for all
finite-projective Weil diagrams. At every generalized ideal
point $I$, with stage codes reindexed along inverse image,
\begin{equation}\label{eq:weils-fibres}
 I^*\mathfrak M\cong j_BM_I.
\end{equation}
Moreover $\mathfrak M^D\cong\mathfrak M\times\mathfrak M$.
For a polynomial resource map
$F:j_{\bar B}\bar H\to j_{\bar B}\bar K$, the universal
constrained locus $S_F=F^{-1}(\mathfrak M)$ has tangent
\begin{equation}\label{eq:universal-normal}
 T_xS_F=\ker\bigl(j_{\bar B}(q)DF_x\bigr),
 \qquad q:\bar K\to\bar K/\mathscr M.
\end{equation}
These comparisons specialize to the corresponding budget fibres.
\end{theorem}
\begin{proof}\leavevmode
\begin{enumerate}[leftmargin=*,align=left]
\item \textbf{Realize the universal family and prove microlinearity.}
The Yoneda and finite-dual-basis proof of \cref{thm:weil} is
internal and applies in $\mathscr S$, giving the functorial
properties and the Weil coefficient identities. The flatness
of $\mathscr M$ from \cref{thm:universal-flat} lets the proof
of \cref{thm:microlinearity} apply to the full projective stages.

\item \textbf{Compute ideal fibres.}
Under \cref{thm:budget-weil}, inverse image at $I$ gives
\[
 \varinjlim_{a\in I}(A\otimes_BM_a)
 \cong A\otimes_B\varinjlim_{a\in I}M_a
 \cong A\otimes_BM_I.
\]

\item \textbf{Verify parameter and stage coherence.}
For a locally varying stage code, the same isomorphism follows
by reindexing its finite matrices; inverse image preserves the
idempotent equations and the balanced tensor presentation.
Both comparison paths along a stage map send $c\otimes m$
to the same substituted elementary tensor. This proves the
parameter and stage coherence of \eqref{eq:weils-fibres}.

\item \textbf{Compute tangent coefficients and the normal kernel.}
The dual-number basis gives the exponential formula.
Finally, exactness identifies $\mathfrak M$ with the kernel
of $j_{\bar B}(q)$; expanding $qF(x+\eps v)$ at a feasible
point gives $\eps qDF_x(v)$. Coefficient uniqueness proves
\eqref{eq:universal-normal}. Inverse image preserves this
kernel and its finite coefficient calculation.
\end{enumerate}
\end{proof}


\subsection{Conservative realization of measured feasibility}
\label{sec:conservative-feasibility}
Realization preserves the measured incidence and its feasibility image.

\begin{theorem}[Conservative computational comparison]
\label{thm:incidence}
Positive exact, signed, and Weil-realized incidences classify the
same measured implementations. At a budget $I$, the natural
map $iK\to j_BK$ sends $f$ to $1\otimes f$, and
\[
 iP\times_{j_BK}j_BM_I\cong iZ_I,
 \qquad \im(iZ_I\to iV)\cong i\CostClass(I).
\]
These comparisons also hold relatively over $\mathscr S$ for
$i_W\mathscr Z\cong i_W\Delta P\times_{j_{\bar B}\bar K}\mathfrak M$.
\end{theorem}
\begin{proof}\leavevmode
\begin{enumerate}[leftmargin=*,align=left]
\item \textbf{Preserve and reflect exact resource membership.}
\Cref{thm:signed-resources} preserves and reflects membership of
an exact positive cost in its marked resource. At a Weil stage
$A$, membership of $1\otimes\Phi(e)$ in $A\otimes_BM_I$
follows from membership in $M_I$. Conversely augmentation
$A\to B$ recovers $\Phi(e)\in M_I$ from a stage witness.

\item \textbf{Identify the incidence subobject and its image.}
Flatness of $A$ makes the scalar-extended marking monic.
Thus the incidence is the same constant subobject at each
stage; images are stagewise. The identical calculation at
each budget $a$ proves the relative assertion.
\end{enumerate}
\end{proof}


Equivalently, the realization square
\begin{equation}\label{eq:cat-conservative-realization}
\begin{tikzcd}[column sep=large,row sep=large]
 iM_I \arrow[r] \arrow[d,hook] & j_BM_I \arrow[d,hook]\\
 iK \arrow[r] & j_BK
\end{tikzcd}
\end{equation}
is cartesian.

\begin{corollary}[Preservation of implemented feasibility and optimality]
\label{sec:formal-implementation-boundary}
\label{prop:formal-optimization-conservative}
Formal principal parts and de Rham coefficient constructions preserve
the original implementation diagram's feasible opens in $\Cplx$ and
their compactness and supercompactness. Thus the optimal-witness and
finite-portfolio equivalences of \cref{sec:feasibility-moduli} apply
with the same implementation fibres and computational environments.
\end{corollary}
\begin{proof}\leavevmode
\begin{enumerate}[leftmargin=*,align=left]
\item \textbf{Recover measured incidence and feasibility by reduction.}
The coefficient constructions act on the resource-coordinate variable.
Reduction recovers the measured incidence by \cref{thm:incidence};
its image and the lattice of opens of $\Cplx$ are unchanged.
\end{enumerate}
\end{proof}

The comparison uses the split Weil extension. Arbitrary coefficient
changes behave differently: tensoring $B\incl K$ with $B/tB$ gives
a nonzero source and zero target.

The finite calculus now differentiates the measured resource families
and their budget constraints. We next classify variations of the
embedded budgets and calculate the successive feasibility obstructions.

\section{Deformations of resource-constrained computation}
\label{sec:resource-deformations}

Ordinary budget families have already been classified by $\Cplx$.
We now ask how an embedded budget can vary over a nilpotent
coefficient extension. This is a different moduli problem.
The earlier point object
$\operatorname{Pt}_{\E}(\Cplx)\cong\Idl(L)$ has constant
realization, with $(i\Idl(L))^D\cong i\Idl(L)$ by
\cref{sec:cohesion}; it carries no infinitesimal variation.
We therefore define deformations directly by extending the
coefficient ring and varying the marked submodule of the
extended fraction algebra. Tangents and lifting obstructions
will describe this functor.


\subsection{The marked deformation functor}
\label{sec:resource-deformation-functor}
For a Weil stage $A=B\oplus\mathfrak n_A$ put
$K_A=K\otimes_BA$ and
$\mathfrak n_{K,A}=K\otimes_B\mathfrak n_A$.
The split augmentation gives this kernel formula. For an idempotent
$e$, the multiplicative group $e+e\mathfrak n_{K,A}$ is formed in
the corner ring $eK_A$; inverses are finite geometric series.


\begin{definition}[The principal embedded-resource deformation functor]
\label{def:principal-deformation-functor}
Fix $M=Bg\incl K$ and its support $e=e_g$.  For a finite-projective
Weil stage $A$, let $\operatorname{Def}^{\mathrm{pr}}_M(A)$ be the
object of marked images $Ax\incl K_A$ with $x\in eK_A$ reducing
to $g$.  Generators determine the same point precisely when their
embedded images agree.  An augmented Weil map $A\to A'$ sends
this point to $A'x'\incl K_{A'}$, where $x'$ is the substituted
generator.
\end{definition}

To verify functoriality, write $x=g(e+z)$ for a unique
$z\in e\mathfrak n_{K,A}$, since $gg^\dagger=e$.
The multiplier $e+z$ is invertible in the corner, and the
finite-projective stage $A$ is flat over $B$.  Multiplication by
$x$ therefore identifies the displayed image with $eA$.  If $Ax=Ay$,
each generator is a scalar multiple of the other; substitution
preserves these two equations, hence preserves equality of images.
Identity and composition hold on generators and descend to the
image quotient.  Tight replacement of $g$ gives the same marked
functor, as proved in \cref{cor:formal-tight-invariance}.

The functor thus classifies nilpotent deformations of an embedded
budget, while $h_{\Cplx}(Z)$ classifies ordinary budget families
over locales. We first calculate the deformation functor at
first order, then use that tangent module to describe its
formal classifying disk.

\subsection{First-order deformations of finite-projective cells}
\label{sec:first-resource-deformations}
At first order, a deformation changes the embedding of a module
while reducing to its original marking.
Let $A_1=B[\eps]/(\eps^2)$ and
$K_1=K\otimes_BA_1\cong K[\eps]/(\eps^2)$.
For a displayed finite-projective $i_M:M\incl K$, a framed
deformation is an $A_1$-linear monomorphism
$\widetilde i:M\otimes_BA_1\incl K_1$ reducing to $i_M$.
The frame is the displayed source $M\otimes_BA_1$.
Forgetting it identifies monomorphisms with the same embedded
image. The theorem calculates this image quotient.

\begin{theorem}[The embedded resource tangent module]
\label{thm:embedded-deformation-tangent}
Every framed first-order embedded deformation has the unique form
\begin{equation}\label{eq:framed-deformation-form}
 i_h(m+\eps n)=m+\eps(n+h(m)),
 \qquad h\in\Hom_B(M,K).
\end{equation}
Two deformations have the same image if and only if
$h-h'$ factors through $M\incl K$.
Consequently the object of image classes is canonically
\begin{equation}\label{eq:resource-moduli-tangent}
 T_M^{\mathrm{emb}}
 \cong \Hom_B(M,K)/\Hom_B(M,M)
 \cong \Hom_B(M,K/M).
\end{equation}
For $M=Bg\cong Be_g$ this reduces to
\begin{equation}\label{eq:principal-resource-tangent}
 T_{Bg}^{\mathrm{emb}}
 \cong e_gK/Bg;
 \qquad T_B^{\mathrm{emb}}\cong K/B.
\end{equation}
These isomorphisms are natural under marked isomorphisms of
finite-projective resource modules and substitution of parameters.
\end{theorem}
\begin{proof}\leavevmode
\begin{enumerate}[leftmargin=*,align=left]
\item \textbf{Classify framed first-order deformations.}
The decompositions $A_1=B\oplus\eps B$ and
$K_1=K\oplus\eps K$ are finite biproducts.  An
$A_1$-linear map reducing to $i_M$ sends $m$ to
$m+\eps h(m)$ for a unique $B$-linear $h:M\to K$;
$A_1$-linearity forces its value at $\eps n$ to be
$\eps n$.  This proves \eqref{eq:framed-deformation-form}.

\item \textbf{Prove that every framed deformation is monic.}
If $i_h(m+\eps n)=0$, comparison of the constant
coefficients gives $m=0$ because $M\incl K$ is monic.
Comparison of the nilpotent coefficients then gives $n=0$.
Thus every $i_h$ is a monomorphism.

\item \textbf{Identify changes of framing.}
If $i_h$ and $i_{h'}$ have the same image, their
monomorphisms factor through each other as inverse
$A_1$-linear automorphisms of $M\otimes_BA_1$.
Reduction of such an automorphism is the identity, since
$i_M$ is monic.  Its formula is necessarily
$m+\eps n\mapsto m+\eps(n+u(m))$ for a
$B$-linear $u:M\to M$.  The equation
$i_{h'}=i_h(1+\eps u)$ yields $h'-h=i_Mu$.

\item \textbf{Identify the image-equivalence quotient.}
Conversely this equation constructs the required
inverse automorphism $1-\eps u$.
The first quotient in \eqref{eq:resource-moduli-tangent}
is therefore the image-equivalence quotient.

\item \textbf{Compute the quotient by finite-projective duality.}
Apply $\Hom_B(M,-)$ to
$0\to M\to K\to K/M\to0$.
The specified finite-free retract supplies a finite dual basis,
so $\Hom_B(M,-)\cong M^\vee\otimes_B-$ is exact: the dual
module is a retract of a finite free module. Consequently the
cokernel is $\Hom_B(M,K/M)$.

\item \textbf{Specialize to principal resources and verify naturality.}
For
$M\cong Be_g$ evaluation at $e_g$ identifies
$\Hom_B(M,K/M)$ with $e_g(K/M)=e_gK/Bg$.
All maps are finite module operations and quotient
universal properties, hence commute with parameter substitution.
\end{enumerate}
\end{proof}


\paragraph{Change of frame and its effective quotient.}
The tangent calculation can be organized by the changes of
frame that leave the image fixed. Write
$T_M=T_M^{\mathrm{emb}}\cong\Hom_B(M,K/M)$,
$\mathcal F_M=\Hom_B(M,K)$, and
$\mathcal G_M=\operatorname{End}_B(M)$.
The additive action is $u\cdot h=h+i_Mu$.
Its compatibility with the embedding is
\[
\begin{tikzcd}[column sep=large,row sep=large]
 M\otimes_BA_1 \arrow[r,"{1+\eps u}","{\cong}"']
   \arrow[d,"{i_{h'}}"']
   & M\otimes_BA_1 \arrow[d,"{i_h}"]\\
 K_1 \arrow[r,equal] & K_1 .
\end{tikzcd}
\qquad h'=h+i_Mu.
\]
The action is free since $i_M$ is monic.  Its effective quotient
is the marked-image object, presented by the coequalizer
\begin{equation}\label{eq:cat-frame-quotient}
\begin{tikzcd}[column sep=large]
 \mathcal G_M\times\mathcal F_M
   \arrow[r,shift left=.6ex,"{\mathrm{act}}"]
   \arrow[r,shift right=.6ex,"{\pi_2}"']
   & \mathcal F_M \arrow[r,two heads] & T_M .
\end{tikzcd}
\end{equation}
Exactness of $\Hom_B(M,-)$ gives the short exact sequence
\[
 0\longrightarrow\operatorname{End}_B(M)
 \xrightarrow{i_M\circ-}\Hom_B(M,K)
 \xrightarrow{q\circ-}\Hom_B(M,K/M)\longrightarrow0 .
\]
The cokernel therefore removes exactly the infinitesimal changes
of frame. Monicity of the marking makes the action free,
and the effective quotient records its image.


For a principal resource $M=Bg$, the first-order image classes in
\cref{def:principal-deformation-functor} coincide with those just
computed.  A frame sends $g$ to $g+\eps h(g)$; $B$-linearity makes
$h(g)$ supported on $e_g$.  Conversely each supported generator
reducing to $g$ gives a frame on $M\otimes_BA_1\cong e_gA_1$.
Equality of images is the same equivalence in both descriptions.
Thus \eqref{eq:principal-resource-tangent} is the first-order value
of the already defined deformation functor.

For cofinite numerical input the class of
$t^{-1}=n+1$ is nonzero in $K/B$; the deformation
\[
 B[\eps](1+\eps t^{-1})\subset K[\eps]
\]
therefore differs from the constant deformation
$B[\eps]\subset K[\eps]$ while reducing to the same
embedded module $B\subset K$.
This exhibits a nonzero moduli tangent already on the
rank-one resource cell.

\subsection{Formal resource cells and their classifying disks}
\label{sec:formal-resource-moduli}
The first-order quotient identifies the intrinsic tangent module.
For a principal resource, finite logarithm and exponential extend
this description to every nilpotence order. We now construct
the quotient and representing formal disk of
\cref{def:principal-deformation-functor}.

\begin{theorem}[The formal classifier of a principal resource cell]
\label{thm:principal-formal-deformations}
\label{thm:linearized-deformations}
For $M=Bg\incl K$, with $e=e_g$, the image classes of principal
$A$-submodules generated by an element reducing to $g$ and supported
in $eK_A$ form the functor
\begin{equation}\label{eq:formal-resource-functor}
 \operatorname{Def}^{\mathrm{pr}}_M(A)
 =\frac{e+e\mathfrak n_{K,A}}{e+e\mathfrak n_A},
 \qquad [u]\longmapsto A(gu)\subseteq K_A.
\end{equation}
Put $T_M=\Hom_B(M,K/M)$. There are canonical natural isomorphisms
\begin{equation}\label{eq:linearized-deformations}
 eK/eB\xrightarrow{\ \cong\ }T_M,\qquad
 \operatorname{Def}^{\mathrm{pr}}_M(A)
       \cong\mathfrak n_A\otimes_BT_M.
\end{equation}
The first map sends a multiplier class $[z]$ to
$(m\mapsto zm\bmod M)$. With $0:1\to iT_M$, the formal disk
\begin{equation}\label{eq:resource-formal-disk}
 \mathfrak D_M=1\times_{iT_M}j_BT_M
\end{equation}
represents the functor: its value at $A$, equivalently the internal
object of maps $D_A\to\mathfrak D_M$, is
$\operatorname{Def}^{\mathrm{pr}}_M(A)$.
\end{theorem}
\begin{proof}\leavevmode
\begin{enumerate}[leftmargin=*,align=left]
\item \textbf{Classify supported deformations by unit cosets.}
The support inverse $g^\dagger=e/g$ satisfies $gg^\dagger=e$.
Thus every supported generator reducing to $g$ is uniquely $gu$
with $u\in e+e\mathfrak n_{K,A}$. Multiplication by $gu$ embeds
$eA$ in $eK_A$: the inclusion $A\incl K_A$ is monic since $A$
is a finite-free retract and regular denominators act injectively
on it. Two such images agree precisely when $u^{-1}v$ is a unit
of $eA$ reducing to $e$, namely an element of $e+e\mathfrak n_A$.
This proves the quotient classification and its functoriality.

\item \textbf{Construct finite logarithm and exponential.}
Choose the displayed bound $r$ with $\mathfrak n_A^{r+1}=0$.
The ideal $e\mathfrak n_{K,A}$ has the same nilpotence bound.
For $x$ in that ideal define in the corner
\[
 \log(e+x)=\sum_{k=1}^r\frac{(-1)^{k+1}}k x^k,
 \qquad \exp(x)=e+\sum_{k=1}^r\frac{x^k}{k!}.
\]

\item \textbf{Identify the multiplicative and additive quotients.}
The rational polynomial identities, reduced modulo terms of degree
greater than $r$, prove that these are inverse and that logarithm
takes products to sums. Increasing the displayed bound changes
neither map. They carry the subgroup over $eA$ onto
$e\mathfrak n_A$, so they induce
\[
 \operatorname{Def}^{\mathrm{pr}}_M(A)
 \cong (eK\otimes_B\mathfrak n_A)/(eB\otimes_B\mathfrak n_A)
 \cong\mathfrak n_A\otimes_B(eK/eB).
\]

\item \textbf{Identify the tangent module intrinsically.}
Here $\mathfrak n_A$ is a displayed finite-free retract, hence
flat. Multiplication by $g$ identifies $eK/eB$ with $eK/Bg$;
the dual-basis calculation of
\cref{thm:embedded-deformation-tangent} identifies this with $T_M$.
Its intrinsic formula is the multiplier map in the statement,
so it is independent of the chosen generator.

\item \textbf{Identify the representing formal disk.}
Finally, augmentation on $A\otimes_BT_M$ has kernel
$\mathfrak n_A\otimes_BT_M$ by the split decomposition of $A$.
This calculates \eqref{eq:resource-formal-disk}; the internal
Yoneda lemma gives its representing property. All polynomial
identities and quotients commute with Weil maps and parameters.
\end{enumerate}
\end{proof}


At every stage the logarithm is an isomorphism of the quotient
presentations.  In the corner supported on $e$, it gives
\begin{equation}\label{eq:cat-logarithmic-quotient}
\begin{tikzcd}[column sep=large,row sep=large]
 e+e\mathfrak n_A \arrow[r,hook] \arrow[d,"{\log}","{\cong}"']
   & e+e\mathfrak n_{K,A} \arrow[r,two heads]
       \arrow[d,"{\log}","{\cong}"']
   & \operatorname{Def}^{\mathrm{pr}}_M(A) \arrow[d,"{\cong}"]\\
 e\mathfrak n_A \arrow[r,hook]
   & e\mathfrak n_{K,A} \arrow[r,two heads]
   & \mathfrak n_A\otimes_BT_M .
\end{tikzcd}
\end{equation}
The upper row is a quotient of multiplicative groups and the lower
row is its additive cokernel presentation.  Naturality of the finite
logarithm and exponential makes this a diagram of Weil-stage functors.


\begin{corollary}[Universal images and finite lifts]
\label{cor:formal-resource-lifts}
For $z\in\mathfrak n_A\otimes_BT_M$, choose a representative
$\widetilde z\in e\mathfrak n_{K,A}$ locally along the quotient
epimorphism. The image $A(g\exp\widetilde z)$ is independent of
that representative and constructs the universal marked family
over $\mathfrak D_M$. Surjective Weil-stage maps induce internal
epimorphisms on these deformation objects. For
$A_r=B[d]/d^{r+1}$ and $M=B$,
\begin{equation}\label{eq:formal-truncated-moduli}
 \operatorname{Def}^{\mathrm{pr}}_B(A_r)
 \cong\frac{1+dK[d]/d^{r+1}}{1+dB[d]/d^{r+1}}
 \cong\bigoplus_{j=1}^r d^j(K/B).
\end{equation}
Reduction to order $r-1$ has kernel $K/B$.
\end{corollary}
\begin{proof}\leavevmode
\begin{enumerate}[leftmargin=*,align=left]
\item \textbf{Descend the universal image through the logarithmic quotient.}
Changing $\widetilde z$ by $e\mathfrak n_A$ multiplies its
exponential by a unit of $eA$, hence preserves the image.
These equal images descend across the quotient epi and commute
with stage maps, giving a submodule of $K_A$ over the formal
classifier.

\item \textbf{Construct lifts along surjective stage maps.}
A surjective augmented algebra map is surjective
on augmentation ideals: subtract the augmentation of a lift.
Tensoring preserves that epimorphism. The last assertion follows
by the displayed direct-sum description.
\end{enumerate}
\end{proof}

\begin{corollary}[Invariance under tight replacement]
\label{cor:formal-tight-invariance}
If $g'\in\Theta(g)$, then $Bg'=Bg$ as marked modules and their
formal classifiers, universal images, and tangent comparisons
are canonically identical.
\end{corollary}
\begin{proof}\leavevmode
\begin{enumerate}[leftmargin=*,align=left]
\item \textbf{Identify the deformation classifier under tight replacement.}
Write $g'=vg$ with $v\in B^\times$ by \cref{thm:orbits}.
The support is unchanged and $A(g'u)=A(gu)$. The multiplier
description of $T_M$ and logarithm identify the same quotient
and representing disk at every stage.
\end{enumerate}
\end{proof}

Logarithm and exponential use finite rational arithmetic at the
displayed nilpotence bound, so they act on the coefficient
presentations. Subgroup membership and equality of image classes
remain the associated quotient predicates, with their separately
specified observation structure.
A deformation of a resource need not preserve a particular
resource constraint. We now vary a polynomial resource value
and its budget together. At first order, the change of the
value modulo the budget must match the change of the budget
at that value.

\subsection{Deformation of incidence with a moving budget}
\label{sec:moving-incidence}
Let $M\incl K$ be a displayed finite-projective resource cell,
$q:K\to Q=K/M$, and let $F:j_BH\to j_BK$ be one of the
polynomial resource maps constructed in
\cref{thm:polynomial-lift,prop:polynomial-laws}.
A reduced feasible point $x:1\to H$ has $F(x)\in M$.
The displayed generators of $M$ supply the principal generator
$g=\sum_i|i_Mr(e_i)|$, where $i_M:M\incl K$ is the marking,
$r:B^n\to M$ is its displayed retraction, and $(e_i)$ is the
standard basis. Thus $M=Bg$ by bounded division, and the preceding
formal classifier applies. Its independence of the generator
was proved in \cref{cor:formal-tight-invariance}.
Define the incidence deformation functor by
\begin{equation}\label{eq:incidence-deformation-functor}
 \mathsf{Inc}_{F,x,M}(A)=
 \left\{(x_A,N_A):
 \begin{array}{l}
 x_A\in A\otimes_BH,\quad (\epsilon_A\otimes1)x_A=x,\\
 N_A\in\operatorname{Def}^{\mathrm{pr}}_M(A),\quad F_A(x_A)\in N_A
 \end{array}\right\}.
\end{equation}
Here $\epsilon_A:A\to B$ is the augmentation and
$N_A\incl K_A$ is the universal marked image of
\cref{cor:formal-resource-lifts}. Weil substitution acts on both
entries and preserves membership by naturality of $F$ and of
that image. This constructs a subfunctor of the formal disk of
$j_BH$ at $x$ times $\mathfrak D_M$. Source frames are forgotten
by equality of embedded images. Its first-order object is denoted
$T_{(x,M)}\mathsf{Inc}_F:=\mathsf{Inc}_{F,x,M}(B[\eps]/\eps^2)$.

\begin{theorem}[The tangent of resource incidence]
\label{thm:moving-incidence}
The first-order incidence deformations are the kernel of the
$B$-linear map
\[
 H\oplus T_M\longrightarrow Q,
 \qquad(v,\theta)\longmapsto qDF_x(v)-\theta(F(x)).
\]
Equivalently,
\begin{equation}\label{eq:moving-incidence}
 T_{(x,M)}\mathsf{Inc}_F
 =\{(v,\theta):qDF_x(v)=\theta(F(x))\}.
\end{equation}
For a fixed budget this gives
\begin{equation}\label{eq:fixed-budget-tangent}
 T_xF^{-1}(j_BM)=\ker(qDF_x).
\end{equation}
\end{theorem}
\begin{proof}\leavevmode
\begin{enumerate}[leftmargin=*,align=left]
\item \textbf{Represent the varying point and budget.}
A first-order point is $x+\eps v$ and its value is
$F(x)+\eps DF_x(v)$ by the Weil coefficient identity.
By \cref{thm:embedded-deformation-tangent}, a budget deformation
is locally represented by $h:M\to K$, with image of
$m+\eps n\mapsto m+\eps(n+h(m))$ and class $\theta=qh$.

\item \textbf{Derive the incidence equation in the quotient.}
Membership of the value is exactly
$DF_x(v)-h(F(x))\in M$. Passing to $Q$ gives
\eqref{eq:moving-incidence}.

\item \textbf{Recover membership and descend the calculation.}
Conversely that equation supplies
the unique element of $M$ representing the displayed difference,
and hence the required membership. Replacing $h$ by a map with
the same class changes neither condition. The calculations
therefore descend from the quotient cover, naturally in all
parameters. The fixed-budget slice is $\theta=0$.
\end{enumerate}
\end{proof}


\paragraph{The incidence deformation complexes.}
The same linearized map governs first-order solutions and
higher-order corrections. We record the fixed and moving
versions as two-term complexes of $B$-modules in degrees $0,1$:
\begin{equation}\label{eq:cat-incidence-complexes}
\begin{aligned}
 \DCx^{\mathrm{fix}}_{F,x,M}
   &= [\,H\xrightarrow{qDF_x}Q\,],\\
 \DCx^{\mathrm{mov}}_{F,x,M}
   &= [\,H\oplus T_M\xrightarrow{L^{\mathrm{mov}}_{(x,M)}}Q\,],\\
 L^{\mathrm{mov}}_{(x,M)}(v,\theta)
   &=qDF_x(v)-\theta(F(x)).
\end{aligned}
\end{equation}
In each complex, $H^0$ is the kernel and $H^1$ the cokernel.
The kernel gives the first-order incidence deformations.
At higher orders, a residual can be cancelled precisely when
its class in the cokernel vanishes; the lift theorems below
then identify the inhabited lift object as a torsor under
$H^0$. Allowing the budget to vary gives the comparison
\begin{equation}\label{eq:cat-fixed-moving-complex}
\begin{tikzcd}[column sep=large,row sep=large]
 H \arrow[r,"{qDF_x}"] \arrow[d,hook,"{v\mapsto(v,0)}"']
   & Q \arrow[d,equal]\\
 H\oplus T_M \arrow[r,"{L^{\mathrm{mov}}_{(x,M)}}"']
   & Q .
\end{tikzcd}
\end{equation}
This compares the tangent and obstruction modules for the
two problems: the vertical inclusion allows the additional
budget direction while preserving the value residual.


\begin{proposition}[Normal directions for a fixed marked budget]
\label{prop:normal}
For any embedded resource $M$ and geometric valuation family
$F:j_BH\to j_BK$, the tangent of $S_F=F^{-1}(j_BM)$ at a
feasible point is $\ker(j_B(q)DF_x)$.
\end{proposition}
\begin{proof}\leavevmode
\begin{enumerate}[leftmargin=*,align=left]
\item \textbf{Realize the fixed resource as a kernel.}
Exact realization identifies $j_BM$ with the kernel of $j_B(q)$.
Exponentiation by the dual-number object preserves this
pullback.

\item \textbf{Compute the tangent kernel at a feasible point.}
Its coefficient formula gives $(qF(x),qDF_x(v))$,
whose zero fibre at a feasible point is the asserted kernel.
This proves \eqref{eq:fixed-budget-tangent} for arbitrary embedded budgets.
\end{enumerate}
\end{proof}

\begin{corollary}[Rigidity of a tight budget containing a fixed cost]
\label{cor:tight-budget-rigidity}
For a fixed implemented cost $c\in K$, the only principal formal
deformation of $Bc$ containing the constant cost $c$ is the
constant deformation. Its first-order version is
$\{\theta\in T_{Bc}:\theta(c)=0\}=0$.
\end{corollary}
\begin{proof}\leavevmode
\begin{enumerate}[leftmargin=*,align=left]
\item \textbf{Translate containment of the fixed cost into a unit equation.}
Write $e=e_c$. A deformation is $A(cu)$ for
$u\in e+e\mathfrak n_{K,A}$. If $c\in A(cu)$, multiplication
by $c^\dagger$ yields $e=au$ for some $a\in eA$.

\item \textbf{Prove triviality of the deformation.}
Thus $a=u^{-1}$ reduces to $e$ and is a nilpotent perturbation
of the unit in $eA$. Its finite geometric inverse belongs to
$eA$, so $u\in e+e\mathfrak n_A$ and its class is trivial.

\item \textbf{Deduce first-order rigidity.}
At first order, $c$ generates $Bc$, so a linear map vanishing
on $c$ is zero.
\end{enumerate}
\end{proof}

\subsection{Higher-order feasibility and lifting obstructions}
\label{sec:feasibility-obstructions}
A feasible deformation through order $r$ need not extend to
order $r+1$. Its lower-order coefficients determine the next
failure of the constraint, and a new coefficient changes that
failure through the linearized incidence map. We calculate
whether this residual can be cancelled.
Use the fixed budget $M$ and polynomial $F$ above, and put
$L_x=qDF_x:H\to Q$. An order-$r$ feasible jet is
\[
 x_r=x+\sum_{j=1}^r d^jv_j\quad\text{over }B[d]/d^{r+1},
 \qquad qF(x_r)=0.
\]
View the same displayed polynomial over $B[d]/d^{r+2}$ and define
\[
 R_{r+1}=[d^{r+1}]qF(x_r)\in Q.
\]
The coefficients $v_j$ are uniquely determined by the jet, so
this is its canonical residual.

\begin{theorem}[The obstruction and the object of lifts]
\label{thm:feasibility-obstructions}
The obstruction to lifting $x_r$ by one order is
\begin{equation}\label{eq:feasibility-obstruction}
 \operatorname{ob}_{r+1}(x_r)
       =[R_{r+1}]\in\operatorname{coker}L_x.
\end{equation}
It vanishes exactly when the internal object of feasible lifts
is inhabited. That object is the fibre
$\{w:H:L_x(w)=-R_{r+1}\}$; when inhabited, it is a torsor
under $\ker L_x$. The construction is natural in the polynomial
data and reduced parameters, and applies relatively over
the structured resource space.
\end{theorem}
\begin{proof}\leavevmode
\begin{enumerate}[leftmargin=*,align=left]
\item \textbf{Compute the next-order correction equation.}
A lift has the unique form $x_r+d^{r+1}w$.
The polynomial Taylor identity gives, modulo $d^{r+2}$,
\[
 qF(x_r+d^{r+1}w)
   =d^{r+1}\bigl(R_{r+1}+L_x(w)\bigr).
\]
Indeed terms quadratic in the added correction have degree
at least $r+2$, and the positive-degree part of $DF_{x_r}-DF_x$
also has that property after multiplication by $d^{r+1}$.
For $r=0$, the same identity is the dual-number formula.

\item \textbf{Identify lifts with the linear fibre.}
Equality of polynomial coefficients proves the fibre assertion.

\item \textbf{Characterize existence by the obstruction class.}
The kernel of $Q\to\operatorname{coker}L_x$ is the regular
image of $L_x$, so zero obstruction means exactly internal
existence of $w$.

\item \textbf{Identify the lift torsor and parameter stability.}
Two lifts differ by a unique element of
$\ker L_x$, and addition by such an element preserves a lift.
This proves the torsor assertion.
All maps are coefficient operations, kernels, and cokernels;
their internal constructions and identities apply in
$\mathscr S$ as well.
\end{enumerate}
\end{proof}


\begin{example}[The second-order obstruction of the measured family]
\label{ex:second-obstruction}
For the family $F(\lambda)=1+t^{-2}\lambda^2$ of
\cref{ex:weighted-resource-family}, take $H=B$, $M=B$, and $x=0$.
Here $L_0=qDF_0=0$, and the first-order jet $\lambda_1=d$ is feasible.
Its next residual is $R_2=[t^{-2}]\in K/B$, which is nonzero because
$(n+1)^2$ is unbounded. Every lift $d+d^2w$ has the same square
modulo $d^3$, so no parameter correction removes the obstruction.
Under the coordinate isomorphism $B\to t^{-1}B$,
$\lambda\mapsto b=t^{-1}\lambda$, this is equivalently the family
$1+b^2$ with first-order direction $t^{-1}$.
\end{example}

\subsection{Complete principal-resource moduli}
\label{sec:complete-principal-resource-moduli}
The principal-resource classifier of
\cref{thm:principal-formal-deformations} applies at every finite
nilpotence order. We now assemble its image classes into a
compatible formal family. The classification is taken at the
level of those classes, so compatible generators need not be
chosen.

\begin{theorem}[Completed principal deformation functor]
\label{thm:complete-principal-moduli}
For a principal marked resource $M=Bg\incl K$ put
$T_M=\Hom_B(M,K/M)$ and $A_r=B[d]/d^{r+1}$.  Then
\begin{equation}\label{eq:complete-principal-moduli}
 \widehat{\operatorname{Def}}^{\mathrm{pr}}_M
 :=\varprojlim_r\operatorname{Def}^{\mathrm{pr}}_M(A_r)
 \ \cong\ d\,T_M[[d]]
 =\prod_{j\ge1}d^jT_M.
\end{equation}
The transition maps are projections onto initial coefficient
segments.  Every finite-order class admits a compatible extension
to all higher orders, and the displayed comparison is natural in
marked principal-resource isomorphisms and parameters.
\end{theorem}
\begin{proof}\leavevmode
\begin{enumerate}[leftmargin=*,align=left]
\item \textbf{Compute each finite deformation stage.}
The augmentation ideal of $A_r$ is the displayed free module
$\mathfrak n_r=dB\oplus\cdots\oplus d^rB$.  Applying the
finite logarithm classification of
\cref{thm:principal-formal-deformations} yields
\[
 \operatorname{Def}^{\mathrm{pr}}_M(A_r)
 \cong\mathfrak n_r\otimes_BT_M
 \cong\bigoplus_{1\le j\le r}d^jT_M.
\]

\item \textbf{Identify the inverse limit with coefficient sequences.}
Its naturality in Weil-stage maps identifies truncation with
deletion of the highest coefficient.  A point of the inverse limit
is therefore exactly a compatible finite-prefix family of
$T_M$-coefficients, equivalently an internal sequence
$(z_j)_{j\ge1}$.  This gives \eqref{eq:complete-principal-moduli}.

\item \textbf{Construct an extension of every finite prefix.}
Extend a given coefficient prefix by zero coefficients to obtain
an explicit compatible sequence of deformation \emph{classes}.
\end{enumerate}
\end{proof}


Principal-resource image deformations are formally unobstructed by
coefficientwise zero extension. The following calculation describes
the quotient that supplies their coefficients.

\begin{proposition}[$B$-linear nonsplitting of the normal quotient]
\label{prop:no-formal-representative-selector}
In the cofinite numerical fibre, the quotient
$q:K\twoheadrightarrow K/B$ does not admit a $B$-linear
section.
\end{proposition}
\begin{proof}\leavevmode
\begin{enumerate}[leftmargin=*,align=left]
\item \textbf{Exhibit nonzero torsion in the normal quotient.}
Let $t(n)=1/(n+1)\in B$.  It is regular in $B$ and
invertible in $K$, while $t^{-1}=n+1\notin B$.  Therefore
$\xi=q(t^{-1})\in K/B$ is nonzero but
$t\xi=q(1)=0$.

\item \textbf{Derive a contradiction from a linear section.}
If a $B$-linear section
$s:K/B\to K$ existed, then
\[
 t\,s(\xi)=s(t\xi)=0.
\]
Multiplication by $t$ is invertible on $K$, so
$s(\xi)=0$, contradicting
$q(s(\xi))=\xi\ne0$.
\end{enumerate}
\end{proof}

We now combine these two calculations. The budget itself admits
formal extension, but its extension must still be compatible
with the resource value. Allowing a budget correction enlarges
the linearized incidence operator and gives the moving-budget
obstruction.

\subsection{Higher-order incidence with a moving principal budget}
\label{sec:moving-higher-order}
Let $H$ be a displayed finite-projective $B$-module, let
$F:H\to K$ be a specified polynomial resource map, and let
$x\in H$ satisfy $F(x)\in M=Bg$.  Write
\[
 Q=K/M,\quad q:K\onto Q,\quad
 T_M=\Hom_B(M,Q).
\]
The first-order moving-incidence operator of
\cref{sec:resource-deformations} is
\begin{equation}\label{eq:moving-incidence-linear}
 L^{\mathrm{mov}}_{(x,M)}:
   H\oplus T_M\longrightarrow Q,\qquad
 (v,\theta)\longmapsto qDF_x(v)-\theta(F(x)).
\end{equation}
This map is intrinsic to the marking $M\incl K$.

\begin{definition}[Feasible moving formal jet]
\label{def:moving-formal-feasible}
At order $r$, a moving formal jet is a valuation coefficient
$x_r=x+\sum_{j=1}^r d^jv_j\in H\otimes_BA_r$
together with a principal deformation image
$M_r\in\operatorname{Def}^{\mathrm{pr}}_M(A_r)$.
It is feasible when the image of $F(x_r)$ in
$(K\otimes_BA_r)/M_r$ is zero.  The object of these pairs is
formed as a pullback of the corresponding polynomial evaluation
and marked-incidence subobjects.  All its equalities are
$\Omega$-valued.
\end{definition}

\begin{theorem}[Moving-budget lifting obstruction]
\label{thm:moving-budget-obstructions}
For a feasible moving jet $(x_r,M_r)$, there is a canonical class
\begin{equation}\label{eq:moving-budget-obstruction}
 \operatorname{ob}^{\mathrm{mov}}_{r+1}(x_r,M_r)
 \in\operatorname{coker}L^{\mathrm{mov}}_{(x,M)}.
\end{equation}
Its vanishing is equivalent to internal inhabitation of the
object of feasible lifts to order $r+1$.  When inhabited, that
object is a torsor under $\ker L^{\mathrm{mov}}_{(x,M)}$.
The constructions are natural in polynomial resource maps,
marked principal modules, and all supplied parameters.  The
fixed-budget theorem \cref{thm:feasibility-obstructions} is the
specialization in which the budget tangent component is zero.
\end{theorem}
\begin{proof}\leavevmode
\begin{enumerate}[leftmargin=*,align=left]
\item \textbf{Represent the budget by a supported logarithm.}
Write $e=e_g$.  A principal image class admits, locally along
the quotient epimorphism of
\cref{thm:principal-formal-deformations}, a supported logarithmic
multiplier
\[
 z_r(d)=\sum_{j=1}^r d^jz_j\in eK\otimes_BdA_r.
\]

\item \textbf{Express feasibility through an ambient unit.}
In the corner algebra $eK\otimes_BA_r$, its finite exponential
$\exp_e(z_r)$ has constant coefficient $e$.  Extend it to a
unit of $K\otimes_BA_r$ by
\[
 U_r(d)=(1-e)+\exp_e(z_r(d)),\qquad
 U_r^{-1}=(1-e)+\exp_e(-z_r(d)).
\]
Then $M_r=U_r(d)\,(M\otimes_BA_r)$, and its membership equation
is equivalent to
\[
 q\bigl(U_r^{-1}F(x_r)\bigr)=0\pmod{d^{r+1}}.
\]

\item \textbf{Define the next residual coefficient locally.}
Regard the displayed polynomials as their zero-extended
representatives modulo $d^{r+2}$, evaluate the finite exponential
to that order, and put
\[
 R_{r+1}=[d^{r+1}]q\bigl(U_r^{-1}F(x_r)\bigr)\in Q.
\]

\item \textbf{Descend the residual through changes of representative.}
This is well-defined locally. Two logarithmic representatives
of the same image class differ by
$b(d)\in deB[d]/d^{r+1}$, by the logarithmic quotient
in \cref{thm:principal-formal-deformations}.
Zero-extending both logarithms and taking their exponentials
modulo $d^{r+2}$ gives
\[
 U'_r=U_r h,\qquad h=(1-e)+\exp_e(b)\in
 (B[d]/d^{r+2})^\times,\qquad h(0)=1.
\]
The residual series is consequently multiplied by $h^{-1}$.
Its coefficients through degree $r$ vanish in $Q$, so the
coefficient of degree $r+1$ is unchanged. This equality on the
kernel pair of the representative epimorphism constructs the
descended map $R_{r+1}$ by the effective-quotient property.

\item \textbf{Represent the valuation and budget corrections.}
A next-order lift adds $d^{r+1}v$ to the valuation and a new
budget tangent coefficient $\theta\in T_M$.  Locally represent
$\theta$ by a supported multiplier $\zeta\in eK$, so its
infinitesimal action on $m\in M$ is
$\theta(m)=\zeta m\bmod M$.

\item \textbf{Compute the linearized correction equation.}
A finite polynomial Taylor expansion
gives
\[
 F(x_r+d^{r+1}v)
 \equiv F(x_r)+d^{r+1}DF_x(v)\pmod{d^{r+2}}.
\]
The inverse supported multiplier changes by
$U_{r+1}^{-1}\equiv U_r^{-1}(1-d^{r+1}\zeta)$ at that order.
Since its constant term is $1$ and $F(x)\in M$, the new
coefficient of the feasibility equation is
\[
 R_{r+1}+qDF_x(v)-\theta(F(x))
 =R_{r+1}+L^{\mathrm{mov}}_{(x,M)}(v,\theta).
\]

\item \textbf{Verify independence of the new representative.}
Changing $\zeta$ by $eB$ has no effect on this equation because
$F(x)\in M$, so it is independent of the new local lift as well.

\item \textbf{Identify the canonical obstruction.}
It follows that a feasible lift exists internally if and only if
$-R_{r+1}$ belongs to the regular image of
$L^{\mathrm{mov}}_{(x,M)}$, equivalently if and only if its
cokernel class vanishes.  This defines the canonical obstruction
\eqref{eq:moving-budget-obstruction}.

\item \textbf{Identify the torsor and parameter coherence.}
The fibre of solutions to a linear inhomogeneous equation is a
torsor under its kernel: subtraction of two solutions gives a
unique kernel element, and addition of such an element preserves
solutions. Every step was a finite
polynomial equation or exact quotient calculation and hence
commutes with parameter substitution.
\end{enumerate}
\end{proof}


\begin{corollary}[Comparison of fixed and moving obstructions]
\label{cor:cat-obstruction-comparison}
For a fixed-budget feasible jet, regarded as a moving jet with
constant budget, the morphism
\eqref{eq:cat-fixed-moving-complex} carries its fixed obstruction
to its moving obstruction.  Explicitly the comparison square is
\[
\begin{tikzcd}[column sep=large,row sep=large]
 Q \arrow[r,two heads] \arrow[d,equal]
   & H^1(\DCx^{\mathrm{fix}}_{F,x,M}) \arrow[d]\\
 Q \arrow[r,two heads]
   & H^1(\DCx^{\mathrm{mov}}_{F,x,M}).
\end{tikzcd}
\]
\end{corollary}
\begin{proof}\leavevmode
\begin{enumerate}[leftmargin=*,align=left]
\item \textbf{Identify the residual for a constant budget.}
The constant budget is represented by $U_r=1$.
Both residual constructions therefore give
$[d^{r+1}]qF(x_r)$ in $Q$.

\item \textbf{Construct the comparison of obstruction classes.}
The inclusion $\operatorname{im}(qDF_x)\subseteq
\operatorname{im}(L^{\mathrm{mov}}_{(x,M)})$ induces the displayed
map of cokernels and proves the assertion.
\end{enumerate}
\end{proof}


\begin{example}[The measured family with a moving budget]
\label{ex:moving-removes-obstruction}
For \cref{ex:weighted-resource-family,ex:second-obstruction}, choose
\[
 M_2=(1+d^2t^{-2})B[d]/d^3.
\]
Then $F(d)=1+d^2t^{-2}\in M_2$. Here $T_B=K/B$ and
$L^{\mathrm{mov}}_{(0,B)}(v,\theta)=-\theta(1)$ is surjective,
so the moving obstruction vanishes. The displayed generator supplies
the feasible second-order budget deformation.
\end{example}

\paragraph{The measured linear budget.}
Multiplication by the full-support germ $g=t^{-1}$ transports
this example to the budget of the doubling program in
\cref{app:complexity-size-example}. In fact
\[
 K/B\xrightarrow{\ \cong\ }K/Bg,\quad [a]\longmapsto[ga],
 \qquad F_g(b)=g(1+b^2).
\]
The map is well-defined and invertible because $g$ is a unit
of $K$, and it carries $B$ onto $Bg$. Thus the fixed
second-order residual is $[gt^{-2}]$, and the moving budget
$(1+d^2t^{-2})Bg[d]/d^3$ contains $F_g(dt^{-1})$.


\subsection{Infinitesimal compatibility of budget diagrams}
\label{sec:flag-deformations}
\begin{proposition}[The tangent of a resource flag]
\label{prop:flag-tangent}
For displayed finite-projective embedded cells $M_0\subseteq M_1$,
write $q_{01}:K/M_0\to K/M_1$. The first-order deformations
of this flag are the pairs
\[
 (\theta_0,\theta_1)\in
 \Hom_B(M_0,K/M_0)\oplus\Hom_B(M_1,K/M_1)
 \quad\text{with}\quad
 q_{01}\theta_0=\theta_1|_{M_0}.
\]
For a finite diagram, impose this equation on every inclusion.
\end{proposition}
\begin{proof}\leavevmode
\begin{enumerate}[leftmargin=*,align=left]
\item \textbf{Derive the first-order inclusion equation.}
Locally lift the classes to $h_i:M_i\to K$ using the exact
finite-dual-basis presentations. The deformed image of $M_0$
lies in that of $M_1$ exactly when
$h_0(m)-h_1(m)\in M_1$ for every $m\in M_0$.
This is the displayed equation.

\item \textbf{Recover the inclusion and descend the condition.}
Its converse constructs the
inclusion by the unique correction in $M_1$; changes of lifts
preserve it. The condition therefore descends and is necessary
and sufficient.

\item \textbf{Extend the calculation to finite diagrams.}
Imposing it for finitely many edges gives the
diagram statement, with commutativity following from the common
ambient marking. This is the formal incidence problem lying
over the classifier of \cref{thm:budget-flags}, not the tangent
of its infinitesimally discrete constant point object.
\end{enumerate}
\end{proof}

\subsection{Complete finite budget-diagram equations}
\label{sec:formal-flags}
A system of resource constraints may include inclusions between
several budgets. Deforming it requires those inclusions to persist.
For a finite diagram $M_i\subseteq M_j$ of principal marked
resources, write $M_i=Bg_i$, put $e_i=e_{g_i}$ and choose local supported
logarithms $z_i(d)$. Use the ambient units
\[
 U_i(d)=(1-e_i)+\exp_{e_i}(z_i(d))\in K[d]/d^{r+2},
 \qquad U_i(0)=1.
\]
Their products with $M_i$ represent the deformations.
Extending by $1-e_i$ makes every inverse below an inverse
in the common ambient coefficient algebra.
An edge remains an inclusion precisely when
\[
 q_j\bigl(U_j(d)^{-1}U_i(d)m\bigr)=0
 \quad(m\in M_i),\qquad q_j:K\to K/M_j.
\]
It is enough to test a displayed finite family of generators.
The coefficient at first order is the existing tangent flag
equation $q_{ij}\theta_i=\theta_j|_{M_i}$.
At degree $r+1$ the new coefficients occur through this same
linear operator, while all lower coefficients determine an
explicit residual.


Intrinsically the finite diagram has the two-term complex
\begin{equation}\label{eq:cat-flag-complex}
 \DCx_{\mathrm{flag}}^\bullet
 =
 \left[
 \bigoplus_i T_{M_i}
 \xrightarrow{\partial_{\mathrm{flag}}}
 \prod_{(i\to j)}
       \Hom_B(M_i,K/M_j)
 \right],
\end{equation}
where
$\partial_{\mathrm{flag}}(\theta)_{ij}
 =q_{ij}\theta_i-\theta_j|_{M_i}$.
The products and sums here are finite.  Each target is an intrinsic
edge quotient module; a displayed generator identifies it with the
corresponding supported quotient used in the coefficient calculation.
The vanishing of one edge component is the commutative square
\[
\begin{tikzcd}[column sep=large,row sep=large]
 M_i \arrow[r,hook] \arrow[d,"{\theta_i}"']
   & M_j \arrow[d,"{\theta_j}"]\\
 K/M_i \arrow[r,"{q_{ij}}"'] & K/M_j .
\end{tikzcd}
\]
Use this differential as $L_{\mathrm{flag}}$ below.
The degree-$(r+1)$ residual on an edge is the $B$-linear map
$m\mapsto[d^{r+1}]q_j(U_j^{-1}U_i m)$, so it has precisely the
target in \eqref{eq:cat-flag-complex}.  Changes of next-order
coefficients act through $\partial_{\mathrm{flag}}$.
Thus $H^0$ organizes compatible tangent directions and $H^1$
the finite-diagram obstructions.


\begin{proposition}[Formal obstruction for a finite budget flag]
\label{prop:formal-flag-obstruction}
A finite diagram of principal embedded-resource inclusions has
a constructed linear first-order compatibility map
$L_{\mathrm{flag}}$ into the finite product of its edge quotient
modules.  Every compatible order-$r$ moving diagram has a
canonical order-$(r+1)$ obstruction in
$\operatorname{coker}L_{\mathrm{flag}}$.  Its vanishing is
internal inhabitation of the compatible lift object; when
inhabited that object is a torsor under
$\ker L_{\mathrm{flag}}$.  Adding polynomial valuation
incidences enlarges the same finite operator by their first-order
rows and gives the analogous obstruction theorem.
\end{proposition}
\begin{proof}\leavevmode
\begin{enumerate}[leftmargin=*,align=left]
\item \textbf{Extract the residual on each inclusion edge.}
For each inclusion edge, conjugate the deformed inclusion by
the target supported unit as displayed above.  Its expression
has zero coefficients through degree $r$.  Extract the next
coefficient.

\item \textbf{Prove independence of representatives.}
Changing logarithmic representatives gives $U'_i=U_i h_i$,
where $h_i$ is an ambient $B[d]$-unit with constant term $1$
and preserves $M_i$. For an edge $i\to j$, commutativity
of the coefficient algebra gives
\[
 q_j\bigl((U'_j)^{-1}U'_i m\bigr)
   =h_j^{-1}h_i\,
      q_j\bigl(U_j^{-1}U_i m\bigr).
\]
The quotient is a $B[d]$-module, and the original expression
has zero coefficients through order $r$.
Since $(h_j^{-1}h_i)(0)=1$, the order-$(r+1)$ residual
is unchanged. This proves descent on both the source and
target representative kernel pairs.

\item \textbf{Identify the flag correction map.}
New coefficients of
source and target units enter linearly and give precisely the
first-order flag-tangent map.

\item \textbf{Assemble the obstruction and lift torsor.}
The finite graph of edges allows
all residual coefficients and maps to be assembled into one
finite product.  The lift equation is therefore
$L_{\mathrm{flag}}(w)=-R_{r+1}$.  Its image, cokernel, and
kernel supply the obstruction, inhabited fibre, and torsor
statements exactly as in the preceding theorem.  Adding
valuation rows uses the same polynomial Taylor calculation.
\end{enumerate}
\end{proof}

The next example combines two measured incidences with a budget
inclusion. Its obstruction will then be compared under changes of the
allowed overheads.

\begin{example}[An obstruction to jointly moving two budgets]
\label{ex:coupled-budget-obstruction}
In the cofinite numerical fibre put $t(n)=1/(n+1)$ and take
\[
 M_0=B\subset M_1=t^{-1}B,\qquad
 F_0(\lambda)=1+t^{-2}\lambda^2,\qquad F_1(\lambda)=t^{-1}.
\]
The first valuation is \cref{ex:weighted-resource-family}; the second
is measured by a program that emits $n+1$ unit-cost idle events and
returns its input. At order one, $\lambda=d$ and the two constant
budgets satisfy both incidences and their inclusion.

Each incidence separately extends to order two: the inner budget may
be $(1+d^2t^{-2})B[d]/d^3$, while the outer budget remains constant.
For a joint lift, put $Q_i=K/M_i$, let $q_{01}:Q_0\to Q_1$, and
write $a=\theta_0(1)\in Q_0$, $b=\theta_1(t^{-1})\in Q_1$ for the
new budget coefficients. A parameter correction $d+d^2v$ has square
$d^2$ modulo $d^3$. Consequently the two incidence equations and the
inclusion equation require
\[
 a=[t^{-2}]_{Q_0},\qquad b=0,\qquad q_{01}a-tb=0.
\]
The last equation uses $\theta_1(1)=t\theta_1(t^{-1})$.
It would imply $t^{-2}\in t^{-1}B$, contrary to the unboundedness
of $t^{-1}$. Thus no joint lift exists.

The full obstruction calculation identifies its intrinsic quotient.
The correction operator is
\[
 L:B\oplus Q_0\oplus Q_1\longrightarrow Q_0\oplus Q_1\oplus Q_1,
 \qquad L(v,a,b)=(-a,-b,q_{01}a-tb).
\]
The surjection $(x,y,z)\mapsto q_{01}x-ty+z$ has kernel
$\im L$: if its value is zero, the triple is $L(0,-x,-y)$.
Hence $\operatorname{coker}L\cong Q_1$, and the residual
$([t^{-2}]_{Q_0},0,0)$ represents
\begin{equation}\label{eq:coupled-budget-obstruction}
 o=[t^{-2}]\in K/t^{-1}B,\qquad \Ann_B(o)=tB.
\end{equation}
Indeed $bo=0$ means $bt^{-2}\in t^{-1}B$, equivalently $b\in tB$.
The obstruction records compatibility of the budget system even
though both individual moving-incidence problems are solvable.
\end{example}


\subsection{Gauge change of infinitesimal resources}
\begin{proposition}[Localization of finite-cell deformation tangents]
\label{prop:gauge-deformation}
For a multiplicative overhead ideal $P$ containing $[1]$
and a finite-projective resource module $M\subseteq K$,
localization induces a canonical isomorphism
\begin{equation}\label{eq:gauge-deformation}
 B_P\otimes_B T_M^{\mathrm{emb}}
 \cong\Hom_{B_P}(B_P\otimes_BM,\,K/B_PM).
\end{equation}
For $M=B$ it sends the tangent quotient $K/B$
to $K/B_P$.  The principal formal-deformation functor
commutes with finite-projective Weil coefficient changes
$A\mapsto B_P\otimes_BA$ at the level of its defining
unit-group quotients.
\end{proposition}
\begin{proof}\leavevmode
\begin{enumerate}[leftmargin=*,align=left]
\item \textbf{Localize the ambient module and normal quotient.}
The coefficient ring $B_P$ is a localization of $B$ at
regular denominators, and $K$ is already the total
quotient ring.  Hence $B_P\otimes_BK\cong K$.
Flat localization sends the quotient $K/M$ to
$K/B_PM$.

\item \textbf{Commute tangent Hom objects with localization.}
A finite-projective module has a finite dual
basis, so its internal hom commutes with scalar extension:
\[
 B_P\otimes_B\Hom_B(M,K/M)
 \cong\Hom_{B_P}(B_P\otimes_BM,K/B_PM).
\]
Apply \cref{thm:embedded-deformation-tangent} to obtain
\eqref{eq:gauge-deformation}.

\item \textbf{Construct the comparison of principal deformation functors.}
For principal cells, tensor
associativity identifies the extended finite Weil algebra
and its nilpotent augmentation ideal.  The embeddings
$B_P\subset K$ and $B\subset B_P$ map the numerator and
denominator unit groups of \eqref{eq:formal-resource-functor}
into their localized counterparts, commuting with all
stage maps and reduction.  This constructs the functorial coefficient-change comparison.
\end{enumerate}
\end{proof}


\begin{theorem}[Overhead localization of resource obstructions]
\label{thm:gauge-obstruction-criterion}
For the fixed, moving, or finite-diagram incidence complexes constructed
above, write $C^\bullet=[E^0\xrightarrow{L}E^1]$. Let
$B_P=S_P^{-1}B$ be an overhead localization from \cref{thm:gauge}.
There are natural identifications
\[
 C_P^\bullet\cong B_P\otimes_BC^\bullet,\qquad
 H^i(C_P^\bullet)\cong B_P\otimes_BH^i(C^\bullet)
 \quad(i=0,1).
\]
At every finite order, the obstruction $o$ of a feasible jet maps to
$o_P=1\otimes o$, and
\begin{equation}\label{eq:overhead-obstruction-criterion}
 o_P=0\quad\Longleftrightarrow\quad
 \exists s\in S_P\; so=0
 \quad\Longleftrightarrow\quad
 S_P\cap\Ann_B(o)\text{ is inhabited}.
\end{equation}
The existential and the lift objects are interpreted internally.
\end{theorem}
\begin{proof}\leavevmode
\begin{enumerate}[leftmargin=*,align=left]
\item \textbf{Compare the localized deformation complexes.}
Localization identifies $B_P\otimes_BK$ with $K$ and carries $K/M$
to $K/B_PM$. Every Hom object in the complexes has a displayed
finite-projective source. Its finite dual basis therefore gives the
scalar-extension identification of \cref{prop:gauge-deformation}.
The finite sums, edge products, derivatives, and evaluation maps in
$L$ commute with these identifications, proving the complex comparison.

\item \textbf{Preserve tangent and obstruction modules.}
Flatness of localization preserves its kernel and cokernel.

\item \textbf{Identify the localized obstruction class.}
Residuals are obtained by finite coefficient extraction from polynomial
valuations and finite supported exponentials. Those operations commute
with scalar extension. Their equality on representative kernel pairs
therefore identifies the localized obstruction with $1\otimes o$.

\item \textbf{Characterize vanishing by denominator annihilation.}
By the fraction-equivalence relation, $o/1=0$ in $S_P^{-1}H^1(C)$
exactly when some $s\in S_P$ annihilates $o$. This gives
\eqref{eq:overhead-obstruction-criterion}; the localized obstruction
theorem identifies its vanishing with the inhabited lift object.
\end{enumerate}
\end{proof}

For \cref{ex:coupled-budget-obstruction}, allowing the polynomial gauge
$B[t^{-1}]$ removes the obstruction, since $to=0$. Polylogarithmic
overheads do not: writing $\lambda(n)=1+\operatorname{bitlength}(n+1)$,
no $\lambda^{-k}$ belongs to $tB$. On $n=2^m$ the ratio
$t^{-1}\lambda^{-k}$ is unbounded, as the binomial lower bound for
$2^m$ of degree $k+1$ shows. Thus the same finite diagram has no lift
after polylogarithmic localization but has a lift after polynomial
localization. The annihilator gives the precise overhead criterion.


The resource deformation problem is now expressed by tangent modules,
incidence complexes, and their successive lift objects. The following
sections construct the geometric function objects and transport
structures associated with these families.

\section{Geometric function objects and differential cohesion}
\label{sec:primary-geometry}
The preceding calculus describes variations of measured resource values
and embedded budgets. To compare functions on their infinitesimal
neighbourhoods with differential forms and transport, we now represent
the coordinate spaces themselves. A nilpotent arrow records both a
family and its reduced boundary. Its two endpoints construct the
adjunctions below, while affine representability identifies its
function objects. Restriction to based probes recovers the coefficient
calculations used in the resource deformation theory.

\subsection{The primary infinitesimal-arrow geometry and its \texorpdfstring{\'etale}{etale} topology}
\label{sec:etale-infinitesimal-site}
A nilpotent thickening has a smaller boundary and a larger
infinitesimal space. Keeping both endpoints produces the primary
geometry. Its restriction to the based stages of \cref{sec:weil}
will recover their resource coefficient calculus.

\begin{definition}[Ordinary algebraic and infinitesimal tests]\label{def:infinitesimal-arrow-category}
Let $\AffB$ be the opposite of the internal category of finitely
presented commutative $B$-algebras.  A presentation is a quotient
$B[x_1,\ldots,x_n]/(f_1,\ldots,f_m)$ by a displayed finite list;
its equations, homomorphisms, and changes of presentation are
formed by internal finite lists, quotients, and algebraic equality.
Write $\operatorname{Spec}A$ for the corresponding object of $\AffB$.
Let $\InfB$ be the internal category whose objects are surjective
maps of those algebras $S\twoheadrightarrow R$ with a nilpotent
kernel, written
\[
 j=(U=\operatorname{Spec}R\hookrightarrow
       T=\operatorname{Spec}S).
\]
Its arrows are \emph{all} commutative squares of the corresponding
affine arrows.  The nilpotence condition is an existential finite
power equation, $\exists n\ge1\;(\ker(S\to R))^n=0$; witnesses enter the arguments by existential elimination.
\end{definition}

Finite polynomial codes and their interpretations define these
categories internally. Surjectivity and nilpotence are bounded power-object predicates. Finite polynomial and finite-projective presentations provide the computationally named tests.

\begin{lemma}[The two affine endpoints]\label{lem:infinitesimal-endpoints}
There are functors $c,d:\InfB\to\AffB$ and
$\iota:\AffB\to\InfB$ with
\[
 c(j)=T,\qquad d(j)=U,\qquad
 \iota(V)=(V\xrightarrow{\id}V),
 \qquad c\dashv\iota\dashv d,\quad c\iota=d\iota=\id.
\]
\end{lemma}
\begin{proof}\leavevmode
\begin{enumerate}[leftmargin=*,align=left]
\item \textbf{Identify maps into the identity arrow.}
A square $j\to\iota(V)$ is uniquely specified by its larger
endpoint map $T\to V$; the smaller map must be its composite with
$U\to T$.

\item \textbf{Identify maps out of the identity arrow.}
A square $\iota(V)\to j$ is uniquely specified by a
map $V\to U$; its larger map is forced.

\item \textbf{Verify naturality and the endpoint identities.}
These give mutually
inverse internal hom-object bijections, natural in both variables;
the two identities follow on identity arrows.
\end{enumerate}
\end{proof}


The two endpoints have different universal properties. Their
adjunctions are expressed by
\[
 \Hom_{\AffB}(c(j),V)\cong\Hom_{\InfB}(j,\iota V),
 \qquad
 \Hom_{\InfB}(\iota V,j)\cong\Hom_{\AffB}(V,d(j)).
\]
Their canonical arrows are the commutative squares
\[
\begin{tikzcd}[column sep=large,row sep=large]
 U \arrow[r,hook] \arrow[d,hook]
   & T \arrow[d,equal]\\
 T \arrow[r,equal] & T
\end{tikzcd}
\qquad
\begin{tikzcd}[column sep=large,row sep=large]
 U \arrow[r,equal] \arrow[d,equal]
   & U \arrow[d,hook]\\
 U \arrow[r,hook] & T .
\end{tikzcd}
\]
These are respectively the arrows $j\to\iota c(j)$ and
$\iota d(j)\to j$. Their opposite directions will determine the
endpoint adjunctions on sheaves.


On $\AffB$, use the \emph{big finitely presented \'etale site}:
a basic covering family consists of a finite presented family of
finitely presented \'etale maps $T_\lambda\to T$ with an effective
jointly covering image.  Such finite families generate the topology
on the finitely presented affine site; covering diagrams are indexed internal families.
On $\InfB$, declare the pullbacks
\[
 U_\lambda=U\times_TT_\lambda\hookrightarrow T_\lambda
\quad(T_\lambda\to T\text{ an \'etale covering})
\]
to be basic covers.

\begin{lemma}[\'Etale lift of finite nilpotent covering diagrams]
\label{lem:etale-nilpotent-lift}
For a finite nilpotent thickening $U\hookrightarrow T$, base
change induces an equivalence between the categories of
finitely presented \'etale $T$-schemes and \'etale $U$-schemes.
It preserves finite intersections and jointly covering families,
and the comparison is natural in commutative squares of thickenings.
Consequently the larger- and smaller-endpoint constructions
carry the stated \'etale covers to \'etale covers.
\end{lemma}
\begin{proof}\leavevmode
\begin{enumerate}[leftmargin=*,align=left]
\item \textbf{Construct lifts across square-zero and nilpotent ideals.}
For a square-zero ideal the standard-\'etale Jacobian
presentation gives uniqueness and existence of lifts: the
invertible Jacobian solves the correction equation by finite
linear algebra.  Iteration over the powers of a nilpotent ideal
provides the finite-order statement.

\item \textbf{Glue local lifts by uniqueness.}
An arbitrary \'etale
presentation is locally standard \'etale; uniqueness on overlaps
and ordinary \'etale descent glue the local lifts and maps.
This is the infinitesimal invariance of the \'etale site
\citep[Theorems 41.15.2 and 41.16.1; Section 59.45]{StacksEtale}.

\item \textbf{Preserve covers, refinements, and descent diagrams.}
The equivalence identifies covering categories, refinements, and \v{C}ech diagrams.  Surjectivity of a family is preserved since
a nilpotent closed immersion induces the same underlying
open-support lattice.  The proof is finite in each nilpotence witness.
\end{enumerate}
\end{proof}

\begin{theorem}[Differential cohesion over the Landau ring]
\label{thm:etale-endpoint-cohesion}
The internal \'etale sheaf categories
\[
 \BaseEt=\ShEt(\AffB;\E),\qquad
 \ThickEt=\ShEt(\InfB;\E)
\]
carry four adjoints
\begin{equation}\label{eq:etale-endpoint-adjunction}
 \iota_!\dashv\iota^*\dashv\iota_*\dashv\iota^!.
\end{equation}
Both endpoint embeddings are fully faithful, and
\[
 (\iota_!F)(U\hookrightarrow T)=F(T),\qquad
 (\iota_*F)(U\hookrightarrow T)=F(U).
\]
The first three functors preserve all internal limits and colimits;
$\iota^!$ preserves all internal limits.  The associated idempotent
modalities
\[
 \mathfrak R=\iota_!\iota^*,\qquad
 \InfShape=\iota_*\iota^*,\qquad
 \mathfrak C=\iota_*\iota^!
\]
obey $\mathfrak R\dashv\InfShape\dashv\mathfrak C$.
The canonical
$\iota_!\to\iota_*$ and $\mathfrak R\to\InfShape$ are not
isomorphisms.
\end{theorem}
\begin{proof}\leavevmode
\begin{enumerate}[leftmargin=*,align=left]
\item \textbf{Identify the presheaf adjoints and full faithfulness.}
On presheaves, left and right Kan extension along
$\iota^{\mathrm{op}}$ are evaluation at $c$ and $d$,
respectively, by \cref{lem:infinitesimal-endpoints}.  Their
restrictions to the identity-arrow objects recover the input
functor, proving full faithfulness.

\item \textbf{Compare the endpoint covering diagrams.}
We spell out the sheafification comparison. At an arrow
$j=(U\hookrightarrow T)$, basic covers and their refinements
are precisely the finite \'etale covering diagrams of $T$.
For the larger endpoint $c$ this is the defining identification.
For the smaller endpoint $d$, \cref{lem:etale-nilpotent-lift}
identifies that category, including its intersections and
refinements, with the corresponding category over $U$.
At an identity arrow, the same diagrams are identity-arrow
covers.

\item \textbf{Commute endpoint functors with sheafification.}
In each case, matching families for a precomposed
presheaf are therefore the corresponding endpoint matching
families. Taking their equivalence classes under refinement
gives natural isomorphisms
\[
 (F\circ c)^+\cong F^+\circ c,\qquad
 (F\circ d)^+\cong F^+\circ d,\qquad
 (H\circ\iota)^+\cong H^+\circ\iota.
\]
 Applying the construction
twice gives the same comparisons for sheafification.
Thus endpoint precomposition and restriction preserve sheaves
and commute with sheafification.

\item \textbf{Verify preservation of limits and colimits.}
Presheaf limits are pointwise,
and sheaf colimits are sheafifications of pointwise colimits;
this proves all the asserted limit and colimit preservation
for $\iota_!,\iota^*,\iota_*$.

\item \textbf{Construct the remaining adjoint.}
The fourth right adjoint is constructed by
the internal end
\[
 (\iota^!H)(T)=
 \operatorname{Nat}_{\ThickEt}(\iota_*a y_T,H),
\]
where $a y_T$ is the sheaf represented by $T$: the internal
Yoneda density formula represents maps out of arbitrary
$\iota_*F$ as maps out of $F$ to this end.  Equalizers and
power objects in the internal sheaf categories construct the
end, including its naturality and adjunction laws.

\item \textbf{Distinguish the two endpoint embeddings.}
Take the affine-line sheaf $F(T)=\Gamma(T,\OO_T)$ and
$j=(\operatorname{Spec}B\hookrightarrow
\operatorname{Spec}(B[\eps]/\eps^2))$.  At $j$ the two
endpoint embeddings have components $B[\eps]/\eps^2$ and
$B$, with canonical map given by reduction.  This is not an
isomorphism, since coefficient extraction sends $\eps$ to
its nonzero coefficient $1$.

\item \textbf{Distinguish the modalities and verify their laws.}
Apply the two associated
modalities to $\iota_!F$ to see that their comparison is
also not an isomorphism.  Finally all adjunction and
idempotency equations follow from the endpoint adjunctions
and their full faithfulness.
\end{enumerate}
\end{proof}

\begin{proposition}[Smooth automorphisms under nilpotent reduction]
\label{prop:smooth-not-rigid}
Over $S=B[\eps]/\eps^2\twoheadrightarrow R=B$,
the distinct smooth $S$-automorphisms of the affine line
$x\mapsto x$ and $x\mapsto x+\eps$ become identical after
reduction.  Thus the functor on categories of smooth schemes
induced by a nilpotent thickening is not faithful; the
proof of \cref{thm:etale-endpoint-cohesion} cannot be
obtained by substituting ``smooth'' for ``\'etale''.
\end{proposition}
\begin{proof}\leavevmode
\begin{enumerate}[leftmargin=*,align=left]
\item \textbf{Exhibit a nontrivial automorphism with identity reduction.}
The inverse of $x\mapsto x+\eps$ is $x\mapsto x-\eps$;
its reduction is the identity.  The two maps differ at
$x$ by the nonzero element $\eps$.  The conclusion follows.
\end{enumerate}
\end{proof}

\subsection{Recovery of algebraic synthetic differential geometry}
\label{sec:primary-sdg}
The primary site gives an affine calculation of function objects.
We use it to show that the coordinate line supports the algebraic
SDG identities needed by the subsequent differential and formal
constructions.
Let $\Delta_{\mathrm{inf}}:\E\to\ThickEt$ be structural inverse
image. For a finitely presented $B$-algebra $A$, write
\[
 F_A(\operatorname{Spec}S)=\Hom_{B\text{-alg}}(A,S),\qquad
 X_A=\iota_!F_A,\qquad
 \mathcal O_{\mathrm{inf}}=X_{B[z]}.
\]
At $j=(\operatorname{Spec}R\hookrightarrow\operatorname{Spec}S)$,
these have values $\Hom_B(A,S)$ and $S$, respectively.
A square from $j$ to $\iota(\operatorname{Spec}A)$ is determined by
its larger-endpoint map, so
$X_A=y\iota(\operatorname{Spec}A)$.  These representables are sheaves:
algebra maps at both endpoints glue along the stated \'etale covers,
and their commutative-square equation is local.  In the tensor formulas
below, scalar extension is along
$\Delta_{\mathrm{inf}}B\to\mathcal O_{\mathrm{inf}}$.

\begin{theorem}[Primary algebraic SDG recovery]
\label{thm:primary-sdg}\label{thm:primary-affine-duality}
For every finitely presented $B$-algebra $A$, canonical evaluation
induces an isomorphism of $\mathcal O_{\mathrm{inf}}$-algebras
\begin{equation}\label{eq:primary-affine-duality}
 \mathcal O_{\mathrm{inf}}
   \otimes_{\Delta_{\mathrm{inf}}B}\Delta_{\mathrm{inf}}A
 \ \xrightarrow{\ \cong\ }\
 \mathcal O_{\mathrm{inf}}^{X_A}.
\end{equation}
Its value at $j$ is $S\otimes_BA$.  Consequently, for every internal
commutative $\mathcal O_{\mathrm{inf}}$-algebra $C$, coordinate
evaluation gives the natural isomorphism
\begin{equation}\label{eq:primary-comprehensive-sdg}
 \underline{\Hom}_{\mathcal O_{\mathrm{inf}}\text{-alg}}
       (\mathcal O_{\mathrm{inf}}^{X_A},C)
 \ \cong\ \operatorname{Spec}_C(A),
\end{equation}
where $\operatorname{Spec}_C(A)$ is the object of solutions in $C$
of a displayed finite presentation of $A$.  In particular, taking
$A=B\otimes_\QQ A_0$ for finitely presented rational $A_0$ proves
the comprehensive algebraic SDG axiom over $\QQ$.

For a rational Weil algebra $W$, put $W_B=B\otimes_\QQ W$ and
$D_W^{\mathrm{all}}=X_{W_B}$.  Then
\begin{equation}\label{eq:primary-weil-kl}
 \mathcal O_{\mathrm{inf}}^{D_W^{\mathrm{all}}}
 \cong\mathcal O_{\mathrm{inf}}\otimes_\QQ W.
\end{equation}
For $W=\QQ[\eps]/\eps^2$, the probe is
$D^{\mathrm{all}}=\{d\in\mathcal O_{\mathrm{inf}}:d^2=0\}$ and
the comparison is the Kock--Lawvere isomorphism
\begin{equation}\label{eq:primary-dual-number-kl}
 \mathcal O_{\mathrm{inf}}\times\mathcal O_{\mathrm{inf}}
 \xrightarrow{\ \cong\ }\mathcal O_{\mathrm{inf}}^{D^{\mathrm{all}}},
 \qquad(a,b)\longmapsto(d\mapsto a+db).
\end{equation}
This infinitesimal object is nontrivial.  The line and every displayed
finite-projective affine module cell are microlinear for rational Weil
diagrams: their augmented Weil-prolongation functors preserve finite
limits in the slice over the cell.  All comparisons are internal,
natural in the displayed algebras and their maps, and uniform in
parameters.  The construction therefore also applies in
$\mathscr S=\Sh_\E(\Cplx)$ with coefficient ring $\bar B$.
\end{theorem}
\begin{proof}\leavevmode
\begin{enumerate}[leftmargin=*,align=left]
\item \textbf{Construct the product with an affine test.}
The product $j\times\iota(\operatorname{Spec}A)$ is the arrow
\[
 \operatorname{Spec}(R\otimes_BA)
 \hookrightarrow\operatorname{Spec}(S\otimes_BA).
\]
Its endpoints are finitely presented.  The displayed quotient is
surjective, and its kernel is the ideal generated by the image of
$\ker(S\to R)$, so it is nilpotent with the same bound.
This is the right-exact scalar-extension formula for a quotient.

\item \textbf{Compute the primary function algebra.}
Products of
represented sheaves and internal Yoneda now give
\[
 \begin{aligned}
 (\mathcal O_{\mathrm{inf}}^{X_A})(j)
 &\cong\underline{\Hom}_{\ThickEt}
        (yj\times y\iota(\operatorname{Spec}A),\mathcal O_{\mathrm{inf}})\\
 &\cong\mathcal O_{\mathrm{inf}}
       (j\times\iota(\operatorname{Spec}A))
 =S\otimes_BA.
 \end{aligned}
\]

\item \textbf{Verify the algebra structure and naturality.}
The evaluation map sends $s\otimes a$ to the function
$\alpha\mapsto s\alpha(a)$, in the corresponding parameter context.
Thus the comparison preserves multiplication and units, and is natural
under test substitution and maps of $A$.

\item \textbf{Establish affine duality.}
If $A=B[x_1,\ldots,x_n]/(f_1,\ldots,f_m)$, the internal
scalar-extension algebra on the left of
\eqref{eq:primary-affine-duality} is the sheafification of the
pointwise coefficient algebras with this presentation over $S$.
The Yoneda calculation has already identified those algebras with
the sheaf $\mathcal O_{\mathrm{inf}}^{X_A}$.  This proves
\eqref{eq:primary-affine-duality}.

\item \textbf{Identify internal algebra maps independently of presentation.}
An algebra map from the
scalar extension to $C$ is exactly a tuple in $C^n$ satisfying
the displayed equations.  The same argument holds in every slice,
so it gives the internal-hom comparison
\eqref{eq:primary-comprehensive-sdg} and its presentation independence.

\item \textbf{Recover the Weil and dual-number coefficient identities.}
Specializing to $W_B$ gives \eqref{eq:primary-weil-kl}.
For dual numbers, the basis $1,\eps$ gives
\eqref{eq:primary-dual-number-kl} with its displayed canonical map.

\item \textbf{Prove nontriviality of the square-zero probe.}
At the dual-number test, the generalized element $\eps$ differs
from zero by coefficient extraction, since $1\ne0$ in $B$.
Thus the square-zero probe is not the zero subobject.

\item \textbf{Compute the probes of a module cell.}
For a displayed finite-projective $B$-module $E$, put
$X_E=X_{\Sym_B(E^\vee)}$.  Finite duality identifies
$X_E(j)=S\otimes_BE$, and the same product calculation gives
\[
 (X_E^{D_W^{\mathrm{all}}})(j)
 \cong(S\otimes_BE)\otimes_\QQ W.
\]

\item \textbf{Prove rational microlinearity and parameter stability.}
Finite limits of augmented rational Weil diagrams are computed
by finite products over $\QQ$ and kernels of linear differences,
retaining augmentation.  Their finite-dimensional rational linear
algebra splits by Gaussian elimination, so tensoring with
$S\otimes_BE$ preserves these augmented limits.  The comparisons
are therefore isomorphisms at each test and in the slice over
$X_E$. All constructions are internal finite presentations,
quotients and universal properties, and hence apply with the
indicated parameters in $\mathscr S$ as well.
\end{enumerate}
\end{proof}

This establishes an algebraic SDG model over the Landau coefficient
ring, in the sense of the comprehensive and Weil axioms of
\citep[Sections I.12, I.16, III.1 and Appendix D]{Kock}.
It supplies the function objects used for formal completion and jets.

\begin{lemma}[Full and zero-boundary infinitesimal probes]
\label{lem:primary-probe-comparison}
For $W_B=B\otimes_\QQ W$, set
\[
 D_W^0=y(\operatorname{Spec}B\hookrightarrow\operatorname{Spec}W_B).
\]
There is a canonical monomorphism
$D_W^0\incl D_W^{\mathrm{all}}$.  It induces isomorphisms
\begin{equation}\label{eq:primary-boundary-kl}
 \mathcal O_{\mathrm{inf}}^{D_W^{\mathrm{all}}}
 \xrightarrow{\ \cong\ }\mathcal O_{\mathrm{inf}}^{D_W^0},
 \qquad
 X_E^{D_W^{\mathrm{all}}}\xrightarrow{\ \cong\ }X_E^{D_W^0}
\end{equation}
for displayed finite-projective affine module cells $X_E$.
The two probes themselves need not be isomorphic.  For dual numbers,
$D^0$ consists of square-zero coordinates whose reduction at the
smaller endpoint is zero.
\end{lemma}
\begin{proof}\leavevmode
\begin{enumerate}[leftmargin=*,align=left]
\item \textbf{Identify the zero-boundary condition.}
At $j$, a $W_B$-point of $S$ belongs to $D_W^0$ precisely when
its composite to $R$ equals the augmentation followed by $B\to R$.
This is a subobject of the full set of $W_B$-points.

\item \textbf{Construct the zero-boundary product test.}
The product of $j$ with the representing zero-boundary arrow has
larger algebra $S\otimes_BW_B$ and smaller algebra $R$.
Its quotient kernel is the sum of the two extended nilpotent ideals,
which is nilpotent by a finite binomial bound.

\item \textbf{Compare the two function objects.}
Yoneda therefore
computes the same larger-endpoint function algebra as in
\cref{thm:primary-sdg}.  The map between the products is the
identity on that larger algebra, proving the first isomorphism;
the identical argument with $S\otimes_BE$ proves the second.

\item \textbf{Distinguish the full and zero-boundary probes.}
At the identity test of $B[\eps]/\eps^2$, the element $\eps$
belongs to the full square-zero probe and has nonzero reduction,
so it does not belong to $D^0$.
\end{enumerate}
\end{proof}

The zero-boundary probes describe displacements relative to the
smaller endpoint of a thickening.  They are the probes used by the
shape relation and its boundary principal parts below.

\subsection{Comparison with the based resource calculus}
\begin{proposition}[Based recovery of the primary coefficient identities]
\label{prop:sdg-based-recovery}
Restriction to $j_C=(\operatorname{Spec}B\hookrightarrow
\operatorname{Spec}C)$, $C\in\Weil_B$, defines a functor
$\operatorname{res}_W:\ThickEt\to\Geo$ with natural identifications
\[
 \operatorname{res}_W\mathcal O_{\mathrm{inf}}\cong\Rline,
 \qquad \operatorname{res}_WX_E\cong j_BE,
 \qquad
 \operatorname{res}_WD_W^{\mathrm{all}}
 \cong\operatorname{res}_WD_W^0\cong D_{C_W}.
\]
For the line and these affine module cells, restriction also identifies
the exponentials in \cref{thm:primary-sdg,lem:primary-probe-comparison}
with the corresponding based Weil exponentials of \cref{thm:weil}.
These comparisons preserve augmentation, Weil substitution and
polynomial coordinate maps.
\end{proposition}
\begin{proof}\leavevmode
\begin{enumerate}[leftmargin=*,align=left]
\item \textbf{Evaluate the primary objects at based tests.}
Evaluation at $j_C$ gives $C$, $C\otimes_BE$ and
$\Hom_B(C_W,C)$, respectively.

\item \textbf{Identify the two probes on based stages.}
Every $B$-algebra map from a
Weil stage to the reduced ring $B$ kills its nilpotent ideal.
The zero-boundary condition is therefore automatic at $j_C$,
identifying the two probes there.

\item \textbf{Recover the coefficient identities and coherence.}
Evaluating their primary
function objects gives $C\otimes_\QQ W$, or
$(C\otimes_BE)\otimes_\QQ W$ for a module cell.  These are
exactly the expressions obtained from \cref{thm:weil} in $\Geo$.
Every map is the same coefficient substitution or elementary-tensor
map, proving the stated coherence.
\end{enumerate}
\end{proof}


\begin{proposition}[Budget-indexed endpoint geometry]
\label{prop:etale-budget-effectivity}
The constructions of \cref{lem:infinitesimal-endpoints,thm:etale-endpoint-cohesion}
may be performed in the budget-diagram category
$\mathscr S=[L,\E]$, yielding endpoint adjunctions on the
associated indexed internal \'etale sheaves.  Their defining
formulas are pointwise in the budget parameter and compatible
with $M_a\incl M_b$.  Transport of an actual name-cover
and tracked operation into this indexed geometry preserves
its equality of denotations, its $\Sigma$-classified domain
and its original positive-observation evaluation.
\end{proposition}
\begin{proof}\leavevmode
\begin{enumerate}[leftmargin=*,align=left]
\item \textbf{Reindex the endpoint geometry over budgets.}
A diagram $a\mapsto M_a$ reindexes every polynomial,
nilpotent-arrow, and endpoint formula using the same ring
$B$; finite algebraic operations and naturality commute
with the inclusions.

\item \textbf{Preserve classified computation and its witnesses.}
The named computation equations are
originally equations of maps and classified spans in $\E$.
Applying the left-exact sheaf realization to their finite
diagrams preserves equality, pullbacks and regular images;
the stored classifier is transported with the span.  Thus
the original $\Sigma$-valued observation remains typed and
its finite witnesses remain tracked.
\end{enumerate}
\end{proof}


The primary function objects now supply coordinate algebras for the
finite neighbourhoods used above. Completing their diagonal ideals
assembles the neighbourhoods of every order.

\section{Formal completion of differential resource geometry}
\label{sec:formal-completion}

The finite calculus measures resource variation to a specified order.
Completion assembles compatible orders and exposes a further resource
question: whether one budget contains every coefficient of the assembled
family. The measured tower in \cref{ex:measured-formal-family} has a
polynomial budget at each finite order and no common polynomial growth
degree. It will also supply the de Rham and crystalline example.

We construct formal neighbourhoods by powers of the diagonal ideal,
first over $B$ in $\E$ and then over the resource parameters in
$\mathscr S=\Sh_\E(\Cplx)$. Principal parts retain both endpoints and
the full Taylor action of coordinate changes. Boundary principal parts
represent pairs with equal smaller-endpoint value on the primary site;
their inverse-limit function objects identify the completed infinitesimal
nerve. Effective finite observations are then obtained from the same
coefficient tower.

\subsection{The diagonal and its finite principal parts}
\label{sec:formal-diagonal}
An order-$r$ neighbourhood retains differences between two points
through degree $r$. Its coordinate algebra is obtained by quotienting
the corresponding power of the diagonal ideal.
Let $E$ be a displayed finite-projective $B$-module, let $V=E^\vee$,
and write $A_E=\Sym_B(V)$ for its coordinate algebra, with the
variance of \cref{prop:dual-coordinates}. Put
\[
 D_E=A_E\otimes_BA_E,\qquad
 \mu_E:D_E\longrightarrow A_E,\qquad
 I_E=\ker\mu_E.
\]
Both tensor-factor inclusions will be retained; they give the source
and target actions on every principal-parts object.

\begin{definition}[Formal diagonal filtration]
\label{def:formal-principal-parts}
For $r\in\N$ construct the $A_E$-algebra
\[
 P^r_{E/B}=D_E/I_E^{r+1}.
\]
The transition maps $P^{r+1}_{E/B}\onto P^r_{E/B}$ define the
pro-algebra $\boldsymbol P_E=(P^r_{E/B})_r$ and its internal
inverse-limit realization
\[
 \widehat P_{E/B}=\varprojlim_rP^r_{E/B}.
\]
The pro-algebra is retained along with its realization.  A
$B$-algebra map $P^r_{E/B}\to C$ is an \emph{order-$r$ formal
neighbourhood point} of the diagonal.
\end{definition}

\begin{theorem}[Polynomial calculation of formal principal parts]
\label{thm:formal-principal-parts}
There are natural filtered $A_E$-algebra isomorphisms
\begin{equation}\label{eq:formal-principal-parts}
 \begin{aligned}
  P^r_{E/B}&\cong
    \bigoplus_{0\le j\le r}A_E\otimes_B\Sym_B^j(V),\\
  \widehat P_{E/B}&\cong
    \prod_{j\ge0}A_E\otimes_B\Sym_B^j(V),\\
  \gr_{I_E}(D_E)&\cong
    \Sym_{A_E}(\Omega^1_{A_E/B}),
    \qquad I_E/I_E^2\cong\Omega^1_{A_E/B}.
 \end{aligned}
\end{equation}
Each finite degree is a finite-projective $A_E$-module; multiplication
on the first line is truncated by total degree and on the second is
coefficientwise finite convolution.  The filtered isomorphisms are
natural under polynomial coordinate substitution with its full Taylor
transport: for $y=p(x)$ the displacement is sent to
$p(x+u)-p(x)$.  On the associated graded this transport is the
symmetric algebra of the cotangent map.  All comparisons also commute
with changes of displayed finite-projective presentation.
\end{theorem}
\begin{proof}\leavevmode
\begin{enumerate}[leftmargin=*,align=left]
\item \textbf{Introduce displacement coordinates.}
The symmetric-algebra universal property gives
$D_E\cong\Sym_B(V\oplus V)$.  Apply the invertible linear
transformation $(v_0,v_1)\mapsto(v_0,v_1-v_0)$, with inverse
$(v_0,d)\mapsto(v_0,v_0+d)$.  In the resulting coordinates,
$\mu_E$ is evaluation at $d=0$.  Hence
\[
 D_E\cong A_E\otimes_B\Sym_B(V),\qquad
 I_E^k=\bigoplus_{j\ge k}A_E\otimes_B\Sym_B^j(V).
\]

\item \textbf{Compute principal parts and finite-projective coefficients.}
The quotient and product formulas follow; each multiplication
coefficient uses finitely many summands.  By the finite-free retract
for $V$, every tensor power $V^{\otimes j}$ is finite-projective.
Since $j!$ is invertible in the rational ring $B$, the averaging
idempotent $\frac1{j!}\sum_{\sigma\in S_j}\sigma$ exhibits
$\Sym_B^j(V)$ as a finite-projective retract of $V^{\otimes j}$.

\item \textbf{Identify the cotangent module and associated graded algebra.}
The free-algebra derivation formula
$\Omega^1_{A_E/B}\cong A_E\otimes_BV$ identifies the first
quotient $I_E/I_E^2$ with the cotangent module, and identifies the
whole graded multiplication with its symmetric algebra.

\item \textbf{Transport linear substitutions and projective retractions.}
The linear changes and symmetric-power maps commute with displayed
linear substitutions and retractions.

\item \textbf{Compute polynomial transport in displacement coordinates.}
For a polynomial map $p$,
the two-vertex algebra map sends $(y_0,y_1)$ to
$(p(x_0),p(x_1))$.  In displacement coordinates this is exactly
\begin{equation}\label{eq:principal-parts-taylor-naturality}
 y=p(x),\qquad v=p(x+u)-p(x).
\end{equation}

\item \textbf{Preserve the displacement filtration.}
The second expression has zero constant term in $u$, so substitution
preserves every power of the displacement ideal and defines compatible
maps on finite quotients and their inverse limits.

\item \textbf{Verify composition and identity coherence.}
For composable
polynomial maps $p,q$, the displacement in the composite is
\[
 q\bigl(p(x)+[p(x+u)-p(x)]\bigr)-q(p(x))
       =q(p(x+u))-q(p(x)),
\]
which proves composition coherence; the identity map gives the unit
equation.

\item \textbf{Identify the associated-graded transport.}
Its degree-one part is $dp_x(u)$, so the associated-graded
map is the symmetric algebra of the induced cotangent map.  These
formulas are expressions of the two-vertex universal property and
therefore descend along every displayed projective retraction.

\item \textbf{Distinguish filtered transport from its linear part.}
For example $p(x)=x^2$ sends $v$ to $2xu+u^2$; retaining only
$2xu$ gives the associated-graded transport, not the filtered one.
This proves polynomial naturality with its specified Taylor action.
\end{enumerate}
\end{proof}


The coordinate diagram below fixes the two actions and reduction.
\begin{equation}\label{eq:diagonal-diagram}
\begin{tikzcd}[column sep=large]
 A_E \arrow[r,shift left=.6ex,"{s^\sharp}"]
     \arrow[r,shift right=.6ex,swap,"{t^\sharp}"]
 & A_E\otimes_BA_E \arrow[r,two heads] \arrow[dr,swap,"{\mu_E}"]
 & P^r_{E/B} \arrow[d,"{\operatorname{aug}_r}"]\\
 && A_E .
\end{tikzcd}
\end{equation}
The augmentation ideal of $P^r_{E/B}$ has $(r+1)$-st power zero.
The quotient $P^0_{E/B}=A_E$, so this filtration begins at the
ordinary resource-coordinate space and adds no independently
specified infinitesimals.

\subsection{Finite formal neighbourhoods on based Weil probes}
\label{sec:formal-realization}
Principal parts describe a neighbourhood over the whole polynomial
cell: their reduction is $A_E$. They need not be finite-projective
over $B$ and hence need not be $B$-Weil stages themselves.
We realize them by evaluating their functor of points on the
already constructed category $\Weil_B^{\mathrm{proj}}$:
\[
 \mathcal D_E^{\le r}(C)
   =\Hom_{B\text{-alg}}(P^r_{E/B},C),
   \qquad C\in\Weil_B^{\mathrm{proj}}.
\]
The quotient map makes this a subobject of $j_BE\times j_BE$.

\begin{theorem}[Completion--Weil comparison]
\label{thm:completion-weil}
There is a canonical isomorphism in $\Geo$,
\begin{equation}\label{eq:completion-weil}
 \varinjlim_r\mathcal D_E^{\le r}
 \ \cong\ j_BE\times_{\Shape(j_BE)}j_BE.
\end{equation}
It is natural in finite-projective polynomial coordinate maps and
is compatible with the source and target projections.  For every
finite ordinal $[q]$, completion of the higher diagonal
$A_E^{\otimes_B(q+1)}\to A_E$ gives a functor of points
naturally identified with the $(q+1)$-fold fibre product of
$j_BE$ over $\Shape(j_BE)$.  These isomorphisms respect every
simplicial operator.
\end{theorem}
\begin{proof}\leavevmode
\begin{enumerate}[leftmargin=*,align=left]
\item \textbf{Factor equal-reduction pairs through finite principal parts.}
Fix a stage $C=B\oplus N$, with the supplied finite nilpotence
bound $N^{s+1}=0$.  By finite duality, a pair of points of $j_BE(C)$
is a pair of $B$-algebra maps $A_E\to C$.  It belongs to the
right-hand side of \eqref{eq:completion-weil} precisely when its
two reductions to $B$ agree.  Equivalently the difference of the
two evaluations of every $v\in V$ belongs to $N$.  The induced map
$A_E\otimes_BA_E\to C$ therefore sends $I_E$ to $N$ and kills
$I_E^{s+1}$; it factors through $P^s_{E/B}$.

\item \textbf{Recover equal reductions from finite factorization.}
Conversely, if the map factors through $P^r_{E/B}$, the images of
the coordinate differences are nilpotent.  Their reductions in $B$
are zero because $B$ is reduced, by
\cref{prop:signed-coefficients}.  Duality then implies equality
of the reduced $E$-points.  The implications hold internally,
uniformly in $C$ and its nilpotence datum.  They identify the
subobjects at every stage, and hence the corresponding functors.

\item \textbf{Extend the comparison to higher diagonals.}
For $q+1$ points use differences with the first point, or
successive differences; both generate the kernel of the higher
multiplication map.  The same finite nilpotence argument identifies
its completed formal points with points having equal reductions.

\item \textbf{Verify simplicial coherence.}
Deleting, repeating, or permuting vertices acts by substitution of
these differences and preserves all powers of the higher diagonal
ideal.  The stagewise identifications therefore commute with faces,
degeneracies, and permutations.
\end{enumerate}
\end{proof}


\subsection{Formal completion in the primary geometry}
\label{sec:formal-etale-completion}
The based calculation records displacements at fixed-base probes.
We now identify the full relation on the primary site, where both
endpoints of a thickening remain visible. For
$A=\Sym_B(E^\vee)$, retain the affine sheaf
$F_A\in\BaseEt$ of \cref{sec:primary-sdg}, with
$F_A(\operatorname{Spec}S)=\Hom_{B\text{-alg}}(A,S)$, and put
\[
 X_A=\iota_!F_A\in\ThickEt,
 \qquad\InfShape X_A=\iota_*F_A.
\]
The shape unit $\eta_{X_A}:X_A\to\InfShape X_A$
evaluated at $j=(\operatorname{Spec}R\hookrightarrow
\operatorname{Spec}S)$ is the restriction map
$\Hom_B(A,S)\to\Hom_B(A,R)$.

\begin{theorem}[\'Etale completion of the full infinitesimal relation]
\label{thm:completion-etale}
For each finite order $r$, define a subobject of
$X_A\times X_A$ by
\begin{equation}\label{eq:boundary-principal-parts}
 \mathcal D_{A,\mathrm{et}}^{\le r}(j)
 =\Hom_{B\text{-alg}}(P^r_{A/B},S)
  \mathop{\times}_{\Hom_{B\text{-alg}}(A\otimes_BA,R)}
      \Hom_{B\text{-alg}}(A,R),
\end{equation}
where the last arrow is the diagonal and the first arrow
reduces both vertices modulo $S\to R$.  These are \'etale
sheaves, they form an increasing union, and the comparison
\begin{equation}\label{eq:completion-etale}
 \varinjlim_r\mathcal D_{A,\mathrm{et}}^{\le r}
 \ \xrightarrow{\ \cong\ }\
 \mathcal D_{X_A}:=
 X_A\times_{\InfShape X_A}X_A
\end{equation}
is an isomorphism of internal \'etale sheaves.  It is natural
under maps of displayed polynomial coordinate algebras,
retains both vertex maps, and identifies completed higher
diagonals with the complete infinitesimal \v{C}ech nerve of
the shape unit, including all face and degeneracy maps.

Restriction to the based probe
$\operatorname{Spec}B\hookrightarrow\operatorname{Spec}C$,
$C\in\Weil_B^{\mathrm{proj}}$, identifies
\eqref{eq:boundary-principal-parts} with the previously
constructed $\Hom_B(P^r_{A/B},C)$, and
\eqref{eq:completion-etale} with \eqref{eq:completion-weil}.
\end{theorem}
\begin{proof}\leavevmode
\begin{enumerate}[leftmargin=*,align=left]
\item \textbf{Factor equal-boundary pairs through finite principal parts.}
At $j$, a pair of $A$-points of $S$ lies in the right side
exactly when its reductions in $R$ agree.  Its algebra map
$A\otimes_BA\to S$ therefore takes the diagonal ideal
$I$ into $J=\ker(S\to R)$.  For a nilpotence witness
$J^N=0$, it annihilates $I^N$, so the pair comes from
$P^{N-1}$.

\item \textbf{Recover boundary equality and identify the sheaves.}
Conversely every element of
\eqref{eq:boundary-principal-parts} has equal reduced
vertices by its displayed fibre-product equation.  There is
no chosen global $N$: these implications prove equality of
subobjects by existential elimination in $\E$.  Pullback of
nilpotent quotients along \'etale maps preserves their power
filtration, so the formulas define sheaves and the stagewise
isomorphism is natural in coverings and refinement.

\item \textbf{Extend the comparison to the completed nerve.}
For $q+1$ vertices use the kernel of
$A^{\otimes_B(q+1)}\to A$.  Equality of all reductions
places every difference into $J$, and the same nilpotence
argument applies.  Face, degeneracy, and permutation maps
are algebra substitutions preserving powers of this ideal;
therefore the comparison is simplicial.

\item \textbf{Recover the based-Weil formula.}
On the based arrow $B\hookrightarrow C=B\oplus N$,
the image of every difference under a $P^r$-point is
nilpotent.  The reduced ring $B$ has no nonzero nilpotents,
so the boundary equality is automatic. These identities give the
based-Weil formula directly at every stage.
\end{enumerate}
\end{proof}


The completed relation is the kernel pair of the primary shape unit:
\[
\begin{tikzcd}[column sep=large,row sep=large]
 \mathcal D_{X_A} \arrow[r,"{t}"] \arrow[d,"{s}"']
   & X_A \arrow[d,"{\eta_{X_A}}"]\\
 X_A \arrow[r,"{\eta_{X_A}}"'] & \InfShape X_A .
\end{tikzcd}
\]
The theorem identifies its finite boundary filtration with
$\varinjlim_r\mathcal D_{A,\mathrm{et}}^{\le r}$ inside this
cartesian square.  Higher fibre products give its full nerve.


\subsection{Represented principal parts and synthetic formal functions}
\label{sec:primary-formal-functions}
Identifying a formal relation is not yet a calculation of its
functions. Representability supplies that calculation: finite
neighbourhoods are primary tests, so Yoneda evaluates their function
objects; compatible truncations then give the functions on the
formal union. Retain the coordinate-line ring object
\[
 \mathcal O_{\mathrm{inf}}=\iota_!F_{B[z]},\qquad
 \mathcal O_{\mathrm{inf}}(\operatorname{Spec}R\hookrightarrow
                         \operatorname{Spec}S)=S.
\]
For $q\ge0$ put
\[
 P_q^r=A^{\otimes_B(q+1)}/I_q^{r+1},\qquad
 I_q=\ker(A^{\otimes_B(q+1)}\to A),\qquad
 k_{q,r}=(\operatorname{Spec}A\hookrightarrow\operatorname{Spec}P_q^r).
\]
The finite-projective polynomial presentation supplies finitely many
generators for $I_q$ and for each of its finite powers.  Thus $P_q^r$
is finitely presented over $B$, and $k_{q,r}$ is an actual object
of $\InfB$.  Write $\mathcal D_{A,q}^{\le r}$ for the higher-vertex
version of \eqref{eq:boundary-principal-parts}, and
$\mathcal D_{X_A}^{(q)}=X_A\times_{\InfShape X_A}\cdots
\times_{\InfShape X_A}X_A$ with $q+1$ factors.

\begin{lemma}[Primary representability of finite boundary principal parts]
\label{lem:primary-principal-representability}
There are natural isomorphisms of primary sheaves
\[
 yk_{q,r}\cong\mathcal D_{A,q}^{\le r},\qquad
 \varinjlim_r yk_{q,r}\cong\mathcal D_{X_A}^{(q)}.
\]
They preserve all vertex maps, truncations, and polynomial coordinate
substitutions.  The representing sheaves are already \'etale sheaves.
\end{lemma}
\begin{proof}\leavevmode
\begin{enumerate}[leftmargin=*,align=left]
\item \textbf{Identify the representing boundary square.}
A square from $j=(\operatorname{Spec}R\hookrightarrow\operatorname{Spec}S)$
to $k_{q,r}$ is precisely a pair of algebra maps
$P_q^r\to S$, $A\to R$ whose composites from
$A^{\otimes_B(q+1)}$ agree.  This is the displayed boundary
fibre product, including equality of all reduced vertices.

\item \textbf{Verify the sheaf condition.}
Algebra
maps into an \'etale covering diagram glue uniquely at both endpoints;
the square equation glues because equality is local.  Hence the
representable is a sheaf.

\item \textbf{Identify the completed union and coherence.}
The union assertion is the higher-diagonal
case of \cref{thm:completion-etale}.  Every comparison is induced by
the same algebra substitution on the two endpoints, which proves
the stated coherence.
\end{enumerate}
\end{proof}

\begin{theorem}[Synthetic functions on the completed primary nerve]
\label{thm:primary-synthetic-cochains}
Let $\underline{\Hom}_{\ThickEt}$ denote the internal $\E$-object
of sheaf maps, and define the ring objects
\[
 \mathcal C_r^q=\mathcal O_{\mathrm{inf}}^{\mathcal D_{A,q}^{\le r}},
 \qquad
 \widehat{\mathcal C}^{\,q}
       =\mathcal O_{\mathrm{inf}}^{\mathcal D_{X_A}^{(q)}}.
\]
Then
\begin{align}
 \underline{\Hom}_{\ThickEt}
   (\mathcal D_{A,q}^{\le r},\mathcal O_{\mathrm{inf}})
   &\cong P_q^r,\label{eq:primary-finite-functions}\\
 \underline{\Hom}_{\ThickEt}
   (\mathcal D_{X_A}^{(q)},\mathcal O_{\mathrm{inf}})
   &\cong\varprojlim_rP_q^r,\label{eq:primary-formal-functions}\\
 \widehat{\mathcal C}^{\,q}&\cong\varprojlim_r\mathcal C_r^q.
 \label{eq:primary-complete-coefficient-sheaf}
\end{align}
At $j=(\operatorname{Spec}R\hookrightarrow\operatorname{Spec}S)$,
these coefficient sheaves have the explicit values
\[
 \mathcal C_r^q(j)\cong S\otimes_BP_q^r,\qquad
 \widehat{\mathcal C}^{\,q}(j)
       \cong\varprojlim_r(S\otimes_BP_q^r).
\]
All comparisons respect the cosimplicial maps, algebra operations,
and polynomial coordinate substitutions with their full Taylor law.
\end{theorem}
\begin{proof}\leavevmode
\begin{enumerate}[leftmargin=*,align=left]
\item \textbf{Compute finite and completed internal function objects.}
Internal Yoneda and \cref{lem:primary-principal-representability} give
$\underline{\Hom}(yk_{q,r},\mathcal O_{\mathrm{inf}})
=\mathcal O_{\mathrm{inf}}(k_{q,r})=P_q^r$.
Mapping out of the internally indexed union gives the inverse-limit
formula, uniformly in any additional $\E$-parameter.

\item \textbf{Compute the completed coefficient sheaf.}
In the sheaf
topos, product with a fixed object preserves colimits.  Applying this
fact to the exponential adjunction proves
\eqref{eq:primary-complete-coefficient-sheaf}.

\item \textbf{Construct the product test.}
The product $j\times k_{q,r}$ has larger-endpoint algebra
$S\otimes_BP_q^r$ and smaller-endpoint algebra $R\otimes_BA$.
The kernel of its quotient is generated by the two extended
nilpotent kernels, whose sum is nilpotent by a finite binomial bound.
Both endpoints are finitely presented, so this is a product in
$\InfB$.

\item \textbf{Compute stagewise coefficients and their limit.}
Yoneda therefore calculates
\[
 (\mathcal O_{\mathrm{inf}}^{yk_{q,r}})(j)
 =\underline{\Hom}(yj\times yk_{q,r},\mathcal O_{\mathrm{inf}})
 =S\otimes_BP_q^r.
\]
Sheaf limits are computed pointwise, giving the completed value.

\item \textbf{Recover internal maps and verify coherence.}
At the terminal test $\iota(\operatorname{Spec}B)$ this recovers
the two internal map objects above.  Algebra substitutions represent
the vertex and coordinate maps, so Yoneda identifies every associated
coherence diagram, including multiplication of functions.
\end{enumerate}
\end{proof}


Restriction of functions along a finite-neighbourhood inclusion is
the truncation map under Yoneda:
\[
\begin{tikzcd}[column sep=large,row sep=large]
 {\underline{\Hom}(yk_{q,r+1},\mathcal O_{\mathrm{inf}})}
   \arrow[r] \arrow[d,"{\cong}"']
   & {\underline{\Hom}(yk_{q,r},\mathcal O_{\mathrm{inf}})}
       \arrow[d,"{\cong}"]\\
 P_q^{r+1} \arrow[r,two heads] & P_q^r .
\end{tikzcd}
\]
Here the hom objects are those of $\ThickEt$ internal to $\E$.
Taking the inverse limit of this natural diagram gives the completed
function comparison, including its cosimplicial structure.


The function calculation must also respect the endpoint structure.
We next compare the endpoint functors with inverse limits of the
coefficient sheaves, including the units, counits, and maps of the
infinitesimal nerve.

\subsection{The endpoint modalities through formal completion}
\label{sec:modal-formal-completion}
\begin{theorem}[Endpoint adjunctions on formal towers and their realizations]
\label{thm:modal-formal-completion}
The endpoint quadruple of \cref{thm:etale-endpoint-cohesion}
acts levelwise on internally $\N^{\mathrm{op}}$-indexed towers,
with the same four adjunctions, full-faithfulness equations, and
idempotent modalities.  Every endpoint functor commutes with
realization of these towers by internal inverse limits.  In particular
\begin{equation}\label{eq:modal-completion-comparison}
 \mathfrak M(\varprojlim_rH_r)
 \cong\varprojlim_r\mathfrak M(H_r),
 \qquad
 \mathfrak M\in\{\mathfrak R,\InfShape,\mathfrak C\}.
\end{equation}
These comparisons commute with all adjunction units and counits,
the modality multiplication and comultiplication, and maps of towers.
They apply to the completed coefficient sheaves
$\widehat{\mathcal C}^{\,q}$ of
\cref{thm:primary-synthetic-cochains}, including their cosimplicial
operations and normalized additive cochains.  Rings and modules
are transported over the corresponding transported coefficient rings.
\end{theorem}
\begin{proof}\leavevmode
\begin{enumerate}[leftmargin=*,align=left]
\item \textbf{Prove limit preservation for the endpoint functors.}
Limits of internal sheaves are computed pointwise.  The endpoint
formulas identify $\iota_!$, $\iota^*$, and $\iota_*$ with
precomposition on sheaves, so each commutes with internally indexed
limits.

\item \textbf{Prove limit preservation for the end-defined adjoint.}
For the fourth functor, its constructed end gives directly
\[
 (\iota^!\varprojlim_rH_r)(T)
 =\underline{\Hom}(\iota_*ay_T,\varprojlim_rH_r)
 \cong\varprojlim_r\underline{\Hom}(\iota_*ay_T,H_r).
\]
This is the limit-preservation equation for $\iota^!$, natural in $T$.

\item \textbf{Lift the adjunctions to towers.}
The componentwise adjunction bijections carry natural transformations
of towers to natural transformations of towers, by naturality of the
unit and counit.  They therefore construct the four tower adjunctions;
their triangle and full-faithfulness equations are componentwise.
Composing the four limit comparisons proves
\eqref{eq:modal-completion-comparison}.

\item \textbf{Verify compatibility with modal structure maps.}
For each unit, counit, or modality structure map, both paths through
the claimed comparison have the same composites with every limit
projection.  The limit universal property proves their equality.
The same argument applies to any map of towers.

\item \textbf{Preserve algebraic and normalized cochain structures.}
Ring operations are
finite-product diagrams, and additive normalization is a finite
intersection of degeneracy kernels in each cochain degree.  All four
functors preserve those limits, so the comparisons preserve these
structures as well.  Apply the result to
\eqref{eq:primary-complete-coefficient-sheaf}; the identities for
cofaces, degeneracies, and coordinate changes are identities of maps
of the finite towers and hence persist in their realizations.
\end{enumerate}
\end{proof}


For the shape modality, compatibility with the unit is the square
\[
\begin{tikzcd}[column sep=large,row sep=large]
 \varprojlim_r H_r \arrow[r,"{\eta}"] \arrow[d,equal]
   & \InfShape(\varprojlim_rH_r) \arrow[d,"{\cong}"]\\
 \varprojlim_r H_r \arrow[r,"{{\varprojlim_r\eta_{H_r}}}"']
   & \varprojlim_r\InfShape H_r .
\end{tikzcd}
\]
The analogous counit squares and their mates express the compatibility
for the other endpoint modalities.


There are two related limiting operations here. Formal spaces are
unions of finite represented neighbourhoods, identified in
\cref{thm:completion-etale}; their function objects are inverse
limits. The theorem applies to the inverse limits of coefficient sheaves.

\begin{proposition}[Completion of geometric coefficients]
\label{prop:geometric-coefficient-completion}
For an internal coefficient tower $P_r$, scalar extension has a
canonical comparison
\[
 S\otimes_B\varprojlim_rP_r
 \longrightarrow\varprojlim_r(S\otimes_BP_r).
\]
It is not always an isomorphism, even for a finitely presented
polynomial algebra $S$.  The completed coefficient sheaf in
\cref{thm:primary-synthetic-cochains} is the right-hand construction.
\end{proposition}
\begin{proof}\leavevmode
\begin{enumerate}[leftmargin=*,align=left]
\item \textbf{Construct the coefficient-completion comparison.}
Tensor the limit projections and use the target limit universal
property to construct the comparison.

\item \textbf{Exhibit the polynomial-series comparison.}
Take $S=B[z]$ and
$P_r=B[u]/u^{r+1}$.  It becomes
\[
 B[[u]][z]\longrightarrow B[z][[u]].
\]

\item \textbf{Prove injectivity and failure of surjectivity.}
It is injective by coefficient extraction.  Every element of the
source has a finite bound on its $z$-degree, whereas
$\sum_{m\ge0}z^mu^m$ has none: its coefficient at any larger
$u$-degree has that larger $z$-degree, and $1\ne0$ in $B$.
Thus the comparison is not surjective.

\item \textbf{Recover the stagewise sheaf calculation.}
The sheaf formula is the pointwise calculation already proved.
\end{enumerate}
\end{proof}

\begin{proposition}[Polynomial lifting, effective descent, and a singular counterexample]
\label{prop:shape-lifting-scope}
For a displayed finite-projective $E$, the unit
$\eta_{X_A}:X_A\to\InfShape X_A$ is an effective
(epimorphic) descent morphism of $\ThickEt$.  This
assertion fails for general finitely presented affine
algebras: the unit of
$X=\iota_!\operatorname{Spec}(B[x]/x^2)$ is not an
epimorphism.
\end{proposition}
\begin{proof}\leavevmode
\begin{enumerate}[leftmargin=*,align=left]
\item \textbf{Prove effective lifting for finite-projective affine cells.}
The functor of points of $A=\Sym_B(E^\vee)$ is
$S\mapsto E\otimes_BS$.  A displayed finite-free retract
for $E$ exhibits
$E\otimes_BS\to E\otimes_BR$ as a retract
of a finite power of the surjection $S\to R$; hence it
is a regular epimorphism in $\E$, uniformly in $j$.
The sheaf unit is thus locally surjective, hence effective
epic.  In an elementary topos every epimorphism is an
effective descent morphism for objects, by the regular
kernel-pair and descent construction.

\item \textbf{Construct the singular lifting problem.}
For the counterexample use
$S=B[t]/(t^3)\twoheadrightarrow R=B[t]/(t^2)$.
The element $t\in R$ is a point of $\operatorname{Spec}(B[x]/x^2)$.

\item \textbf{Show that the lift fails locally.}
Any lift to $S$ would have the form $t+at^2$, but
$(t+at^2)^2=t^2\ne0$ in $S$; the equation cannot hold
even after a faithfully flat \'etale extension of $S$.
Thus the associated component of the unit is not locally
surjective.  The obstruction is algebraic, not a failure
of an implementation name.
\end{enumerate}
\end{proof}

\begin{proposition}[A natural nonpolynomial Weil transformation]
\label{prop:nonpolynomial-reduction}\label{rem:algebraic-vs-full}
The reduction--constant-extension transformation
$R=\epsilon_{\Rline}\eta_{\Rline}:\Rline\to\Rline$
is natural in projective Weil stages, satisfies $R_B=\id_B$,
and annihilates every nilpotent coefficient.  It is not represented
by any polynomial in $B[x]$.  Hence the algebraic higher-diagonal
cochains form a proper subcalculus of all maps of Weil functors.
\end{proposition}
\begin{proof}\leavevmode
\begin{enumerate}[leftmargin=*,align=left]
\item \textbf{Verify naturality of reduction.}
At a stage $C=B\oplus N$ the map is $R_C(b+n)=b$ under
constant extension.  Naturality is the equation
$u\eta_C=\eta_{C'}u$ for augmented stage morphisms $u:C\to C'$.

\item \textbf{Force a polynomial representative to be the identity.}
Suppose a polynomial $p(x)=\sum_{j=0}^m b_jx^j$ induced $R$.
Evaluating at $C=B$ yields $p(b)=b$ for all $b\in B$, in
particular at $m+1$ distinct rational constants.  Subtracting
$x$ gives a polynomial of degree at most $\max(m,1)$; evaluating
at one more rational constant if necessary supplies that many
more than its degree.  The Vandermonde determinant is a product
of nonzero rational differences and thus a unit of $B$, so all
coefficients of $p-x$ vanish.  Hence $p=x$.

\item \textbf{Contradict polynomiality on a nilpotent coefficient.}
At the existing
dual-number stage $B[\varepsilon]/\varepsilon^2$, however,
$R(\varepsilon)=0$ while $p(\varepsilon)=\varepsilon\ne0$.
The inequality holds because $1,\varepsilon$ are independent
components of the underlying free $B$-module and $1\ne0$ in
the rational-germ ring.  This is a contradiction.

\item \textbf{Place the example in the two geometric settings.}
The map
$R$ is therefore a nonpolynomial morphism of the based-Weil geometry
$\Geo$. On the primary site,
\cref{thm:primary-synthetic-cochains} identifies the completed
coordinate cochains with all functions on the represented formal
nerve with values in $\mathcal O_{\mathrm{inf}}$.
\end{enumerate}
\end{proof}

At each finite order, the coefficient constructions commute with
the stated gauge and input changes. Passing to an infinite
coefficient object raises a further question: can one denominator
serve all orders? The next comparison records precisely where
finite-order transport and completed transport differ.

\subsection{Formal-order coherence, gauges, and exact coefficients}
\label{sec:formal-coherence}
\begin{proposition}[Finite-order coefficient change]
\label{prop:formal-base-change}
For a localization $B\to B_P$ constructed by an overhead gauge,
put $A_P=B_P\otimes_BA_E$.  For every $r$, the canonical map
\[
 B_P\otimes_BP^r_{E/B}
 \xrightarrow{\ \cong\ }P^r_{A_P/B_P}
\]
is an isomorphism of augmented algebras.  It commutes with the
source and target actions, the differential comparison, truncation,
and all higher-diagonal substitution maps.
\end{proposition}
\begin{proof}\leavevmode
\begin{enumerate}[leftmargin=*,align=left]
\item \textbf{Localize the diagonal ideal and its powers.}
Localization is flat by \cref{thm:gauge}.  It preserves the
kernel of multiplication $\mu_E$, and hence carries its ideal to
the kernel of the base-changed multiplication.  Products of a
finite number of ideal elements commute with tensor extension, so
it carries $I_E^{r+1}$ onto the $(r+1)$-st power of the new ideal.

\item \textbf{Compare finite principal parts and their structure maps.}
Flatness preserves the resulting quotient sequence.  Every map in
the comparison is induced by the universal algebra maps and finite
multiplication, so the stated diagrams commute.
\end{enumerate}
\end{proof}

For the completed algebra the finite comparisons induce
\begin{equation}\label{eq:formal-gauge-comparison}
 B_P\otimes_B\widehat P_{E/B}\longrightarrow
 \varprojlim_r(B_P\otimes_BP^r_{E/B}).
\end{equation}
The following measured family identifies the resource distinction
between its two sides.

\begin{example}[Measured finite programs and a completed resource family]
\label{ex:measured-formal-family}
Work in the cofinite numerical instance and put
$t(n)=1/(n+1)$, $q=t^{-1}$, and $R=B[t^{-1}]$.
Use an idle invocation carrying a natural-number tag; it leaves the
returned input unchanged, and its tag is retained in the history.
NNO recursion constructs a finite description $e_r$ which, for each
$0\le j\le r$, emits $(n+1)^j$ copies of the event tagged $j$ and
then returns its input. Exponentiation and the finite concatenation
are implemented by the retained iteration constructors. The map
$r\mapsto e_r$ is itself given by finite code construction.

Give event $j$ the weight $z^j$ and other invocation events weight zero.
For each nonnegative rational $z$,
the exact measured cost and its germ polynomial are
\begin{equation}\label{eq:measured-geometric-series}
 T_{r,z}(n)=\sum_{j=0}^r(n+1)^jz^j,
 \qquad F_r(z)=\sum_{j=0}^rq^jz^j.
\end{equation}
The polynomial coefficient extension is calculated in the exact signed
algebra. The histories have the same returned value, and
\cref{thm:measured-jet-adequacy} identifies their measured coefficients
with geometric Taylor coefficients. Every coefficient of $F_r$ belongs
to $t^{-r}B$, since $q^j=t^{-r}t^{r-j}$ for $j\le r$.
Truncation gives $F_{r+1}\bmod z^{r+1}=F_r$, so these polynomials
determine
\[
 F(z)=\sum_{j\ge0}q^jz^j\in R[[z]],\qquad (1-qz)F(z)=1.
\]
The measured objects are the finite programs $e_r$ and their costs;
$F$ is the compatible formal generating function of their coefficients.
\end{example}

\begin{proposition}[Finite polynomial budgets without a common degree]
\label{prop:measured-completion-overhead}
The canonical map
\begin{equation}\label{eq:infinite-denominators}
 B[[z]][t^{-1}]\longrightarrow B[t^{-1}][[z]]
\end{equation}
is injective. Its image consists exactly of coefficient families
contained in one budget $t^{-N}B$ for some $N\in\N$.
The measured formal family $F$ of \cref{ex:measured-formal-family}
lies outside this image, although every finite truncation lies in it.
\end{proposition}
\begin{proof}\leavevmode
\begin{enumerate}[leftmargin=*,align=left]
\item \textbf{Characterize series with a common denominator.}
A fraction $t^{-N}b(z)$ has all its coefficients in $t^{-N}B$.
Conversely, a coefficient family with that common budget becomes a
unique $B$-coefficient series after multiplication by $t^N$, giving
the required fraction.

\item \textbf{Prove injectivity and exclude the displayed series.}
Regularity of $t$ proves injectivity
coefficientwise. If $F$ had such a denominator, its coefficient at
$j=N+1$ would give $t^Nq^{N+1}=t^{-1}\in B$, a contradiction.

\item \textbf{Recover the finite-order budget bounds.}
The finite budget calculation in \cref{ex:measured-formal-family}
proves the last assertion.
\end{enumerate}
\end{proof}

Thus localization and completion distinguish an order-dependent
polynomial growth degree from a single degree valid for every
coefficient. The constants in each coefficient's Landau estimate may
depend on that coefficient. Formal order is also distinct from merely
quotienting by asymptotically vanishing errors, as the next calculation
shows.

\begin{proposition}[The Landau vanishing ideal is idempotent]
\label{prop:vanishing-idempotent}
For every proper-filter rational germ context,
\[
 \mathfrak m^2=\mathfrak m,\qquad
 \varprojlim_r B/\mathfrak m^{r+1}\cong B/\mathfrak m.
\]
Consequently completion at the original small-error ideal cannot
replace the diagonal-adic formal-order tower.
\end{proposition}
\begin{proof}\leavevmode
\begin{enumerate}[leftmargin=*,align=left]
\item \textbf{Construct a rational square-root bound.}
Let $f\in\mathfrak m$.  There is a uniform rational operation
$q\mapsto s(q)$ on nonnegative rationals satisfying
\[
 s(0)=0,\qquad q\le s(q)^2\le4q\quad(q>0),
\]
obtained by locating a positive rational between successive
powers of two using finite rational comparison and doubling.

\item \textbf{Factor a vanishing germ.}
Apply it pointwise to $|f|$ and write $g=s(|f|)$.  Define
$h=f/g$ on the nonzero support, zero elsewhere.  These are
rational profiles and respect eventual equality.  One has
$f=gh$ and $|h|\le g$ pointwise.

\item \textbf{Prove that both factors vanish.}
For every $k\ge1$, the
smallness bound $4k^2|f|\le1$ on an eventual extent yields
$kg\le1$ there, and hence also $k|h|\le1$.  Thus
$g,h\in o(1)=\mathfrak m$.  So $f\in\mathfrak m^2$.

\item \textbf{Deduce idempotence and the constant quotient tower.}
The opposite inclusion follows from ideality.  All positive
powers of the idempotent ideal equal $\mathfrak m$, and the
inverse system of quotients is constant at $B/\mathfrak m$.
\end{enumerate}
\end{proof}

\subsection{Formal coordinates on indexed resource cells}
\label{sec:formal-budget-index}
We now apply formal coordinates to the universal resource family.
The budget classifier remains $\Cplx$; the new objects are
neighbourhoods of its resource cells.
At a principal budget $a\in L$, the module is $\mathscr M(a)=M_a$.
Over the regular-epimorphic cover supplying a representative $g$,
it is the displayed finite-projective module $Bg\cong Be_g$.
Construct $A_a=\Sym_B(M_a^\vee)$ and its principal parts locally,
then descend along the representative kernel pair. An inclusion
$M_a\incl M_b$ induces the contravariant coordinate map
\[
 A_b\longrightarrow A_a.
\]
Free algebras, diagonal ideals, finite quotients, and completed
higher diagonals are functorial in these coordinate maps.
The coordinate algebras therefore vary contravariantly with the
budget order, while their represented formal spaces vary covariantly.
This gives the geometric budget diagram with the required variance.

The descent depends only on the marked module. Equal growth
classes give the same embedded image $Bg\subseteq K$, and
finite-projective duality and principal parts are intrinsic to
that image. Their local presentations consequently agree on
kernel pairs and descend by effective quotients in $\E$.


\begin{proposition}[Preservation of the global resource classifiers]
\label{prop:formal-preserves-classifiers}
The functors and formal systems of this section leave unchanged
the classifying equivalence
$\Loc_\E(Z,\Cplx)\cong\Sub_{\Delta_ZB}(\Delta_ZK)$,
the universal marked module, all finite resource-diagram classifiers,
and the feasibility classifier $\mathfrak F=\mathbb S^{\Cplx}$.
The formal differential objects are constructed \emph{over} their
already classified parameters.
\end{proposition}
\begin{proof}\leavevmode
\begin{enumerate}[leftmargin=*,align=left]
\item \textbf{Preserve the indexed classifying data under completion.}
The classifiers are defined from $B,K,L$ and their marked subobjects.
Formal completion acts on the polynomial coordinate algebras over
these parameters. The finite algebraic constructions commute with
parameter pullback, so the original classifying maps and universal
properties continue to classify the indexed data.
\end{enumerate}
\end{proof}

\subsection{Observed formal towers and their effective points}
\label{sec:formal-effective-points}
A compatible family of formal coefficients also needs an observation
structure. We generate it from the finite-order observations and
apply the effective-point construction of
\cref{sec:standalone-observations}. The generator test of
\cref{lem:standalone-tests} then reduces the limit comparison
to the observations at each finite stage.

\begin{theorem}[Effective points of an observed inverse tower]
\label{thm:effective-formal-towers}
For an internal tower $(D_r,\mathcal O_r)$ in $\ED$, form the
underlying compatible-family poset $D_\infty=\varprojlim_rD_r$.
Its pointwise directed suprema and the observation structure
\begin{equation}\label{eq:formal-observation-closure}
 \mathcal O_\infty
 =\cl_\Sigma\{\pi_r^*U:r\in\N,\ U\in\mathcal O_r\}
\end{equation}
construct its inverse limit in $\ED$.  There is a canonical
isomorphism over $D_\infty$,
\begin{equation}\label{eq:effective-formal-limit}
 \Pt D_\infty\cong\varprojlim_r\Pt D_r.
\end{equation}
It is uniform in parameters and natural in maps of observed towers.
On the generating observations its positive evaluation is
\[
 \operatorname{Obs}_{D_\infty}(x,\pi_r^*U)
 =\operatorname{Obs}_{D_r}(\pi_rx,U).
\]
\end{theorem}
\begin{proof}\leavevmode
\begin{enumerate}[leftmargin=*,align=left]
\item \textbf{Construct directed suprema coordinatewise.}
For an inhabited directed ideal of compatible families, take its
lower image under each projection.  This is an inhabited directed
ideal in $D_r$; take its supremum there.
Scott continuity of the transition maps makes these suprema a
compatible family.  Their coordinatewise universal property proves
that this family is the supremum in $D_\infty$.

\item \textbf{Construct the observation closure and limit property.}
Every projection
is Scott-continuous.  Its pulled-back observations are therefore
Scott open; their $\Sigma$-observation closure remains Scott open
by the closure construction of \cref{sec:standalone-observations}.
For a cone of observation-preserving maps, inverse image preserves
each generating observation and hence its closure.  This proves
the limit universal property in $\ED$.

\item \textbf{Define the forward effective-point comparison.}
The effective projections give the forward comparison in
\eqref{eq:effective-formal-limit}.

\item \textbf{Construct the profile on generating observations.}
Conversely work with a
parameterized element $(x_r)_r$ of the limit on the right.  Evaluation
on the internally indexed family $\sum_r\mathcal O_r$ is the single
map $(r,U)\mapsto\operatorname{Obs}_{D_r}(x_r,U)$ into $\Sigma$.

\item \textbf{Descend and extend the observation profile.}
If two generators have the same pullback to $D_\infty$, these
evaluations agree in $\Omega$ and hence in $\Sigma$, since
$\Sigma\incl\Omega$ is monic.  They therefore descend to the
image family of generating observations in
\eqref{eq:formal-observation-closure}.  The testing-generators
lemma \cref{lem:standalone-tests} extends this profile to the whole
observation closure.  The effective-point pullback then gives a
unique point of $\Pt D_\infty$ over $(x_r)_r$.

\item \textbf{Verify the isomorphism and parameter coherence.}
The two maps are inverse over the compatible-family object, proving
the isomorphism.  All maps are constructed in the parameter context;
uniqueness in the effective-point pullback proves naturality and
coherence.
\end{enumerate}
\end{proof}


\begin{construction}[Canonical observed exact coefficient towers]
\label{constr:observed-coefficient-towers}
For an exact coefficient object $C$ of $\E$, write
$C_{\mathrm{disc}}$ for its discrete poset and set
$\mathsf D(C)=\F(C_{\mathrm{disc}})$, using the canonical
observed free completion of \cref{sec:standalone-observations}.
An inhabited directed lower ideal in a discrete poset is a singleton.
Thus the underlying dcpo of $\mathsf D(C)$ is canonically $C$,
and its observations are precisely the predicates represented by
maps $C\to\Sigma$, viewed in $\Omega$.  In particular
\[
 \Pt\mathsf D(C)\cong C.
\]
Every exact arrow $f:C\to C'$ preserves these observations by
precomposition and defines $\mathsf D(f)$.  This fixes the finite
observation structure by the canonical free construction.

Apply this construction to $C=P_q^r$ and its truncation maps,
and then use \cref{thm:effective-formal-towers}.  It constructs
an observed dcpo $\widehat{\mathsf D}^{\,q}$ with
\begin{equation}\label{eq:effective-primary-cochains}
 \Pt\widehat{\mathsf D}^{\,q}
 \cong\varprojlim_r\Pt\mathsf D(P_q^r)
 \cong\varprojlim_rP_q^r
 \cong\underline{\Hom}_{\ThickEt}
       (\mathcal D_{X_A}^{(q)},\mathcal O_{\mathrm{inf}}).
\end{equation}
The same construction applies to finite normalized coefficient
objects and their compatible-family limit.  The final comparison
is \cref{thm:primary-synthetic-cochains}.
\end{construction}

\begin{proposition}[Finite-output observation and name transport]
\label{prop:formal-finite-output-tracking}
An internal map between constructed coefficient towers, or between
their limits whose order-$r$ output is given uniformly by a finite
coefficient map at a constructed order $b(r)$, preserves the
generated observations.  It therefore induces the corresponding
map of effective points.  Finite arithmetic, polynomial Taylor
substitution, truncation, and the coface maps have this property;
the de Rham differentials and completed homotopies constructed
below have the same property by their stated finite coefficient
formulas.

In the effective and relative-realizability instances,
the retained exact arithmetic name maps implement these operations
on their projective name presentations, uniformly in the order
and supplied parameters.  A uniformly implemented compatible input
family consequently gives uniformly implemented finite output
observations of the resulting formal datum.
\end{proposition}
\begin{proof}\leavevmode
\begin{enumerate}[leftmargin=*,align=left]
\item \textbf{Transport finite-output observations.}
For a generating output observation $\pi_r^*U$, its pullback is
$\pi_{b(r)}^*(U\circ f_r)$, where $f_r$ is the displayed finite
coefficient map.  Precomposition is an allowed observation in the
canonical exact coefficient structure.  For a map of general
observed towers the same assertion is its finite-stage
observation-preservation equation.  \Cref{lem:standalone-tests}
extends the result to all generated observations; its effective-point
functor then supplies the claimed map.

\item \textbf{Compose the retained arithmetic name maps.}
Each finite formula is a composition of the retained exact arithmetic,
finite-list, rational-coefficient, and matrix operations.  Their
name maps compose by \cref{lem:standalone-syntax}.

\item \textbf{Construct one uniform operation on finite outputs.}
The order bound $b(r)$ and the finite lists of coefficient operations
are constructed by NNO recursion, so the index is an input to one
uniform operation.

\item \textbf{Verify products, descent, and dependent name transport.}
A map defined on several coefficient objects
uses their finite product presentation.  Passing from presentations
to denotations preserves its equation; dependent inputs use the
projective refinement of \cref{prop:standalone-name-lifting} before taking a tracker.  This proves the finite-output comparison on names.
\end{enumerate}
\end{proof}

The limit classifies compatible effective points with the generated
observations. A uniform compatible family of names carries its
tracker; projective presentations of the limit are supplied by
\cref{prop:standalone-name-lifting}.

We can now use the formal neighbourhoods and their identified
function objects to build cochains on the infinitesimal nerve.
The next section compares those cochains with algebraic differential
forms.

\section{Synthetic and algebraic de Rham resource geometry}
\label{sec:derham}

Resource variations have differential coefficients, and compatible
changes of those coefficients require transport. De Rham forms express
the infinitesimal data; curvature tests the compatibility of a proposed
connection, while a stratification records coherent transport at every
formal order. The event instrumentations of
\cref{prop:instrumentation-curvature} exhibit this compatibility test.
The completed measured family of \cref{ex:measured-formal-family}
provides a nonzero logarithmic class and an explicit flat transport
(\cref{prop:measured-logarithmic-transport}).

Functions on the completed infinitesimal nerve and algebraic K\"ahler
forms give two descriptions of this differential information. We
construct a chain-homotopy equivalence between them over the exact
resource coefficients. The geometric function calculation is
\cref{thm:primary-sdg,thm:primary-synthetic-cochains}; its effective
finite observations are
\cref{constr:observed-coefficient-towers,prop:formal-finite-output-tracking}.
The algebraic precedents are
\citep{GrothendieckDeRham,BerthelotOgus}, and the finite rational
contraction uses a bar--Koszul calculation related to \citep{HKR}.

Relative forms differentiate valuation parameters while fixing $B$.
Absolute forms over $\QQ$ also differentiate resource coefficients;
the logarithmic connections on $t^{-k}B$ in \cref{prop:log-connection}
use that second calculus. The primary site supplies the function
objects for the comparison, and based-Weil restriction recovers their
finite polynomial coefficient formulas.

\subsection{Completed higher-diagonal cochains}
\label{sec:completed-cochains}
A cochain evaluates a function on a finite sequence of
infinitesimally related points. The completed higher diagonals
provide these configurations and their coordinate algebras.
Fix $E$ finite-projective, $V=E^\vee$, and $A=\Sym_B(V)$.
For $q\ge0$ let
\[
 D_q=A^{\otimes_B(q+1)},\quad
 I_q=\ker(D_q\xrightarrow{\mu_q}A),\quad
 \widehat C^q(A/B)=\varprojlim_rD_q/I_q^{r+1}.
\]
By \cref{thm:primary-synthetic-cochains} there is a natural
cosimplicial algebra identification
\[
 \widehat C^q(A/B)
 \cong\underline{\Hom}_{\ThickEt}
       (\mathcal D_{X_A}^{(q)},\mathcal O_{\mathrm{inf}}).
\]
These limits are taken in internal $B$-modules or algebras;
\eqref{eq:primary-complete-coefficient-sheaf} gives their
coefficient-sheaf realization. To calculate their maps, write
the successive coordinate differences as
$u_i=x_i-x_{i-1}$, $1\le i\le q$. The finite-projective
change of variables gives
\begin{equation}\label{eq:completed-nerve-coordinates}
 \widehat C^q(A/B)
 \cong A[[u_1,\ldots,u_q]],
\end{equation}
where each $u_i$ denotes a block of variables dual to $E$ and
the completion is in the total difference degree.  The degree
$q=0$ algebra is $A$.

A simplicial map of the higher diagonal sends a vertex to a
selected or repeated vertex, hence sends each difference to a finite
sum of successive differences.  This preserves the powers of
$I_q$ and constructs the cosimplicial algebra
\[
 \widehat C^\bullet(A/B):\Delta\longrightarrow\CAlg_B(\E).
\]
Its normalization and differential are
\begin{equation}\label{eq:normalized-completed-derham}
 N_{\mathrm{inf}}^q(A/B)
  =\bigcap_{i=0}^{q-1}\ker(s_i:\widehat C^q\to\widehat C^{q-1}),
 \qquad \delta=\sum_{i=0}^{q+1}(-1)^id_i.
\end{equation}
By the cosimplicial identities $\delta^2=0$ and
$\delta N^q\subseteq N^{q+1}$.  These are equations of maps,
proved by pairing the two terms corresponding to any two
deleted vertices.

\begin{proposition}[Explicit completed coface formula]
\label{prop:completed-cofaces}
For a normalized $q$-cochain $f(x;u_1,\ldots,u_q)$,
\begin{align}
 \delta f(x;u_1,\ldots,u_{q+1})
 &=f(x+u_1;u_2,\ldots,u_{q+1})\notag\\
 &\quad+\sum_{i=1}^q(-1)^i
  f(x;u_1,\ldots,u_i+u_{i+1},\ldots,u_{q+1})\notag\\
 &\quad+(-1)^{q+1}f(x;u_1,\ldots,u_q).
 \label{eq:formal-cofaces}
\end{align}
Normalization consists exactly of series having positive degree
in each successive difference block $u_i$.  All displayed
operations are defined by finite polynomial substitution at each
fixed homogeneous degree and commute with truncation.
\end{proposition}
\begin{proof}\leavevmode
\begin{enumerate}[leftmargin=*,align=left]
\item \textbf{Compute the coface substitutions.}
Deleting the first vertex changes the base point from $x$ to
$x+u_1$; deleting the $i$-th interior vertex replaces the two
adjacent differences by their sum; deleting the last vertex removes
the last difference.  This proves the formula.

\item \textbf{Identify normalized cochains.}
Repeating one
vertex sets one of its adjacent differences to zero.  Being in
all codegeneracy kernels therefore means that every monomial
contains a positive contribution from each difference block.

\item \textbf{Extend the formulas to completed cochains.}
The substitutions respect the total-degree filtration, and each
resulting coefficient is a finite sum, proving the completion
assertion.
\end{enumerate}
\end{proof}

The normalized-completed comparison is summarized by the commuting diagram
\begin{equation}\label{eq:derham-comparison-diagram}
\begin{tikzcd}[column sep=huge]
N^\bullet_{\mathrm{inf}}(A/B)
 \arrow[r,"{\rho}"] \arrow[d,"{\delta}"']
& \Omega^\bullet_{A/B} \arrow[d,"{d}"] \\
N^{\bullet+1}_{\mathrm{inf}}(A/B)
 \arrow[r,"{\rho}"']
& \Omega^{\bullet+1}_{A/B}.
\end{tikzcd}
\end{equation}
The maps will be constructed and compared in
\cref{thm:de-rham-comparison}.

The normalized complex has an Alexander--Whitney cup product,
\begin{equation}\label{eq:formal-cup}
 (f\smile g)(x;u_1,\ldots,u_{p+q})
 =f(x;u_1,\ldots,u_p)
  g(x+u_1+\cdots+u_p;u_{p+1},\ldots,u_{p+q}).
\end{equation}
Associativity is reassociation of the same ordered vertices.
The coface calculation gives
$\delta(f\smile g)=\delta f\smile g+(-1)^p f\smile\delta g$.


\subsection{Relative forms and the infinitesimal comparison}
\label{sec:derham-comparison}
We now compare the normalized cochains with the algebraic
de Rham complex. The K\"ahler module of \cref{sec:geometry}
already supplies its degree-one forms. Define
\[
 \Omega^p_{A/B}=\bigwedge_A^p\Omega^1_{A/B},
 \qquad d:A\to\Omega^1_{A/B}
\]
and extend $d$ as the unique graded derivation of degree one
with $d(da)=0$.  The alternating Leibniz identities give $d^2=0$.
For a polynomial resource cell,
\begin{equation}\label{eq:polynomial-de-rham-forms}
 \Omega^\bullet_{A/B}
 \cong A\otimes_B\bigwedge_B^\bullet V.
\end{equation}
This construction is internal to the signed rational algebraic
sector and imposes no completion on the quantitative scalar
semiring $S$.

The comparison extracts the first-order contribution in each
successive displacement and alternates it. For a normalized
$q$-cochain $f$, a temporary finite free presentation gives
the coefficient formula
\begin{equation}\label{eq:formal-hkr-map}
 \rho_q(f)=
 \sum_{a_1,\ldots,a_q}
 [u_1^{a_1}\cdots u_q^{a_q}]f\;
 dx^{a_1}\wedge\cdots\wedge dx^{a_q}.
\end{equation}
The exterior quotient makes this alternating tensor formula
independent of the temporary basis. At degree zero it is
$\rho_0=\id_A$. The theorem proves that this map retains the
cohomological information of the completed cochains.

\begin{lemma}[Homogeneous additive-cochain calculation and natural contraction]
\label{lem:additive-bar-koszul}
Let $V$ be a displayed finite-projective module over the rational
commutative ring $B$.  In homogeneous difference degree $m$, put
\[
 C^{q,m}_0(V)=
 \bigoplus_{\substack{m_1+\cdots+m_q=m\\m_i\ge1}}
 \Sym_B^{m_1}(V)\otimes_B\cdots\otimes_B\Sym_B^{m_q}(V)
\]
with $C^{0,0}_0=B$ and all other degree-zero summands absent.
The degree-preserving component $\delta_0$ of
\eqref{eq:formal-cofaces} has cohomology
\[
 H^q(C^{\bullet,m}_0(V))\cong
 \begin{cases}\bigwedge_B^mV,&q=m,\\0,&q\ne m.\end{cases}
\]
There are \emph{constructed}, $B$-linear homogeneous operators
$i_m,p_m,h_m$ of cochain degrees $0,0,-1$, respectively,
with target complex $\bigwedge_B^mV[-m]$ (and $B[0]$ at $m=0$),
such that
\begin{equation}\label{eq:homogeneous-contraction}
 \begin{gathered}
 p_mi_m=1,\quad \delta_0i_m=0,\quad p_m\delta_0=0,\\
 \delta_0h_m+h_m\delta_0=1-i_mp_m,\qquad
 h_m^2=h_mi_m=p_mh_m=0.
 \end{gathered}
\end{equation}
They are natural under linear maps of displayed finite-projective
modules, and their coefficients are finite rational expressions
constructed uniformly in $m$.
\end{lemma}
\begin{proof}\leavevmode
\begin{enumerate}[leftmargin=*,align=left]
\item \textbf{Identify the homogeneous bar--cobar complex.}
The differential is the reduced cobar differential of the graded
symmetric coalgebra on $V$, with coproduct induced by
$u\mapsto u\otimes1+1\otimes u$.  In homogeneous degree $m$ this
is a finite complex.  Its graded dual is the reduced bar complex
of $\Sym_B(V^\vee)$: the divided-power dual is identified with
the symmetric algebra using invertible factorials.

\item \textbf{Compute its cohomology by the bar--Koszul comparison.}
The Koszul
resolution on a displayed free basis of $V^\vee$ resolves its
augmentation $B$; exactness follows by tensoring the one-variable
free resolutions.  The normalized bar resolution resolves the
same augmentation by its extra-degeneracy contraction.  Their
antisymmetrization comparison is a quasi-isomorphism of free
resolutions (the usual bar--Koszul calculation \citep{HKR});
tensoring with $B$ and taking finite graded duals gives the
stated cohomology, concentrated at $q=m$.

\item \textbf{Present the finite permutation complexes.}
We construct the homotopy on finite permutation modules.  Give the set of ordered block
decompositions of $\{1,\ldots,m\}$ its decidable finite
permutation presentation.  In degree $q$ the induced permutation
modules
\[
 P_m^q=\bigoplus_{\substack{m_1+\cdots+m_q=m\\m_i\ge1}}
 \QQ[S_m]\otimes_{\QQ[S_{m_1}\times\cdots\times S_{m_q}]}\QQ
\]
carry the bar--cobar differential by splitting adjacent blocks
into their shuffles with the signs of $\delta_0$.  With right/left
permutation conventions matched, symmetrization gives a natural
isomorphism
\[
 C^{q,m}_0(V)\cong
 V^{\otimes_Bm}\otimes_{\QQ[S_m]}P_m^q.
\]
The alternating representation $\mathrm{sgn}_m$ gives the
corresponding presentation of $\bigwedge_B^mV$.
All these representations are finite-dimensional rational vector
spaces with specified permutation bases.

\item \textbf{Identify the alternating cohomology representation.}
To identify its cohomology equivariantly, evaluate the free-module
bar--Koszul calculation at $V=\QQ^m$ and take the multidegree
$(1,\ldots,1)$ in its $m$ labelled generators.  The resulting
summand of $C^{\bullet,m}_0(V)$ is precisely $P_m^\bullet$,
and the corresponding summand of $\bigwedge^mV$ is the sign
representation $\mathrm{sgn}_m$.  Thus the calculation identifies
the cohomology of $P_m^\bullet$ as $\mathrm{sgn}_m$ in degree
$m$ and zero elsewhere, equivariantly under permutations.

\item \textbf{Construct the projection and antisymmetrization.}
Take $p:P_m^\bullet\to\mathrm{sgn}_m[-m]$ to be the
exterior quotient and take $i$ to be normalized antisymmetrization:
\[
 i(v_1\wedge\cdots\wedge v_m)=\frac1{m!}
 \sum_{\sigma\in S_m}\operatorname{sgn}(\sigma)
 v_{\sigma(1)}\otimes\cdots\otimes v_{\sigma(m)}.
\]
These equivariant chain maps satisfy $pi=1$. The cohomology
calculation above shows that $K_m^\bullet=\ker p$ is acyclic.

\item \textbf{Construct equivariant splittings of the acyclic kernel.}
Write $B_m^q=\operatorname{im}(\delta_0:K_m^{q-1}\to K_m^q)$.
In every degree there is an exact sequence of finite rational
representations
\[
 0\longrightarrow B_m^q\longrightarrow K_m^q
 \xrightarrow{\delta_0}B_m^{q+1}\longrightarrow0.
\]
Lexicographically ordered rational Gaussian elimination constructs
a linear section $s_q:B_m^{q+1}\to K_m^q$ of this last
surjection. Replace it by the finite average
\[
 \bar s_q=\frac1{m!}\sum_{\sigma\in S_m}
       \sigma s_q\sigma^{-1}.
\]
Since $\delta_0$ is equivariant,
$\delta_0\bar s_q=1$ still holds. Consequently
$K_m^q=B_m^q\oplus\bar s_q(B_m^{q+1})$ is an
equivariant direct-sum decomposition.

\item \textbf{Define the contracting homotopy.}
For $z\in P_m^q$ set
\[
 h_q(z)=\bar s_{q-1}
       \bigl(z-i pz-\bar s_q\delta_0z\bigr).
\]
The argument of $\bar s_{q-1}$ belongs to $B_m^q$:
it lies in $\ker p$ and its differential is zero.

\item \textbf{Verify the homotopy equation and side conditions.}
Now $\delta_0h_qz=z-ipz-\bar s_q\delta_0z$, while
$h_{q+1}\delta_0z=\bar s_q\delta_0z$.
Thus $\delta_0h+h\delta_0=1-ip$.
The image of $h_q$ is the complementary summand
$\bar s_{q-1}(B_m^q)$, on which the next $h$ vanishes;
also $hi=0$ and $ph=0$. This proves all the side conditions
in \eqref{eq:homogeneous-contraction}. The $m=0$ case is
$i=p=1$, $h=0$.

\item \textbf{Transport the contraction naturally to coefficient modules.}
Finally apply $V^{\otimes m}\otimes_{\QQ[S_m]}(-)$.
The homotopy equations remain equations of finite matrices,
valid over every rational $B$; all operations are natural in
$V$ because $S_m$ acts by permuting tensor factors.  If $V$
is a retract of a finite free module, the same functorial maps
restrict along its displayed idempotent.

\item \textbf{Construct the family uniformly in the degree.}
The lexicographic
elimination and finite average are algorithms on $m$ and its
finite permutation lists, so the entire family is constructed
by NNO recursion, not by selecting one splitting for each $m$.
\end{enumerate}
\end{proof}

\begin{theorem}[Completed infinitesimal de Rham comparison]
\label{thm:de-rham-comparison}
For every displayed finite-projective $B$-module $E$, the maps
\eqref{eq:formal-hkr-map} assemble into a natural cochain map
\[
 \rho:N^\bullet_{\mathrm{inf}}(\Sym_B(E^\vee)/B)
  \longrightarrow\Omega^\bullet_{\Sym_B(E^\vee)/B}.
\]
It is a \emph{chain-homotopy equivalence}, with a parameter-uniform
homotopy constructed on compatible truncations.  In particular it
is a quasi-isomorphism, natural under polynomial substitution and
in internal resource parameters.  Through \cref{thm:primary-synthetic-cochains}, this is the
comparison from normalized synthetic functions on the full primary
infinitesimal nerve to algebraic relative forms.  It carries the
Alexander--Whitney product to the exterior product and induces
an isomorphism of graded cohomology rings.
\end{theorem}
\begin{proof}\leavevmode
\begin{enumerate}[leftmargin=*,align=left]
\item \textbf{Decompose normalized completed cochains by degree.}
Put $V=E^\vee$, $A=\Sym_B(V)$ and use the total difference-degree
filtration.  Normalization and the finite symmetric-power
decomposition of the previous lemma identify
\[
 N^q_{\mathrm{inf}}(A/B)
 \cong\prod_{m\ge q} A\otimes_BC^{q,m}_0(V).
\]

\item \textbf{Separate the homogeneous differential and perturbation.}
Split the coface differential as
$\delta=\delta_0+\Delta$.  The first part keeps the base point
$x$ fixed and has homogeneous difference degree zero.  The
perturbation $\Delta$ comes only from expanding the coefficients
of $f(x+u_1;\ldots)$ around $x$: it raises total difference degree
by at least one.

\item \textbf{Assemble the homogeneous contraction.}
The contraction operators of
\cref{lem:additive-bar-koszul}, extended over $A$ and assembled
as product maps $i,p,h$, obey
\[
 pi=1,\quad\delta_0h+h\delta_0=1-ip,\quad
 h^2=hi=ph=0,\quad p=\rho.
\]
The target is $\Omega^\bullet_{A/B}$ with zero homogeneous
differential before perturbation.

\item \textbf{Construct the completed perturbation maps.}
In every finite truncation,
$h\Delta$ and $\Delta h$ are nilpotent because $\Delta$ raises
filtration.  Therefore the following sums are finite modulo
any specified difference degree, and define internal maps of
the completed product:
\begin{equation}\label{eq:completed-perturbation-maps}
 \begin{aligned}
 J&=(1+h\Delta)^{-1}i
   =\sum_{k\ge0}(-h\Delta)^k i,\\
 H&=h(1+\Delta h)^{-1}
   =\sum_{k\ge0}(-1)^k h(\Delta h)^k,\\
 p'&=p(1+\Delta h)^{-1},\qquad
 d' =p\Delta J.
 \end{aligned}
\end{equation}

\item \textbf{Verify the perturbation identities coefficientwise.}
These are the finite-geometric-series formulas of the basic
perturbation lemma, applied on each finite truncation
\citep{MillerRahmati}.  Substitution into
$(\delta_0+\Delta)^2=0$ and the homogeneous contraction
equations gives
$\delta J=Jd'$, $p'\delta=d'p'$,
$p'J=1$, and $\delta H+H\delta=1-Jp'$.
The equations agree under truncation and hence hold in the
internal inverse limit. This proves the perturbation identities coefficientwise.

\item \textbf{Identify the transferred projection.}
Filtration identifies their transferred terms exactly.
The map $p$ extracts difference degree equal to cochain degree;
$h$ preserves difference degree but lowers cochain degree, so
$p\Delta h=0$.  Thus $p'=p=\rho$.

\item \textbf{Identify the transferred de Rham differential.}
Terms in
$p\Delta(h\Delta)^ki$ with $k\ge1$ have difference degree
strictly above target cochain degree, and hence vanish.  Only
$p\Delta i$ survives in $d'=p\Delta J$.
For a polynomial coefficient $a$, the component of degree one
of $a(x+u)-a(x)$ is $da$; applying alternating projection to
the remaining linear difference coefficients gives precisely
\[
 d'=p\Delta i=d_{\mathrm{dR}}.
\]

\item \textbf{Establish the completed homotopy equivalence.}
Consequently the actual completed maps satisfy
\begin{equation}\label{eq:completed-derham-retraction}
 \boxed{\quad
 \rho J=1,\qquad
 \delta J=Jd_{\mathrm{dR}},\qquad
 \rho\delta=d_{\mathrm{dR}}\rho,\qquad
 \delta H+H\delta=1-J\rho.\quad}
\end{equation}
This proves the asserted homotopy equivalence, including
injectivity and surjectivity on all cohomology objects.

\item \textbf{Verify finite-output construction and naturality.}
All operators in a fixed output difference degree depend on only
finitely many lower coefficients.  Their defining algorithms use
finite rational matrices and sums and commute with parameter
substitution.  In particular the compatible infinite corrections
are produced by the NNO and the product universal property.
The comparison $\rho$ commutes with polynomial coordinate substitutions by the chain rule and the universal alternating quotient. The auxiliary maps $H,J$ have the linear naturality established by their finite coefficient construction.

\item \textbf{Verify compatibility with cup and wedge products.}
Finally the degree-$(p+q)$ component of
\eqref{eq:formal-cup} linear in every difference block is the
concatenation of those of $f$ and $g$.  Alternation gives
$\rho(f\smile g)=\rho(f)\wedge\rho(g)$, as claimed.
\end{enumerate}
\end{proof}


For a polynomial coordinate map $p:E\to F$, naturality of the
comparison itself is the commutative square of complexes
\[
\begin{tikzcd}[column sep=large,row sep=large]
 N^\bullet_{\mathrm{inf}}(A_F/B) \arrow[r,"{\rho_F}"]
   \arrow[d,"{p^*}"']
   & \Omega^\bullet_{A_F/B} \arrow[d,"{p^*}"]\\
 N^\bullet_{\mathrm{inf}}(A_E/B) \arrow[r,"{\rho_E}"']
   & \Omega^\bullet_{A_E/B}.
\end{tikzcd}
\]
The finite homogeneous contraction constructs its homotopy equivalence;
the full Taylor substitution supplies the vertical cochain map.


\subsection{Relative formal Poincar\'e contraction}
\label{sec:formal-poincare}
The preceding comparison relates two descriptions of forms.
A second calculation checks the local formal neighbourhood
itself by contracting its relative displacement directions.
For $n$ finite free formal variables put
$R=A[[u_1,\ldots,u_n]]$ and define the continuous relative forms
\[
 \widehat\Omega^p_{R/A}
   =R\otimes_A\bigwedge_A^p\Bigl(\bigoplus_{i=1}^nA\,du_i\Bigr).
\]
The differential is termwise formal differentiation, which is
well-defined because each output coefficient is a finite sum.

\begin{theorem}[Formal Poincar\'e lemma]
\label{thm:formal-poincare}
The inclusion $A\to\widehat\Omega^\bullet_{R/A}$ is a chain-homotopy
equivalence.  Its relative de Rham cohomology is $A$ in degree
zero and zero in positive degrees.  The result is natural in
finite-projective formal displacement modules.
\end{theorem}
\begin{proof}\leavevmode
\begin{enumerate}[leftmargin=*,align=left]
\item \textbf{Compute the Euler contraction identity.}
Let $\mathsf E=\sum_i u_i\partial/\partial u_i$ and let
$\iota_{\mathsf E}$ be contraction.  The graded Leibniz rule
gives
\[
 d\iota_{\mathsf E}+\iota_{\mathsf E}d
  =\mathcal L_{\mathsf E}.
\]

\item \textbf{Define the contraction on positive weights.}
On homogeneous coefficient degree $m$ and form degree $p$ the
Lie derivative is multiplication by $m+p$.  For $m+p>0$ define
$h=(m+p)^{-1}\iota_{\mathsf E}$, and set $h=0$ on degree-zero
constants.  Rational invertibility and coefficientwise
finite arithmetic make this an internal operator on the
completed product.

\item \textbf{Verify the deformation retraction and projective descent.}
Its identity is
\[
 dh+hd=\id-\operatorname{const}\operatorname{ev}_0.
\]
This exhibits the required deformation retraction.  The Euler
operator and the homogeneous decomposition are intrinsic to
symmetric powers and commute with finite-projective retractions.
\end{enumerate}
\end{proof}

\subsection{Crystalline stratifications and integrable resource connections}
\label{sec:crystals}
Differential coefficients can also be transported between
infinitesimally related points. A stratification records that
transport and its composition law; an integrable connection
records its first-order differential data. We construct their
comparison for a displayed finite-projective module $N$ over
$A=\Sym_B(E^\vee)$. Pull $N$ back along the two coefficient
actions on principal parts:
\[
 s_r^*N=P^r_{A/B}\otimes_{A,s_r}N,
 \qquad
 t_r^*N=P^r_{A/B}\otimes_{A,t_r}N.
\]
\begin{definition}[Formal resource stratification]
\label{def:formal-stratification}
With $s_r(a)=a(x)$ and $t_r(a)=a(x+u)$, a formal
stratification in the source-to-target convention is a family
of $P^r_{A/B}$-linear isomorphisms
\[
 \varepsilon_r:s_r^*N\xrightarrow{\ \cong\ }t_r^*N
\]
compatible with truncation and equal to the identity on the
reduced diagonal.  At every finite completed triple diagonal,
writing $\varepsilon_{xy}$ for transport from the fibre over
$x$ to the fibre over $y$, its cocycle equation is
\[
 \varepsilon_{yz}\,\varepsilon_{xy}=\varepsilon_{xz}.
\]
The morphisms are $A$-linear maps commuting with this family.
\end{definition}

An integrable relative connection is a $B$-linear map
\[
 \nabla:N\to N\otimes_A\Omega^1_{A/B}
\]
satisfying the Leibniz identity
$\nabla(an)=a\nabla(n)+n\otimes da$ and whose graded
extension satisfies $\nabla^2=0$.  Curvature is an $A$-module map.

\begin{theorem}[Connections are formal stratifications]
\label{thm:connection-stratification}
For displayed finite-projective $E$ and finite-projective
$A=\Sym_B(E^\vee)$-modules $N$, integrable $B$-connections
on $N$ and formal stratifications of $N$ form naturally
equivalent categories.  The equivalence commutes with
polynomial substitution, parameter substitution, and finite-order
coefficient base change.
\end{theorem}
\begin{proof}\leavevmode
\begin{enumerate}[leftmargin=*,align=left]
\item \textbf{Derive the iterated covariant Leibniz formula.}
Write $P^r=P^r_{A/B}$ and $u=y-x$.  We first construct the
transport in the direction forced by its semilinearity.  In a
finite-free coordinate presentation write
$\nabla n=\sum_i\nabla_i(n)\otimes dx^i$.  The Leibniz
identity and flatness say
\[
 \nabla_i(an)=\partial_i(a)n+a\nabla_i(n),\qquad
 [\nabla_i,\nabla_j]=0.
\]
The finite multi-index Leibniz formula follows by induction:
\[
 \nabla^\alpha(an)=
 \sum_{\beta\le\alpha}\binom\alpha\beta
 \partial^\beta(a)\nabla^{\alpha-\beta}(n).
\]

\item \textbf{Construct the finite Taylor transport.}
Define a $P^r$-linear map by its values on $t_r$-generators,
\begin{equation}\label{eq:connection-taylor}
 \begin{aligned}
 T_r:t_r^*N&\longrightarrow s_r^*N,\\
 T_r(1\otimes_{t_r}n)
  &=\sum_{|\alpha|\le r}\frac{u^\alpha}{\alpha!}
           (1\otimes_{s_r}\nabla^\alpha n).
 \end{aligned}
\end{equation}

\item \textbf{Verify balancing and reduction.}
Indeed finite reindexing of the Leibniz sum gives
\[
 T_r(1\otimes_{t_r}an)
 =a(x+u)T_r(1\otimes_{t_r}n)
 =t_r(a)T_r(1\otimes_{t_r}n),
\]
so the defining formula respects the $t_r$-balanced relation.
Its reduction modulo $I$ is the identity.

\item \textbf{Prove the transport cocycle.}
On the finite completed triple diagonal use vertices $x,y,z$,
with differences $u=y-x$ and $v=z-y$.  Expanding the iterated
covariant derivatives and using their commutation, the ordinary
multinomial identity gives
\[
 T_{x\leftarrow y}\,T_{y\leftarrow z}=T_{x\leftarrow z},
 \qquad T_{x\leftarrow x}=1.
\]
Each coefficient is a finite sum, since total difference degree
is truncated.

\item \textbf{Construct the inverse stratification.}
Swapping the vertices gives the reverse transport;
the cocycle with $z=x$ proves that $T_r$ is invertible.  Define
\[
 \varepsilon_r=T_r^{-1}:s_r^*N\xrightarrow{\cong}t_r^*N.
\]
The displayed identity for $T$ now yields the required
source-to-target cocycle, while reduction and truncation are
immediate.  In particular no arbitrary lifting of an inverse is
needed.

\item \textbf{Extract the connection from first-order transport.}
Conversely, given $\varepsilon$, set $T=\varepsilon^{-1}$.
Working in $s_1^*N$, define $\nabla n$ by
\[
 T_1(1\otimes_{t_1}n)-(1\otimes_{s_1}n)
 \quad\bmod I^2\ \in\ (I/I^2)\otimes_A N
  \cong\Omega^1_{A/B}\otimes_A N.
\]

\item \textbf{Verify Leibniz and zero curvature.}
The equality $T_1(1\otimes_{t_1}an)
=t_1(a)T_1(1\otimes_{t_1}n)$ gives
$\nabla(an)=a\nabla(n)+n\otimes da$ after taking the
first difference coefficient.  The alternating coefficient of
$u_iv_j$ in the triple cocycle is
$[\nabla_i,\nabla_j]$, so the cocycle gives zero curvature.

\item \textbf{Recover every higher Taylor coefficient.}
Differentiating the same triple cocycle once in the final
nilpotent displacement recursively gives the higher coefficient
rule $T_{\alpha+e_i}=(\alpha_i+1)^{-1}\nabla_iT_\alpha$.
It recovers uniquely the finite Taylor polynomial
\eqref{eq:connection-taylor} because all positive integers are
invertible in $B$.

\item \textbf{Descend the construction and verify naturality.}
For a displayed projective retract, replace the temporary free
basis by a finite dual basis of $E$; constant directional
derivations commute, and their symmetric tensor expressions
factor through the displayed idempotent.  Every identity above
is an identity of finite symmetric-power morphisms.  The
construction therefore descends to the finite-projective cell,
commutes with coordinate substitutions and coefficient base
change, and is functorial on connection-preserving module maps.

\item \textbf{Verify that the constructions are inverse.}
The two procedures are mutually inverse by the first-order
coefficient and the unique higher-coefficient recurrence.
\end{enumerate}
\end{proof}


\subsection{Measured instrumentation, curvature, and coherent transport}
\label{sec:instrumentation-transport}
Two changes of event measurement can preserve returned values while
their order changes the resource data. Use the idle events $a,b$ and
let $N=(N_a,N_b)^{\mathsf T}$ be the corresponding count profiles of
a retained program. For natural numbers $u,v$, the finite substitutions
\[
 \sigma_u(a)=a,\quad \sigma_u(b)=ba^u,
 \qquad \tau_v(a)=ab^v,\quad\tau_v(b)=b
\]
are implemented by inserting the indicated idle blocks. Other
primitives are unchanged. Recursion on finite descriptions performs
the replacement under sequencing, branching, and iteration. Since the
inserted blocks are finite and leave the state unchanged, the finite
return comparison identifies the return domains and values. The
returning histories are the displayed free-monoid substitutions.

Writing count vectors as columns, their exact actions are
\[
 U_u=I+uP,\qquad V_v=I+vQ,\qquad
 P=\begin{pmatrix}0&1\\0&0\end{pmatrix},\quad
 Q=\begin{pmatrix}0&0\\1&0\end{pmatrix}.
\]
For instance $\sigma_u$ changes $(N_a,N_b)$ to
$(N_a+uN_b,N_b)$. These polynomial laws extend to resource coefficients.
The count profiles lie in a common principal cell, so the two-component
calculation is realized on a finite-projective resource bundle.

\begin{proposition}[Instrumentation curvature and ordered transport]
\label{prop:instrumentation-curvature}
The two instrumentations have mixed-jet discrepancy
\begin{equation}\label{eq:instrumentation-mixed-jet}
 U_\eps V_\delta-V_\delta U_\eps
 =\eps\delta[P,Q],\qquad
 [P,Q]=\begin{pmatrix}1&0\\0&-1\end{pmatrix}
\end{equation}
over $B[\eps,\delta]/(\eps^2,\delta^2)$.
On the coefficient bundle $B[u,v]^2$, the connection prescribing
these constant infinitesimal updates,
\[
 \nabla_0=d-P\,du-Q\,dv,
\]
has curvature $[P,Q]\,du\wedge dv$ and admits no formal stratification
inducing it. The ordered family $G(u,v)=V_vU_u$ instead constructs
the integrable connection
\[
 \nabla_G=d-(dG)G^{-1}
\]
and the formal transport
$\varepsilon_{x,y}=G(y)G(x)^{-1}$. Its associated crystalline coaction
is given by \cref{prop:crystalline-jet-coaction}.
\end{proposition}
\begin{proof}\leavevmode
\begin{enumerate}[leftmargin=*,align=left]
\item \textbf{Compute the measured mixed coefficient.}
Expand the two products and compare the coefficient of $\eps\delta$.
The signed difference is \eqref{eq:instrumentation-mixed-jet}.
These are measured coefficients by the count formulas and
\cref{thm:measured-jet-adequacy}.

\item \textbf{Compute curvature and exclude a constant-rule stratification.}
The differential of
$-P\,du-Q\,dv$ is zero, and its exterior square is
$[P,Q]\,du\wedge dv$, which is nonzero in the free coefficient bundle.
The connection--stratification equivalence therefore rules out a
stratification with these constant update rules.

\item \textbf{Construct the polynomial frame and flat connection.}
Since $P^2=Q^2=0$, the inverse of $G$ is
$(I-uP)(I-vQ)$, an actual polynomial matrix. For a column $h$,
$\nabla_G(Gh)=G\,dh$; applying the graded connection again gives
zero curvature.

\item \textbf{Verify the stratification cocycle.}
The transport $G(y)G(x)^{-1}$ is the identity on the
diagonal, is compatible with every truncation, and satisfies
\[
 G(z)G(y)^{-1}\,G(y)G(x)^{-1}=G(z)G(x)^{-1}.
\]
It therefore gives the required stratification.

\item \textbf{Identify the ordered update rules.}
Its $v$-update is
$Q$, while its $u$-update is $V_vPV_v^{-1}$; at the origin these are
the original two updates. The conjugation records their specified
order and supplies the compatible transport.
\end{enumerate}
\end{proof}

This example distinguishes the mixed response of two implemented
instrumentations from their common returned computation. It also
exhibits, on the same resource coefficients, a curvature obstruction
and a constructed flat transport. Formal stratifications act on the
primary infinitesimal relation by \cref{prop:crystalline-jet-coaction};
based restriction gives the displayed Taylor matrices.

\subsection{Global resource coefficients, absolute forms, and a nonzero class}
\label{sec:derham-global}
The construction is uniform in $E$, its displayed idempotent,
and its coefficient maps. It therefore applies at each budget
to $A_a=\Sym_B(M_a^\vee)$. Budget transitions act by polynomial
coordinate restriction, and \cref{thm:de-rham-comparison}
identifies the induced maps on cochains and forms.
This assembles the relative complexes and crystalline transport
over the resource cells. The parameter space is still $\Cplx$,
with its previously specified points and positive observations.

\begin{theorem}[Formal principal parts and de Rham comparison after localization]
\label{thm:localized-formal-derham}
Let $A=\Sym_B(V)$ for a displayed finite-projective $V$, let
$f\in A$, and put $A_f=A[f^{-1}]$.  Write $D_q=A^{\otimes_B(q+1)}$
and $I_q=\ker(D_q\to A)$, with analogous notation $I_{q,f}$
over $A_f$.  At every $q,r$ there are natural isomorphisms
\begin{equation}\label{eq:localized-higher-diagonals}
 \frac{A_f^{\otimes_B(q+1)}}{I_{q,f}^{r+1}}
 \ \cong\ A_f\otimes_A\frac{D_q}{I_q^{r+1}}.
\end{equation}
They respect every vertex and truncation map.  Consequently
normalized completed coordinate cochains of $A_f$ are
chain-homotopy equivalent to $\Omega^\bullet_{A_f/B}$.
The connection--stratification equivalence and finite Taylor
transport extend to finite-projective $A_f$-modules.
\end{theorem}
\begin{proof}\leavevmode
\begin{enumerate}[leftmargin=*,align=left]
\item \textbf{Construct inverses at every vertex.}
Write $f_i$ for the image of $f$ at vertex $i$ in the
higher-diagonal algebra.  Modulo $I_q^{r+1}$, the element
$f_i-f_0$ is nilpotent of exponent $r+1$.  After inverting
$f_0$, the inverse of every other $f_i$ is constructed by
\[
 f_i^{-1}=f_0^{-1}\sum_{k=0}^r
  \left(-\frac{f_i-f_0}{f_0}\right)^k
  \pmod{I_q^{r+1}}.
\]

\item \textbf{Identify the localized principal-parts algebra.}
Thus the universal algebra obtained by inverting all the
$f_i$ at finite order is precisely the one obtained by
inverting $f_0$.  Since localization commutes with quotients
and with finite powers of their defining ideals, the two
universal properties give the inverse maps of
\eqref{eq:localized-higher-diagonals}.

\item \textbf{Verify simplicial and truncation coherence.}
The finite inverse
is unique, so the comparison commutes with all simplicial
vertex maps and truncation.

\item \textbf{Extend the homogeneous contraction over the localization.}
The normalized completed complex is consequently the inverse
limit of its finite localizations
$A_f\otimes_A(D_q/I_q^{r+1})$.  In difference coordinates its
homogeneous additive complex is
$A_f\otimes_B C^{\bullet,m}_0(V)$.  The maps of
\cref{eq:homogeneous-contraction} extend over $A_f$.

\item \textbf{Construct the localized de Rham retraction.}
Expanding the inverse $f(x+u)^{-1}$ by the displayed finite
geometric formula produces a perturbation strictly raising
formal difference degree.  The same finite-order perturbation
construction \eqref{eq:completed-perturbation-maps} therefore
gives a deformation retraction onto
$\Omega^\bullet_{A_f/B}$, whose transferred differential is
the universal localized derivation
$d(f^{-1})=-f^{-2}df$.  Localization is performed at each finite diagonal order before taking the limit.

\item \textbf{Extend connection transport over the localization.}
The Taylor transport of \cref{thm:connection-stratification}
uses the same finite inverse identities; it therefore extends
and satisfies its unit and cocycle equations over $A_f$.
\end{enumerate}
\end{proof}

\begin{proposition}[A logarithmic relative de Rham class]
\label{prop:logarithmic-derham}
On $A=B[x,x^{-1}]$, relative to $B$,
\[
 H^0_{\mathrm{dR}}(A/B)\cong B,\qquad
 H^1_{\mathrm{dR}}(A/B)\cong B\cdot[d\log x],
\]
and the higher cohomology objects vanish.  The formal diagonal
has $P^r_{A/B}\cong A[u]/(u^{r+1})$, and $d\log x=dx/x$
is represented by the completed normalized cocycle
\[
 \ell(x,y)=\log\left(1+\frac{y-x}{x}\right)
 =\sum_{j\ge1}\frac{(-1)^{j+1}}j
       \left(\frac{y-x}{x}\right)^j.
\]
\end{proposition}
\begin{proof}\leavevmode
\begin{enumerate}[leftmargin=*,align=left]
\item \textbf{Compute the Laurent de Rham cohomology.}
Termwise relative differentiation gives
$d(b_jx^j)=jb_jx^{j-1}dx$.  Every $j\ne0$ is invertible
in $B$ and the derivative has no $x^{-1}dx$ coefficient.
This computes the two nonzero cohomology objects; the
cotangent module is free of rank one, so higher forms vanish.

\item \textbf{Compute the finite and completed principal parts.}
Put $y=x+u$.  At finite order $y$ has the explicitly
constructed inverse
\[
 y^{-1}=x^{-1}\sum_{k=0}^r(-u/x)^k\pmod{u^{r+1}}.
\]
It follows from \cref{thm:localized-formal-derham} that
$P^r_{A/B}\cong A[u]/(u^{r+1})$ and
$\widehat P_{A/B}\cong A[[u]]$.

\item \textbf{Construct the logarithmic cocycle and identify its class.}
The formal logarithm
$\ell(x,y)$ is therefore a legitimate completed coordinate
cochain.  Its finite logarithm identity
$\ell(x,y)+\ell(y,z)=\ell(x,z)$ holds modulo each higher
diagonal power, so $\delta\ell=0$.  Its first difference
coefficient is $x^{-1}$; therefore
$\rho_1(\ell)=dx/x$, as claimed.
\end{enumerate}
\end{proof}

\begin{proposition}[A measured logarithmic class and its transport]
\label{prop:measured-logarithmic-transport}
For the measured family of \cref{ex:measured-formal-family}, retain
$R=B[t^{-1}]$, $q=t^{-1}$, and put
\[
 A=R[z,(1-qz)^{-1}],\qquad f(z)=(1-qz)^{-1},\qquad
 \omega=\frac{q}{1-qz}\,dz.
\]
Relative to $R$, the form $\omega$ is closed and has nonzero de Rham
class. Its completed diagonal representative is
\begin{equation}\label{eq:measured-log-cocycle}
 \ell(z,z+u)
 =\log\frac{1-qz}{1-q(z+u)}
 =\sum_{j\ge1}\frac1j\left(\frac{qu}{1-qz}\right)^j.
\end{equation}
The rank-one connection $\nabla=d-\omega$ is integrable, $f$ is
horizontal, and its source-to-target stratification is multiplication by
\begin{equation}\label{eq:measured-rank-one-transport}
 \tau(z,z')=\frac{f(z')}{f(z)}=\frac{1-qz}{1-qz'}.
\end{equation}
The Taylor coefficients of $f$ at zero are exactly the measured
coefficients $q^j$ of the family $F_r$.
\end{proposition}
\begin{proof}\leavevmode
\begin{enumerate}[leftmargin=*,align=left]
\item \textbf{Identify the nonzero logarithmic de Rham class.}
Since $q$ is a unit of $R$, $x=1-qz$ identifies $A$ with
$R[x,x^{-1}]$ and $\omega$ with $-dx/x$.
\Cref{prop:logarithmic-derham}, applied over $R$, identifies its
nonzero class. The same relative calculation is available over $B$:
localization gives $\Omega^1_{R/B}=0$. Alternatively
$A=B[z,(t(t-z))^{-1}]$, a principal open of a polynomial $B$-cell.

\item \textbf{Construct the measured logarithmic cocycle.}
At every finite diagonal order, $qu/(1-qz)$ is nilpotent, so
\eqref{eq:measured-log-cocycle} is a finite logarithm there. Its
truncations are compatible. Multiplication of the ratios in
\eqref{eq:measured-rank-one-transport} gives the additive cocycle
identity for their formal logarithms, and the first difference
coefficient is $q/(1-qz)$. Thus the de Rham comparison sends $\ell$
to $\omega$.

\item \textbf{Verify integrability and the horizontal frame.}
The relative exterior square vanishes in this one-coordinate algebra,
and $d\omega=0$, giving integrability. Direct differentiation gives
$df=\omega f$. The unit $f$ is a horizontal frame, so transport of
its scalar coefficients is multiplication by $f(z')/f(z)$.

\item \textbf{Verify transport coherence and the measured expansion.}
These ratios satisfy the identity and triple cocycle equations at
every finite order. Expanding $(1-qz)^{-1}$ at zero gives $q^j$,
which agrees with \eqref{eq:measured-geometric-series}.
\end{enumerate}
\end{proof}

The de Rham class concerns the absence of an algebraic additive
primitive for $\omega$. The displayed multiplicative transport exists
and will give the crystalline jet coaction in
\cref{ex:measured-crystalline-family}. Its horizontal generating
function retains the complete coefficient family whose resource bounds
were calculated in \cref{prop:measured-completion-overhead}.


The absolute algebraic de Rham complex on these cells has an Euler contraction.
\begin{theorem}[Absolute Euler contraction and the Laurent class]
\label{thm:absolute-euler}
For every displayed finite-projective $V$ over the rational ring
$B$, inclusion and zero-section augmentation construct a natural
chain-homotopy equivalence
\[
 \Omega^\bullet_{B/\QQ}
 \ \simeq\ \Omega^\bullet_{\Sym_B(V)/\QQ}.
\]
For $L=B[x,x^{-1}]$, there is a natural chain-homotopy
equivalence
\[
 \Omega^\bullet_{L/\QQ}\ \simeq\
 \Omega^\bullet_{B/\QQ}\oplus
 \Omega^{\bullet-1}_{B/\QQ}\cdot d\log x,
\]
where the shifted summand has the corresponding shifted
differential.  Consequently
\[
 H^k_{\mathrm{dR}}(L/\QQ)
 \cong H^k_{\mathrm{dR}}(B/\QQ)
 \oplus H^{k-1}_{\mathrm{dR}}(B/\QQ)\,[d\log x].
\]
\end{theorem}
\begin{proof}\leavevmode
\begin{enumerate}[leftmargin=*,align=left]
\item \textbf{Decompose the absolute de Rham complex by weight.}
The symmetric algebra has its canonical weight grading
$A=\bigoplus_{m\ge0}\Sym_B^m(V)$.  Give the universal
differential $da$ the same weight as homogeneous $a$.
The relations defining $\Omega^1_{A/\QQ}$ and its exterior
powers are homogeneous, so the absolute de Rham complex is a
direct sum by nonnegative weight.  Its weight-zero subcomplex
is exactly $\Omega^\bullet_{B/\QQ}$, with inclusion and
projection induced by $B\to A\to B$.

\item \textbf{Construct the Euler contraction.}
The Euler derivation $\mathcal E$ is zero on $B$ and acts as
multiplication by $m$ on $\Sym_B^m(V)$.  Its interior contraction
on forms obeys the Cartan identity
\[
 d\iota_{\mathcal E}+\iota_{\mathcal E}d
   =\mathcal L_{\mathcal E}.
\]

\item \textbf{Contract the positive-weight summands.}
On weight $m>0$ the right side equals $m$.  Define
$h=m^{-1}\iota_{\mathcal E}$ on this weight and zero on
weight zero.  Since each algebraic form has finitely many
weights, $h$ is an actual map of internal $B$-modules.  It
gives $dh+hd=1-\iota_B\epsilon_B$, proving the first
homotopy equivalence.  The grading and Euler map are
intrinsic and commute with finite-projective retractions.

\item \textbf{Contract the nonzero Laurent weights.}
For $L=B[x,x^{-1}]$, give $x$ weight $1$, $x^{-1}$ weight
$-1$, and $dx$ weight $1$.  The same Euler contraction
kills every nonzero integer weight, since every nonzero
integer is a unit of $B$.  The weight-zero complex is
\[
 \Omega^\bullet_{B/\QQ}\oplus
 \Omega^{\bullet-1}_{B/\QQ}\wedge\frac{dx}{x}.
\]

\item \textbf{Identify the remaining complex and its cohomology.}
The absolute differential preserves this direct sum and
$d(dx/x)=0$.  More explicitly, a chain isomorphism from the
standard cochain shift (whose differential is $-d_B$) sends
$\alpha$ of degree $k-1$ to
$(-1)^{k-1}\alpha\wedge d\log x$ in degree $k$.
Hence the same Cartan identity supplies the asserted deformation
retraction and its cohomology calculation.
\end{enumerate}
\end{proof}

Relative formal cochains use $B$ as base. The absolute retractions
above use $\QQ$ and apply to the numerical scale differentials and
logarithmic connections of \cref{sec:scales}.


\subsection{The smooth-groupoid Artin sector}
\label{sec:artin-sector}
A finite-projective resource-coordinate module may be described in
local frames. Changes of frame act on its coordinate expression and
on every polynomial resource map defined on it. To organize these
isomorphisms and their stabilizers, we use algebraic stacks in
groupoids on the same big \'etale site.
An internal groupoid-valued presheaf is constructed as a finite-limit
category object in $\E$. Descent requires its values on an
\'etale \v{C}ech nerve to give the groupoid of compatible objects
and arrows, and stackification glues these finite covering data.
This is a further sector alongside the $\E$-valued sheaf categories
of \cref{thm:etale-endpoint-cohesion}.

\begin{construction}[Smooth Artin quotients]
\label{constr:artin-quotient}
For a displayed smooth affine groupoid $R\rightrightarrows U$
of finitely presented $B$-schemes, with smooth source and target,
let $[U/R]$ be the stack of local $U$-objects and their
$R$-valued change-of-trivialization arrows, glued along
\'etale covering diagrams.  Equivalently, on a test $T$
its objects are groupoid-principal $R$-bundles $P\to T$
with equivariant map $P\to U$, and its arrows are their
equivariant isomorphisms. The atlas $U\to[U/R]$
has kernel-pair groupoid $R$ and is smooth and surjective in
the ordinary algebraic-stack sense.  The construction includes
the full simplicial nerve $U_\bullet$ with
$U_0=U$, $U_1=R$ and $U_q=R\times_U\cdots\times_U R$.
\end{construction}
\begin{proof}\leavevmode
\begin{enumerate}[leftmargin=*,align=left]
\item \textbf{Verify descent for the torsor quotient.}
The torsor cocycle on a covering family is finite diagram data;
refining two coverings compares them by their pullbacks.
The sheaf condition for $U$ and $R$ supplies effectiveness of
object and arrow descent, and torsor descent makes the
displayed assignment a stack.

\item \textbf{Recover the groupoid nerve and smooth atlas.}
Pullback of the universal
torsor along $U\to[U/R]$ reconstructs the groupoid arrows
and its full nerve.  Smoothness and surjectivity of the atlas
follow by pulling back to a torsor trivializing cover; this
is the standard ordinary smooth-groupoid presentation theorem
\citep[Sections 94.12 and 94.17]{StacksArtin}.
\end{enumerate}
\end{proof}


\begin{theorem}[Endpoint cohesion on groupoid-valued \'etale stacks]
\label{thm:etale-stacky-endpoints}
Let $\StEt^{(1)}(\AffB;\E)$ and
$\StEt^{(1)}(\InfB;\E)$ denote the 2-categories of
internal groupoid-valued \'etale stacks, with pseudonatural
maps and modifications.  The endpoint construction lifts to
a string of 2-adjoints
\[
 \iota_!\dashv\iota^*\dashv\iota_*\dashv\iota^!,
\]
with the same larger- and smaller-endpoint formulas as
\eqref{eq:etale-endpoint-adjunction}.  On a representable
affine $T$, the rightmost functor is determined by
\[
 (\iota^!\mathcal H)(T)
 =\operatorname{PsNat}(\iota_*y_T,\mathcal H),
\]
the groupoid of pseudonatural transformations and their
invertible modifications.  The two endpoint embeddings remain
nonisomorphic, witnessed by the discrete affine-line stack.
The full Artin quotient stacks from
\cref{constr:artin-quotient} occur as objects of the base
stack 2-category.
\end{theorem}
\begin{proof}\leavevmode
\begin{enumerate}[leftmargin=*,align=left]
\item \textbf{Lift the endpoint adjunctions to groupoid-valued presheaves.}
The endpoint universal properties of
\cref{lem:infinitesimal-endpoints} apply to pseudonatural
transformations of groupoid-valued presheaves, yielding the
first three 2-adjunctions and full faithfulness by the two
identities $c\iota=d\iota=1$.

\item \textbf{Compare endpoint descent and stackification.}
The \'etale lift of
\cref{lem:etale-nilpotent-lift} identifies not only
covering families but their complete descent groupoids;
therefore endpoint precomposition commutes with stackification
by gluing effective descent data.  It preserves the
2-colimits of stack diagrams and consequently has a right
2-adjoint.

\item \textbf{Construct the remaining right 2-adjoint.}
To construct it, express a groupoid-valued presheaf by the
Grothendieck construction on its representable objects and
arrows, with their specified invertible coherence.  Its
2-colimit is the original presheaf by the bicategorical
Yoneda calculation.  Applying $\iota_*$ to that presentation
and mapping to $\mathcal H$ identifies the resulting groupoid
of maps with the pseudonatural transformations into the
displayed rightmost formula.

\item \textbf{Descend the adjunction and distinguish the embeddings.}
Stackification gives the same
adjunction on \'etale stacks.  The affine-line example in
\cref{thm:etale-endpoint-cohesion} consists of discrete
groupoids, so the endpoint embeddings are also distinct here.
\end{enumerate}
\end{proof}

\begin{proposition}[Infinitesimal symmetries of resource frames]
\label{prop:artin-gl}
The finitely presented group scheme
\[
 \mathrm{GL}_{n,B}=
 \operatorname{Spec}B[x_{ij},\det(x_{ij})^{-1}]
\]
is smooth, and its quotient $B\mathrm{GL}_{n,B}=[\operatorname{Spec}B/\mathrm{GL}_{n,B}]$
is a displayed Artin stack.  Its atlas has self-intersection
$\mathrm{GL}_{n,B}$.  Over any square-zero algebra $R[\eps]/\eps^2$,
its infinitesimal stabilizer at the identity is
\[
 \ker\bigl(\mathrm{GL}_n(R[\eps]/\eps^2)
           \to\mathrm{GL}_n(R)\bigr)
 =\{1+\eps M:\ M\in M_n(R)\},
\]
with multiplication $(1+\eps M)(1+\eps N)=1+\eps(M+N)$.
In particular the displayed atlas is not \'etale for $n>0$.
\end{proposition}
\begin{proof}\leavevmode
\begin{enumerate}[leftmargin=*,align=left]
\item \textbf{Represent invertible matrices by a smooth principal open.}
Localization at the determinant represents invertible
matrices.  The Jacobian criterion makes the principal open in
a finite affine space smooth over $B$.

\item \textbf{Compute the infinitesimal group law.}
A matrix congruent
to the identity modulo a square-zero ideal is uniquely
$1+\eps M$, with inverse $1-\eps M$.  Multiplying and
using $\eps^2=0$ gives the stated group law.

\item \textbf{Identify the nonzero tangent and quotient atlas.}
The group
scheme has a nonzero relative tangent module, so its atlas
cannot be \'etale.  Its torsor quotient and atlas are
\cref{constr:artin-quotient}.
\end{enumerate}
\end{proof}


\begin{theorem}[The infinitesimal relation on presented Artin groupoids]
\label{thm:artin-stacky-formal-relation}
For the displayed quotient stack $\mathcal X=[U/R]$ of
\cref{constr:artin-quotient}, define on every infinitesimal
arrow $j=(\operatorname{Spec}R_0\hookrightarrow\operatorname{Spec}S_0)$
the groupoid
\[
 \mathcal D_{\mathcal X}(j)
 =\mathcal X(S_0)
   \mathop{\times}^{\mathrm{h}}_{\mathcal X(R_0)}
   \mathcal X(S_0).
\]
An object is a pair of $S_0$-families with a specified
isomorphism between their reductions to $R_0$; arrows are
pairs of isomorphisms compatible with that boundary
isomorphism.  The assignment is an \'etale stack on
$\InfB$.  Its iterated homotopy fibre products give a
simplicial infinitesimal groupoid, with unit and composition
induced by identity and composition of the specified
boundary isomorphisms.  On the ordinary affine sheaf
$F_A$ this groupoid reduces to the fibre-product
relation of \cref{thm:completion-etale}, and consequently
has its explicit finite principal-parts filtration.
\end{theorem}
\begin{proof}\leavevmode
\begin{enumerate}[leftmargin=*,align=left]
\item \textbf{Construct the functorial infinitesimal groupoid relation.}
The restriction $\mathcal X(S_0)\to\mathcal X(R_0)$
is functorial in commutative squares of thickenings.
Homotopy pullbacks of groupoids may be constructed as
the groupoids of triples $(x,y,\alpha)$ with
$\alpha:x|_{R_0}\cong y|_{R_0}$; their morphisms are
compatible pairs.  This construction proves functoriality.

\item \textbf{Verify descent and the simplicial operators.}
On an \'etale covering diagram of $j$, each
endpoint gives a covering diagram by
\cref{lem:etale-nilpotent-lift}.  Torsors and their
isomorphisms descend along those covers, so the groupoids
at both endpoints satisfy descent; finite homotopy pullbacks
preserve the required limit diagrams.  Composition and
inversion of the boundary isomorphisms obey the groupoid
laws, giving all the simplicial operators.

\item \textbf{Recover the affine principal-parts relation.}
For a represented ordinary affine scheme $F_A$, its point
groupoids are discrete.  An isomorphism between their
reductions is precisely equality, so the displayed
homotopy pullback is the ordinary kernel pair
$F_A(S_0)\times_{F_A(R_0)}F_A(S_0)$.  Apply the finite
nilpotence proof of \cref{thm:completion-etale} for its
principal-parts filtration.  For a general Artin quotient the boundary isomorphism remains part of the groupoid data.
\end{enumerate}
\end{proof}

\subsection{\'Etale descent and the smooth Artin groupoid de Rham interface}
\label{sec:artin-derham}
The smooth-groupoid sector extends the coefficient calculation
to presentations with stabilizers. To retain that information
we use the full simplicial atlas of the Artin presentations
from \cref{constr:artin-quotient}. The following construction
provides its internal coefficient diagram over $B$.
The broader de Rham formalism for ordinary Artin stacks is
the cited theory of \citep{BehrendDR,StacksArtin}.

\begin{construction}[De Rham cochains of a displayed smooth affine groupoid]
\label{constr:artin-derham}
Let $R\rightrightarrows U$ be a smooth affine groupoid over $B$
whose nerve has finitely presented smooth affine terms $U_q$;
put $A_q=\OO(U_q)$.  The face maps of this full simplicial
scheme induce a cosimplicial diagram of relative differential
graded algebras
\[
 [q]\longmapsto\Omega^\bullet_{A_q/B}.
\]
Define its cochains by the total complex of the normalized
cosimplicial complex, with the de Rham differential $d$ and
the alternating coface differential $\partial$; explicitly
\[
 \operatorname{DR}_{\mathrm{grp}}([U/R]/B)
 =\operatorname{Tot}_{q\ge0}
           N^q\Omega^\bullet_{A_\bullet/B},
 \qquad D|_{(q,p)}=\partial+(-1)^q d.
\]
Normalization is the intersection of degeneracy kernels.
In each fixed total degree only finitely many terms occur,
so $D$ and all its signs are maps of internal $B$-modules.
For polynomial and principal-localization terms, the natural
comparison of \cref{thm:de-rham-comparison,thm:localized-formal-derham}
replaces each relative-form complex with its completed formal
coordinate-cochain complex.  Its simplicial naturality gives
a morphism of the corresponding total complexes.
\end{construction}
\begin{proof}\leavevmode
\begin{enumerate}[leftmargin=*,align=left]
\item \textbf{Verify the total differential.}
All simplicial maps are maps of $B$-algebras and commute
strictly with universal derivations.  The simplicial
identities imply $\partial^2=0$; naturality gives
$d\partial=\partial d$, so the displayed signed total
differential squares to zero.

\item \textbf{Normalize the simplicial complex.}
Degenerate summands split off
by the standard finite normalization identities.

\item \textbf{Compare completed infinitesimal cochains with the total de Rham complex.}
The
comparison $\rho$ from completed normalized infinitesimal
cochains to forms is natural under polynomial substitution
and algebraic localization, hence is compatible with all
cofaces wherever those local charts apply.  Levelwise
chain homotopy equivalence then gives a total-complex
quasi-isomorphism by the first-quadrant finite-total-degree
filtration.
\end{enumerate}
\end{proof}

\begin{proposition}[\'Etale refinement and infinitesimal stabilizer retention]
\label{prop:artin-derham-stabilizers}
The above cosimplicial coefficient construction is functorial
under maps of displayed smooth groupoids and under \'etale
refinement of their affine nerve terms.  The augmentation
of an \'etale \v{C}ech refinement is a quasi-isomorphism of
the corresponding total complexes of forms.  For the classifying
stack $B\mathrm{GL}_{n,B}$, the nerve is
\[
 U_q=(\mathrm{GL}_{n,B})^q,
\]
and its first nonconstant term has differential module
$\Omega^1_{\mathrm{GL}_{n,B}/B}$; its cotangent at the unit
contains $\mathfrak{gl}_n(B)^\vee$.  Hence this construction
retains the nonzero infinitesimal stabilizers which would be
lost by using only the atlas $U_0=\operatorname{Spec}B$.
\end{proposition}
\begin{proof}\leavevmode
\begin{enumerate}[leftmargin=*,align=left]
\item \textbf{Construct the map induced by a groupoid morphism.}
A morphism of groupoids induces a map of whole nerves,
so it induces every coface and degeneracy map and hence a
map of the total complexes.

\item \textbf{Apply differential base change to an \'etale refinement.}
Under an \'etale affine
refinement $A\to A'$, the vanishing of
$\Omega^1_{A'/A}$ and the flat base-change formula give
$\Omega^p_{A'/B}\cong A'\otimes_A\Omega^p_{A/B}$.

\item \textbf{Prove the refinement quasi-isomorphism.}
Faithfully flat \v{C}ech descent is exact on each displayed
coefficient module; the spectral sequence of the resulting
first-quadrant double complex proves the claimed
quasi-isomorphism.

\item \textbf{Identify the stabilizer tangent retained by the nerve.}
For the last assertion the groupoid
nerve has $U_0=\operatorname{Spec}B$,
$U_1=\mathrm{GL}_{n,B}$ and the indicated higher fibre
products.  Differentiating the matrices $1+\eps M$ of
\cref{prop:artin-gl} identifies the identity tangent with
$M_n(B)$ and its dual with the identity cotangent module.
The full simplicial diagram therefore records this
information although $\Omega^1_{B/B}=0$ at the zero-simplicial
atlas.
\end{enumerate}
\end{proof}

For ordinary smooth Artin stacks, presentation independence follows
from \citet{BehrendDR}. The internal nerve construction above gives
functoriality under displayed groupoid maps, invariance under the
stated \'etale refinements, and the completed-polynomial comparison.

\subsection{Effective observations of differential coefficients}
\label{sec:derham-effectivity}
The observation structures of
\cref{thm:effective-formal-towers,constr:observed-coefficient-towers}
apply to the differential coefficient objects. Their finite formulas
give the following tracking result.

\begin{proposition}[Positive tests and completed differential data]
\label{prop:derham-positive-tests}
The finite coefficient operations defining $P^r$, normalized
cochains, $d$, truncation, and formal Taylor transport commute
with substitution of all computational parameters.  Their
realizations on supplied effective coefficient names are
uniformly tracked when the underlying exact arithmetic names
are supplied.  Their semantic equalities, closed-form predicates,
and cohomology classes remain predicates and quotients in $\E$;
their classifiers take values in $\Omega$.
\end{proposition}
\begin{proof}\leavevmode
\begin{enumerate}[leftmargin=*,align=left]
\item \textbf{Implement the finite differential operations.}
The finite operations are compositions of the exact scalar
constructors, finite lists, sums, products, and finite matrix
identities implemented in \cref{app:names-codes}. The
tracked structural maps are closed under these
finite compositions and substitute uniformly in the ambient
parameters.  The cochain differentials and truncated Taylor
maps are given by the same finite expressions at each order,
proving the first two assertions.

\item \textbf{Identify the logical status of coefficient equations.}
Equality and closedness are
kernel and equalizer conditions on coefficient maps; their
classifiers are generally $\Omega$-valued.

\item \textbf{Verify finite coefficient dependence of the completed maps.}
For the normalized coefficient tower, use the canonical discrete
observed presentations of \cref{constr:observed-coefficient-towers}.
Truncation, $\delta$, and $\rho$ have finite coefficient dependence;
the maps $J,H$ of \eqref{eq:completed-perturbation-maps} have the
same property because the perturbation raises difference degree.
Each fixed output coefficient therefore uses finitely many of its
geometric-series terms.  Formal displacement differentiation uses
one additional input degree.  The order bounds and rational finite
matrices are uniform in the indices.

\item \textbf{Obtain effective-point and name comparisons.}
Applying
\cref{prop:formal-finite-output-tracking} proves preservation of
the generated observations and the corresponding effective-point
and finite-output name comparisons for the completed maps.
\end{enumerate}
\end{proof}

We now have formal functions, their comparison with differential
forms, and transport of algebraic coefficients along the completed
relation. The next section uses that same relation to describe
jets: sections over infinitesimal neighbourhoods and the operators
determined by their finite coefficients.

\section{Crystalline jets and formal differential operators}
\label{sec:jets}
A resource jet collects the finite response coefficients of a valuation
or a section of a resource bundle. \Cref{thm:measured-jet-adequacy}
identifies these coefficients with Taylor-valued execution measurements.
Principal-parts correspondences now specify their geometric re-centring
and composition. Their full compatible system constructs the jet
comonad, and flat resource transport gives its crystalline coalgebras.
The family of \cref{ex:measured-formal-family} will realize all these
comparisons explicitly in \cref{ex:measured-crystalline-family}.

The shape-based jet construction follows the synthetic perspective of
\citet{KS}. Here the primary endpoint model supplies represented
algebraic function objects and exact resource coefficients. Finite
principal parts identify fibrewise-linear jet operators with algebraic
differential operators. Effective descent identifies the full coalgebra
category on the stated cells, while the connection--stratification
comparison supplies its finite-projective algebraic sector.

\subsection{Finite principal-parts correspondences}
\label{sec:jet-formal-correspondences}
For a finite-projective coordinate cell $E$, write
$A_E=\Sym_B(E^\vee)$ and $X_E=X_{A_E}=\iota_!F_{A_E}$
in the primary infinitesimal-arrow category.
Write $\mathcal D_E^{\le r}$ for the \'etale sheaf
\eqref{eq:boundary-principal-parts}; its restriction to based
Weil probes is the previously constructed finite diagonal
of \cref{sec:formal-realization}.  The two projections
\[
 s_r,t_r:\mathcal D_E^{\le r}\rightrightarrows X_E
\]
form the order-$r$ formal correspondence.  At $r=0$ both
projections are the identity; for $r>0$ they retain
nilpotent differences through order $r$.


The correspondence acts on sections by an adjoint pair.
Dependent sum carries data along it, and dependent product
collects the corresponding sections:
\begin{equation}\label{eq:cat-finite-correspondence-adjoints}
 \mathbb T_{X_E}^r=(t_r)_!s_r^*
 \quad\dashv\quad
 \mathbb J_{X_E}^r=(s_r)_*t_r^*.
\end{equation}
Here $(t_r)_!$ is dependent sum and $(s_r)_*=\Pi_{s_r}$ is
dependent product.  The adjunction follows by composing
$(t_r)_!\dashv t_r^*$ and $s_r^*\dashv(s_r)_*$.
The common underlying correspondence is
\[
\begin{tikzcd}[column sep=large]
 X_E & \mathcal D_E^{\le r}
   \arrow[l,"{s_r}"'] \arrow[r,"{t_r}"] & X_E .
\end{tikzcd}
\]
Composing two correspondences pulls back over their common
endpoint. The next calculation maps that composite to the
order-$(r+s)$ neighbourhood; taking mates will give the jet
comparison in the reverse direction.


\begin{theorem}[Addition of formal order]
\label{thm:jet-order-addition}
Composition of adjacent formal displacements defines natural maps
\[
 \mathcal D_E^{\le r}\times_{X_E}\mathcal D_E^{\le s}
 \longrightarrow\mathcal D_E^{\le r+s},
\]
compatible with order truncation, unit, associativity, and
interchange of the two vertices.  Applying dependent sums gives
$\mathbb T_{X_E}^s\mathbb T_{X_E}^r\to\mathbb T_{X_E}^{r+s}$.
Their right-adjoint mates give the filtered jet comparison maps
\begin{equation}\label{eq:jet-order-addition}
 \mathbb J_{X_E}^{r+s}Y
 \longrightarrow
 \mathbb J_{X_E}^{r}\mathbb J_{X_E}^{s}Y,
 \qquad
 \mathbb J_{X_E}^rY=\Pi_{s_r}(t_r^*Y).
\end{equation}
At order zero $\mathbb J^0_{X_E}\cong\Id$.  The displayed
comparisons obey the filtered comonad associativity and counit
identities.
\end{theorem}
\begin{proof}\leavevmode
\begin{enumerate}[leftmargin=*,align=left]
\item \textbf{Bound the order of a composite displacement.}
At any nilpotent arrow, two successive differences $u,v$ have
nilpotence bounds $u^{r+1}=0$ and $v^{s+1}=0$ in the sense that
all their monomials of the indicated total degrees vanish.  Every
monomial of total degree $r+s+1$ in $u+v$ contains at least
$r+1$ factors from $u$ or at least $s+1$ factors from $v$.
Thus $u+v$ has order at most $r+s$.

\item \textbf{Verify units, associativity, and vertex interchange.}
This defines the displayed
map, naturally in the stage.  The unit is displacement zero,
associativity is $(u+v)+w=u+(v+w)$, and vertex interchange
sends $u$ to $-u$.  These polynomial equations commute with
all truncation maps.

\item \textbf{Compose the finite correspondences.}
Put
$Q_{r,s}=\mathcal D_E^{\le r}\times_{t_r,X_E,s_s}\mathcal D_E^{\le s}$.
The cartesian square
\[
\begin{tikzcd}[column sep=large,row sep=large]
 Q_{r,s} \arrow[r,"{\pi_2}"] \arrow[d,"{\pi_1}"']
   & \mathcal D_E^{\le s} \arrow[d,"{s_s}"]\\
 \mathcal D_E^{\le r} \arrow[r,"{t_r}"'] & X_E
\end{tikzcd}
\]
gives dependent-sum Beck--Chevalley.  Consequently the composite
$(t_s)_!s_s^*(t_r)_!s_r^*$ is the action of the composite
correspondence.  The map $Q_{r,s}\to\mathcal D_E^{\le r+s}$
therefore induces
$\mathbb T_{X_E}^s\mathbb T_{X_E}^r\to\mathbb T_{X_E}^{r+s}$.

\item \textbf{Construct the jet comparison and its coherence.}
Taking right-adjoint mates under
\eqref{eq:cat-finite-correspondence-adjoints} gives
\eqref{eq:jet-order-addition}.  Mates preserve the equations
between these composites, proving the filtered associativity
and unit identities.
\end{enumerate}
\end{proof}


Write $\delta_{r,s}:\mathbb J^{r+s}\to\mathbb J^r\mathbb J^s$.
The mate of associativity is the commutative square
\begin{equation}\label{eq:cat-filtered-jet-coherence}
\begin{tikzcd}[column sep=large,row sep=large]
 \mathbb J^{r+s+t} \arrow[r,"{\delta_{r+s,t}}"]
   \arrow[d,"{\delta_{r,s+t}}"']
   & \mathbb J^{r+s}\mathbb J^t
       \arrow[d,"{\delta_{r,s}\mathbb J^t}"]\\
 \mathbb J^r\mathbb J^{s+t}
   \arrow[r,"{\mathbb J^r\delta_{s,t}}"']
   & \mathbb J^r\mathbb J^s\mathbb J^t .
\end{tikzcd}
\end{equation}
Order-zero factors give the counit identities. Thus the
filtered comonad equations are the mates of the associativity
and unit equations for formal correspondences. The addition
of orders records the bound on a composite displacement.


The corresponding algebra maps of principal parts are the duals
of the finite correspondence compositions.  They retain the two
$A_E$-actions; replacing their target by an unaugmented symmetric
power would discard the source--target information needed for
substitution.

\subsection{Formal disks and jets in the primary geometry}
\label{sec:jet-etale-primary}
For $X\in\ThickEt$ let $\eta_X:X\to\InfShape X$ be
the unit of the middle endpoint adjunction, and put
\[
 \mathcal D_X=X\times_{\InfShape X}X,
 \qquad s,t:\mathcal D_X\rightrightarrows X.
\]
Because $\ThickEt$ is an internal sheaf topos, its slices
admit dependent sums and products
$(\eta_X)_!\dashv\eta_X^*\dashv(\eta_X)_*$.
Define the formal-disk monad and jet comonad by
\[
 \mathbb T_X^\infty=\eta_X^*(\eta_X)_!,
 \qquad\mathbb J_X^\infty=\eta_X^*(\eta_X)_*.
\]

\begin{theorem}[Full jet and precise coalgebra descent]
\label{thm:jets}
The displayed functors are a monad and a comonad on
$\ThickEt/X$, with
\[
 \mathbb T_X^\infty\dashv\mathbb J_X^\infty,
 \qquad \mathbb J_X^\infty Y\cong\Pi_s(t^*Y).
\]
For every $X$ there is a canonical comparison functor
\[
 \ThickEt/\InfShape X\longrightarrow
 (\ThickEt/X)^{\mathbb J_X^\infty}.
\]
It is an equivalence if $\eta_X$ is an effective epimorphism,
in particular for the polynomial finite-projective cells
$X=X_A$ of \cref{prop:shape-lifting-scope}.
\end{theorem}
\begin{proof}\leavevmode
\begin{enumerate}[leftmargin=*,align=left]
\item \textbf{Construct the monad, comonad, and their adjunction.}
The adjunctions construct the monad and comonad, with
multiplication and comultiplication given by the unit and
counit mates.  The triangle identities prove both unit and
counit equations and the associativity/coassociativity equations.
Successive slice adjunctions give
\[
 \Hom_X(\eta_X^*(\eta_X)_!E,F)
 \cong\Hom_{\InfShape X}((\eta_X)_!E,(\eta_X)_*F)
 \cong\Hom_X(E,\eta_X^*(\eta_X)_*F).
\]

\item \textbf{Identify the formal-pair correspondence and comultiplication.}
Dependent-product Beck--Chevalley for the defining pullback
of $\mathcal D_X$ gives the correspondence formula
$\Pi_s(t^*Y)$.  Comparing the two groupoid action diagrams
determines the jet comultiplication by re-centering along
composable formal pairs.

\item \textbf{Construct the descent comparison.}
The comparison functor sends $H\to\InfShape X$ to its
pullback $\eta_X^*H$ with the canonical unit-induced
coaction.

\item \textbf{Prove comonadicity under effective descent.}
If $\eta_X$ is an effective epimorphism, it is
an effective descent morphism in the sheaf topos: its entire
\v{C}ech kernel-pair diagram realizes the target, and
objects with compatible descent action are precisely its
pullbacks.  The same \v{C}ech action is the coalgebra
structure of $\eta_X^*(\eta_X)_*$, proving comonadicity
in this case.  The effective-epimorphism calculation on
polynomial cells is \cref{prop:shape-lifting-scope}.
\end{enumerate}
\end{proof}

The jet counit evaluates a formal-disk section at its centre.
Comultiplication re-centres it along composable nilpotent
pairs. A section $\sigma:X\to Y$ prolongs by
$X\cong\mathbb J_X^\infty X\to\mathbb J_X^\infty Y$,
and coKleisli composition uses that comultiplication.
This is the crystalline differential-operator calculus of
the primary endpoint geometry. Restricting its relations to
$\operatorname{Spec}B\hookrightarrow\operatorname{Spec}C$
recovers the based-Weil calculation of \cref{thm:cohesion}.


\subsection{Primary jet coefficients and algebraic principal parts}
\label{sec:primary-jet-coefficients}
To calculate the jets of an algebraic bundle, we realize its
coefficient module as an affine bundle sheaf. For a displayed
finite-projective $A$-module $N$, the fibre at $x:A\to S$
is $S\otimes_{A,x}N$. Equivalently, for a displayed idempotent
$e_N\in M_k(A)$, put
\[
 Y_N^0(\operatorname{Spec}S)
 =\{(x,z):x:A\to S,\ z\in S^k,\ e_N(x)z=z\},
 \qquad Y_N=\iota_!Y_N^0\longrightarrow X_A.
\]
This bundle is represented by
$\operatorname{Spec}\Sym_A(N^\vee)$ over $\operatorname{Spec}A$;
in particular it satisfies \'etale descent.  The following
calculation takes place in the primary slice throughout.

\begin{theorem}[Primary finite jets are algebraic principal parts]
\label{thm:primary-jet-coefficients}
Let $j=(\operatorname{Spec}R\hookrightarrow\operatorname{Spec}S)$
and let $x:A\to S$ be a generalized point of $X_A$.
Define
\[
 Q_{j,x,r}=S\otimes_{A,x,s}P^r_{A/B},\qquad
 k_{j,x,r}=(\operatorname{Spec}R\hookrightarrow
                         \operatorname{Spec}Q_{j,x,r}),
\]
where $Q_{j,x,r}\to R$ is reduction after the augmentation to $S$.
Then $k_{j,x,r}$ is a primary test and represents the pullback
$yj\times_{X_A,s_r}\mathcal D_A^{\le r}$.  For the bundle $Y_N$,
\begin{equation}\label{eq:primary-jet-coefficients}
 (\mathbb J_{X_A}^rY_N)(j,x)
 \cong Q_{j,x,r}\otimes_{A,t}N.
\end{equation}
Here the left side denotes the fibre over the generalized point $x$.
Consequently, with source $A$-action,
\begin{equation}\label{eq:primary-jets-principal-module}
 P^r_{A/B}(N):=P^r_{A/B}\otimes_{A,t}N,
 \qquad
 \mathbb J_{X_A}^rY_N\cong Y_{P^r_{A/B}(N)}.
\end{equation}
For $A=B[x_1,\ldots,x_n]$ and $N=A\otimes_BM$ with $M$
displayed finite-projective over $B$, the fibre formula is
\[
 (\mathbb J_{X_A}^rY_N)(j,x)
 \cong\prod_{|\alpha|\le r}(S\otimes_BM).
\]
All comparisons preserve truncation and the full Taylor
coordinate-substitution law of \cref{thm:formal-principal-parts}.
\end{theorem}
\begin{proof}\leavevmode
\begin{enumerate}[leftmargin=*,align=left]
\item \textbf{Construct the representing nilpotent test.}
The algebra $Q_{j,x,r}$ is finitely presented over $B$ by its
finite tensor-product presentation.  Let $L$ be its displacement
ideal and $J=\ker(S\to R)$.  Then $L^{r+1}=0$ and $Q/L=S$.
If $J^a=0$, the kernel $K=\ker(Q\to R)$ satisfies $K^a\subseteq L$,
so $K^{a(r+1)}=0$.  The quotient to $R$ is surjective.
Thus $k_{j,x,r}$ is an object of $\InfB$.

\item \textbf{Identify the pullback of finite principal parts.}
The endpoint pullback of $j$ and
$k_{1,r}=(\operatorname{Spec}A\hookrightarrow\operatorname{Spec}P^r)$
over $\iota(\operatorname{Spec}A)$ has larger algebra $Q$ and
smaller algebra $R\otimes_AA=R$.  It is exactly $k_{j,x,r}$.
The representing property of \cref{lem:primary-principal-representability}
therefore proves the asserted sheaf pullback.

\item \textbf{Compute the primary jet fibre.}
Dependent-product
adjunction and Yoneda give
\[
 \underline{\Hom}_{X_A}(yj,\Pi_{s_r}t_r^*Y_N)
 \cong\underline{\Hom}_{X_A}(yk_{j,x,r},Y_N)
 \cong Q\otimes_{A,t}N,
\]
where the map from $yk_{j,x,r}$ to $X_A$ on the right is its
target vertex.  This proves \eqref{eq:primary-jet-coefficients}.

\item \textbf{Identify the represented principal-parts bundle.}
Tensor associativity identifies this fibre with
$S\otimes_{A,x}P^r_{A/B}(N)$.  The module $P^r_{A/B}(N)$
is a retract of a finite power of $P^r$, which is finite-projective
as a source $A$-module; hence its bundle is defined above and the
fibrewise comparison proves \eqref{eq:primary-jets-principal-module}.

\item \textbf{Recover monomial coefficients and coherence.}
For free coordinates, the truncated monomials in $u=y-x$ are a
finite basis of $Q$ over $S$, giving the product formula.
Every map is induced by a representing pullback or algebra substitution;
these universal properties prove all coherence assertions.
\end{enumerate}
\end{proof}


The coefficient algebra of a jet at $x:A\to S$ is the pushout
\[
\begin{tikzcd}[column sep=large,row sep=large]
 A \arrow[r,"{s}"] \arrow[d,"{x}"']
   & P^r_{A/B} \arrow[d]\\
 S \arrow[r] & Q_{j,x,r}.
\end{tikzcd}
\]
The source action fixes the base point; the target action
$A\xrightarrow{t}P^r_{A/B}\to Q_{j,x,r}$ pulls back the bundle.
Its fibre is therefore $Q_{j,x,r}\otimes_{A,t}N$.
Keeping the two actions visible explains the tensor convention
and the Taylor evaluation at the target.


At the universal target point of $k_{j,x,r}$, a polynomial
section $f$ is evaluated by $f(x+u)$.  The finite binomial identity
therefore gives its primary Taylor coefficients directly:
\begin{equation}\label{eq:primary-polynomial-taylor}
 f(x+u)=\sum_{|\alpha|\le r}
       \frac{u^\alpha}{\alpha!}\partial^\alpha f(x)
       \pmod{(u)^{r+1}}.
\end{equation}
The symmetric-power expression descends along displayed projective
retracts.  Thus primary polynomial Taylor transport is proved on the
representing test, independently of any preservation of dependent
products by restriction to based-Weil probes.

\begin{theorem}[Algebraic differential operators and primary jets]
\label{thm:algebraic-primary-operators}
For $A$-modules $N,M$, let $\operatorname{Diff}_{A/B}^{\le r}(N,M)$
be the $B$-linear maps whose $(r+1)$-fold iterated commutators with
multiplication by elements of $A$ vanish.  There is a natural
isomorphism
\begin{equation}\label{eq:algebraic-primary-operators}
 \operatorname{Diff}_{A/B}^{\le r}(N,M)
 \cong\Hom_A(P^r_{A/B}\otimes_{A,t}N,M),
\end{equation}
where the source $A$-action is used on the tensor product.
For displayed finite-projective $N,M$, these are exactly the
fibrewise-linear primary geometric maps
$\mathbb J^r_{X_A}Y_N\to Y_M$.  Their action on sections is
composition with jet prolongation, and their composition agrees
with the order-additive principal-parts comparisons.
\end{theorem}
\begin{proof}\leavevmode
\begin{enumerate}[leftmargin=*,align=left]
\item \textbf{Translate a linear operator into a tensor map.}
For a $B$-linear $D:N\to M$, define the source-$A$-linear map
$\widetilde D:A\otimes_BN\to M$ by
$\widetilde D(a\otimes n)=aD(n)$.

\item \textbf{Identify iterated commutators with diagonal powers.}
The diagonal generator $\Delta b=1\otimes b-b\otimes1$ acts on
$A\otimes_BN$ through its natural $A\otimes_BA$-action, and
\[
 \widetilde D\bigl(\Delta b\,(a\otimes n)\bigr)
       =a\bigl(D(bn)-bD(n)\bigr).
\]
Induction identifies a product of $r+1$ such generators with the
corresponding iterated commutator.  They generate $I^{r+1}$,
so the commutator condition is precisely that $\widetilde D$
annihilate $I^{r+1}(A\otimes_BN)$.

\item \textbf{Prove the principal-parts universal property.}
The quotient is
$P^r_{A/B}\otimes_{A,t}N$ by the right-exact balanced tensor
presentation.  Conversely any map from that quotient gives
$D(n)=\widetilde D(1\otimes n)$ and satisfies the same condition.
These inverse constructions prove \eqref{eq:algebraic-primary-operators},
including its universal order-$r$ map $n\mapsto1\otimes n$.
This is the internal principal-parts universal property
\citep[Tag~09CH]{Stacks}.

\item \textbf{Identify fibrewise-linear primary bundle maps.}
For finite-projective coefficients, affine-bundle Yoneda identifies
fibrewise-linear bundle maps with module maps: on the symmetric
coordinate algebras, fibrewise linearity selects the degree-one
maps, and finite duality reverses them to the asserted module maps.
Full faithfulness of $\iota_!$ and
\cref{thm:primary-jet-coefficients} now give the primary comparison.

\item \textbf{Identify the universal prolongation.}
On a section $n$, its prolongation evaluates $n$ at the universal
target vertex, hence is $1\otimes n$.

\item \textbf{Compare operator composition and its order bound.}
For two operators, pulling
back to successive universal vertices and then composing evaluates
the same two maps on those vertices.  The algebra substitution
representing composition of the finite correspondences is the
order-additive map of \cref{thm:jet-order-addition}.  Its dependent-product
mate thus gives exactly the same composite operator and the
bound $r+s$, with all associativity equations following from the
three-vertex substitution.
\end{enumerate}
\end{proof}


The operator classification is the following factorization:
every order-$r$ operator is obtained uniquely by applying an
$A$-linear map after the universal $B$-linear jet map.
\[
\begin{tikzcd}[column sep=large,row sep=large]
 N \arrow[r,"{j^r}"] \arrow[dr,"{D}"']
   & P^r_{A/B}(N) \arrow[d,dashed,"{\overline D}"]\\
 & M .
\end{tikzcd}
\]
Here $j^r(n)=1\otimes n$, and $\overline D$ is the unique
$A$-linear arrow with this composite, using the source $A$-action
on principal parts.  The triangle is interpreted in $B$-modules,
with the stronger linearity of its dashed leg specified.


\subsection{Finite formal coefficients and exact Taylor formulas}
Put
\[
 W_{n,r}=\QQ[d_1,\ldots,d_n]/(d_1,\ldots,d_n)^{r+1},
 \qquad D_{n,r}=D_{W_{n,r}},\qquad
 D_{n,\infty}=\bigcup_rD_{n,r}.
\]
This is the based-Weil specialization of the primary construction.
Its polynomial coefficient formulas agree with
\cref{thm:primary-jet-coefficients}; the proposition below also treats
arbitrary $B$-modules and arbitrary natural maps in the based-Weil
category, on that category's own domain.
At a based stage $C$, the last union consists of the nilpotent
$n$-tuples. A tuple with zero augmentation lies in the
nilpotent augmentation ideal and hence in some $D_{n,r}(C)$.
Conversely a nilpotent tuple has zero augmentation because
$B$ is reduced. Therefore
\begin{equation}\label{eq:formal-affine}
 \mathcal D_{\Rline^n}\cong\Rline^n\times D_{n,\infty},
 \qquad (x,y)\longmapsto(x,y-x).
\end{equation}
For the finite relation let $s_r(x,d)=x$ and $t_r(x,d)=x+d$,
and define $\mathbb J_{\Rline^n}^rE=\Pi_{s_r}(t_r^*E)$.

\begin{proposition}[Finite jet coefficients]\label{prop:jet-coefficients}
For every $B$-module $M$, the original based-Weil
coefficient functor, considered in its own slice category,
has the explicit finite coefficients
\[
 \mathbb J_{\Rline^n}^r(\Rline^n\times j_BM)
 \cong
 \Rline^n\times
 \prod_{\substack{\alpha\in\N^n\\|\alpha|\le r}}j_BM.
\]
For a geometric map $F:j_BM\to j_BN$ and directions
$v_1,\ldots,v_n$, its finite Taylor formula is
\[
 F\left(x+\sum_i d_iv_i\right)
 =
 \sum_{|\alpha|\le r}
 \frac{d^\alpha}{\alpha!}
 D^{|\alpha|}F(x)
 (\underbrace{v_1,\ldots,v_1}_{\alpha_1},\ldots,
  \underbrace{v_n,\ldots,v_n}_{\alpha_n})
 \quad(d\in D_{n,r}).
\]
The term of degree zero is $F(x)$.
\end{proposition}
\begin{proof}\leavevmode
\begin{enumerate}[leftmargin=*,align=left]
\item \textbf{Identify finite jets with monomial coefficients.}
The formal-affine trivialization reduces the first assertion
to $(j_BM)^{D_{n,r}}$ and its monomial coefficient basis in
\eqref{eq:weil-coefficients}.

\item \textbf{Compute coefficients using square-zero substitutions.}
For one direction, write its uniquely represented polynomial
as $\sum_{j=0}^r a_jd^j$.
Substitute $d=\eps_1+\cdots+\eps_r$ in the Weil algebra
with $\eps_i^2=0$. Iterating the first-order expansion
identifies the coefficient of $\eps_1\cdots\eps_j$
with $D^jF(x)(v,\ldots,v)$.
In the polynomial expression this coefficient is $j!a_j$.

\item \textbf{Compare truncations and divide by factorials.}
The substitution into $j$ square-zero variables is defined
on the order-$j$ truncation, and truncation identifies its
coefficients with the first $j$ coefficients of the order-$r$
expression. Rational factorials are invertible.

\item \textbf{Recover symmetric mixed derivatives.}
For several directions use disjoint groups of square-zero
variables. The same coefficient comparison gives
$\alpha!$ and the displayed mixed derivative.
Symmetry proved above makes the formula independent of
the order within each group.
\end{enumerate}
\end{proof}

The directed union in \eqref{eq:formal-affine} gives
\[
 \mathbb J_{\Rline^n}^\infty E
 \cong\varprojlim_r\mathbb J_{\Rline^n}^rE.
\]
Indeed maps out of its union are compatible families of
maps out of the finite disks; the dependent-product
adjunction gives the assertion in the slice.
All countable coefficient objects here are internal
dependent products indexed by the NNO.

\subsection{Exact, signed and completed coefficient comparison}
For a positive exact algebra $A$, the truncated polynomial
algebra $A[d]/d^{r+1}$ has underlying module $A^{r+1}$ and
finite convolution multiplication. Hence
\[
 \Gr(A[d]/d^{r+1})
 \cong(\Gr A)[d]/d^{r+1}.
\]
Finite biproducts, multiplication, and the derivative
comparisons give the same finite coefficient compatibility
with $\Ext$. Several variables use their finite monomial
lists. The comparisons commute with truncation and every
substitution represented by these polynomial presentations.
Finite Weil coefficient modules also commute with gauge and
filtered coefficient change by tensor associativity.

At every finite order these comparisons preserve the measured coefficients
of \cref{thm:measured-jet-adequacy}. At infinite order the distinction is
\eqref{eq:infinite-denominators}: the measured family of
\cref{ex:measured-formal-family} has a polynomial budget at each finite
order but no single polynomial budget for all its coefficients.

\subsection{Formal-order completion and its de Rham interface}
\label{sec:jet-derham-interface}
\begin{proposition}[Algebraic crystalline bundles give primary jet coalgebras]
\label{prop:crystalline-jet-coaction}
Let $N$ be a displayed finite-projective module over
$A=\Sym_B(E^\vee)$, represented by $e_N\in M_k(A)$.
Use the affine bundle $Y_N^0\to F_A$ and its primary realization
$Y_N=\iota_!Y_N^0\to X_A$ constructed in
\cref{sec:primary-jet-coefficients}.
Every formal stratification on $N$ gives naturally a
$\mathbb J_{X_A}^\infty$-coalgebra structure
\[
 \gamma_N:Y_N\longrightarrow\mathbb J_{X_A}^\infty Y_N
\]
in the primary \'etale infinitesimal-arrow geometry.
Morphisms of stratifications induce coalgebra morphisms.
Its restriction to based $B$-Weil arrows is exactly the earlier
coefficient-functor bundle and Taylor transport.
\end{proposition}
\begin{proof}\leavevmode
\begin{enumerate}[leftmargin=*,align=left]
\item \textbf{Construct the coefficient bundle and its sheaf.}
The displayed idempotent equation makes $Y_N^0$ functorial
in $S$; its fibre is $S\otimes_{A,\xi}N$.  Since this is
represented by an affine bundle over $\operatorname{Spec}A$,
it satisfies \'etale descent and defines the claimed sheaf.

\item \textbf{Pull back finite stratification transport.}
For a nilpotent arrow $j=(R\leftarrow S)$ and a pair of
$S$-points $x,y$ agreeing in $R$, \cref{thm:completion-etale}
gives a finite $r$ and an algebra map $P^r_{A/B}\to S$.
Pull back the $P^r$-linear stratification to obtain the
fibre isomorphism
\[
 \eps_{x,y}:(Y_N)_x\xrightarrow{\cong}(Y_N)_y.
\]

\item \textbf{Extend transport to the full relation.}
Compatibility with principal-parts truncation shows that
this transport is independent of the exhibited bound;
it defines a map $s^*Y_N\to t^*Y_N$ on the entire
infinitesimal relation.

\item \textbf{Construct the jet coaction and verify its laws.}
Transpose through
$s^*\dashv\Pi_s$ to obtain $\gamma_N$.
The identity pair produces the coalgebra counit equation;
the triple-diagonal stratification cocycle produces
$\delta\gamma_N=\mathbb J(\gamma_N)\gamma_N$.

\item \textbf{Verify functoriality and based restriction.}
For a morphism of stratifications $f$, the equation
$f_y\eps_{x,y}=\eps'_{x,y}f_x$ proves the induced
jet-coalgebra square.  At a based arrow the matrix fibres
and finite transport evaluate on $S=C$ exactly as in the
based-Weil construction.
\end{enumerate}
\end{proof}

\begin{lemma}[Crystalline restriction to a principal open]
\label{lem:principal-open-jets}
Let $A=\Sym_B(E^\vee)$ be a displayed finite-projective polynomial
coordinate algebra and $h\in A$. For its principal open $A_h$, the
finite principal-parts coefficients, formal stratifications, and
crystalline jet coactions have the same comparisons after restriction
to $X_{A_h}$. In particular the shape unit of $X_{A_h}$ is an effective
epimorphism, and a finite-projective $A_h$-module with integrable
connection gives a primary jet coalgebra. Finite and full jet functors
restrict along $X_{A_h}\incl X_A$ by the corresponding slice
Beck--Chevalley maps.
\end{lemma}
\begin{proof}\leavevmode
\begin{enumerate}[leftmargin=*,align=left]
\item \textbf{Lift invertibility through a nilpotent quotient.}
At a nilpotent arrow $S\onto R_0$, an element of $S$ is a unit
exactly when its image in $R_0$ is a unit. To construct its inverse,
lift an inverse modulo the nilpotent kernel; the product differs
from $1$ by a nilpotent element, whose finite geometric series
corrects the lift.

\item \textbf{Identify the restricted infinitesimal relation.}
Thus, for two infinitesimally related $A$-points,
$h$ is invertible at one endpoint exactly when it is invertible at
the other. The full infinitesimal relation over $X_{A_h}$ is therefore
both the source and the target pullback of the relation over $X_A$.
The same statement holds at each finite diagonal order, with the
finite inverse formula of \cref{thm:localized-formal-derham}.

\item \textbf{Construct the jet restriction comparisons.}
These cartesian squares give the slice dependent-product
Beck--Chevalley maps for the finite and full jet correspondences.

\item \textbf{Identify principal-parts bundles and their coactions.}
The finite coordinate algebras are
$A_h\otimes_AP^r_{A/B}$ by
\cref{thm:localized-formal-derham}; their represented bundle
calculation is the finite-dual-basis calculation of
\cref{thm:primary-jet-coefficients}. For a finite-projective
$A_h$-module its displayed idempotent has entries in $A_h$ and
defines the same affine-bundle presentation on this open. A
stratification therefore gives its full-relation transport and its
jet coaction by the transposition of
\cref{prop:crystalline-jet-coaction}.

\item \textbf{Prove effective lifting and restriction coherence.}
Finally, an $A_h$-point of $R_0$ first gives an $A$-point. The
polynomial lifting construction of \cref{prop:shape-lifting-scope}
lifts it locally to $S$; the lifted value of $h$ is a unit by the
preceding inverse construction. Hence the lift factors through $A_h$.
This proves that its shape unit is an effective epimorphism and
establishes the stated descent comparison by \cref{thm:jets}.
All inverse and transport formulas are compatible with truncation
and parameter substitution, proving the claimed restriction laws.
\end{enumerate}
\end{proof}


\begin{theorem}[Formal-order completion of the primary crystalline jets]
\label{thm:completed-jet-comparison}
For $X_E=X_{A_E}$ a finite-projective polynomial resource-coordinate
sheaf in the primary \'etale infinitesimal-arrow geometry,
\begin{equation}\label{eq:completed-jet-comparison}
 \mathbb J_{X_E}^\infty Y\cong
 \varprojlim_r\mathbb J_{X_E}^rY
\end{equation}
naturally in every $Y\to X_E$.  The order-additive
comparison maps of \cref{thm:jet-order-addition} construct the
comultiplication of the jet comonad of
\cref{thm:jets}.  Its jet-coalgebra descent equivalence
holds here by \cref{prop:shape-lifting-scope}.  For finite-projective algebraic coefficients it is the completed
primary principal-parts bundle
\[
 \mathbb J_{X_A}^\infty Y_N
 \cong\varprojlim_rY_{P^r_{A/B}(N)},
\]
with finite coefficients calculated by
\cref{thm:primary-jet-coefficients}.  Algebraic stratifications
construct its crystalline jet coactions by
\cref{prop:crystalline-jet-coaction}.
\end{theorem}
\begin{proof}\leavevmode
\begin{enumerate}[leftmargin=*,align=left]
\item \textbf{Express full jets as a limit of finite jets.}
By \cref{thm:completion-etale}, the full relation is the
internally increasing union
$\mathcal D_{X_E}=\varinjlim_r\mathcal D_E^{\le r}$.
For every object $Z\to X_E$, the universal property of the
union gives compatible maps on its finite stages:
\[
\begin{aligned}
 \Hom_{X_E}(Z,\mathbb J_{X_E}^\infty Y)
 &\cong\Hom_{\mathcal D_{X_E}}(s^*Z,t^*Y)\\
 &\cong\varprojlim_r
       \Hom_{\mathcal D_E^{\le r}}(s_r^*Z,t_r^*Y)\\
 &\cong\Hom_{X_E}
       \left(Z,\varprojlim_r\mathbb J_{X_E}^rY\right).
\end{aligned}
\]

\item \textbf{Identify the natural isomorphism and comultiplication.}
Internal Yoneda proves the natural isomorphism.  Addition
of displacements and ordinary composition of formal pairs
are the same maps in the union, so their dependent-product
mates identify the limit comultiplication with the
comultiplication of \cref{thm:jets}.

\item \textbf{Recover crystalline descent and based coefficients.}
The effective descent
and crystalline coaction clauses are the explicit constructions
cited in the statement.  Restriction to based $B$-Weil
probes recovers the original finite principal-parts maps by the
displayed stagewise coefficient formulas.
\end{enumerate}
\end{proof}

\begin{example}[The measured family as a crystalline resource line]
\label{ex:measured-crystalline-family}
Continue with $R=B[t^{-1}]$, $q=t^{-1}$, and
$A=R[z,(1-qz)^{-1}]=B[z,(t(t-z))^{-1}]$ from
\cref{prop:measured-logarithmic-transport}.
The principal-open comparison of \cref{lem:principal-open-jets}
realizes the line $N=A$ and its integrable connection $d-\omega$ in
the primary resource geometry. Its stratification
$\tau(z,z')=(1-qz)/(1-qz')$ induces the coaction
$Y_N\to\mathbb J_{X_A}^{\infty}Y_N$.

The horizontal section $f=(1-qz)^{-1}$ has finite jet
\begin{equation}\label{eq:measured-crystalline-taylor}
 f(z+u)=\sum_{j=0}^r\frac{q^j}{(1-qz)^{j+1}}u^j
 \pmod{u^{r+1}}.
\end{equation}
This follows by multiplying the finite geometric-series inverse of
$1-qu/(1-qz)$ by $f(z)$. At $z=0$ it is precisely $F_r(u)$,
so \cref{thm:measured-jet-adequacy} identifies this geometric jet
with the measured coefficients of the program $e_r$.
The complete jet is the compatible family of these finite expressions.
Its fibre at zero has coefficients in $R[[u]]$ and is outside
$B[[u]][t^{-1}]$ by \cref{prop:measured-completion-overhead}.

Re-centring from $z$ to $z'$ and then to $z''$ multiplies the
transport factors to $\tau(z,z'')$. Transposition is exactly the
coalgebra coassociativity equation. The logarithm of the same factor
is the completed cocycle \eqref{eq:measured-log-cocycle}, whose de Rham
image is $\omega$. Thus measured coefficients, completed functions,
forms, and crystalline transport are calculated on one resource family.
\end{example}


\begin{corollary}[Finite principal-parts differential operators]
\label{cor:finite-differential-operators}
For a finite-projective polynomial resource cell, an order-$r$
operator on a geometric bundle $Y\to X_E$ with target bundle
$Z\to X_E$ is a map
$\mathbb J_{X_E}^rY\to Z$.  The maps
\eqref{eq:jet-order-addition} compose an order-$s$ operator after
an order-$r$ operator into an operator of order at most $r+s$.
The corresponding complete operators are maps
$\mathbb J_{X_E}^\infty Y\to Z$ and compose by the jet comonad.
Polynomial resource operators have their previously constructed
finite Taylor coefficients under this correspondence.
\end{corollary}
\begin{proof}\leavevmode
\begin{enumerate}[leftmargin=*,align=left]
\item \textbf{Construct composition using order addition.}
The first assertion is the definition via the finite formal
disk adjunction.  For $\alpha:\mathbb J^rY\to Z$ and
$\beta:\mathbb J^sZ\to W$, compose
\[
 \mathbb J^{r+s}Y\longrightarrow\mathbb J^s\mathbb J^rY
 \xrightarrow{\mathbb J^s\alpha}\mathbb J^sZ
 \xrightarrow{\beta}W.
\]
The comparison is \eqref{eq:jet-order-addition} with its indices
interchanged if necessary.

\item \textbf{Verify finite and completed composition laws.}
The finite associativity theorem
proves the operator composition laws; the completed comonad gives
the same argument without an order bound.

\item \textbf{Compare polynomial and fibrewise-linear algebraic operators.}
Primary polynomial agreement is the representing-test Taylor
calculation \eqref{eq:primary-polynomial-taylor} and
\cref{thm:primary-jet-coefficients,thm:polynomial-lift}.
For fibrewise-linear algebraic operators, the universal-property
identification and its composition coherence are
\cref{thm:algebraic-primary-operators}.  Their based-Weil
specializations agree with \cref{prop:jet-coefficients}.
\end{enumerate}
\end{proof}

The completed relation now has both a modal and an algebraic
description. Where the shape unit is an effective epimorphism,
including polynomial finite-projective cells, its jet coalgebras
are identified with $\ThickEt/\InfShape X$.
Within the algebraic coefficient sector,
\cref{thm:connection-stratification} identifies finite-projective
stratifications with integrable $B$-connections.
These comparisons use the same infinitesimal relation but
have different domains; an arbitrary semantic jet coalgebra
need not be an algebraic connection.

The measured family now has an operational coefficient calculation, a
completed function, a de Rham cocycle, and crystalline transport. The
remaining instances relate this geometry to numerical clocks, input
sizes, and the logarithmic connections of asymptotic scales.

\section{Effective instances and geometric calculations}
\label{sec:instances}
\label{sec:recovery}

This section connects the constructions to two kinds of examples.
The numerical execution instance supplies measured costs, input-size
comparisons, and finite resource certificates. The later geometric
calculations exhibit normal resource directions, asymptotic scales,
and logarithmic connections. Each example keeps its specified
implementation family and coefficient context, so its conclusion
can be read directly against the general classification.

We begin with an operational realization of the measured span.
The machine configuration and update map determine its histories;
the first-hit clock then identifies the numerical cost of a
returning execution.

\subsection{The numerical execution instance}
\begin{construction}[Numerical configurations and measured return]
\label{constr:numerical-execution}
In the effective numerical fibre, use the finite register-program syntax
of \cref{prop:standalone-numerical-cover}: increment, zero test with decrement, jumps, and halt.
Finite instruction lists have primitive-recursive codes; invalid codes
denote a permanently looping program. A configuration consists of its
program code, instruction pointer, and finite register data, with the
Boolean halted flag. Reading an instruction and updating these finite
data constructs the total primitive-recursive maps
\[
 h:\mathsf{Conf}\to2,\qquad
 s:\mathsf{Conf}\to\mathsf{Conf},\qquad
 \operatorname{init}:\N\times\N\to\mathsf{Conf},
\]
and an output map $o:h^{-1}(1)\to\N$. Extend $s$ by the identity on
halted configurations. These are the numerical code and configuration
objects specified by these finite instruction clauses.

Define an equation morphism with additive natural cost by
\[
 b:\mathsf{Conf}\to L_{\Sigma,\N}(\N+\mathsf{Conf}),\qquad
 b(c)=\begin{cases}
 \eta(0,\mathsf{inl}\,o(c)),&h(c)=1,\\
 \eta(1,\mathsf{inr}\,s(c)),&h(c)=0.
 \end{cases}
\]
Finite case analysis constructs this map; the discreteness lemma makes
the indicated total data maps computational algebra morphisms.
Then
\[
 \operatorname{ev}_{\mathrm{time}}
   =b^{\ddagger_{\N}}\operatorname{init}:
   \N\times\N\to\LL(\N\times\N)
\]
returns the number of updates before the halted state and the output.
The cost is zero for an initially halted configuration. Tagging each
branch by its current configuration instead constructs
\[
 \operatorname{ev}_H:\N\times\N\to
       \LL(\mathsf{Conf}^*\times\N)
\]
by \cref{thm:weighted-return,thm:history-valuation}. Its word includes
the initial, intermediate, and final configurations. The valuation that
assigns $1$ to a nonhalted configuration and $0$ to a halted one recovers
$\operatorname{ev}_{\mathrm{time}}$.

A fixed total primitive-recursive occupancy function
$\operatorname{occ}:\mathsf{Conf}\to\N$, defined from the chosen finite
register representation, folds these same recorded configurations by
maximum. The product monoid records time and that peak occupancy on one
execution. Other finite measurement expressions built from the recorded
fields, constants, addition, maximum, and finite folds have total
interpretations $\nu:\mathsf{Conf}^*\to\N$ by finite recursion.
Here the expression is selected before forming its implementation image.

The first halting stage defines the partial evaluator of
\cref{prop:standalone-numerical-cover}. Its finite run is exactly
the corresponding return component of \cref{lem:standalone-return},
so erasure has the same value and domain. The regular epimorphism
$\N\onto\mathsf{PC}=\J\mathsf{FG}$ of that proposition supplies
the numerical code presentation. The added time and occupancy values are
the folds just defined. The
oracle-relative instance uses the same relative code construction and
records the declared oracle-invocation events.
\end{construction}

\begin{proposition}[The bounded numerical clock]\label{prop:numerical-clock}
The NNO recursion $\operatorname{run}(e,x,k)=s^k(\operatorname{init}(e,x))$
constructs the monotone Boolean clock
\[
 \alpha(e,x,k)=h(\operatorname{run}(e,x,k)).
\]
Its first-hit domain is the return domain of the preceding execution,
and its unique first-hit index is the time value of
$\operatorname{ev}_{\mathrm{time}}$. Write its domain and index as
\[
 D_\alpha=\{(e,x):\exists k:\N,\ \alpha(e,x,k)=1\}\incl\N\times\N,
 \qquad\tau_\alpha:D_\alpha\to\N.
\]
\end{proposition}
\begin{proof}\leavevmode
\begin{enumerate}[leftmargin=*,align=left]
\item \textbf{Establish monotonicity and define the first-hit test.}
The halted state is absorbing, so the clock is monotone. Define the
Boolean predicate
\[
 \operatorname{First}_\alpha(e,x,k)=
 \alpha(e,x,k)\wedge\bigwedge_{j<k}\neg\alpha(e,x,j).
\]

\item \textbf{Construct the unique first hit.}
Finite induction through a successful stage produces a least one, and
linearity of the NNO order proves uniqueness.

\item \textbf{Identify the classified return domain and its cost.}
The corresponding
subobject therefore projects isomorphically onto
$\{(e,x):\exists k\,\alpha(e,x,k)=1\}$, whose classifier lies in
$\Sigma$. On its $k$th summand there are exactly $k$ update branches
followed by the zero-cost return branch. This is $R_k$ of
\cref{thm:weighted-return}, with the same output and cost $k$.
\end{enumerate}
\end{proof}

\subsection{Numerical cost adequacy}
\begin{theorem}[Numerical cost adequacy]
\label{thm:numerical-cost-adequacy}
In the numerical instance, the measured evaluator erases to the
partial evaluator $\varphi(e,x)$ of \cref{prop:standalone-numerical-cover}, with exactly the same
domain. Its partial time map $T(e,x)$ has that domain and satisfies
\begin{equation}\label{eq:numerical-cost-adequacy}
 T(e,x)=k\quad\Longleftrightarrow\quad
 \operatorname{First}_\alpha(e,x,k)=1.
\end{equation}
The graph predicate on the right is primitive recursive, uniformly
in $e,x,k$. The bounded-return predicate
\[
 \mathsf{Within}(e,x,b)=\alpha(e,x,b)
\]
is primitive recursive and is true exactly when the execution returns
with time at most $b$. In the oracle-effective instance these same
claims hold relative to its fixed oracle.
\end{theorem}
\begin{proof}\leavevmode
\begin{enumerate}[leftmargin=*,align=left]
\item \textbf{Identify erasure and the return domain.}
The first-hit definition of $\varphi$ and \cref{lem:standalone-return}
identify erasure and its domain; the weighted-return comparison
retains the same components.

\item \textbf{Construct the bounded first-hit test.}
The primitive-recursive
configuration step gives primitive-recursive bounded iteration
$\operatorname{run}$ by NNO recursion. Thus $\alpha$ and its finite
first-hit test are primitive recursive.

\item \textbf{Identify the numerical output and measured cost.}
On the exact-return component
$R_k$, the retained history contains $k$ update events of weight $1$
and a final event of weight $0$; its fold is $k$ and its returned
output is $\varphi(e,x)$. Conversely a first hit at $k$ belongs to
that component, by \cref{prop:numerical-clock}. This proves
\eqref{eq:numerical-cost-adequacy}.

\item \textbf{Prove the bounded-time equivalence and oracle version.}
A true $\alpha(e,x,b)$ has a
unique first hit at some $k\le b$ by bounded finite search; a first
hit at $k\le b$ makes $\alpha(e,x,b)$ true by monotonicity. The
oracle proof uses the same bounded simulation relative to the fixed
oracle-query instruction of \cref{prop:standalone-numerical-cover}.
\end{enumerate}
\end{proof}


On the returning domain, the measured and unmeasured evaluators
are related by the commuting triangle
\[
\begin{tikzcd}[column sep=large,row sep=large]
 D_\alpha \arrow[r,"{{(T,\varphi)}}"] \arrow[dr,"{\varphi}"']
   & \N\times\N \arrow[d,"{\pi_2}"]\\
 & \N .
\end{tikzcd}
\]
Both partial maps use the same domain mono
$D_\alpha\incl\N\times\N$ of program--input pairs.
Thus the measured span is obtained from an exact operational
domain-and-value comparison.


\begin{corollary}[The numerical time is a Blum measure]
\label{cor:blum-measure}
For the register-program enumeration of
\cref{prop:standalone-numerical-cover}, the numerical time
$T_e(x)$ is a Blum complexity measure: its domain is the domain
of $\varphi_e$, and the relation $T_e(x)=k$ is decidable uniformly
in $e,x,k$. The same statement holds relative to a fixed oracle.
\end{corollary}
\begin{proof}\leavevmode
\begin{enumerate}[leftmargin=*,align=left]
\item \textbf{Verify the two defining measure properties and their relative form.}
Domain equality and the primitive-recursive first-hit test are
\cref{thm:numerical-cost-adequacy}. They are the two defining
properties of an abstract complexity measure in
\citep{Blum1967MachineIndependentComplexity}.
Oracle-relative bounded simulation gives the relative test.
\end{enumerate}
\end{proof}

This corollary concerns the constructed numerical time; an arbitrary
valuation of histories need not have decidable exact-cost fibres.
For example, assigning zero to every event leaves exact cost zero
on precisely the halting domain.

The classical external interpretation of the numerical resource
space also gives a geometric reading of Blum speedup.
Blum's theorem provides computable functions for which every
implementation admits an asymptotic improvement, with a prescribed
recursive speedup. Choosing an unbounded multiplicative improvement
rules out an $O$-optimal implementation. In the ordinary external
growth locale, \cref{thm:optimization-compactness} then identifies its feasibility
open as non-supercompact: an optimal implementation would generate
that open and contradict speedup.
The numerical measurement throughout this instance counts the updates
of its displayed register machine. Other valuations use their own
history folds, as in \cref{thm:history-valuation}.

Measured execution gives a profile on actual inputs. Complexity
bounds expressed in an input size must be pulled back along the
chosen size map. Keeping that map explicit distinguishes, for
example, polynomial bounds in numerical value from polynomial
bounds in binary length.

\subsection{Input sizes and effective bound profiles}
\begin{proposition}[Size-indexed bounds]
\label{prop:size-indexed-bounds}
An implemented total size map $\ell:X\to\N$ and a displayed total
numerical bound expression $b_\beta:\N\to\N$ construct the raw
profile $g_\beta=b_\beta\ell$ and its exact germ on the supplied
proper-filter context of actual inputs. They therefore construct the
principal admitted budget $\downarrow[g_\beta]$. For numerical inputs,
\begin{equation}\label{eq:size-bounded-clock}
 \forall x:\N\;\alpha(e,x,g_\beta(x))=1
 \quad\Longleftrightarrow\quad
 e\in P\ \text{ and }\ \forall x:\N\;T_e(x)\le g_\beta(x).
\end{equation}
Every such pointwise bound gives $\Phi(e)\in B g_\beta$ and hence
the corresponding upper-feasibility judgment for $p(e)$.
\end{proposition}
\begin{proof}\leavevmode
\begin{enumerate}[leftmargin=*,align=left]
\item \textbf{Construct the size-indexed profile and admitted budget.}
Composition of the retained size and bound descriptions computes the
raw profile; rational inclusion and the germ quotient give its exact
class. The principal admitted-ideal constructor and its name
implementation from \cref{prop:standalone-numerical-cover} give the budget. Its pairing is
$u\mapsto u([g_\beta])$, so admission is explicit.

\item \textbf{Identify the bounded clock with the time inequality.}
For each $x$,
\cref{thm:numerical-cost-adequacy} identifies the bounded clock with
return and the stated time inequality.

\item \textbf{Deduce incidence and the represented-input version.}
Quantifying over $x$ proves
\eqref{eq:size-bounded-clock}; a global bound is eventual with constant
$1$, giving the incidence. For represented inputs the same profile
construction is on $X=C_X$, with no quotienting of their costs.
\end{enumerate}
\end{proof}

Finite-word length and binary bit length are constructed by finite
recursion on the chosen representations. Thus a polynomial in numerical
value uses $\ell(n)=n$, whereas a polynomial in binary input length
uses $\ell(n)=\operatorname{bitlength}(n)$, with the encoding's declared
convention at zero. These are different profiles. On higher-type or
redundantly named inputs the pointwise formulation does not require
finite size fibres or an enumeration of all names of a given size.
An extensional size map is pulled back along $q_X$ when available;
an encoding size may instead depend on the name. Computing a bound
for purposes of checking it also has its own recorded execution cost.

\subsection{A numerical program and two size conventions}
\label{app:complexity-size-example}

The following program makes the distinction concrete. Work in the
register-machine environment of
\cref{constr:numerical-execution,thm:numerical-cost-adequacy}.
Place the natural input $n$ in an input register and $0$ in an
output register. The instruction scheme uses the existing
zero-test with decrement, increment, jump, and halt instructions:
\[
 \begin{array}{ll}
 \mathsf{loop}:&\text{if the input register is zero, go to }\mathsf{halt};\\
 &\text{otherwise decrement the input register;}\\
 &\text{increment the output register twice; jump to }\mathsf{loop};\\
 \mathsf{halt}:&\text{return the output register}.
 \end{array}
\]
Its finite code $e_{2\cdot}$ has denotation
$v(n)=2n$.  For each input $n$ there are exactly $n$ nonzero loop
passes followed by the final zero branch.  Counting the fixed finite
instruction block gives fixed integers $a>0$ and $b\ge0$, independent
of $n$, with
\begin{equation}\label{eq:app-double-cost}
 T_{e_{2\cdot}}(0)=b,\qquad
 T_{e_{2\cdot}}(n+1)=T_{e_{2\cdot}}(n)+a,
 \qquad T_{e_{2\cdot}}(n)=an+b.
\end{equation}
For the literal four-update loop consisting of a test/decrement,
two increments, and a jump, with one final zero test and zero cost for
the halted state, $a=4$ and $b=1$.  If the fixed encoding expands one
of these commands into several primitive updates, its finite expansion
merely changes $a,b$ and not the growth class.
The clock adequacy theorem identifies each first-hit index with this
update count, so \eqref{eq:app-double-cost} is a statement about
\emph{measured execution}.
As $a>0$,
\[
 \kappa(e_{2\cdot})=[n+1].
\]

Write
\[
 I_{\mathrm{lin}}=\downarrow[n+1],\qquad
 I_{\mathrm{poly,val}}
   =\downarrow\{[(n+1)^k]:k\in\N\}.
\]
Both are actual ideal points of $\Cplx$, and the latter is admitted
in the overt-NNO instance: its positive pairing with a monotone
$u:L\to\Sigma$ is $\exists k\,u([(n+1)^k])$, which factors through
$\Sigma$ by \cref{prop:standalone-dominance}.  The displayed code
gives concrete witnesses
\begin{equation}\label{eq:app-double-probes}
 I_{\mathrm{lin}}\models\mathcal U_v,
 \qquad I_{\mathrm{poly,val}}\models\mathcal U_v.
\end{equation}
Indeed $\Phi(e_{2\cdot})\in B(n+1)$, and this same cost belongs to
the larger polynomial module.  The first assertion tests the
principal linear budget; the second tests the nonprincipal ideal of
all polynomial bounds in \emph{numerical input value}.

Now put $\lambda(n)=1+\operatorname{bitlength}(n+1)$, as in
\cref{sec:transport}, and define the different polynomial budget
\[
 I_{\mathrm{poly,bit}}
     =\downarrow\{[\lambda^k]:k\in\N\}.
\]
For every fixed $k$, the inequality
$n+1\le c\lambda(n)^k$ fails on arbitrarily large inputs, as can
be seen along $n=2^m$ from exponential growth against fixed
polynomials in $m$.  Hence
\[
 \kappa(e_{2\cdot})\notin I_{\mathrm{poly,bit}}.
\]
The displayed code therefore meets the polynomial budget for numerical
value and fails it for binary length. The problem's feasible open is
the union over its implementation fibre, as in
\cref{thm:universal-cost}.


\paragraph{The numerical implementation over a budget family.}
We can also evaluate this measured code on the nonconstant budget
family of \eqref{eq:app-generic-budget-family}.
There $b:\mathbb S\to\Cplx$ uses $a_0=[1]$, $a_1=[n+1]$,
and distinguished open $u$.
The one-description presentation containing $e_{2\cdot}$ has
denotation image $\{v\}$ and feasible open $U_{a_1}$.
By \eqref{eq:app-nonconstant-probe}, its pullback is $u$:
the code is feasible exactly where the budget becomes linear.
In the full implementation fibre this same witness proves
$u\le b^*\mathcal U_v$; equality additionally depends on the
other implementations of $v$.

\subsection{Closed resource certificates and semantic incidence}
\label{sec:resource-certification}
Finite proof certificates provide another way to witness a resource
assertion. We use the checker and ordinary-context interpretation of
\cref{app:certificates}, in a fixed observational environment.
This gives a concrete comparison with the cost-aware logics of
\citep{GueneauChargueraudPottier2018Fistful,NiuSterlingGrodinHarper2022Calf}:
the assertions and their proofs are descriptions in the stated
language.

Fix a total bound expression $g_\beta(x)$ built from implemented
numerical constructors, composition, and primitive recursion.
The parameter $\beta$ records its constructor data and may include
a fixed implemented size expression. Constructor quotation represents
$e$ and $\beta$ in the language; arbitrary semantic cost predicates
are not supplied with such quotation.

The bounded simulation and first-hit tests are fixed implemented data
programs. Ordinary identity types over $2$ therefore define
\begin{align}
 \mathsf{Bound}(e,\beta)
   &:=\Pi_{x:\N}\Id_2(\alpha(e,x,g_\beta(x)),1),
       \label{eq:bound-type}\\
 \mathsf{Total}(e)
   &:=\Pi_{x:\N}\Sigma_{k:\N}
       \Id_2(\operatorname{First}_\alpha(e,x,k),1).
       \label{eq:total-type}
\end{align}
For cofinite numerical input define the Boolean program
\[
 \mathsf{Tail}(e,\beta,a,N,x)
   :=(x<N)\vee\alpha(e,x,(a+1)g_\beta(x))
\]
and the ordinary type
\begin{equation}\label{eq:landau-bound-type}
 \mathsf{Land}(e,\beta):=
 \mathsf{Total}(e)\times
 \Sigma_{a:\N}\Sigma_{N:\N}\Pi_{x:\N}
       \Id_2(\mathsf{Tail}(e,\beta,a,N,x),1).
\end{equation}
All of these descriptions use the existing ordinary sums, products,
and Boolean identity types; the notation abbreviates these finite descriptions.

\begin{proposition}[The interpreted resource assertions]
\label{prop:resource-assertion-interpretation}
In the concrete model, the support of $\mathsf{Total}(e)$ is precisely
$P(e)$. The supports of $\mathsf{Bound}$ and $\mathsf{Land}$ are
respectively
\[
 \{(e,\beta):e\in P,\ \forall x\ T_e(x)\le g_\beta(x)\},
 \qquad
 \{(e,\beta):e\in P,\ \Phi(e)\in B g_\beta\}.
\]
Here support means the image of the interpreted type's projection to
its parameter object, and $g_\beta$ in the second expression denotes
its exact germ.
\end{proposition}
\begin{proof}\leavevmode
\begin{enumerate}[leftmargin=*,align=left]
\item \textbf{Interpret totality using the unique first hit.}
At a returning input there is exactly one first-hit index. The
first-hit subobject therefore projects isomorphically onto the return
domain. Unique choice supplies the dependent function in
\eqref{eq:total-type} exactly on $P$;

\item \textbf{Interpret the global bound.}
Boolean identity is interpreted by its equality
subobject, so \cref{thm:numerical-cost-adequacy} proves the first bound
formula.

\item \textbf{Extract eventual domination from a tail witness.}
In \eqref{eq:landau-bound-type}, $\mathsf{Total}$ retains
termination also at the finitely many inputs below $N$. Above $N$,
the Boolean tail test is exactly
$T_e(x)\le(a+1)g_\beta(x)$. Such a witness is an eventual domination
witness.

\item \textbf{Construct a tail witness from eventual domination.}
Conversely a rational domination coefficient is bounded
above by a positive integer $a+1$, and a cofinite extent supplies
its cutoff $N$. These data give the tail factor. The totality factor
was already constructed uniquely on $P$. Taking supports proves
the second formula, uniformly in the data parameters.
\end{enumerate}
\end{proof}

For either ordinary type description $A(e,\beta)$, use the ordinary
inhabitance tag of the mixed checker. Constructor substitution defines
the closed-instance boundary $j^{\mathrm{cl}}_A(e,\beta)$: its
described ordinary and linear telescopes are both empty, and its type
is $A(\bar e,\bar\beta)$. Define, with the pointwise closed-instance
convention of \cref{app:certificates},
\begin{equation}\label{eq:resource-closed-certificates}
 \begin{aligned}
 C^{\mathrm{cl}}_{\mathcal T}[A;e,\beta]
   &:=\Sigma_{d:\mathsf{Tree}}
       \Id_2(\mathsf{Acc}_{\mathcal T,c}
                    (d,j^{\mathrm{cl}}_A(e,\beta)),1),\\
 \mathbb B^{\mathrm{cl}}[A;e,\beta]
   &:=F_d C^{\mathrm{cl}}_{\mathcal T}[A;e,\beta].
 \end{aligned}
\end{equation}
The ordinary boundary tag uses its ordinary checker; the empty linear
boundary is retained as in the mixed presentation. The environment
$c$, included among the fixed supplied parameters, is not quoted as
constructor data. The ambient variables $e,\beta$ index closed
instances; they are not hypothetical premises in the described proof.

\begin{theorem}[Resource-certificate soundness and finite replay]
\label{thm:resource-certificate-soundness}
The interpreted supports of the ordinary certificate families
$C^{\mathrm{cl}}_{\mathcal T}[\mathsf{Bound};e,\beta]$ and
$C^{\mathrm{cl}}_{\mathcal T}[\mathsf{Land};e,\beta]$ lie in the
respective resource incidences of
\cref{prop:resource-assertion-interpretation}. Their denotation images
therefore lie in $\operatorname{Upper}_{g_\beta}$.
The comparison commutes with constructor-data substitution and with
the verified proof transformations of \cref{thm:standalone-checker}.

The numerical first-hit and bounded-return facts have uniform finite
replay: a true test at supplied numerical data constructs a closed
certificate of its Boolean identity at the exact stated root.
The corresponding finite-source resource certificates, and finite
positive-probability witnesses for closed bound certificates, also
map into these semantic incidences.
\end{theorem}
\begin{proof}\leavevmode
\begin{enumerate}[leftmargin=*,align=left]
\item \textbf{Apply soundness to an accepted closed certificate.}
An accepted ordinary root is an actual derivation of the instantiated
type with both described telescopes empty. Apply the concrete
soundness theorem \ref{thm:standalone-checker}.
The interpretation inhabits that instantiated type; its support is
the incidence calculated in
\cref{prop:resource-assertion-interpretation}. This is a semantic
soundness comparison, not a reflection rule in the described theory.

\item \textbf{Identify the feasibility image and compiler stability.}
For bound certificates it factors through $P$ and then through the
upper-bound subobject of $P$; composition with $p$ gives the asserted
denotation image. Soundness is uniform in the retained data, and the
closed-instance instantiation program commutes with their substitution.
Verified proof compilers preserve their specified roots, so the same
soundness calculation applies to their outputs.

\item \textbf{Construct a checked derivation from the numerical trace.}
For finite replay, run the actual bounded numerical test, retaining
its finite configuration trace. On its true branch the verified
fixed-program calculation construction of \cref{lem:standalone-replay} constructs
the derivation of the corresponding Boolean identity. Its exact-root
invariant gives the required closed boundary. It branches on the
computed Boolean result; no runtime acceptance proof or arbitrary
dependent value is quoted.

\item \textbf{Extract the exact return history.}
For a bounded-return test, bounded search
also extracts the unique first hit and its finite return history.

\item \textbf{Forget publication data while preserving acceptance.}
A finite-source resource certificate of \cref{app:certificates} retains the
accepted tree along with its publication data; forgetting that data
gives the ordinary certificate above.

\item \textbf{Extract an accepted tree from positive publication weight.}
A finite witness
$\Theta^{\mathrm{cl}}(q;A;e,\beta)$, with rational $q>0$, contains
a finite list of accepted trees with total weight exceeding $q$.
The list cannot be empty, so finite inspection extracts one accepted
tree. These explicit maps followed by the preceding support comparison
give the remaining assertions. They retain the same closed boundary
and supplied environment.
\end{enumerate}
\end{proof}

Finite replay is complete for the stated finite numerical
observations. The universally quantified $\mathsf{Bound}$ and
$\mathsf{Land}$ assertions retain a different interpretation:
sound finite proofs witness their truth, but truth is not
identified with finite provability. The free quantitative
certificate module likewise has no rule evaluating an arbitrary
certificate into its described type. Contextual and pointwise
closed certificate families keep their respective meanings.
Search has a positive return domain, whereas the universal
assertions have the quantified supports calculated above.
Thus the relation between proof-source probability and execution
time is carried by finite witnesses and semantic incidence.

\subsection{Exact overhead constructors}
An overhead bound must account for the events introduced by its
implementation. The basic operation is concatenation of recorded
words. In the fixed finite-description signature, include the
idle primitive $\operatorname{idle}=\eta_1:1\to\LL1$ with event
$a=\iota_{\mathrm{idle}}(*,*)$.
For a supplied total nonnegative integer $b(x)$, finite recursion
constructs $\operatorname{rep}(a,b(x))$ by that many invocations.
Prefixing it gives the history transformation
\[
 D_b(H)=\operatorname{rep}(a,b(x)){+\!+}H,
 \qquad D_bD_c=D_{b+c}
\]
on a fixed input. For the selected valuation its exact formula is
\begin{equation}\label{eq:delay-measurement}
 T_{\nu,D_b e}(x)=\nu(\operatorname{rep}(a,b(x)){+\!+}H_e(x)).
\end{equation}
Here $D_b e$ denotes the retained execution description that prefixes
these events. For a monoid fold this is
$\overline w(\operatorname{rep}(a,b(x)))\,T_{\overline w,e}(x)$.
For additive time with $w(a)=1$ it is $b(x)+T_e(x)$; for a maximum fold
it is the maximum of the recorded prefix and original occupancies.
The equality follows from the fold clauses, so no additive formula is
assigned to an arbitrary measurement $\nu$.

If $b(x)$ must itself be computed, the execution of that computation is
prefixed too. In a numerical code translation, the actual dispatch and
loop instructions of its finite wrapper also contribute their recorded
events. Thus \eqref{eq:delay-measurement} describes the specified execution
word; a concrete wrapper's bound is obtained by evaluating the valuation
on all of its events. The cost of executing a code transformation and
the execution cost of its output program have different inputs and are
kept separate.

For the bounded clock in \cref{prop:numerical-clock}, a numerical
reindexing is also available. Define
\[
 \beta(e,x,k)=\begin{cases}
 0,&k<b(x),\\
 \alpha(e,x,k-b(x)),&k\ge b(x).
 \end{cases}
\]
Its return extent is unchanged, and its first-hit index satisfies
\begin{equation}\label{eq:delay-firsthit}
 \tau_\beta(e,x)=b(x)+\tau_\alpha(e,x).
\end{equation}
This follows by the two Boolean cases and uniqueness of the first hit.
For a nonnegative integral polynomial $P(x,c)$, the inflationary map
$r_x(c)=c+P(x,c)$ gives the monotone Boolean clock
\[
 \beta_P(e,x,k)=\bigvee_{j\le k}
   (\operatorname{First}_\alpha(e,x,j)\wedge r_x(j)\le k).
\]
A first hit at $j$ becomes a first hit at $r_x(j)$, since $j\le r_x(j)$
and the original first hit is unique. Conversely a hit contains the
original returning witness. This proves equality of extents and the
exact reindexed clock value. The procedure computing this Boolean
predicate has its own execution history; its runtime is not inferred
from the reindexed clock value.

\subsection{Polynomial valuations and compiler comparisons}
For $F(x)=\sum_{j=0}^d b_jx^j$ and a marked resource $M$,
\[
 F(M)\subseteq Bb_0+\sum_{j=1}^d(Bb_j)M^j,\qquad
 DF(x,v)=\sum_{j=1}^d j b_jx^{j-1}v.
\]
The first follows from resource multiplication, and the
second from the differential comparison theorem.
Substituting polynomials commutes with their exact,
quantitative, signed, and geometric interpretations.

Let $\gamma:P\to P'$ be a finite implemented translation with
the measured comparison obtained from its compiled histories,
$T_{\gamma(e)}\le a_eT_e+b_e$ eventually, where
$a_e\in\Qplus$ is a supplied numerical constant (more generally a
supplied nonnegative bounded coefficient) and $b_e$ is an exact
nonnegative remainder profile.
An exact bound $b_e\in O(T_e)$ gives
$T_{\gamma(e)}\in O(T_e)$.
A reverse measured comparison gives $\Theta$.
An exact remainder
$T_{\gamma(e)}-T_e\in o(T_e)$ gives $\sim$.
These implications follow from \cref{prop:landau-arithmetic}
with the implementation parameter retained.


Description measurements use the same calculus with their
own exact profiles. The prefix-free word family $1^n0$ has
length $n+1$. Prefixing a fixed word of length $c$ gives
length $n+1+c\sim n+1$, with error constant $c$.
The execution cost of prefixing is calculated in \cref{app:certificates}.

\subsection{Numerical and oscillating examples}
\begin{example}[Leading coefficients and errors]
Let $g=n+1$. Then $n+2\sim g$, whereas $2g\in\Theta(g)$
and $2g\not\sim g$.
The first difference is $1\in o(g)$; the second is $g$,
which fails its smallness inequality at tolerance $1/2$.
In the orbit description the first multiplier is $1+t$,
and the second is the unit $2$ outside $1+\mathfrak m$.
\end{example}

\begin{example}[Polynomial, logarithmic, and exponential orders]
For integers $0\le r<s$, $(n+1)^r\in o((n+1)^s)$,
by the direct cutoff for $(n+1)^{s-r}$.
The profile $(n+1)\lambda$ is in
$\widetilde\Theta(n+1)$ and in $\omega(n+1)$.
Every fixed polynomial belongs to $o(2^n)$ by the
finite binomial estimate of \cref{sec:scales}.
The parity examples of \cref{sec:landau} give comparisons
which have no single limiting ratio.
\end{example}

\begin{example}[Tight versus asymptotically exact variation]
At a reference resource $g$, first-order variation within
the tight-bound orbit has direction module $Bg$.
Variation within its asymptotic-equivalence class has
direction module $\mathfrak m g$ by
\eqref{eq:orbit-tangents}. For $g=t^{-1}$ these are
$t^{-1}B$ and $t^{-1}\mathfrak m$.
The direction $g$ is allowed by the first and excluded
by the second, whereas the constant direction $1=tg$
belongs to both. This is the differential form of the
leading-coefficient distinction.
\end{example}

The next calculation returns to the geometric interpretation of
resource constraints. It exhibits a direction excluded by a
bounded budget and shows how allowing a chosen overhead removes
that normal obstruction.

\subsection{An explicit normal resource direction}
For cofinite input put $H=t^{-1}B$ and take a fixed
implementation of bounded cost $c=\Phi(e)\in B$.
The valuation family $F_e(b)=c+b$ has budget-$B$ tangent
$B\subset H$ at $b=0$. The direction $t^{-1}$ gives a
nonzero class in $H/B$, since $n+1$ is not bounded.
After localization at $t$, this quotient becomes zero
by \cref{thm:gauge}. This computes an explicit normal
resource obstruction and its change under overhead inversion.

A numerical asymptotic scale organizes resource coefficients by order of decay. Absolute differentials describe logarithmic connections along this scale.

\subsection{Asymptotic expansions, scale filtrations, and connections}
\label{sec:scales}
\subsubsection{The numerical scale and its remainders}
Work in the cofinite numerical fibre and put $t(n)=1/(n+1)$.
It is a regular element of $B$, belongs to $\mathfrak m$,
and is not a unit: its reciprocal $n+1$ is unbounded.
Thus its powers form a strict asymptotic scale,
$t^{r+1}\in o(t^r)$.

\begin{definition}[Asymptotic expansion]\label{def:expansion}
For a rational sequence $(a_j)_{j\in\N}$, the formal expansion
relation (distinct from relative-error equivalence $\sim$) is written
\[
 f\mathrel{\sim_t}\sum_{j\ge0}a_jt^j
\]
and means the family of finite remainder statements
\[
 f-\sum_{j=0}^r a_jt^j\in o(t^r)\qquad(r\in\N).
\]
No value for an infinite sum is part of this definition.
More generally a displayed scale $(\varphi_j)$ with
$\varphi_{j+1}\in o(\varphi_j)$ defines its expansion by
the corresponding remainders. This is the finite-remainder
use of asymptotic expansion \citep{DLMF}.
\end{definition}

\begin{proposition}[Coefficient uniqueness and arithmetic]
\label{prop:asymptotic-expansion}
The rational coefficients of a $t$-expansion are unique.
Expansions add and multiply with the usual finite
coefficient formulas. On the subobject of expansions with
$a_0\in\QQ^\times$, their germs are units of $B$ and their
inverses have the recursively determined reciprocal expansion.
\end{proposition}
\begin{proof}\leavevmode
\begin{enumerate}[leftmargin=*,align=left]
\item \textbf{Exclude nonzero rational constants from the vanishing ideal.}
A nonzero rational constant cannot belong to $o(1)$:
choose $k$ with $k\abs a>1$; its required eventual extent
would be empty, contrary to properness.

\item \textbf{Prove coefficient uniqueness by induction.}
Suppose coefficients agree below degree $r$ in two
expansions of $f$. Subtract their order-$r$ remainder
equations and divide by the everywhere nonzero $t^r$.
Then $a_r-b_r\in o(1)$, hence $a_r=b_r$.
Induction proves uniqueness.

\item \textbf{Verify addition of expansions.}
Addition follows from the small-error sum law.

\item \textbf{Compute the product expansion.}
For multiplication fix order $r$. The order-zero estimates
make both germs bounded. Multiplying the two order-$r$
remainders by bounded germs gives errors in $o(t^r)$.
Terms of the finite product with degree greater than $r$
belong to $t^{r+1}B\subseteq o(t^r)$.
The remaining coefficients are the finite convolution.

\item \textbf{Construct the reciprocal expansion.}
For $a_0\ne0$, the order-zero expansion gives
$f=a_0(1+a)$ with $a\in\mathfrak m$. Hence $f$ is a unit
by \cref{thm:jacobson}. Set $b_0=a_0^{-1}$ and
$b_r=-a_0^{-1}\sum_{j=1}^r a_jb_{r-j}$.
At each finite order, convolution gives
$f\sum_{j=0}^r b_jt^j=1+o(t^r)$.
Multiply the remainder by the bounded inverse of $f$.
This is the reciprocal expansion.
\end{enumerate}
\end{proof}

An expansion need not determine the exact germ.
The rational profile $2^{-n}$ is nonzero but belongs to
$o(t^r)$ for every $r$. Indeed for $k=r+1$ and $n\ge2k$,
\[
 2^n\ge \binom{n}{k}\ge (n/2)^k/k!,
\]
so $(n+1)^r2^{-n}$ is bounded by a rational constant
times $1/n$ eventually. Each tolerance has an explicit
integer cutoff. Thus its expansion has all coefficients
zero, while its exact resource remains visible.

\subsubsection{Filtered coefficients and differentiation}
The filtration $F^rB=t^rB$ records a specified scale of
error. It is finer than declaring an error merely vanishing.
For example
$b(n)=1/(1+\lfloor\sqrt{n+1}\rfloor)$ belongs to $\mathfrak m$
but not to $tB$: square inputs make $b/t$ arbitrarily large.
The floor square root is a bounded finite search, so this
is an exact rational profile.

\begin{proposition}[Scale boundary]\label{prop:scale-boundary}
There are natural isomorphisms and exact sequences
\[
 t^rB/t^{r+1}B\cong B/tB,\qquad
 \gr_{tB}B\cong(B/tB)[z],
\]
\[
 0\longrightarrow M/tM\xrightarrow{\ t\ }M/t^2M
 \longrightarrow M/tM\longrightarrow0
\]
for every embedded resource $M\subseteq K$.
\end{proposition}
\begin{proof}\leavevmode
\begin{enumerate}[leftmargin=*,align=left]
\item \textbf{Identify the associated graded algebra by cancellation.}
Multiplication by $t^r$ is injective on $K$ and hence on $B$
and $M$. It identifies the quotient in degree $r$ with
$B/tB$. Products of these identifications add degrees,
giving the polynomial graded algebra.

\item \textbf{Verify exactness of the scale sequence.}
In the displayed sequence, the first map sends $[m]$ to
$[tm]$. If $tm\in t^2M$, cancellation gives $m\in tM$.
Its image is exactly the kernel of reduction modulo $tM$.
Surjectivity of the last map is quotient surjectivity.
\end{enumerate}
\end{proof}

For a finite polynomial parameter algebra $A=\Sym_B(E)$
and an ideal $J\subseteq B$, relative differentiation gives
\[
 d_{/B}(JA)\subseteq J(A\otimes_BE).
\]
Indeed $d_{/B}(bp)=b\,d_{/B}p$ for $b\in J$.
This proves coefficientwise stability of remainders in
$t^rA$ and $\mathfrak m A$ under polynomial parameter
differentiation. It supplies an explicit differential
calculus for resource-valued asymptotic expansions.

The nilpotent Taylor parameter $d$ of \cref{sec:jets} and
the regular asymptotic scale $t$ have different algebraic
roles. Relative Weil differentiation fixes $t$ as a
coefficient. Absolute coefficient differentiation is
constructed next.

\subsubsection{Logarithmic differentials and pole resources}
Let
\[
 L_t=(\Omega^1_{B/\QQ}\oplus B\ell)/(dt-t\ell),
\]
and write $d:B\to L_t$ for the universal differential
followed by this quotient.
A $B$-linear map $L_t\to N$ is exactly a pair
$(\delta,\omega)$ with $\delta:B\to N$ a $\QQ$-derivation
and $\delta(t)=t\omega$. This follows directly from the
generators and the one displayed relation.

For $k\ge0$ let $M_k=t^{-k}B\subseteq K$.
\begin{proposition}[Logarithmic connections]\label{prop:log-connection}
The formula
\begin{equation}\label{eq:connection}
 \nabla_k(bt^{-k})=t^{-k}\otimes(db-kb\ell)
\end{equation}
defines an $L_t$-connection on $M_k$.
It is compatible with $M_k\incl M_{k+1}$ and with
$M_k\otimes_BM_l\cong M_{k+l}$.
Its residue in the generator $t^{-k}$ is $-k$.
\end{proposition}
\begin{proof}\leavevmode
\begin{enumerate}[leftmargin=*,align=left]
\item \textbf{Verify the connection and Leibniz identity.}
The coefficient $b$ is unique because $t$ is regular.
The derivation identity gives
$\nabla_k(am)=a\nabla_k(m)+m\otimes da$.

\item \textbf{Verify compatibility with inclusions.}
Under the inclusion, $bt^{-k}=tb\,t^{-(k+1)}$.
Substitute $d(tb)=t\,db+b\,dt$ and $dt=t\ell$ into
the right-hand formula at $k+1$; the result is exactly
the formula at $k$.

\item \textbf{Verify tensor compatibility.}
For a product
$bt^{-k}\cdot ct^{-l}$, the Leibniz rule for $d(bc)$
adds the two logarithmic terms, giving the tensor
compatibility.

\item \textbf{Construct the residue map.}
Define $\operatorname{res}:L_t\to B/tB$ by $db\mapsto0$
and $\ell\mapsto1$. The relation maps to zero, so this
is well-defined and factors through $L_t/tL_t$.

\item \textbf{Compute the residue of the resource line.}
Applying it to the coefficient in
\eqref{eq:connection} gives $-k\bar b$.
Under the generator identification of $M_k/tM_k$ with
$B/tB$, this is multiplication by $-k$.
\end{enumerate}
\end{proof}

\begin{corollary}[Residues distinguish logarithmic resource lines]
\label{cor:logarithmic-residue-invariance}
In the cofinite numerical fibre, the residue of $\nabla_k$ is $-k$
in every bounded-unit frame of $M_k=t^{-k}B$. The logarithmic resource
lines $(M_k,\nabla_k)$ and $(M_l,\nabla_l)$ are isomorphic as
$L_t$-connections exactly when $k=l$. Tensor product adds their
residues.
\end{corollary}
\begin{proof}\leavevmode
\begin{enumerate}[leftmargin=*,align=left]
\item \textbf{Show that the residue is independent of the generator.}
Every generator of the free rank-one module $M_k$ has the form
$ut^{-k}$ with $u\in B^\times$. Its connection form is
$u^{-1}du-k\ell$. Since $\operatorname{res}(du)=0$, its residue
is $-k$.

\item \textbf{Use residues to distinguish the indices.}
An isomorphism $M_k\to M_l$ has the form
$t^{-k}\mapsto ut^{-l}$ with $u\in B^\times$.
Compatibility with the connections gives
$du+(k-l)u\ell=0$; applying residue yields
$(k-l)\bar u=0$ in $B/tB$. This quotient is nonzero because $t$
is not a unit, and every nonzero rational remains a unit in it.
As $\bar u$ is also a unit, $k=l$.

\item \textbf{Verify the converse and tensor rule.}
Conversely the identity is an
isomorphism when the indices agree. Tensor compatibility follows
from \cref{prop:log-connection}, and its residues are $-(k+l)$.
\end{enumerate}
\end{proof}


\subsubsection{Principal parts and geometric realization}
Put $A_t=B\ltimes L_t$ and use the two algebra maps
\[
 \lambda(b)=(b,0),\qquad \rho(b)=(b,db).
\]
For a $B$-module $M$ define
$P_t^1(M)=A_t\otimes_{B,\rho}M$, with its left
$B$-module structure induced by $\lambda$.
The right-$B$ sequence
$0\to L_t\to A_t\to B\to0$ splits through $\rho$.
Tensoring therefore gives a left-$B$ exact sequence
\begin{equation}\label{eq:principal-parts}
 0\to L_t\otimes_BM\to P_t^1(M)\to M\to0.
\end{equation}
Its canonical $\QQ$-linear map $j(m)=1\otimes m$ obeys
$j(bm)=b\,j(m)+db\otimes m$.
Consequently a connection $\nabla:M\to L_t\otimes_BM$
gives the $B$-linear splitting $j-\nabla$, and a splitting
$s$ gives the connection $j-s$. These formulas are
mutually inverse. Exact full faithfulness of $j_B$
realizes the sequence and this correspondence inside
$\Geo$. This constructs the differential resource bundle
with its actual extension class and splittings.


\section{Conclusion}
\label{sec:conclusion}
The construction gives a differential geometry of computational
resources with a common exact coefficient interpretation. Measured
execution determines cost profiles and implementations; embedded
resource modules and their locale classify budgets and feasibility.
Taylor-valued history instrumentation identifies finite geometric jets
with measured response coefficients. Resource deformation complexes
then describe constrained variations, including compatibility between
moving budgets and the effect of permitted overheads on obstructions.

The primary infinitesimal geometry represents the function objects
needed for completion, de Rham cochains, and crystalline jets.
The comparison with algebraic forms, the equivalence between
connections and stratifications, and the principal-parts description
of differential operators retain the same resource parameters and
effective finite observations.

The examples exercise these structures together. A measured family has
polynomial bounds at every finite order without a common polynomial
degree for its complete coefficient family. Its generating function
carries a logarithmic de Rham class and a crystalline resource line.
Event instrumentations exhibit a mixed-jet curvature defect and an
explicit coherent ordered transport. Coupled budgets can obstruct a
joint lift despite separate feasibility, and logarithmic residues
distinguish asymptotic pole orders. These calculations connect the
geometric foundations to the measured resources from which they were
constructed.

\section*{Acknowledgements}
\paragraph{Use of AI-assisted tools.} Generative AI systems were used during the research and development of this manuscript. Responsibility for the mathematical content, arguments, and conclusions rests with the author.

\paragraph{Funding.}
This research received no financial support.
\paragraph{Competing interests.}
The author declares no competing interests.

\appendix

\section{Admitted presentations and parameter coherence}
\label{app:presentations}

This appendix supplies the admitted presentation theorem
\cref{thm:standalone-presentation} used to construct the continuous
linear and algebraic objects in the main text.
The external dcpo method is due to Jung--Moshier--Vickers
\citep{JungMoshierVickers2008PresentingDcpos}.
We give its internal selected-ideal construction, including the
parameter comparisons needed by tensors and differentials.

For a preorder take its effective antisymmetrization; lower subsets
and monotone maps into posets descend uniquely. Thus it suffices to
write $P$ for a poset. For $C\subseteq P\times\J P$, let $L_C(P)$
be the object of lower subsets $Z\subseteq P$ satisfying
\[
 (p,I)\in C,\quad I\subseteq Z\quad\Longrightarrow\quad p\in Z.
\]
Intersections define a closure $c_C$ on lower subsets.
The object $L_C(P)$ is a complete lattice, with joins given by
closure of unions, and $\eta_C(p)=c_C(\downarrow p)$ is monotone.
A complete lattice is a $\J$-algebra by restriction of its joins.

\begin{construction}[Generated computational subalgebra]
\label{constr:standalone-generated}
Intersect all subobjects of $L_C(P)$ containing $\eta_C(P)$
and closed under images of their intrinsic admitted ideals.
Denote the intersection by $G$.
For $I\in\J G$, its image in each such subobject is admitted,
and its ambient supremum lies in that subobject. Hence it lies
in $G$. The restricted supremum defines a $\J$-algebra on $G$,
because the least-upper-bound characterization proves its unit
and associativity laws. Set $\pres(P,C)=G$ and $j=\eta_C$ with
this restricted codomain.
\end{construction}

The use of intrinsic weights matters: an ideal admitted on the ambient
complete lattice and supported in a subobject has not thereby been
shown admitted on that subobject.
For a cover $(p,I)\in C$, the ambient join of $\eta_C(x)$,
$x\in I$, contains $I$, so its cover closure contains
$p$ and $\eta_C(p)$. The lower image of $I$ under $j$ is admitted.
It follows that $j$ respects every cover.

\begin{lemma}[Extension and uniqueness]\label{lem:standalone-extension}
Restriction along $j$ gives the order isomorphism
\eqref{eq:admitted-presentation-interface}.
\end{lemma}
\begin{proof}\leavevmode
\begin{enumerate}[leftmargin=*,align=left]
\item \textbf{Embed the target in its lattice of admitted-closed subsets.}
For a target algebra $A$, let $\operatorname{Cl}_{\J}(A)$
be the complete lattice of lower subsets closed under every
admitted ideal of $A$ that they contain.
The order embedding $a\mapsto\downarrow a$ preserves admitted
suprema. Indeed, the least such closed lower subset containing
an admitted ideal $I$ is $\downarrow\sup I$:
closure forces $\sup I$ and then its lower set, while this
principal lower set is itself closed.

\item \textbf{Define the extension on closed lower subsets.}
For a monotone cover-respecting $f:P\to A$, define
\[
 f^\sharp:L_C(P)\longrightarrow\operatorname{Cl}_{\J}(A),
 \qquad f^\sharp(Z)=\bigvee_{p\in Z}\downarrow f(p).
\]

\item \textbf{Verify cover compatibility and join preservation.}
If $W\in\operatorname{Cl}_{\J}(A)$, then $f^{-1}W$ is lower
and cover-closed. In fact, $f(I)\subseteq W$ implies
$\sup\J f(I)\in W$, and the cover inequality puts $f(p)$
in $W$. Consequently closing a union in $L_C(P)$ introduces
no additional image outside any upper bound $W$.
This proves that $f^\sharp$ preserves all joins and that
$f^\sharp\eta_C(p)=\downarrow f(p)$.

\item \textbf{Restrict the extension to the generated algebra.}
The embedded copy of $A$ in $\operatorname{Cl}_{\J}(A)$
is closed under intrinsic admitted suprema. Its inverse image
under $f^\sharp$ is therefore a computational subalgebra
containing every generator, so contains $G$.
Restricting and inverting the order embedding gives
$\bar f:G\to A$, with $\bar f j=f$.

\item \textbf{Prove uniqueness and the order isomorphism.}
For two such maps the equality locus is admitted-closed and
contains all generators, hence is all of $G$.
The same argument with the inequality locus proves the order
isomorphism, including monotonicity of extension.
\end{enumerate}
\end{proof}

\begin{lemma}[Canonical presentations]\label{lem:standalone-canonical}
Every algebra $(A,a)$ has the canonical presentation on its
underlying poset with covers $(a(I),I)$ for $I\in\J A$.
Its presented algebra is canonically $A$.
\end{lemma}
\begin{proof}\leavevmode
\begin{enumerate}[leftmargin=*,align=left]
\item \textbf{Construct the comparison and prove continuity.}
The identity respects the covers, giving $h:\pres(A,C_A)\to A$.
The cover gives $j(aI)\le\sup\J j(I)$; monotonicity gives
the reverse inequality. Thus $j$ is an algebra morphism.

\item \textbf{Verify the inverse equations and naturality.}
We have $hj=\id_A$ on generators, and $jh=\id$ by uniqueness
of extension. The same uniqueness proves naturality.
\end{enumerate}
\end{proof}

\begin{lemma}[Products and operations]\label{lem:standalone-product-generation}
Finite products of generated algebras are generated by the products
of their generators. Cover-compatible finitary operations extend,
and their equations and inequalities extend.
\end{lemma}
\begin{proof}\leavevmode
\begin{enumerate}[leftmargin=*,align=left]
\item \textbf{Generate the product from pairs of generators.}
Let $A,B$ be generated by $G_A,G_B$, and let $D$ be the
least admitted-closed subobject of $A\times B$ containing
$G_A\times G_B$. For $a\in G_A$, the section
$\{b:(a,b)\in D\}$ is admitted-closed, because pairing with
a fixed element preserves inhabited admitted suprema.
It contains $G_B$, hence equals $B$.

\item \textbf{Complete the finite-product induction.}
For arbitrary $b\in B$, the other section now contains
$G_A$ and is admitted-closed, hence equals $A$.
Thus $D=A\times B$. Induction proves the finite-product assertion.

\item \textbf{Define the extended operation.}
For a monotone operation $\omega:P^n\to P$ respecting covers
separately, put
\[
 \widetilde\omega(Z_1,\ldots,Z_n)
 =c_C\downarrow\{\omega(p_1,\ldots,p_n):p_i\in Z_i\}.
\]

\item \textbf{Prove separate continuity.}
Cover compatibility says that the inverse image of a closed lower
subset under a section at generators is cover-closed.
Applying the closure universal property in each variable proves
that $\widetilde\omega$ preserves joins separately.

\item \textbf{Deduce joint continuity and restrict to the generated algebra.}
For a directed family in a finite product, common upper bounds
identify the iterated coordinate suprema with the joint supremum.
Thus separate continuity is joint continuity. Product generation
now shows that $\widetilde\omega$ restricts to $G^n\to G$.

\item \textbf{Extend the algebraic laws.}
Equalities and inequalities between the extended term maps have
admitted-closed loci. They contain the tuples of generators when
the original laws hold, so contain all tuples. This also handles
laws with repeated variables.
\end{enumerate}
\end{proof}

For an algebraic signature with several sorts, use one such lattice
per sort. The same section proof extends each typed finite-arity
operation. Finite term objects exist by encoding well-founded finite
trees using finite lists and natural-number recursion.
Closing covers under one-hole contexts is an internal relation on
these terms; its weights remain admitted by functoriality of $\J$.
Equations are imposed by inequalities in both directions before
antisymmetrization.

For the tensor presentation of \cref{prop:standalone-linear}, the
relations are precisely additivity, zero, and
$(sm)\otimes n=m\otimes(sn)$, together with the two domain-continuity
and scalar-continuity covers. A map out of it is therefore a
continuous balanced bilinear map. Iterated tensors represent the
same multilinear maps, so associativity and symmetry are canonical.
The pentagon, triangles, and symmetry equations hold on elementary
tensors and then everywhere by
\cref{lem:standalone-product-generation}.
A relative tensor adds the balance relations for its coefficient
algebra. A quotient module adds its displayed equations.
For differentials these are additivity, scalar linearity, and the
Leibniz relations; the presentation consequently represents
continuous derivations. This constructs the signed quotients used in the differential sector.

\begin{proposition}[Parameter coherence]\label{prop:standalone-parameters}
The presentation, free-module, tensor, and quotient universal
properties above hold in each parameter slice. Their induced
comparison maps commute with substitution, and their coherence
diagrams commute.
\end{proposition}
\begin{proof}\leavevmode
\begin{enumerate}[leftmargin=*,align=left]
\item \textbf{Construct the slice substitution comparisons.}
Every carrier is a specified subobject of a power object, every
cover retains its parameters, and every extension is determined by
a unique factorization. For a substitution $\sigma:\Delta\to\Gamma$, pullback between
the slices of $\E$ preserves finite limits and regular images.
The slice power-object and dependent-product Beck--Chevalley
comparisons give
\[
 \sigma^*\mathcal P_\Gamma(P)\cong
 \mathcal P_\Delta(\sigma^*P),\qquad
 \sigma^*(Y^X)\cong(\sigma^*Y)^{\sigma^*X}.
\]
Here all objects and exponentials on the left are in $\E/\Gamma$,
and those on the right are in $\E/\Delta$.

\item \textbf{Transport admission and its monad structure.}
The same comparisons transport the predicates for inhabitedness,
directedness, lower closure and the complete upper-test profile.
They therefore identify
$\sigma^*(\J^\Gamma P)\cong\J^\Delta(\sigma^*P)$,
including its unit, multiplication and admitted cover relation.

\item \textbf{Transport cover closure and generated presentations.}
Universal quantification over the slice power object transports
the intersections defining cover closure and the generated
subalgebra. Hence the resulting presentation comparison
\[
 \sigma^*\pres^\Gamma(P,C)\cong
 \pres^\Delta(\sigma^*P,\sigma^*C)
\]
sends each pulled-back generator to the corresponding generator.

\item \textbf{Verify extension and substitution coherence.}
The extension universal property supplies its inverse and shows
that extension commutes with this comparison. Identity and composite substitutions, tensor
coherence, and substitution through extension agree on those
generators and hence everywhere. These comparisons are formed in
parameter slices of the fixed ambient topos.
\end{enumerate}
\end{proof}


\section{Support, representatives, and finite denominator witnesses}
\label{app:support}
The resource calculations use formulas on representatives and
then descend them to germs. This appendix records the support
operations, common principal stages, and finite denominators
that justify those steps. The pointwise operations are rational
formulas with decidable cases:
\[
 e_g(x)=
 \begin{cases}1,&g(x)\ne0,\\0,&g(x)=0,\end{cases}
 \qquad
 q(f,g)(x)=
 \begin{cases}f(x)/g(x),&g(x)\ne0,\\0,&g(x)=0.\end{cases}
\]
Changing either representative on a noneventual complement
changes these formulas only there. Intersecting the
representative-equality extents proves that they descend.
Under a domination witness $f\preceq g$, the equation
$q(f,g)g=f$ holds on the same eventual comparison extent.

For a finite list $f_1,\ldots,f_r$ in a signed resource,
the nonnegative germ $g=\sum_i\abs{f_i}$ belongs to that
resource and dominates the list. Each $f_i$ is a bounded
multiple of $g$. This is the common principal stage used
to verify finite tensor relations and exactness. The
stage is constructed from the finite list; no simultaneous
selection of principal generators for arbitrary submodules
is involved.

A quotient argument with local representatives is an
internal argument over the effective epimorphism from
profiles to germs. To define a function on the quotient
we prove its formula invariant under the equivalence
relation. To prove existence of a multiplier, we construct
it over a representative and then use the image predicate
of its defining equation. This proves internal existence.

Similarly, a finite list of localization equations has a
common denominator given by the product of its denominators.
For equations involving a fixed scale $t$, take a common
maximum exponent. These finite operations justify the
flatness and finite-cover arguments.
The common-denominator condition for an infinite coefficient
family is retained as an existential statement about one
exponent, as in \eqref{eq:infinite-denominators}; it is not
obtained by taking a maximum of infinitely many witnesses.


\section{Finite Weil duality and coefficient coherence}
\label{app:weil}

The linear adjoints and coefficient-change comparisons in the
main text use finite-projective duality. We record the transpose
formulas and check their independence of the chosen dual basis,
then verify compatibility with stage and coefficient maps.
Let $C$ be a finite projective $B$-algebra, $F$ a $C$-module,
and $M$ a $B$-module. For a dual basis $(c_i,\phi_i)$,
the transpose of a balanced map
$h:C^\vee\otimes_CF\to M$ is
\[
 \widetilde h(x)=\sum_i c_i\otimes h(\phi_i\otimes x).
\]
For $\psi\in C^\vee$,
\[
 (\psi\otimes1)\widetilde h(x)
 =h\left(\sum_i\psi(c_i)\phi_i\otimes x\right)
 =h(\psi\otimes x).
\]
The same identity with the right $C$-action on $C^\vee$
shows $\widetilde h(cx)=c\widetilde h(x)$.
It also proves independence of the basis: every contraction
with a functional has the specified value, and finite
duality detects equality in $C\otimes_BM$.
This gives the inverse equations used in
\eqref{eq:linear-coend}.

For a stage map $u:C\to C'$, contracting the naturality
equation of a family $v_C:F(C)\to C\otimes_BM$ with
$\phi\in(C')^\vee$ gives
\[
 h_C((\phi u)\otimes x)
 =
 h_{C'}(\phi\otimes F(u)x).
\]
These are exactly the coequalizer relations defining
$\Lamlin$. Thus its adjunction is natural in both the
coefficient module and the geometric module.

For a coefficient map $\beta:B\to B'$ and a rational
Weil presentation $W$, tensor associativity gives
\[
 (B'\otimes_{\QQ}W)\otimes_{B'}(B'\otimes_BM)
 \cong
 B'\otimes_B\bigl((B\otimes_{\QQ}W)\otimes_BM\bigr).
\]
On a pure tensor both sides have the same coefficient,
Weil basis element, and resource element.
The diagrams for two coefficient maps, two Weil
substitutions, truncation, and interchange therefore
commute on these generators. This proves the finite
coefficient coherence used for gauge and input transport.


\section{History substitution and numerical clock regrading}
\label{app:clocks}

Two distinct comparisons are used in the execution examples:
reassociating a recorded history preserves its fold, while
regrading a numerical clock changes its index. The following
calculations establish their composition laws and keep their
respective measurements explicit.
\subsection{Coherence of the retained resource valuation}
For a monoid $(C,e,\mu)$, an event map $w:A\to C$, and returning words
$s,t,u\in A^*$, the recursive fold satisfies
\[
 \overline w(s{+\!+}t)=\overline w(s)\overline w(t),\qquad
 \overline w((s{+\!+}t){+\!+}u)
 =\overline w(s{+\!+}(t{+\!+}u)).
\]
The first equation is induction on $s$; the second follows from word
associativity and monoid associativity. Empty words give the two unit
equations. Each interface-cut transformation in
\cref{thm:weighted-elgot} preserves the flattened word, so these formulas
prove its resource coherence. Input substitution precomposes the state,
event, and valuation maps in these same formulas; induction on the
finite descriptions proves its compatibility with composition and
recursive execution.

For a finite-history measurement $\nu$ that is not a monoid fold, the
assembled cost remains $\nu(H_{\mathrm{assembled}})$. No monoidal law
is deduced for $\nu$. A richer finite accumulator may instead construct
its stated law. For the sequential time convention, an overhead
execution of cost $H(e,x)$ returning a waiting amount $b(e,x)$ contributes
$H(e,x)+b(e,x)+T_e(x)$ together with every recorded dispatch transition.
This is the sum of the actual component words' time folds.

For two verified measured comparisons
$T_2\le aT_1+b$ and $T_3\le cT_2+d$, on the intersection of their
eventual extents,
\[
 T_3\le caT_1+(cb+d).
\]
The comparison retains the same implementation parameter and the
nonnegative coefficients. For general monotone growth transformations,
their composition is the one in \cref{constr:measured-compilers}.

\subsection{The bounded numerical calculation}
For the numerical clock $\alpha(e,x,k)$ and a supplied inflationary
regrading $r_x:\N\to\N$ with its finite Boolean tests, set
\[
 R_r\alpha(e,x,k)=\bigvee_{j\le k}
    (\operatorname{First}_\alpha(e,x,j)\wedge r_x(j)\le k).
\]
This is monotone and Boolean by its finite formula. A first hit of
$\alpha$ at $j$ becomes the first hit of $R_r\alpha$ at $r_x(j)$;
inflationarity puts $j$ inside the bounded search and uniqueness excludes
another contributor. A hit of the new clock contains a first-hit witness
of the original, so their classified extents agree.

For inflationary $r_x,s_x$, the two clocks
\[
 R_s(R_r\alpha)=R_{s\circ r}\alpha
\]
are true at exactly the indices at least $s_x(r_x(j))$ when the original
first hit is $j$. Existence and uniqueness of that $j$ prove the equality
in every parameter context. The identity regrading gives the original
monotone clock. Input substitution acts on the same bounded formula,
so commutes with regrading. This is a clock-index calculation for the
constructed numerical instance; measured compiler costs are supplied
by the retained histories and their valuations above.

\section{Name presentations, finite descriptions, and numerical codes}
\label{app:names-codes}
The implementation construction requires names with exact
domain comparisons, finite typed descriptions, and a numerical
program presentation. This appendix supplies those three
interfaces from the realizability and exact-completion
constructions of
\citep{Hyland1982EffectiveTopos,BirkedalOosten2002,VanOosten2008Realizability}.
Costs are subsequently assigned by the execution histories
and valuations of the main text.

\begin{proposition}[Names and tracked lifts]\label{prop:standalone-name-lifting}
Every object in the displayed realizability models has a regular-epimorphic
partitioned-assembly presentation. On a retained projective source,
an arrow has a tracked lift through a target name cover.
Partial arrows have implementations with the same classified domain.
Finite products and substitutions admit partitioned refinements;
equations hold after realization and on common refinements.
\end{proposition}
\begin{proof}\leavevmode
\begin{enumerate}[leftmargin=*,align=left]
\item \textbf{Lift total arrows through projective presentations.}
Use the projective cover supplied by the exact-completion
presentation \citep[Corollaries 10.6.4--10.6.5]{Birkedal2000TypesComputability}.
For $q:C\onto X$, $q':C'\onto Y$, and $f:X\to Y$,
pull back $q'$ along $fq$. The resulting regular epimorphism onto
the projective $C$ splits. Composing its section with the other
projection gives $\widehat f:C\to C'$ and
$q'\widehat f=fq$. Fullness of the partitioned-assembly embedding
supplies its single permitted tracker.

\item \textbf{Construct a projective presentation of the hit domain.}
For a partial arrow, pull its domain back to $C$ and lift its
$\Sigma$-classifier through $2^\N\onto\Sigma$. Let $\alpha_c$
be the resulting total stream. Its hit assembly
$H=\{(c,n):\alpha_c(n)=1\}$ is partitioned and projective.
The semantic value on $H$ therefore lifts to $C'$.

\item \textbf{Implement the partial lift by first-hit search.}
On input $c$, search for the least hit and apply that lift to
$(c,n)$. Search terminates exactly on the original domain;
the returned name realizes the original value. It defines
a classified partial arrow by the functional first-hit relation.

\item \textbf{Verify the domain-sensitive implementation.}
The first-hit search is implemented by represented natural-number
recursion with delayed recursive calls. Every value used after
a hit has its stated domain witness, so no test for arbitrary
pca application definedness is used.

\item \textbf{Construct dependent lifts and common refinements.}
Finite products of partitioned presentations are partitioned.
For a dependent pullback, cover the actual pullback by a partitioned
assembly and perform the same construction there. Two presentations
have the common refinement obtained by covering their fibre product
over the denotation object.

\item \textbf{Verify coherence after realization.}
Identity and composition of retained
lifts satisfy their squares by composition; different selected lifts
need only agree after realization. This is precisely the coherence
asserted, and requires no functorial selection of one name per
denotation.
\end{enumerate}
\end{proof}

For a modest object one can use its usual partial-equivalence-relation
presentation on ambient names: the cover consists of the names in
the domain of the PER, and its kernel pair is the PER.
An extensional tracker descends through this pair; a morphism lifts
by the preceding proposition. The same statement for a relation
retains a name for its membership proof. An inclusion of represented
relations lifts after pulling back their proof-name covers, so
internal witness transformations used by finite algorithms can
be implemented on those actual presentations.

\begin{construction}[Typed finite descriptions]\label{constr:standalone-syntax}
For a fixed finite typed computational signature, let $\mathcal S$
contain its constructor tags, arities, source and target sorts,
and supplied parameter slots. A relative raw tree is a finite
prefix list of constructor nodes, with its finite binding and
child-position data, together with the tuple of supplied literal
labels required at its leaves. The valid-tree condition is the
finite stack-parsing condition: each node consumes one pending place
and opens the number specified by its tag; the stack is empty
exactly at the end. Thus the carrier is a subobject of a sum of
finite products of label objects and numerical data. It exists
using the NNO and finite lists and has a partitioned presentation
in the realizability models.

Type annotations record the types at every node, including the
types of its operation slots. The well-typed subobject is defined
by the constructor typing equations. Semantic equalities of
arbitrary literal labels are not used as Boolean grammar tests.
An environment object $\Theta$ retains those labels and operation
parameters. Its well-typed descriptions over $\theta:\Theta$
form the objects $E_{X,Y,\theta}$.
\end{construction}

\begin{lemma}[Interpretation and substitution]\label{lem:standalone-syntax}
Finite recursion constructs the interpretation of the displayed
descriptions in the computational function objects, and tracked
interpretations on their name presentations. Interpretation
commutes with finite typed substitution. Pure structural maps,
composition, branching, sequential pairing, and recursive execution
have actual constructors.
\end{lemma}
\begin{proof}\leavevmode
\begin{enumerate}[leftmargin=*,align=left]
\item \textbf{Interpret the finite description tree.}
A leaf is its retained parameter or primitive; a node applies
the designated constructor morphism to the interpretations of
its children. Products and exponentials of
\cref{thm:standalone-domain}, lifted strength, and
\cref{prop:standalone-recursion} give those morphisms.
A recursion node denotes the least solution of its body, rather
than a recursive definition of the tree interpreter itself.
Evaluate the finite tree from its leaves to its root. This
constructs its semantic interpretation with the stated type.

\item \textbf{Implement the constructors on names.}
On names, retain the primitive implementations and their
classified-domain squares from
\cref{prop:standalone-name-lifting}. Composition runs the first
implementation and then the second on the returned name;
branching inspects its finite tag; a recursive execution runs
its body until its exit tag is returned. The exact finite-return
construction supplies the latter partial map.

\item \textbf{Prove domain and value adequacy.}
Induction on
the number of transitions identifies its domain and returned
denotation with the semantic iteration. Structural induction
then proves the same for the whole finite description.

\item \textbf{Verify identity and composite substitution.}
Substitution replaces slot leaves by their designated trees and
shifts bound indices by the depth of the enclosing context.
At a leaf the interpretation equation is its definition.
At a compound node it follows from the equation for the children
and the same constructor morphism. This proves identity and
composite substitution, including the iteration node.

\item \textbf{Retain parameters and construct code-output descriptions.}
The algorithms traverse only finite trees; arbitrary supplied
names remain labels. A code-output extension of the signature
uses these same trees as data sorts, and its identity and
composition constructors provide the compiler descriptions
used in \cref{constr:measured-compilers}.
\end{enumerate}
\end{proof}

For pure syntax, all literal numerical data are constructor strings
over a finite alphabet. Prefix serialization and stack parsing
therefore give total primitive-recursive maps
$\operatorname{ser}:\mathsf{Tree}\to\N$ and
$\operatorname{parse}:\N\to\mathsf{Tree}$, with
$\operatorname{parse}\operatorname{ser}=\id$ (invalid words parse
to a fixed dummy tree). Enumerating finite words up to a given
length gives finite lists $E(s)$ containing every tree at some
stage. Consequently pure trees are overt.
This assertion concerns pure descriptions, not a numerical
enumeration of arbitrary relative literal carriers.

\subsection{The numerical program presentation}
Let $\mathsf{FG}$ be the decidably ordered poset of finite functional
graphs in $\N\times\N$, ordered by inclusion. Its empty graph is
least. An ideal of $\mathsf{FG}$ determines a partial function by
taking its union, and every partial function determines its ideal
of finite subgraphs. Directedness proves functionality and
reconstructs every finite subgraph; these operations are inverse.

\begin{lemma}[Recognizing admitted finite-information ideals]
\label{lem:standalone-searchable}
For an overt poset $P$ with $\Sigma$-classified order, an ideal
belongs to $\J P$ exactly when its membership predicate is
uniformly $\Sigma$-valued.
\end{lemma}
\begin{proof}\leavevmode
\begin{enumerate}[leftmargin=*,align=left]
\item \textbf{Recover uniform membership from admitted pairings.}
The positive upper test $u_p(q)=(p\le q)$ has pairing
$\chi_I(u_p)\leftrightarrow I(p)$, by lower closure and reflexivity.
Admission therefore gives membership uniformly in $p$.

\item \textbf{Construct admitted pairings from membership.}
Conversely, uniform membership and an upper test have
$\Sigma$-valued conjunction; overtness of $P$ gives their
existential pairing uniformly in the test.
\end{enumerate}
\end{proof}

\begin{proposition}[Numerical codes and extensional denotation]
\label{prop:standalone-numerical-cover}
In $\Eff$, finite register programs construct a partial evaluator
$\varphi(e,x)$ and a regular epimorphism
\[
 q:\N\onto\mathsf{PC}=\J\mathsf{FG},\qquad
 q(e)=\{g:g\subseteq\operatorname{graph}(\varphi_e)\}.
\]
Its kernel pair is extensional equality of partial functions.
Specialization and sequential composition have total recursive
code constructors. The construction relativizes to a fixed oracle.
\end{proposition}
\begin{proof}\leavevmode
\begin{enumerate}[leftmargin=*,align=left]
\item \textbf{Construct the register-program enumeration.}
Encode finite instruction lists with instructions for increment,
zero-test with decrement, jumps, and halt. A configuration stores
the program, instruction pointer, and finite register array.
The total primitive-recursive transition reads one instruction,
updates that finite data, and leaves a halted configuration fixed.
Invalid codes are assigned a fixed looping configuration.
Primitive recursion constructs bounded runs. Minimizing the
first halting stage and reading the output defines $\varphi$.

\item \textbf{Implement the partial-recursive constructors.}
Conversely, zero, successor, projections, composition, primitive
recursion, and minimization have register-program implementations:
use finite register blocks for intermediate values, counted loops
for recursion, and an increasing search loop for minimization.
Thus this is the ordinary partial-recursive enumeration underlying
$K_1$ \citep{VanOosten2008Realizability}.

\item \textbf{Recognize finite subgraphs and prove admission.}
For fixed $e,g$, the total stream whose stage-$s$ value tests
whether every pair in $g$ is verified by an $s$-step run of $e$
has hit support $g\subseteq\operatorname{graph}(\varphi_e)$.
The empty graph is accepted immediately.
The preceding lemma proves that $q(e)$ is admitted, uniformly in $e$.

\item \textbf{Construct a uniform numerical code from a name.}
To prove epicity, cover $\J\mathsf{FG}$ by a partitioned assembly.
Its uniform membership classifier on that cover times $\mathsf{FG}$
lifts to a total stream $\theta(a,g,k)$.
On input $x$, search all graph codes and hit stages by finite
diagonals until a hit is found for a graph containing $(x,y)$,
and return $y$. Directedness makes different possible discoveries
agree; every pair in the represented partial function is eventually
found. The search is one program with the supplied point name $a$
as parameter, hence compiles uniformly into a numerical code.
Its composite with $q$ is the original cover, proving that $q$
is regular epic.

\item \textbf{Identify extensional equality and its name comparison.}
Equality of the finite-subgraph ideals is exactly equality of
their unions. Uniformly, a finite convergence witness for one
code can be transported to an extensionally equal code by
running the latter on the finitely many named inputs until
the same outputs appear. This gives the equality comparison
on names as well as the semantic kernel-pair equation.

\item \textbf{Implement parameter specialization.}
Specialization prefixes instructions loading the fixed numerical
parameter and then runs the original code on its paired input.

\item \textbf{Implement composition and oracle relativization.}
Composition renames the scratch registers of two finite programs,
runs the first, copies its output to the second input block,
and runs the second. These are total recursive operations on
instruction lists; their return domains and values follow from
the two executions. An oracle-query instruction uses the same
finite syntax and finite witnessed transitions relative to that
oracle, so all comparisons relativize.
\end{enumerate}
\end{proof}

\begin{lemma}[Diagonal nontermination is not positive]
\label{lem:standalone-halting}
In the numerical effective model, the complement of
$H=\{e:\varphi_e(e)\downarrow\}$ is not $\Sigma$-classified.
\end{lemma}
\begin{proof}\leavevmode
\begin{enumerate}[leftmargin=*,align=left]
\item \textbf{Recognize diagonal termination positively.}
$H$ is classified by its bounded-run hit stream.

\item \textbf{Derive a total decision from positive nontermination.}
A $\Sigma$-classifier of its complement lifts on the projective
numerical presentation to a total recursive hit stream.
Interleaving the two searches would decide $H$.

\item \textbf{Obtain the diagonal contradiction.}
The register program which loops on input $e$ when that decision
says $e\in H$, and halts otherwise, has an index $d$ by
\cref{prop:standalone-numerical-cover}. Evaluating at $d$
contradicts either Boolean answer. This proves the claim.
\end{enumerate}
\end{proof}

The numerical and relative name comparisons above concern
denotations and exact partial domains. Costs are supplied by
the retained instruction events and valuations in
\cref{sec:retained-execution,constr:numerical-execution}.
Neither an extensional quotient nor projective lifting identifies
the costs of its different representatives.


\section{The finite resource-certificate calculus}
\label{app:certificates}
The certificate comparison in \cref{sec:resource-certification}
uses a finite language whose proofs can be checked and interpreted.
We construct it from the indexed adjunction of
\cref{prop:standalone-lnl}, specifying ordinary and linear
families before the checking and replay procedures.
As in \citep{NiuSterlingGrodinHarper2022Calf,GrodinNiuSterlingHarper2024Decalf},
costs and behavioral assertions have separate semantic roles.
The annotation on a linear judgment is its augmentation grade;
an execution-time interpretation requires the measured comparison.

\subsection{Ordinary and linear families}
Ordinary types over $\Gamma$ are the split cartesian families of
\cref{prop:standalone-lnl}. The universe-free ordinary syntax has
$0,1,2,\N$, finite lists and trees, dependent $\Pi,\Sigma,\Id$,
and Boolean type case. Its terms have variables, the corresponding
constructors and eliminators, and function symbols for finite
descriptions of total data programs. Such programs are built from
constructors, composition, Boolean case, and parameterized structural
recursion, together with a finite supplied family of total observation
maps $a_i(c,-)$; $c$ remains an environment parameter.
The defining equations for a program are read from its finite
description, so their recognition does not require executing an
arbitrary program.
The scalar object $S$ is also an ordinary sort, with rational constants
and its specified finite-operation symbols. Scalar conversions require
displayed equality derivations. Arbitrary semantic scalars are neither
constructor data nor assumed to have decidable equality.

The ordinary rules are formation, introduction, elimination, their
computation equations, congruence, symmetry and transitivity of
equality, and conversion by a displayed equality derivation.
For $\Pi$ these are dependent abstraction, application, and their
beta and eta equations; for $\Sigma$, dependent pairing and
elimination with the corresponding equations; for $\Id$, reflexivity
and identity elimination with its reflexivity equation.
The data sorts have their constructor and structural-recursion
rules. Boolean type case has
$\operatorname{If}(0;A,B)=A$ and
$\operatorname{If}(1;A,B)=B$, including substitution.
There is no equality-reflection or semantic-truth rule.

These rules are interpreted by composition and dependent adjoints in
$\E/\Gamma$, and diagonals for identity. To specify split Boolean
case, decompose a parameter object by the two fibres of its Boolean
term and glue the two displays. Choose the constant-zero case to
be its first display and the constant-one case to be its second;
over an initial base both are the identity display. The unique
cartesian comparisons give the substitution equations.
This external normalization introduces no Boolean test for equality
of semantic types.

For the linear fragment put
$G_bM=\{m:\aug_M(m)=b\}$, and interpret
\[
 \Gamma;\Delta\vdash_b t:M,\qquad
 \Delta=x_1:M_1,\ldots,x_k:M_k,
\]
by a continuous linear map $\bigotimes_iM_i\to M$ satisfying
$\aug_Mt=b\,\aug_{\otimes_iM_i}$.
All displayed linear types and $b$ depend only on the ordinary
context. The empty tensor is $I=(S,\id)$.
The linear type formers are
\[
 I,\quad 0_\ell,\quad M\otimes N,\quad M\multimap N,\quad
 F_dA,\quad \Sigma^\ell_{a:A}M(a),\quad \Pi^\ell_{a:A}M(a),
\]
and $G_bM$ is an ordinary type.

The dependent linear sum is presented by generators $[a,m]$,
linearity in $m$, and the continuity covers in each fibre;
its augmentation is $\aug[a,m]=\aug_{M(a)}m$.
The dependent product is the continuous submodule
\begin{equation}\label{eq:standalone-linear-product}
 \left\{(r,m):r:S,\ m:\prod_{a:A}UM(a),\
       \forall a\;\aug_{M(a)}m(a)=r\right\},\qquad
 \aug(r,m)=r.
\end{equation}
The scalar coordinate is retained even if $A$ is empty.
Consequently the empty sum is $0_\ell$ and the empty product is $I$.
Maps out of the sum are the corresponding families of maps;
maps into the product are the corresponding families with one
common augmentation. These facts prove the dependent adjunctions.
Tensor Frobenius sends $[a,m]\otimes n$ to $[a,m\otimes n]$,
with inverse given by the opposite formula. The presentation theorem
and pointwise product formula prove Beck--Chevalley and coherence.
Defining these operations componentwise on split cartesian families
gives their substitution comparisons.

The following rule schemes specify the linear fragment; every
conversion mentioned in them is a supplied finite equality tree.
A variable and the unit introduction have grade $1$.
Tensor introduction uses disjoint linear contexts and multiplies
grades. Linear abstraction preserves the grade; application
uses disjoint contexts and multiplies grades.
Tensor and unit elimination use the same context splitting and
grade multiplication. There is the unique unit-grade discard
$\epsilon_M:M\to I$ and the initial-type eliminator from $0_\ell$.
Weakening is its explicit elaboration
$\operatorname{weak}_x(t)=\text{let }*=\epsilon_Nx\text{ in }t$.
There is no unrestricted contraction of a linear variable.
Addition uses the same input context in both premises and adds
their grades; scalar multiplication multiplies the grade;
the zero map has grade $0$.

For the ordinary--linear interface, an ordinary $a:A$ gives
$[a]:F_dA$ at empty linear context and grade $1$.
A term $t:M$ at empty linear context and grade $b$ gives
$\operatorname{up}_b(t):G_bM$, with inverse
$\operatorname{down}_b$.
Eliminating $F_dA$ uses its generator and linear universal property:
from a major premise of grade $c$ and a body of grade $b$,
the result has grade $cb$; the other context, result type,
and body grade are pulled back from the context before introducing
the bound ordinary variable. Dependent linear sum introduction
uses a supplied ordinary index and the grade of its component;
its eliminator has the same independence and multiplication rules.
Dependent product abstraction preserves a grade pulled back from
the outer context, and evaluation at an ordinary index preserves it.
Include their beta and eta equations, typed exchange, congruence,
the semimodule equations with their displayed grades, and identity
transport into linear families with its reflexivity equation.
These rules are exactly the adjunction, tensor, and dependent
universal-property maps just constructed.

\begin{lemma}[Graded substitution]\label{lem:standalone-graded-substitution}
Ordinary substitution and graded linear substitution are finite
operations on derivations. Substituting a grade-$c$ term for a
linear variable in a grade-$b$ judgment yields grade $bc$.
They have finite derivations of their typed identity and composition
equations, and preserve concrete interpretation.
\end{lemma}
\begin{proof}\leavevmode
\begin{enumerate}[leftmargin=*,align=left]
\item \textbf{Define substitution at binders and variables.}
Use de Bruijn indices and recurse simultaneously on formation,
typing, and equality trees. At a binder, shift indices and substitute
in the type annotations and conversion derivations.
At a variable, insert the supplied derivation.

\item \textbf{Preserve grades at split-context nodes.}
At a split-context node, transform the unique premise containing
the substituted variable and insert the associativity and
commutativity equations for the grades.

\item \textbf{Preserve additive and zero grades.}
At addition, transform both premise trees and insert
$(b_1+b_2)c=b_1c+b_2c$; at zero use $0c=0$.

\item \textbf{Execute the discarded argument and retain its grade.}
At weakening of the substituted variable the output is
\[
 (\operatorname{weak}_x t)\langle u/x\rangle
   =\text{let }*=\epsilon_Mu\text{ in }t.
\]
The discarded argument is executed as the stated linear morphism;
its grade $c$ remains, and unit elimination has grade $cb$,
converted to $bc$.

\item \textbf{Handle the remaining rule templates.}
At weakening of another variable rebuild
that node after the recursive substitution. The scalar and
dependent eliminator cases use their displayed scalar and
substitution equations. These are finite rule templates.

\item \textbf{Prove the substitution equations.}
The identity and composition equations follow by the same
induction. At weakening, unit beta, eta, and commuting conversions
identify the iterated discards; at addition use distributivity
and congruence. These are equality derivations, not literal
equalities of different proof trees.

\item \textbf{Verify preservation of interpretation.}
The interpretation of each step is composition, evaluation,
or the relevant dependent universal property, which proves
preservation of interpretation.
\end{enumerate}
\end{proof}

\subsection{Annotated checking and calculation replay}
Pure sort, type, term, context, judgment, and derivation descriptions
are finite constructor trees. A derivation node stores its exact
root, rule tag, premise trees, context splits, and all conversion
proofs. Distinct tags distinguish ordinary and linear judgments,
formation, typing, and equality. An observation leaf stores
$(i,\vec x,b)$ with constructor data $\vec x$ and Boolean $b$;
its acceptance locus is the result of actually testing
$a_i(c,\vec x)=b$.
The environment $c$ is supplied to the checker and is not quoted.

\begin{theorem}[Finite checker and concrete soundness]
\label{thm:standalone-checker}
There is a total finite data program
$\mathsf{Check}_{\mathcal T,c}(d,j):\mathsf{Tree}\times\mathsf{Tree}\to2$
whose accepted trees are exactly the annotated derivations at the
stated root $j$, on their observation loci.
Accepted derivations have the stated interpretations in the
indexed ordinary and graded linear model.
Renaming, ordinary substitution, graded linear substitution,
application, and their finite compositions are implemented proof
transformations with verified exact roots.
\end{theorem}
\begin{proof}\leavevmode
\begin{enumerate}[leftmargin=*,align=left]
\item \textbf{Check the root and inference annotations.}
Parse both trees and check that the stored root of $d$ is literally
$j$. At a node compare its annotations with the finite rule schema,
check the context splits and binder indices, and recursively check
its proper premise subtrees.

\item \textbf{Check supplied equality and program equations.}
A conversion node must supply an
equality derivation; a scalar-equation leaf must be an instance of
a semiring law or a proved equation from its defining finite
program. It never tests equality of arbitrary scalar denotations.
The program-symbol scheme checks its finite source description
and the appropriate defining equation.

\item \textbf{Check observations and establish totality.}
An observation leaf makes its designated total observation call
and compares with the recorded Boolean. Malformed trees return
zero. This is structural recursion with total observation calls,
so it is a total data program.

\item \textbf{Prove checker exactness.}
Induction on a derivation proves acceptance. Conversely induction
on an accepted tree reconstructs each inference from its checked
premises and exact annotations. At an observation leaf the actual
Boolean test supplies its claimed validity locus. This proves
checker exactness.

\item \textbf{Interpret the rules and their grades.}
For soundness, interpret contexts by comprehension, ordinary
substitution by pullback, ordinary rules by their dependent
universal properties, and linear rules by the maps above.
At an application or eliminator, augmentation grades multiply;
at addition they add. The weakening case is composition with
discard and retains its grade.

\item \textbf{Verify equality, observations, and semantic soundness.}
The semimodule and adjunction
equations interpret the equality rules; observation leaves are
valid precisely on their tested loci. Rule induction proves
soundness, simultaneously with equality and substitution.

\item \textbf{Verify the proof-transformer invariants.}
The explicit transformations of
\cref{lem:standalone-graded-substitution} rebuild these same finite
rule templates with their exact annotations; node induction
proves their checker-correctness invariants.
\end{enumerate}
\end{proof}

\begin{lemma}[Fixed-program calculation]\label{lem:standalone-replay}
For a fixed total data program $f:X\to D$ on the displayed
constructor data sorts there is a total finite program
$\mathsf{calc}_f(x)$ constructing an accepted closed derivation of
\[
 \Id_D\bigl(\mathsf f(\bar x),\overline{f(x)}\bigr)
\]
at that exact root, with the supplied environment retained.
The root and acceptance identities have finite proofs in the
same calculus.
\end{lemma}
\begin{proof}\leavevmode
\begin{enumerate}[leftmargin=*,align=left]
\item \textbf{Quote constructor data and initialize structural induction.}
Constructor quotation is recursion on the data value; it does
not quote an arbitrary semantic element or the environment.
Induct on the finite description of $f$.
For a constructor, use its defining equation and reflexivity.

\item \textbf{Construct the certificate for composition.}
For composition, recursively calculate the intermediate data,
then use congruence and transitivity.

\item \textbf{Construct the certificate for branching.}
For a Boolean case, calculate its selector and use the equation
for the selected branch.

\item \textbf{Construct the certificate for recursion.}
For structural recursion, recurse on the immediate smaller data,
apply the defining successor equation, and combine the induction
certificates.

\item \textbf{Construct the certificate for an observed answer.}
The observation case calls $a_i$ and emits the
leaf with its actual answer; its check is reflexive.

\item \textbf{Verify exact roots and checker acceptance.}
Each node template has its exact annotated root, and its checker
equation reduces to the comparisons of those annotations and
the premise checks. Induction gives both invariants as finite
derivations using the data eliminators.

\item \textbf{Apply the construction to fixed proof-processing programs.}
This also applies to a fixed checker or proof-transformer program:
it proves the Boolean calculation of that program, not the
conclusion of an arbitrary accepted derivation by a reflection rule.
\end{enumerate}
\end{proof}

An exact judgment and an inhabitance boundary have distinct tags.
At an ordinary boundary $(\gamma;a)$, extract the term from an
ordinary typing root and check that exact root in context $\gamma$.
At a linear boundary $(\gamma;\delta;m;b)$, do the same while
retaining the linear context and grade. A mismatching tag returns
zero. Denote these total tests, together with exact-root checking,
by $\mathsf{Acc}_{\mathcal T,c}(d,j)$.
For the resource types, constructor substitution of $e,\beta$
constructs $j_A^{\mathrm{cl}}(e,\beta)$ with both described contexts
empty. The ambient parameters still index the family of these
closed instances. The ordinary certificate family and its free
quantitative object are therefore exactly
\[
 C^{\mathrm{cl}}_{\mathcal T}[A;e,\beta]
 =\sum_{d:\mathsf{Tree}}
   \Id_2(\mathsf{Acc}_{\mathcal T,c}
              (d,j_A^{\mathrm{cl}}(e,\beta)),1),
 \qquad
 \mathbb B^{\mathrm{cl}}[A;e,\beta]=F_d C^{\mathrm{cl}}_{\mathcal T}[A;e,\beta].
\]
Their supports and denotation images are ordinary regular images.

A finite-source certificate retains such an accepted tree, a finite
source-program description, its finite generation trace, and its
declared publication fields. Its forgetful map returns the tree.
A finite positive-probability witness
$\Theta^{\mathrm{cl}}(q;A;e,\beta)$ retains a finite list of such
accepted trees, their nonnegative rational weights, and the
checked rational inequality that their sum exceeds $q>0$
(and the source trace establishing those weights when specified).
Its list is nonempty; returning its first entry gives the
certificate map used in the soundness theorem.
Neither construction requires a probability-to-time identification.

Finally, finite search through the lists $E(s)$ of pure trees
has support $\exists d\,\mathsf{Acc}_{\mathcal T,c}(d,j)=1$,
by totality of the checker and exhaustiveness of that enumeration.
The support is $\Sigma$-classified, uniformly in $j,c$.
This identifies certificate search with certificate existence.


\begin{thebibliography}{99}
\raggedright

\bibitem[Abramsky and Jung(1994)]{AbramskyJung1994DomainTheory}
Abramsky, S. and Jung, A. (1994).
Domain Theory.
In \emph{Handbook of Logic in Computer Science}, vol. 3. Oxford University Press, pp. 1--168.
\url{https://www.cs.ox.ac.uk/publications/publication1830-abstract.html}.

\bibitem[Aschenbrenner et~al.(2017)]{AschenbrennerVanDenDriesVanDerHoeven2017Transseries}
Aschenbrenner, M., van den Dries, L. and van der Hoeven, J. (2017).
\emph{Asymptotic Differential Algebra and Model Theory of Transseries}.
Princeton University Press.
\href{https://doi.org/10.1515/9781400885411}{doi:10.1515/9781400885411}.

\bibitem[Avigad and Donnelly(2004)]{AvigadDonnelly2004ONotation}
Avigad, J. and Donnelly, K. (2004).
Formalizing {O} Notation in {Isabelle/HOL}.
In \emph{Automated Reasoning: IJCAR 2004}, Lecture Notes in Computer Science, vol. 3097. Springer, pp. 357--371.
\href{https://doi.org/10.1007/978-3-540-25984-8_27}{doi:10.1007/978-3-540-25984-8\_27}.

\bibitem[Behrend(2005)]{BehrendDR}
Behrend, K. (2005).
On the de Rham cohomology of differential and algebraic stacks.
\emph{Advances in Mathematics} \textbf{198}, 583--622.
\href{https://doi.org/10.1016/j.aim.2005.05.025}{doi:10.1016/j.aim.2005.05.025}.

\bibitem[Berthelot and Ogus(1978)]{BerthelotOgus}
Berthelot, P. and Ogus, A. (1978).
\emph{Notes on Crystalline Cohomology}.
Princeton University Press.

\bibitem[Birkedal(2000)]{Birkedal2000TypesComputability}
Birkedal, L. (2000).
\emph{Developing Theories of Types and Computability via Realizability}.
Electronic Notes in Theoretical Computer Science, vol. 34.
Elsevier.
\href{https://doi.org/10.1016/S1571-0661(05)80642-5}{doi:10.1016/S1571-0661(05)80642-5}.

\bibitem[Birkedal and van Oosten(2002)]{BirkedalOosten2002}
Birkedal, L. and van Oosten, J. (2002).
Relative and modified relative realizability.
\emph{Annals of Pure and Applied Logic} \textbf{118}(1--2), 115--132.
\href{https://doi.org/10.1016/S0168-0072(01)00122-1}{doi:10.1016/S0168-0072(01)00122-1}.

\bibitem[Blum(1967)]{Blum1967MachineIndependentComplexity}
Blum, M. (1967).
A Machine-Independent Theory of the Complexity of Recursive Functions.
\emph{Journal of the ACM} \textbf{14}(2), 322--336.
\href{https://doi.org/10.1145/321386.321395}{doi:10.1145/321386.321395}.

\bibitem[Blute et~al.(2020)]{BCS}
Blute, R.~F., Cockett, J.~R.~B., Pacaud Lemay, J.-S. and Seely, R.~A.~G. (2020).
Differential categories revisited.
\emph{Applied Categorical Structures} \textbf{28}, 171--235.
\href{https://doi.org/10.1007/s10485-019-09572-y}{doi:10.1007/s10485-019-09572-y}.

\bibitem[Blute et~al.(2016)]{BLWON}
Blute, R., Lucyshyn-Wright, R.~B.~B. and O'Neill, K. (2016).
Derivations in codifferential categories.
\emph{Cahiers de Topologie et G\'eom\'etrie Diff\'erentielle Cat\'egoriques}
\textbf{57}(4), 243--279.
\url{https://cahierstgdc.com/index.php/volumes/volume-lvii-2016/}.

\bibitem[Dal~Lago and Hofmann(2011)]{DalLagoHofmann2011RealizabilityComplexity}
Dal~Lago, U. and Hofmann, M. (2011).
Realizability models and implicit complexity.
\emph{Theoretical Computer Science} \textbf{412}(20), 2029--2047.
\href{https://doi.org/10.1016/j.tcs.2010.12.025}{doi:10.1016/j.tcs.2010.12.025}.

\bibitem[Danner et~al.(2015)]{DannerLicataRamyaa2015CostSemantics}
Danner, N., Licata, D.~R. and Ramyaa (2015).
Denotational cost semantics for functional languages with inductive types.
In \emph{Proceedings of the 20th ACM SIGPLAN International Conference on Functional Programming}. Association for Computing Machinery, pp. 140--151.
\href{https://doi.org/10.1145/2784731.2784749}{doi:10.1145/2784731.2784749}.

\bibitem[Escard{\'o} and Knapp(2017)]{EscardoKnapp2017}
Escard{\'o}, M.~H. and Knapp, C.~M. (2017).
Partial Elements and Recursion via Dominances in Univalent Type Theory.
In \emph{CSL 2017}, LIPIcs \textbf{82}, 21:1--21:16.
\href{https://doi.org/10.4230/LIPIcs.CSL.2017.21}{doi:10.4230/LIPIcs.CSL.2017.21}.

\bibitem[Giordano and Luperi~Baglini(2016)]{GiordanoLuperiBaglini2016AsymptoticGauges}
Giordano, P. and Luperi~Baglini, L. (2016).
Asymptotic gauges: Generalization of {Colombeau} type algebras.
\emph{Mathematische Nachrichten} \textbf{289}(2--3), 247--274.
\href{https://doi.org/10.1002/mana.201400278}{doi:10.1002/mana.201400278}.

\bibitem[Gomes and Selman(2001)]{GomesSelman2001AlgorithmPortfolios}
Gomes, C.~P. and Selman, B. (2001).
Algorithm portfolios.
\emph{Artificial Intelligence} \textbf{126}(1--2), 43--62.
\href{https://doi.org/10.1016/S0004-3702(00)00081-3}{doi:10.1016/S0004-3702(00)00081-3}.

\bibitem[Grodin et~al.(2024)]{GrodinNiuSterlingHarper2024Decalf}
Grodin, H., Niu, Y., Sterling, J. and Harper, R. (2024).
{Decalf}: A Directed, Effectful Cost-Aware Logical Framework.
\emph{Proceedings of the ACM on Programming Languages} \textbf{8}(POPL), 273--301.
\href{https://doi.org/10.1145/3632852}{doi:10.1145/3632852}.

\bibitem[Grothendieck(1966)]{GrothendieckDeRham}
Grothendieck, A. (1966).
On the de Rham cohomology of algebraic varieties.
\emph{Publications Math\'ematiques de l'IH\'ES} \textbf{29}, 95--103.

\bibitem[Gu{\'e}neau et~al.(2018)]{GueneauChargueraudPottier2018Fistful}
Gu{\'e}neau, A., Chargu{\'e}raud, A. and Pottier, F. (2018).
A Fistful of Dollars: Formalizing Asymptotic Complexity Claims via Deductive Program Verification.
In \emph{Programming Languages and Systems: ESOP 2018}, Lecture Notes in Computer Science, vol. 10801. Springer, pp. 533--560.
\href{https://doi.org/10.1007/978-3-319-89884-1_19}{doi:10.1007/978-3-319-89884-1\_19}.

\bibitem[Hochschild et~al.(1962)]{HKR}
Hochschild, G., Kostant, B. and Rosenberg, A. (1962).
Differential forms on regular affine algebras.
\emph{Transactions of the American Mathematical Society}
\textbf{102}, 383--408.

\bibitem[Hyland(1982)]{Hyland1982EffectiveTopos}
Hyland, J.~M.~E. (1982).
The Effective Topos.
In \emph{The L. E. J. Brouwer Centenary Symposium}, Studies in Logic and the Foundations of Mathematics, vol. 110. North-Holland, pp. 165--216.
\href{https://doi.org/10.1016/S0049-237X(09)70129-6}{doi:10.1016/S0049-237X(09)70129-6}.

\bibitem[Jung et~al.(2008)]{JungMoshierVickers2008PresentingDcpos}
Jung, A., Moshier, M.~A. and Vickers, S.~J. (2008).
Presenting Dcpos and Dcpo Algebras.
\emph{Electronic Notes in Theoretical Computer Science} \textbf{218}, 209--229.
\href{https://doi.org/10.1016/j.entcs.2008.10.013}{doi:10.1016/j.entcs.2008.10.013}.

\bibitem[Kawamura and Cook(2012)]{KawamuraCook2012Operators}
Kawamura, A. and Cook, S. (2012).
Complexity Theory for Operators in Analysis.
\emph{ACM Transactions on Computation Theory} \textbf{4}(2), article 5.
\href{https://doi.org/10.1145/2189778.2189780}{doi:10.1145/2189778.2189780}.

\bibitem[Khavkine and Schreiber(2027)]{KS}
Khavkine, I. and Schreiber, U. (2027).
Synthetic geometry of differential equations:
I. Jets and comonad structure.
\emph{Journal of Geometry and Physics} \textbf{231}, article 106025.
\href{https://doi.org/10.1016/j.geomphys.2026.106025}{doi:10.1016/j.geomphys.2026.106025}.
Preprint \href{https://arxiv.org/abs/1701.06238}{arXiv:1701.06238}
(2017; revised 26 October 2026).

\bibitem[Knuth(1976)]{Knuth}
Knuth, D.~E. (1976).
Big Omicron and big Omega and big Theta.
\emph{ACM SIGACT News} 8(2), 18--24.
\href{https://doi.org/10.1145/1008328.1008329}
{doi:10.1145/1008328.1008329}.

\bibitem[Kock(2006)]{Kock}
Kock, A. (2006).
\emph{Synthetic Differential Geometry}, second edition.
London Mathematical Society Lecture Note Series 333.
Cambridge University Press.

\bibitem[Luperi~Baglini and Giordano(2017)]{LuperiBagliniGiordano2017CategoryColombeau}
Luperi~Baglini, L. and Giordano, P. (2017).
The category of {Colombeau} algebras.
\emph{Monatshefte f{\"u}r Mathematik} \textbf{182}(3), 649--674.
\href{https://doi.org/10.1007/s00605-016-0990-1}{doi:10.1007/s00605-016-0990-1}.

\bibitem[Mac~Lane and Moerdijk(1992)]{MacLaneMoerdijk}
Mac~Lane, S. and Moerdijk, I. (1992).
\emph{Sheaves in Geometry and Logic: A First Introduction to
Topos Theory}.
Springer, Chapter IX (Localic Topoi).

\bibitem[Manuell(2022)]{ManuellSpectrum}
Manuell, G. (2022).
\emph{Quantalic spectra of semirings}.
\href{https://arxiv.org/abs/2201.06408}{arXiv:2201.06408}.

\bibitem[Miller and Rahmati(2024)]{MillerRahmati}
Miller, C. and Rahmati, H. (2024).
Transferring algebra structures on complexes.
\emph{Journal of Homotopy and Related Structures}
\textbf{19}, 561--596.
\href{https://doi.org/10.1007/s40062-024-00356-3}
{doi:10.1007/s40062-024-00356-3}.

\bibitem[NIST(2026)]{DLMF}
National Institute of Standards and Technology (2026).
\emph{Digital Library of Mathematical Functions},
Section 2.1, Definitions and elementary properties.
\url{https://dlmf.nist.gov/2.1}. Accessed 28 October 2026.

\bibitem[Niu et~al.(2022)]{NiuSterlingGrodinHarper2022Calf}
Niu, Y., Sterling, J., Grodin, H. and Harper, R. (2022).
A cost-aware logical framework.
\emph{Proceedings of the ACM on Programming Languages} \textbf{6}(POPL), article 9.
\href{https://doi.org/10.1145/3498670}{doi:10.1145/3498670}.

\bibitem[Niu et~al.(2024)]{NiuSterlingHarper2024CostSensitiveSDT}
Niu, Y., Sterling, J. and Harper, R. (2024).
Cost-sensitive computational adequacy of higher-order recursion in synthetic domain theory.
\emph{Electronic Notes in Theoretical Informatics and Computer Science} \textbf{4}.
\href{https://doi.org/10.46298/entics.14732}{doi:10.46298/entics.14732}.

\bibitem[Schellekens(1995)]{Schellekens1995SmythCompletion}
Schellekens, M. (1995).
The {Smyth} Completion: A Common Foundation for Denotational Semantics and Complexity Analysis.
\emph{Electronic Notes in Theoretical Computer Science} \textbf{1}, 535--556.
\href{https://doi.org/10.1016/S1571-0661(04)00029-5}{doi:10.1016/S1571-0661(04)00029-5}.

\bibitem[Smyth(1977)]{Smyth1977EffectiveDomains}
Smyth, M.~B. (1977).
Effectively given domains.
\emph{Theoretical Computer Science} \textbf{5}(3), 257--274.
\href{https://doi.org/10.1016/0304-3975(77)90045-7}{doi:10.1016/0304-3975(77)90045-7}.

\bibitem[Stacks Project Authors(2026a)]{Stacks}
Stacks Project Authors (2026a).
\emph{The Stacks Project}.
Tags \href{https://stacks.math.columbia.edu/tag/05UT}{05UT}
(filtered colimits of flat modules),
\href{https://stacks.math.columbia.edu/tag/00NX}{00NX}
(finite projective modules),
\href{https://stacks.math.columbia.edu/tag/09CH}{09CH}
(finite-order differential operators), and
\href{https://stacks.math.columbia.edu/tag/01X9}{01X9}
(standard-open \v{C}ech exactness).
Accessed 28 October 2026.

\bibitem[Stacks Project Authors(2026b)]{StacksEtale}
Stacks Project Authors (2026b).
\emph{The Stacks Project}, \S\S 41.15--41.16
(infinitesimal invariance and the functorial characterization
of \'etale morphisms), and \S 59.45 (topological invariance
of the small \'etale site).
\url{https://stacks.math.columbia.edu/tag/04DY}.
Accessed 28 October 2026.

\bibitem[Stacks Project Authors(2026c)]{StacksArtin}
Stacks Project Authors (2026c).
\emph{The Stacks Project}, \S 94.12 (algebraic stacks),
\S 94.17 (algebraic stacks associated to smooth groupoids),
and \S 37.38 (smooth morphisms and \'etale local sections).
\url{https://stacks.math.columbia.edu/tag/02BK}.
Accessed 28 October 2026.

\bibitem[van Oosten(2008)]{VanOosten2008Realizability}
van Oosten, J. (2008).
\emph{Realizability: An Introduction to its Categorical Side}.
Elsevier.

\bibitem[van Oosten and Simpson(2000)]{VanOostenSimpson2000SyntheticDomains}
van Oosten, J. and Simpson, A.~K. (2000).
Axioms and (counter)examples in synthetic domain theory.
\emph{Annals of Pure and Applied Logic} \textbf{104}(1--3), 233--278.
\href{https://doi.org/10.1016/S0168-0072(00)00014-2}{doi:10.1016/S0168-0072(00)00014-2}.

\bibitem[Vickers(1989)]{Vickers}
Vickers, S. (1989).
\emph{Topology via Logic}.
Cambridge Tracts in Theoretical Computer Science 5.
Cambridge University Press.

\end{thebibliography}
\end{document}

